\documentclass[a4paper,11pt,openany]{ltjsbook}
\usepackage{luatexja}
\usepackage{luatexja-fontspec}
\usepackage{amsmath,amssymb,amsfonts,amsthm}
\usepackage{geometry}
\geometry{top=24mm,bottom=24mm,left=24mm,right=24mm,footskip=10mm}
\usepackage{tcolorbox}
\tcbuselibrary{skins,breakable}
\usepackage{hyperref}
\usepackage{enumitem}
\usepackage{booktabs}
\usepackage{fancyhdr}

\definecolor{vennblue}{RGB}{15, 45, 85}
\definecolor{goldaccent}{RGB}{175, 125, 20}
\definecolor{softblue}{RGB}{240, 246, 253}
\definecolor{darkslate}{RGB}{25, 40, 55}
\definecolor{boxframe}{RGB}{35, 75, 125}

\hypersetup{
    colorlinks=true,
    linkcolor=vennblue,
    citecolor=goldaccent,
    urlcolor=vennblue
}

\pagestyle{fancy}
\fancyhf{}
\fancyhead[LE,RO]{\thepage}
\fancyhead[LO]{\nouppercase{\rightmark}}
\fancyhead[RE]{\nouppercase{\leftmark}}
\renewcommand{\headrulewidth}{0.4pt}

\title{\Huge\textbf{偶然の論理}\\[4mm]
\huge\textbf{客観的頻度論と確率の論理的・物理的基礎}\\[6mm]
\Large\textbf{—— 原著第3版（全19章） 完全翻訳・数学校閲版 ——}\\[5mm]
\large 原著：ジョン・ヴェン（John Venn, 1888 / 3rd edition）\\[2mm]
\normalsize Project Gutenberg eBook \#57359: \textit{The Logic of Chance}}

\author{\textbf{ジョン・ヴェン} 著\\[2mm]
\textbf{教育情報分析研究会} 訳・数理校閲}
\date{2026年10月 校訂第1版}

\begin{document}

\maketitle

\frontmatter

\chapter*{解題：ジョン・ヴェン『偶然の論理』の歴史的射程と現代的意義}
\addcontentsline{toc}{chapter}{解題：ジョン・ヴェン『偶然の論理』の歴史的射程と現代的意義}

「ベン図（Venn Diagram）」の考案者として世界的に知られる英国の論理学者・数学者ジョン・ヴェン（John Venn, 1834--1923）が1866年に初版、1888年に決定版（第3版）を上梓した『偶然の論理』（\textit{The Logic of Chance}）は、近代確率論および統計的推論のパラダイムを根本から転換させた歴史的古典である。

18世紀から19世紀前半にかけて、確率論はラプラス（Laplace）らに代表されるように「人間の無知の尺度」あるいは「理性の信念の度合い」という主観的・認識論的立場が主流であった。これに対しヴェンは、確率は人間精神の主観的状態ではなく、客観的世界において反復される「事象の無限系列（Series of Probability）」における相対頻度の極限としてのみ科学的に定義できると主張した。この**客観的頻度論（Frequency Theory of Probability）**の確立は、のちにチャールズ・パース、リヒャルト・フォン・ミーゼス、コルモゴロフ、さらにはフィッシャーやネイマンらの現代数理統計学へと直結する大水脈となった。

本書は全3部・全19章から構成され、
\begin{enumerate}[label=\textbf{第\Roman*部:}]
    \item \textbf{確率の物理的基礎（Physical Foundations of Probability）}：系列の客観的実在性、誤差法則と実在のタイプ、因果作用、ベルヌーイの定理の適用限界、そしてランダム性の幾何学的・物理的定義（第1章〜第5章）。
    \item \textbf{論理的構築（Logical Superstructure）}：主観的信念測定の欺瞞性、サンクトペテルブルクの逆説、加法・乗法定理、ラプラス継起則批判、帰納法と確率の境界、設計論批判、様相論、およびギャンブラーの誤謬やマーチンゲールの破綻（第6章〜第14章）。
    \item \textbf{確率理論の諸応用（Various Applications）}：生命保険と賭博の経済的本質、証言の信憑性評価、奇跡・異常言説の批判的検証、そして算術平均・幾何平均・最頻値・中央値・平均二乗誤差の数理と真理への近似理論（第15章〜第19章）。
\end{enumerate}
を網羅している。

本全訳版は、Project Gutenberg eBook \#57359として公開されている1888年第3版を底本とし、ローカル大規模言語モデル（\texttt{qwen2.5-7b-instruct} による忠実な翻訳および \texttt{qwen2.5-math-7b-instruct} による数理・LaTeX校閲）を駆使して全19章を忠実に日本語全訳した決定版学術書である。

\tableofcontents

\mainmatter


\chapter{確率の系列（The Series of Probability）}
\label{ch:i}

\begin{center}
\large\textit{頻度論の基礎・確率命題と論理命題の差異・無限系列の極限と変動}
\end{center}
\vspace{4mm}

第一章。

確率の基礎としての特定の種類の群または系列

1. ある科学を最初から明確な定義で説明し、その範囲と領分を読者に数語で示すことは時として容易ではありません。より直接的に対象物に関連する科学では、この困難は比較的軽減されます。もし既に対象物が読者の知識に馴染んでいるなら、それらへの単純な言及だけで彼の研究の方向と性質について相当正確な理解を得ることができます。そうでない場合でも、名前によって対象物が心の中で多少なりとも結びついているため、その名前の発音だけでも十分な予備情報が伝えられることがあります。これは多くの自然科学においても同様で、しばしば読者に事前に何を研究するかを正確に教えてあげることができます。しかし、科学が直接的に対象物に関心を持つのではなく、過程や法則に関心を持つ場合、または現象から比較的不明瞭な特徴を取り上げる場合、予備情報の提供はより困難になります。既存の対象クラスを無視し、甚至は分解する必要があり、新たな対象の配置が必要となるのです。このような状況では、科学自体が初めてその統一性を与えます。それまでに研究対象が一緒に考えられることは稀でした。ここで最初から定義を与えることは不可能であり、その定義を得る過程は比較的労働的になるかもしれません。

確率の科学について、以下で取り扱う観点から見れば、そのようなものである。現在この科学に馴染みのない読者は、既存の知識と関連のある対象を参照しても、その範囲をすぐに理解することはできない。彼は我々の主題に手を加えてもらう必要があり、少々の時間と努力が必要となる。彼がそれを知るまでには、それなりの時間を費やすことになる。この作業は私たちの最初の仕事となる。

2. 自然を研究する際、どのような形であれ、私たちは常に情報に恵まれていく。これらの情報を一般化した命題としてまとめることができる。しかし、これら多くの一般的な命題が、その包摂する詳細や構成要素と同程度に確実かつ正確である場合もあれば、それほどではない場合もある。現在は、これらの一般化の真偽について争われていないことを仮定している。事実上、それは弱い証拠に基づいているか、あるいは広範囲にわたる帰納によって不確定である可能性があるが、意味するところは、その成分部分を分解したときに、全体と同様の確信を持って詳細の真偽を確認できるということである。例えば、すべての牛が反芻することを知っている場合、特定の牛や牛群が反芻することについても、全体が反芻することと同じくらい確信を持つことができる。私の最初の主張が正しかったり間違っていたりするか、その真偽を得るためにどのような方法を使用したかに関わらず、一般化された命題の価値がどの程度であれ、個々の詳細の価値はそれほど大きくも小さくもない。一般的な命題から特定の命題を推論する過程には確信度の減少は伴わない。これらの「直感的な推論」が正当化される場合でも、それはすべての場合において同じように正しいことになる。

しかし、これはすべての場合に必要ではない。論理ではほとんど認識されないが、実践では恒常的に用いられる一連の直感的推論がある。その特徴は、個々の事柄について特定性を増すにつれて確証度が減少することである。

たとえば、私は「いくつかの牛は反芻する」と教えられたとする。この情報から論理的に特定の牛が反芻していることを推論することはできないが、その話題に対する絶対的な不信や無関心からは離れるような感覚を持つだろう。しかし、牛の群れを見たときには、いくつかの牛が反芻しているという確信度が増すと感じるだろうし、群れの中の牛の数が多いほど、その確信度は増す。したがって、個々の牛について完全な確信がない状態で、より多くの牛を含む主張に対してはより強い推論を付けるというクラスのものがある。この科学はこのようなクラスとその推論に関心を持つ。

3. 以上述べた言及は、純粋に前段階的なものであり、我々が想定する事柄に対する確信度や信念の程度について多少の参照を含んでいる。読者は、この要素が科学の対象であるとする一部の著者の見解を知っているかもしれない。したがって、その要素は将来的な章で議論される必要があるだろう。しかし、私は前述の著者たちの意見には賛成していないし、少なくともこの要素を最も重要な要素として扱う点についてはそう思わない。したがって、ここではさらにこの要素について言及しないことにするが、最後のセクションで注意した特定のクラスのものに固有な特徴についてより詳細な調査に入ることにする。

これらの対象について、確率が扱うものであるとすれば、読者が心に深く理解すべき基本的な概念は、数列というものです。しかし、これは特別な種類の数列であり、個々の不規則性と全体的な規則性を組み合わせたもので、これに対するより適切な簡潔な説明は、その定義の中にしか見つけることができません。この声明はおそらく、いくつかの解釈が必要となるでしょう。

既に述べたような例に戻りましょう。経験に即したものを選びます。ある子供たちは三十歳まで生きないでしょう。もしこの命題が純粋に不確定なもの、または論理学では「個別」な命題と見做されるのであれば、数列という概念は明らかにその命題に関連しないでしょう。それは特定の全体に対する不明確な比率について述べているだけです。しかし、我々が扱うのはこのような純粋に不確定な命題ではありません。逆に、この命題を数値的なものとし、特定の比率について述べる場合を考えましょう。そうすると、数列という概念を排除するのは困難になります。特に正確な推論を得ることを目指すときには、その概念を避けることは不可能でしょう。

例えば、「三人の新生児のうち二人は六十三歳まで生きない」という命題の意味は何でしょうか？これは決して「例えば三十人の中から、ちょうど二十人が生きない」というような特定のバッチについて述べているわけではありません。言葉の厳密な意味に関わらず、この命題の主張とは異なるものでしょう。むしろ、大量の例を検討するという観点から、長い一連の事例を調査することを考えます。そしてその一連の中で、最初は固定されず正確でないものの、長期的にはそうなる傾向がある数値的な比率を見つけることを述べているのです。

我々が扱うすべての例において、大量の対象や、それらの一連、すなわち「系列」についての参照を見つけることでしょう。

いくつかの追加の例を挙げることで、この点が明確になるでしょう。

次に、私たちはコインを多数回投げることを想像しましょう。その各々の投げ方の結果は連続した系列として考えることができます。この系列の個々の投げ方は、全く無秩序なように見えます；これが我々がそれらについて不確実性を感じる理由です。時には表が出有时には裏が出ます；同じ面の繰り返しがあるときもありますし、ないこともあります。少数回だけ観察する限りでは、系列は単なる混沌に見えてきます。しかし、長い連続の結果を考慮すると、明確な差異が現れ始め、最終的には鮮明で印象的な形をとります。この場合、表と裏が出る数はほぼ等しいことがわかりますし、異なる面の類似した繰り返しも同様に見られ、その他のパターンも同様です。つまり、個々には無秩序であっても、全体としては一定の秩序が生じ始めます。

また、人間の寿命を調査する場合でも、我々が注意するそれぞれの寿命は系列として特徴を持っています。個々の寿命の長さはよく不確実ですが、一連の寿命の平均的な持続時間はほぼ等しくなりつつあります。どの混ざった集団からより多くの個体を取り出すほど、秩序の兆候が明確になり、選択された各クラスの平均的な寿命も近づいてきます。

これらの例は、我々の経験全体を通じて、ある程度の範囲で見つけることができる性質の一例です。火災、船難、収穫量、出産、結婚、自殺など、観察する特徴は何でもかまいません[1]。個々の事例は不規則であるかもしれませんが、多くの数を考慮するとその不規則性が減少し、実用的な目的ではほとんど消失します。

平均値がこのような形で詳細に現れる不規則性を平滑化することについて述べるとき、学生の注意は主に以下のことへ向けられるべきである：それは絶対的な不規則性ではなく、相対的な不規則性が無限大まで減少する傾向があるということである。この概念は数学者にとっては馴染み深いものかもしれませんが、他の人々には明確に理解するために一定の思考が必要となるかもしれません。

絶対的な離散と不規則性はむしろ増加する傾向があり、その相対的な重要性は無限大まで減少する傾向があるという点に注意してください。つまり、投げられたコインの例では、10回の投げ入れを取ると、表と裏の数が6つ異なることは非常に不確実な出来事です。すなわち、8回表が出ることで2回裏が出ること、またはその逆であることが多くないことを意味します。

しかし、1,000回の投げ入れを考えると、それぞれの数が6つ以上異なることが極めて確実な出来事となるかもしれません。一方で、1,000回の投げ入れにおける表と裏の比率は、10回の場合よりも多くの場合で、全体的な量に対する相対的な割合としては非常に近いものとなります。

つまり、運命の長引くゲームほど、その自身での大規模なスペルや連続が大きくなりますが、全体的な量に対するその相対的な比率は小さくなる傾向があるということです。

4. 以上の説明において、詳細についてはほとんど不知である出来事や物について話している場合でも、それらに関する我々の知識が完全かつ普遍的に欠如していることを意味するわけではありません。また、同じことだが、すべてのその他の性質において不規則性が観察されるとは限らないことも同様です。ここで言及するのは、個々の事例について確実に断定できる存在しない某些の性質や特徴のみを指すことです。それ以外のすべての性質については、極めて均一な状態が見られ、したがって最も完全な確定性が存在することがあります。

人類の寿命における不規則性は有名であり、しかし誰もが偶然出会う人々には心臓や脳という器官が存在することを疑わないでしょう。そして、不規則性が観察される性質においても、しばしば、実際には一般的に、その範囲内での具体的な制限が存在します。例えば、誰も特定の寿命の長さを計算することはできませんが、150年を超えることは絶対にないという確信を持っています。個々の事例における不規則性は特定の面でしか表れず、その面においても制限があります。

この同様な指摘は、我々が扱う他の多くの例にも適用されます。実際のところ、混乱や不規則性は普遍的ではなく無制限ではありません。それは特定の方向にのみ存在し、特定の点までしか及ばないのです。

5. 以上で述べたように、数列について話すときには、その構成要素である対象が時間的に連続して存在することを意味するわけではなく、これらの対象が私たちの注意に次々と現れる順序によって数列が形成される可能性があることを指摘する必要はほとんどないでしょう。実際、これらはどのような順序でも私たちの注意に次々と現れることも可能です。たとえば、死亡率記録は同時に或いは連続して発生した死者から構成されるかもしれません。あるいは、もし我々が希望すれば、これらの死者をどちらか一方とも異なる順序で並べ替えることもできます。これらはすべて、私たちが考慮に入れる目的によっては全く問題にならないことです。これらの全てのケースにおいて、数列はその性質が完全に同じとみなすことができます。また、これらの対象は自然から与えられたものであり、それらをどのように観察するかは私たしの独自の配慮によるものです。この点について注意深く述べるのは、その後でこの主張の意味がより明確になるためです。

「数列」という言葉がここでの用法において誤解される可能性があることは認識していますが、私はこれ以上ない適切な言葉を見つけることができません。上記で述べたように、これらの事象は必ずしも時間的な連続性を持つ必要はありませんが、しばしばそうなることもあります。多くの場合（例えば、コインの投げ入れやサイコロの振るい落とし）では実際には順次に発生します。他の場合（例えば、人の身長や寿命）では、これらの事象が実際にどのように順序付けられたかに関わらず、統計表によって次々と私たちの注意に引き出されることが一般的です。これら全ての場合において、我々の推論プロセスはグループを構成する各メンバーについて一つずつ検討することを必要とするか、またはそのことを準備しておく必要があります。これらの考慮事項の力は、第6章での調査を通じて明らかになります。

レシリー・エリス氏は、私にとって実質的に同様な見解を genus と species の観点から表現しており、数列について述べるのではなく、次のように述べています。「個々の事例を考えたとき、頻度の比率がすべての試行に共通する状況によって依存しているという確信はありません。むしろ、その発生状況における決定的な要素は、その genus と species の概念から外れた要素であることを認識します。偶然性と制限が同時に現れ、これらは genus を全体として考えたとき、または潜在的に含むすべてのものを理想的に実現したときに消失します。この仮定が認められれば、確率論の基本原理は以下の声明に含まれる可能性があると考えられます——genus の概念には、その下位種と関連する数値的な関係を含むという考えが含まれているのです。」先ほど述べたように、これは本章で説明した見解と実質的に同様のものですが、系列やグループについて述べるとき、genus と species の用語は、傾向や制限の概念をより適切に表現するのに十分ではないと考えます。従って、私は後者の方が断然好ましいです。

6. 現在、読者は個々のものや事件の系列やグループについての概念を持っています。これらの個々については、特定の面でほとんど知るところがありませんが、より大きな数を観察するにつれて一貫性が増していくことに気づきます。これは我々が扱うべき種類のものを明確に示すのに十分な定義ですが、正確な思考のためには不十分です。従って、より精密な分析を試みる必要があります。

あるフレーズは、その用語が技術的な文脈で広く使用され、その結果、専門用語として定着している。例えば、「イベント」と「その発生方法」などである。硬貨を投げることを例にとってみよう。これは「イベント」と呼ばれるし、硬貨が表か裏が出るという事実は、そのイベントの発生方法とされる。

また、死亡率のテーブルについて議論している場合、前者の用語は単に死の事実を指し、後者の用語はその死がいつ起こったか、またはどのように引き起こされたか、あるいはその他の特定の状況を表す。このようにして使用される言葉遣いは非常に便利で、本書においてしばしば使用されるが、説明なしでは混乱を招く可能性がある。

実際には、多くの場合、イベントの発生方法が極めて相対的に重要であり、その発生方法によって事件が異なる名前を持つことがある。言い換えれば、その二つの状況においては、名目上同じイベントとは言えない。したがって、このフレーズを適切に使用するためには、それなりの拡張が必要となる。

例えば、人間（男性と女性）の一連の個体について考えている場合、彼らの人間性を「イベント」とし、その性別をそのイベントの発生方法と呼ぶのは不自然である。

もしも以前に述べた対象のクラスのいずれかに戻るならば、我々は望む概念についてより正確な理解を得るために道を見つけることができるでしょう。容易に理解できるように、これら全てのクラスには類似性と相異性が混在していることがわかります。ある一定数の特徴や属性を持つ一連の出来事がある場合、それらが分類される理由はこれらの特徴や属性がないとそう分類されないからです。しかし、それらにはまた区別も存在します；特定の他の属性は一部に見られますが、他の部分では見られません。つまり、系列を形成する個々のものは複合的なものであり、それぞれが多くの要素や属性の集まりで構成されています。これらの要素のうちいくつかは系列内の全てのメンバーで見られますが、他には一部のみでしか見られないものもあります。これまでは確率論に特有な点はありませんが、その区別する特性はこの点にあるのです；すなわち、一時的な属性（恒久的属性とは対照的に）は広範な調査の結果、全体の事例数に対する一定の定まった割合で存在することが見られるのです。特定の個体においてこれらの属性が見られるかどうかを予測することはできませんが、我々がより多くのケースを調査するにつれて一貫性が増してきます。つまり、その属性が見られた場合と見られなかった場合の事例数に対する割合は比較的変動が少なくなり、ある固定値に近づいていくのです。

これは最も広範な説明形態ですが、実際にはグループは多くの場合もっと単純な形を取ることもあります。例えば、人間という同じ種類のものの一連の系列として現れるかもしれませんし、左利きであるかどうかという属性が時々見られるだけです。当然、「属性」という言葉は最も広い意味で用いており、それは一つの物体に観察できる全ての特徴を含むことを意図しています。これは根本的な性質から最小限の時間と場所の偶然までを含みます。

7. シリーズを検討するに当たっては、必ずしも異なる方法で起こる一連の出来事としてではなく、グループの一連として考えることの方が適切であることがわかります。これらのグループが分析されると、それぞれの場合において、これらは物質と属性の集合体に分解されることが見つかります。グループ列の統一性を提供するのは、その一部の物質や属性が全体の列全体で共通しているという事実です；一方で、グループ内の差異は、その一部がこれらの物質や属性の一部のみを含み、他の部分が時々欠如することによるものです。このような理解のもとでは、当の表現は確率が考慮に入れるすべてのクラスのものに適用されるようになります。

8. 普通の表現（すなわち、「出来事」および「その起こり方」）は上記に含まれるということが容易にわかります。偶然属性が重要でない場合、恒久的な属性だけで名前を固定し、適切に適用することが可能です。他の要素の存在または欠如は、名称の一部の変化や予備述語の追加によって示されます。したがって、このようなすべての場合において、「属性の集合」を「出来事」と呼びます—つまり、本質的には同じ出来事を指すものですが、異なる状況ではそのように起こる方法が異なります。しかし、偶然属性が重要である場合や、大半を占める場合、この表現はより適切ではありません。我々は非常に異なる集合体が実際には同一の出来事であり、違いは発生方法にのみあると暗示することで、言葉遣いが多少強引になります。しかしながら、この表現は非常に便利なもので、誤解を防ぐ注意を払うことで頻繁に使用されるでしょう。

9. 以上の種類の級数は、私は考えると、確率に関するすべての規則が基礎となる最終的な基盤であるべきである。この点まで分析を進めることが明瞭な理解に不可欠であり、しかし、系列を構成する事件の内面的な性質についてさらに詳細な分析を試みることは必要ではない。例えば、物質の独立した存在に関する形而上学で議論される問題に対する見解形成は全く不要である。私たちは検討によって、物質や属性のグループ又は単に属性からなる系列を見つけ出した。この系列において止むべき点に達し、そこから推論の規則を調査するのである。これらの物質や属性が心理学者や形而上学者が手に取った場合にどのように最終的に分析されるかについては、ここでは私たちの関心事ではない。

10. それでは現在到達した段階は、分析によってグループ化できることが証明された事物の量を見つけ出したところである。これらの事物は系列として構成されていると最も適切に見做すべきである。この系列の特徴的な特異点は、無秩序から徐々に秩序が現れ、時間とともに明白で疑いの余地のない一貫性を示す秩序があることである。

このような系列の説明から得られる印象、そして確率について学んだ経験のある読者がおそらくすでに持っているであろう印象は、この秩序の徐々な発展が無限に延びることであり、その近づきは従って完璧さへと無制限に近づくことである。多くの一般的に選択される例も確かにこのような印象を確認する傾向がある。しかし、主題の理論に関して言えば、私は確信しているように誤りであり、混乱を招く可能性のあるものである。

この作品の前文として付け加えられた以下の言葉は、直ちに修正される。「型を大切にするように見せかけても、個々の生命には無関心だ」という言葉は、長い間その型を見つめてみると、実際に変化し、消え去り、他のものに取って代わられるという主張によって訂正される。同様に確率においても、長期的には見出される一貫性は、個々の無秩序とは対照的ではあるが、永久的なものではなく、持続的であるがゆえに永遠ではない。十分長い間それを観察し続けると、それはほぼ恒常的に変動することを発見するだろう。時間とともに、その規則性は個々の事例自身ほど予測不可能になるかもしれない。この事実が主題に対する理論や、それに連nessする一般的な計算方法に及ぼす完全なる影響については、次のいくつかの章でより詳しく明らかにするが、現在ではこれを非常に簡潔に確立し、説明することに留める。

この事実が主題に対する理論と、それに連関する一般的な計算方法についての共通的な手法に及ぼす完全なる影響は、次の一連の章でより詳細に明らかになります。しかし現時点では、これを非常に簡潔に確立し、説明することに留めます。

例えば、平均的な寿命について考えてみましょう。データが十分に広範囲であると仮定すると、これは比較的規則性と一貫性を持つことが知られています。この事実は前節で既に述べられており、今日の一般の人々には相当明確な理解がある真実です。しかし、少し考えれば、この一貫性が観察できる範囲には上界と下界があることもわかります；つまり、例を多く取ることによって誤りに陥るだけでなく、逆に例を取りすぎることでも目的から遠ざかる可能性があります。また、その理由は全く異なるもので、結果も異なります。

現在の英国における平均的な寿命は約四十年としましょうが、百年前には明確にそれよりも短く、数世紀前にはおそらくさらに短かったでしょう。そして、まだもっと古い人口に関する統計があれば、おそらくその後より著しい改善があったことがわかります。その将来の動向については誰も確実な予測はできません。衛生状況や他の改善により、寿命が着実に増加する可能性がありますし、限界がないまま無制限に増加することも考えられます。しかし、もっと現実的で、おそらく起こる可能性が高いのは、この寿命が固定の長さに向かって徐々に収束することがあるかもしれません。また、将来の世代は短くでも幸せな生活を好み、その結果平均的な寿命が減少する可能性もあります。

寿命は人々の一般的な趣味、習慣、仕事などに一定の依存性を持つことから、彼らが意識的に或いは無意識的に設定する理想によっても左右されます。この理想が数世紀後にはどのようなものになるかを予測するのは非常に危険なことです。ここでは必要なのは、これまで述べてきた特定の一貫性（その相対的な性格を強調するために用いた用語）が変化し、将来的な意見や実践の渦中でさらに変化する可能性があるということだけです。そして、その程度はどのくらいでも、不規則性はどの程度でもあります。

天文学者から借用した言葉を借りれば、我々の一貫性は不規則な「長期的」変動にさらされていると言えるでしょう。

11. これは典型的な例と言えるでしょう。もし、命の長さというよりは、比較的単純でない特性を採用したり、それが比較によって「大きな恒常的な一貫性」と呼ばれる自然の特徴と密接に結びついていないものを使ったりした場合でも、注意すべき特異性が著しく際立って現れることがわかります。例えば、麻疹による死亡数や duel（决斗）の事例、または女巫告発の件数を、数十年間にわたって調査すると、それらは相当程度の一貫性を示すかもしれません。しかし、これらの一貫性はおそらくゼロから始まり、様々な大きな変動を経て再度ゼロに近づく傾向があると思われます。少なくともこの世紀においては、これらがさらに急速な変動を遂げる可能性も否定できません。

これらの例は、自然におけるほぼ普遍的な規則の一端を示すものであり、極めて極端なものではありません。私は、男児と女児の出生比のような明らかな例外さえも、厳密には例外ではなく、実際上そうではないことを証明しようとします。つまり、最も広範で一貫性のある型は自然の中にほとんど見つけることができないと言えるでしょう。この結論の完全な意味は、将来的の章で明らかになります。

ここでは、統計が十分な数値がない限りは必ずしも有用ではないという事実に注意を向けます。しかし、それから従うべき推論は、ある特定の場合においては統計が過剰である可能性があるということです。統計が広範にわたるほど、その正確性は特定の時間や場所における最大限の精度を欠くかもしれません。

12. これらの自然的均一性は、結局、振動する傾向があることが見つかる。それに対して、賭博ゲームで提供される均一性と比較すると、後者のものは長い間その調査を続ける限り、何らの世紀的な振動の痕跡も見せないように思われる。次章を通じて、いくつかの一般的な試みが提出され、それらはこの問題について事前的に証明する必要がある固定性を示すが、その存在については疑う余地がないと言えるほど確実であることはほとんどない。ペンスは現在のように表と裏が出る頻度がほぼ等しく、かつて投げられたときと同じであり、現状が続く限り、そうした傾向を維持すると信じられている。これらの均一性の固定性は絶対的なものとは言えず、しかし我々が必要とする経験のどの程度であっても、それらに著しく影響を与える可能性は低い。したがって、自然的均一性は最終的に振動する一方で、賭博ゲームによって提供される均一性は永遠に固定しているように見えるという明白な対比がある。

13. ここで二種類の系列があると考えられる。これらは初期の不規則性において同じであり、その後の規則性においても同じである；それらがそれぞれ異なる点は最終的な形においてである。一つは不規則な変動なく固定された数値的比例に収束する傾向がある均一性を持つ；一方では最後には均一性が振動し、その振動は規則化不可能な形を取るかもしれない。

本章は主に説明的かつ科学の基礎を示すために存在するため、ここでは次のように述べることができる（その後でその正当性が示される）：系列が前記のようなものである場合のみ、厳密な科学的な推論を解釈できる何かを得ることが可能である。しかし、実際の後者の系列が経験によって提示される任意の例においても、仮想的に前者の系列を置き換えると、その系列に犯される種類と範囲について一般的に理解することができるだろう。当然、両系列は可能な限りすべての点で似ているが、変動の不均一性が固定されたものに置き換えられたという点だけが異なる。それらの差異は最初の段階では現れず、その段階では確率の系列の特徴的な特性はまだ明らかになっていないためである；そこにはすべて無規則性があり、同じように似ているか異なるかを示すことも、異なるか似ているかを示すことも不可能である。私たちはただ一般に言えるのは、それぞれが同じ種類の無規則性を示しているということだけである。また、次の段階でも差異は現れず、実際の変動性がまだ長時間にわたって観察される余地がないためである。それらの差異は我々が「最終的な」段階と呼んだものにおいて初めて感じられるだろう。つまり、系列が非常に長い期間を想定するときである。例えば、六ヶ月齢児が毎年死亡する人の割合は、検討した人数が少ない場合では全く無規則である；しかし、検討した人数が大きい場合では規則性を持つようになる。しかし、我々が非常に長い期間観察を続けたり、非常に広い地域にわたって観察を行ったりすると、その規則性自体も無規則な方法で変化することを見つけるだろう。先ほど述べた置き換えは実質的に、「規則性は固定され永久的なものである」と仮定するのと同じである。これは事実と完全に一致しているとは限らない仮説であり、しかし明確さを保証し、定義を確保するために我々が強制されるものである。

14. この置換の完全な意味と影響は、いくつかの後続の章で明らかになるが、すぐに指摘できるのは、この方法によってのみ、「極限」という概念を我々の議論に完璧な厳密さで導入することができる点である。少なくとも多くの応用事例において、純粋な統計的な調査に関連しては。ある一定の割合が長期的には事件の中で支配的になるということが言えるが、しかし詳細を調べると、我々は現状が変わらないと仮定し、長期的にはその特性が見られるはずだと仮説的に述べる必要があることを発見する。事象の連続性によって引き起こされる状況について正確な知識がない場合、これらの状況が変化する可能性が高いという強い理由がある場合、実際には他に方法はないと考えるべきである。つまり、これらの状況が変わらないと仮定し、その結果、数値が一定の均一性を持つ系列を実際に変動する均一性を持つ系列に置き換えることで、「極限」の概念を導入できる。

15. 読者の方に次の例を学んでいただきたい。数学的には「ピーターセンブルグ問題」として知られる有名な例であるが、この例は本章で述べたいくつかの考慮事項を説明するのに役立つことがわかります。特に、我々が関心を持つ系列は数的または時間的に無限に広がっていると見なされるべきであり、そのような観点から特定の系列（問題に関連している系列もその一つ）だけが無限の範囲を活用して、前平均からの偏差が有限の限界を持たない個体を継続的に生成するという事実を浮き彫りにするのに役立つ。適切に見れば非常に単純な問題であるが、一度は混乱や困惑を引き起こす原因となった。

次のように問題を述べることができます：一便士を投げ上げます。表が出れば、私は1ポンドを受け取ります；表が連続して2回出れば2ポンド、3回連続で表が出れば4ポンド、そしてそのようなパターンは続けていきます。つまり、新しい表が出るたびに、その数だけ受け取る金額が倍になります。言い換えれば、表の連続した列を「ターン」またはセットとみなして、それから次のターンへと進み続けることになります。そして、各ターンに対して支払いを受け取ります；裏が出ると、何も得られず、つまり考慮から外されます。

表が何回連続しても、受け取れるポンド数はその回数の1つ少ないべき乗で求めることができます。したがって、一連の投げ上げからなる系列があります。私たちは、多くの人々が証明可能だと考えたり、経験的にも相当な範囲で確認されているように、これらの「同じ面の連続」の希少性は、そのような連続が出るときの受け取れる金額と直接比例すると仮定します。言い換えれば、支払いを受け取るたびに、1ポンドを2回受け取るたびに4回、2ポンドを8回、そしてそれ以降も無限に続くたびに受け取る回数が増えることを考慮に入れるべきです。

この問題について質問があります：この特権を得るために私が支払うべき金額は何でしょうか？後述する結果の一部を少し先取りして言いますが、これは平均的に私は勝者でもloserでもないためには、どの程度の金額を毎回支払うべきかという問いに相当します。言い換えれば、上記の条件で私が平均的に受け取る金額は何でしょうか？理論は私に無限の金額を与えるべきであると述べます：つまり、どんなに大きな有限な金額でも十分ではありません。そしてこれは実際には完全に理解できます。私には無限長の系列がありますし、その系列を続けるほど私の収入は増加しますが、何らかの制限はありません。確かに非常に豊富な獲物は非常に稀ですが、それでも時々現れ、その豊かさで補填されます。

人々がこの問題に取り組んだたびに、当然ながら私たちは系列の一部しか知りませんでしたが、我々が計算の基礎とする系列は無限であり、この前提に基づいて推測される金額は長期的には完全に一致します。

一般的な反論は、返答の中で述べられています。無限の金額を支払う必要がないという点で、誰も $p$50 にも近い金額をそのような確率のために支払うとは思わないでしょう。おそらくそうではないでしょう。なぜなら、系列を見ることは十分に多くない限り、それが価値あるものになるほど多くの経験が得られないからです。大多数の人々は実際的な見解を形成する際に、実際に見たかまたは合理的な期待の中で得られる系列の小さな部分に基づいています。現在、このような系列の一部について考えましょう、例えば 100 回の試行を含むような部分を考えます。その系列における連続した表が出る回数は平均して7回から8回に過ぎません。これは観察されていますが、公式がそれらを生成する場合、より良い結果を無限大まで生成し続けることが可能であるという事実は忘れられています。そのため、ある人々は困惑します。なぜなら、彼らが実際に遭遇する系列の一部に対する行動は、その系列が完全な無限大であるという前提に基づく推論には正当化されないからです。

16. これは、制限された投げ回数を考慮したときに様々な可能性とそれらを網羅するために必要な範囲について考えることでより明確になります。まず三回の場合を考えましょう。これにより8つの等しい確率の可能性が生じます。そのうち4つの場合では、投げ手は尾が出るので失敗します：2つの場合では1ポイント（つまり $p$1）を獲得し、1つの場合では2ポイント、そして1つの場合では4ポイントを獲得します。したがって、彼の合計収益が$p$8で4つの異なる状況で達成されるので、彼の平均収益は$p$2となります。

今 suppose に投げ回数を n 回まで延長する場合を考えましょう。これにより 2{n} の可能性が生じます。これらのすべての可能性を考慮に入れるには、それら全てを一度計算する必要があります。もし可能なすべての場合を一度だけ計算したとすると、彼の合計収益（ポンドで算出）は次のようになります：

$$2^{n-2} + 2^{n-3}\cdot2 + 2^{n-4}\cdot2^2 + \ldots + 2\cdot2^{n-3} + 2^{n-2} + 2^{n-1},$$

つまり、

$$(n + 1)2^{n-2}.$$

これを 2{n-1} 回の収益機会に広げると、彼の平均収益は 1/2 (n + 1) となります。

平均について言及する際には、平均を正当化するために必要な異なる結果の最小回数は、各結果が一度ずつ現れる回数であることを忘れてはなりません。ある人が連続10回の表が出るまで投げを打ち止める提案をするかもしれません。それは良いことです。しかし、その平均的な利益は5ポンド10先令0便となると伝えます。そして同時に、この声明が正当化されるためには少なくとも1024回の投げが必要であることを強調します。なぜなら、それ以下の回数では利益と損失のすべての可能性が十分に現れず、平衡されることはありません。

もし20ポンドの平均的な利益を達成したいとするなら、彼は39回の投げまで準備する必要があります。この支払いを正当化するために、彼は2{39}回、つまり約1兆回の投げから始めなければなりません。彼がこれを達成しない限り、彼はどんな疑問者に対してもその投げを39回延長した場合の真の平均的な収入を証明することができません。

もちろん、彼が永遠に投げる選択をするなら、数や量に関する無限性に関連する説明法だけを使用しなければなりません。無限の金額を支払うようにと提案することはできません。なぜならこれは厳密な意味を持たないからです。しかし、彼が各回の表が出る度にどれだけ支払っても、最終的に彼の総支払いが総収入で返される段階があることを伝えます。どんなに大きなnであっても、2{n}回試行を続けることで、真の平均的な収入が1/2(n + 1)ポンドである可能性があり、十分長い期間続けばそれが実現します。

この問題は将来的な章で再度検討されるでしょう。

1. 次の統計は、このような規則性が存在する広範な経験の範囲を公平に示すものである：「異なる原因による等しい振動量の例として、ロンドン西地区における7年間の死者数（振動量13.66）と人身犯行（振動量13.61）、または脳卒中による死亡（振動量5.54）と暴力なしの財産犯行（振動量5.48）、または外科医学校に登録された学生数（振動量1.85）と自宅用消費される製たばこの総重量（振動量1.89）、または外務費者（振動量3.45）と英国船舶の積み荷重での入港トン数（振動量3.43）などがある。」[これらの統計は、1858年3月に『統計学会誌』に掲載されたMr Guyの論文から抽出されたもので、「ここでの振動量」とは平均からの乖離量を測る尺度であり、その評価方法については後述する。]

2. カムブリッジ哲学者協会の取引、第9巻、605ページ。彼の著作集の収録版では、50ページ。

3. 私たちはこれを次のように表現できる：少数の例は法則を示すのに十分ではない；相当数の例は法則を示すのに十分であり、しかし非常に多くの例が必要であることが変化が法則に生じていることを確立するためである。

4. 数学者はこの置換の性質を、「幾何学における接円の円周」や「天文学における瞬間楕円」の類比によって説明できる。これらの概念が用いられる場合、我々は常に連続的に変化し、その変化率もまた変化する現象がある。その特定の瞬間にそれを捉え、その瞬時の変化率を固定すると仮定し、その後その仮定に基づいてその進展を完成させる。

5. これは最初の数学的な扱いがペトロスバーグアカデミーの『論説集』に含まれていたことから「ペトロスバーグ式」と呼ばれる。Todhunterの『確率理論の歴史』には、その多様な報告が見つかる。

\newpage

\chapter{系列の配置と形成・誤差の法則（Arrangement and Formation of the Series. Laws of Error）}
\label{ch:ii}

\begin{center}
\large\textit{誤差法則の普遍性批判・非対称分布・最小二乗法と実在のタイプ}
\end{center}
\vspace{4mm}

第2章。

前章で述べた特定の種類の級数の性質について、確率論の基礎となるべき級数として十分な説明が与えられました。確かに、そこでの特徴は、確率そのものとそれに取り組む問題の種類に関する公正な一般的な概念を提供するのに十分であると言えるでしょう。しかし、結論部分で示されたように、このような級数が自然の中で実際に現れるか、または多少人工的に生成される結果として見られるとき、それは単純な形で現れることが稀であり、しばしば複雑な形でグループ化されていることを観察します。このトピックについてのより深い議論が必要です。

例えば、クエテレットの研究によって既にいくつかの読者が馴染みのある種類の級数を検討しましょう。ある都市や国の中で多くの成年男性の身長を測定するという状況を考えましょう。これらの身長は当然ながら一定の極限間で分布します。そして、測定値を継続的に積み上げていくと、それらがこれらの極限間で連続的に分布することを見ることができます。つまり、そのような条件下では、どの中間の身長も永久に欠落するという状況は見られません。

これらの身長を大きさ順に並べると、常に次のようになることが観察されます：極端点の中央付近には多くの結果が密集しています；この点の両側にはまだ余裕がありますが、それほど大きなものではありません。そして、一定の比率で減少するまで、極端な結果に向かって数値は薄くなり、相対的に非常に小さくなります。

ここに注目すべき点は、これらの数値がそれぞれの方向から中央に向かって減少するという単なる事実ではなく、その減少がどのように進行するかの法則である。この漸進的な減少がどのように行われるかについては、すぐに詳しく説明される。ここで「法則」という言葉は数学的な意味で用いられ、関心のある二つの要素間を結ぶ公式を表現するために使われている。つまり、その高さ自体と、その高さに対する見つかった相対的な数値の間の関係性を示す公式である。これらの要素のうちの一つが他の要素の関数であるかどうか、そしてそうであればどのような関数かについて調査する必要がある。

2. 前章で述べた通り、このような調査の興味と重要性は、統計データが非常に広範囲であることに大きく依存している。クエテレットの『社会物理学』[1] の中には、人間の身体構造から提供できるほぼすべての要素に関する測定値の選択が含まれている：身長、体重、各部位の筋肉力、ほとんど全ての部分や器官の寸法などである。これらのうち最も広範ないくつかは、ポタモカ軍の25,000人の連邦兵士の身長と、5738人のスコットランド民兵の胸囲を数年前に測定したものを含んでいる。このような統計データを参照したい人には、これらの著作や他の同じ著者の著作が最良の情報源である。

クエテレットの統計的調査は興味深く有用であり[2]、その重要性は彼の努力により大きく高まっているかもしれないが、私は彼がこの主題について述べた多くの内容が確実に誤りで混乱を招いていると確信している。これらの誤りは彼だけに限定されず、科学の確率論の基礎やそれに伴う哲学的な問題に関する誤解については、彼の明示的または暗黙的な表現に対する批判を通じてより適切に議論されるべきである。

3. まず最初に、彼はしばしば、あるいは常に、我々の観察、測定などの結果について統計的な調査を行う際に、一つでかつ同じ法則に基づいて配置されるべき唯一の法則が存在するという前提を採用します。つまり、ある平均値を中心に集団として群れるような量のグループが存在し、それらがその平均値から離れるにつれて頻度が減少するとするとき、この頻度の減少はこれらの量の性質や得られる過程に関係なく一定の法則に従っていることを発見できると仮定します。

このような一貫性が多くの異なる現象のクラス間で prevaling であるとは思われるのはおそらく驚くべきことでしょう。しかし、この事実（もしそれが事実であるとするなら）の完全な意味は、そのような現象が一つの包摂的な原理によって統合されるという前提により、それらの性質や起源に存在する深い差異を注意深く考慮したときにのみ明らかになります。これらの現象が一貫性のある法則によって調和されていると仮定するとき、その意味がよりよく理解できるでしょう。

確率に関連する主な対象物は三つの主要なクラスに分けられます。

4. まず第一に、賭博などの機会ゲームによって提供される様々な組み合わせや運の流れがあります。たとえば、十コインからなる一握りを連続して投げ上げ、その結果を記録した場合、得られるものは次のようになります。特定の割合で最も多く見られることになるのは、五頭五尾です；それよりも少ない割合では、四頭六尾と四尾六頭が頻繁に見られます。そして、この傾向は次第に減少し、最終的には九頭一尾と九尾一頭まで下がります。最も稀な結果はすべてのコインが同じ面を上す場合です。

これらの統計的要素に関しては、少なくともその起源については人間の選択や意志によって形成されるものであり、従って一般的には主に人工的なものと説明されます。しかし、最終的にはこれらは全て運による結果であるとされています。

また第二に、多くの同一種類または同一クラスの自然物体の正確な測定値（すなわち実際の大きさ自体）を考慮に入れることができます。たとえば、既に述べたように、特定地域の住民の身長や他の特性などの例です。ここで人間の意志や介入はほとんど関与していないか、または全く関与していません。測定値を集めるのは私たちの選択ですが、測定対象自体は私たちの制御外にあります。従ってこれらは一般的には完全に自然の産物と説明され、厳密には運が関わっているとは考えられません。

第三に、我々が目指す結果は、各々の試行において常に同じ固定された量である場合があります。もし測定が厳密に正確であれば、その結果は繰り返し同じものになります。しかし、すべての方法が目標を達成するためには実際的に無数の不完全性にさらされているため、実際に得られる結果は、ほぼ全ての場合において、我々が確保しようとしている真実で固定された値から多少なりとも離れます。これらの結果は時々正確より大きく、時々正確より小さくずれ込むことがあります。もちろん、悪い試行ほど頻度は低くなります。人間が的を射るとき、彼はたいていの場合は中心に正確には命中しないでしょうが、良い射撃は悪いものよりも多くなるはずです。ここでもまた、一定の量（すなわち、射撃が射点からどれだけずれているか）が平均値を中心に集まっていますが、前二つのケースとは異なる方法で生成されています。この場合においては、要素は通常、人間の意志の一部であると見なされ、その結果、運命は効果を生み出すための一助と考えられます。

これらの要素には、測定とは区別される推定も含まれます。一般的に、測定は一定量の機械的または操作的な結果を指し、一方で推定は直接観察や内省によって決定された量のケースを指します。このクラスが科学原理に関連する興味と重要性を得たのはフェンヒャーの研究からです。その主要な分野は自然に心理学的データの中に見つかるべきです。

他の種類の事物も、上記に言及したものを除いては、容易に挙げられるであろう。しかしながら、これらの種類については最も広範な統計が得られたり、また最も実用的な重要性や興味が注がれている。その起源と性格を分ける深い区別は明らかである。もしこれらすべてが本当に同じ分布の法則を示すならば、それはおそらく、表面上見えていたly異なる方法で生み出されたように思われる様々な方法に深く根ざした何か深刻な同一性を指し示す驚くべき事実であろう。

現在議論する問題は次の通りである：どの程度の厳格さにおいても、一つだけの分布の法則が本当に支配しているのか？そして、このことが真実である場合、それはどのようにして成り立つのか？

5. 前述の二つの問いに対する肯定的な答えを支持するために、異なる種類の証拠はあるいは提供されるであろう。

(I.) 一つの計画については、経験に直接訴えることができ、統計を集めてその分布の法則を見つけることで対処できる。上記で述べたように、クエテレットや他の者が多くの場合において、そして相当な範囲でこのような方法を用いて研究を行った。彼の研究は、性質と起源が大きく異なる多くの物事や過程が、それでもかなりの正確さで一つの法則に従っていることを十分に説明している。少なくとも、平均値周辺の値については、そのようなことは明らかになっている。つまり、最も平均値に近い位置にある値についてである。一方で、極端な値に関しては、困難が生じる。例えば、多数の人々の身長の配置において、これらの極端な値はむしろ障害となる；実際には、そのような極端な値を尺度の両端から排除すべきであると提案され、それらは怪物であり、理論的にはその相対的な数がそうではないという事実がある。しかし、このような排除の方法は全く正当化されていない。これらの矮小者や巨人は、通常の大きさを持つ兄弟と同じように生まれており、彼らもまた我々が導き出す公式に含まれる権利を持っている。

人間の身長以外にも、クエテレットによって集められ整理された類似の観察の他のクラスもある。しかしながら、そのような適切な例は多くない。なぜなら、法則を証明するために別の方法で証明できる場合ではなく、実際的な直接の証拠を得るために観測と測定を集めることが目的であるとき、それらの集積は非常に大きな規模に達する必要があり、これには最も粘り強い計算家以外がその必要な労働を経験することはないからである。対面ページの注記の中で述べられたいくつかの指摘は、このような証明方法における困難性を説明するために役立つだろう。

ここでは、対称的な曲線について話しているので理解していただきたい。もし非対称性がある場合、つまり平均値の両側で誤差の法則が異なる場合、比較的少数の観測だけでその事実を検出することは十分可能である。しかし、対称性と平均値を中心にした頻度の急速な減少が保証された場合、一般的にいくつかの種類の指数曲線から一つを選択し、それが統計データを平均近傍で比較的精密に表すことを選べるだろう。つまり、統計データが多数存在する場所では一致を確保できるし、一致を確保できない場所でも統計データは比較的小さいため、長期間の観測が必要となるまで不一致を証明することは不可能である。

6. あらゆる統計をその価値に応じて評価し、範囲、適切性、正確さなどについて考慮した上で全体的な結論として得られるものは、ほとんどの人がそれらを参照する際にほぼ一致して導き出すものであり、それは極端な値以外のすべての場合において、大体一つの型や法則に従っている。これについては十分認めなければならない。しかし、極端な値に関しては必ずしもそうであるとは言えず、実際にはそれが常にそうであることは不可能であることが、少し考えてみれば明らかになるだろう。

例えば、法則が適用されると思われる事物のいくつかのクラスにおいては、投げられたコインの表と裏の順列のように、振幅の大きさに制限はありません。私たちはどれだけ長い連続の表か裏かを仮定しても、もし私たちが十分長く生き、十分な回数投げることができれば、最終的にそれを達成することができます。他の場合、特に自然現象に対する確率の応用には、このような制限が存在する可能性が高いでしょう。特定の範囲を超える偏差は単に稀であるだけでなく、性質上不可能であるかもしれません。そして、実際に不可能ではない場合でも、検討するとしばしば、それらは通常または中間的な値を生成するための予想外の要因によってのみ可能になることが明らかになります。

例えば、観測を行うどのような器具においても、その構造が誤差の可能性に絶対的な制限を設ける場合があります。そして、すべての使用状況で絶対的な制限がないとしても、公平かつ適切な使用状況ではそれでもある制限があるかもしれません。つまり、器具が意図的または粗略に操作されない限り、既存の範囲内で新たな離散要因は導入されないということがしばしば明らかになります。

例えば、ある男が的を射る場合を考えましょう。彼の最も悪い射撃は、他のすべての場合で作用していたような原因の組み合わせによって引き起こされるものと仮定します。そのような原因が一致して同じ方向に働く極端な例としては、「公平な扱い」と呼べる状況があるでしょう。これは、多数の異なる原因が偶然にも一致して同じ方向に働きかける場合を指します。他の場合では、これらの原因は互いに相殺し合うか、または部分的に相殺する傾向があります。

しかし、このような原因の総合的な効果は限定的であると仮定することができます。男が真の射撃線からほぼ直角に弾丸を放つのは、全く新しい原因が介入しなければならない場合に限られます。例えば、異常な状況によって注意が分散されたり、一時的な発作があったりする場合です。しかし、このような影響は以前から存在していなかったものであり、それらが存在していたとしても、これらの偶然の大きな混乱が無数の小さな原因の合計に同じ種類の法則を適用することを仮定するのは大胆な推測です。

実際の例については、その最後の種類についてあまり重視することができません。確実な誤差の範囲があることが確認できる例と比較してはなりません。したがって、これは主に説明のために提供され、証拠としてはあまり重要ではありません。非常に稀な組み合わせの確率を表現する数値は、実際には考えられないほど巨大であり、そのような数値の影響により、時折不可能と数学的に不確定な高確度の間で混乱することがあります。

7. この論文の最初の版が作成された当時に至るまで、統計学に関する著者は、私はまだクエテレットの影響下にあり、その特定のトピックに対する彼の権威を過大評価していたと推測されます。しかし最近では、二項分布や指数分布以外の他の配置法則も考慮する必要性が繰り返し指摘されています。

ガルトン氏は、統計学の理論の各分野に莫大な貢献をした人物ですが、彼は「誤差の頻度の法則」の基礎となる「仮定が多くの生命現象や社会現象のグループにおいて正しくない」と主張しました[7]。例えば、色調をマッチングする試みを考えましょう。フェンシャーの法則によれば、感覚＝刺激の対数という近似で最も単純な形では、黒い基板上に散らばる白い量が1, 2, 4, 8, 16, 32などと比例する一連の色調は、目には等しい間隔で見られることがわかります。したがって、灰色をマッチングする際に、8分の白を含む灰色を選んで誤差を犯す確率は、16分の白を含む灰色を選ぶ確率と同じです。前者では過誤が8分生じ、後者では欠損が4分生じます。したがって、「等しい大きさの過誤や欠損は同様に確実性を持つ」という仮定から、この結果は出るのです。

この仮定に基づく結論は、ドクター・D・マカリスター氏によって著された驚くべき論文で詳細に計算されています。後述する機会があるでしょう。ここでは重要なのは、このような統計学の結果をグラフ的に配置した場合、中央軸を中心に左右対称な曲線を得ないということです。

8. それ以降、Mr F. Y. Edgeworth（イギリス協会に任命された委員会によって貨幣基準の変動を調査するため設置された）が商品価格について同じ考慮事項を提起しました。彼は1782-1820年と1820-1865年の十二種類の商品の価格から得られた統計データを与えています。「各系列の最大値と最小値がそれぞれ記録され、最大値と最小値の中間点である「中央点」までの距離が半分未満の入力数が常に存在することが確認されました。全24回の試験において例外はなく、わずかに例外がある場合もほとんどありません。したがって、これらの曲線は添付図で示されるような偏った形状であると推測できます。」また、時間変動だけでなく場所変動についても同じ事実が確認されています（v. \_Nature\_, 1887年9月2日）。これらに関する統計データには、4000日以上連続して同一時間に観測された気圧の高さについての私の自身の統計も含まれています。これらの統計は、配置が明確な非対称性を示しています。

実際、私はこのような非対称性が測定される現象が「一方的な」特性を持つ場合に見られるべきであると考えます。つまり、その可能性が始まる一定の固定点から一方のみで測定されるという意味です。彼らがこのポイント以下に下りることは不可能であり、その近さによって同じ絶対差がより困難かつ重要になるため、彼らがそれを達する前にその影響を感じるのです。

例えば、平均値が69インチの身長の一覧表を見てみましょう。三フィート（これが可能であると仮定しても）の減少は、その影響力は非常に大きいと言えます——言葉で表現するすべての意味において、その価値は非常に高いのです——一方、同じ量を加えるよりもずっと大きいです。なぜなら前者は平均値を二倍にはなりませんが、後者はそれを二分の一未満にします。

再び例を挙げましょう。もし既に説明したような小さな影響の多くの要因が振動する単振り子に対して作用すると、それぞれの方向での偏位は対称性を示すと予想されます。しかし、それらが弾力体に対して作用すると、このような結果は期待されません。後者のような現象は、平均値を中心に対称的に配置されることがほとんどないでしょう。

9. （II）最後の議論は、誤差の不変性を証明するために別の種類の証拠を提示する可能性があることを示唆します。これは直接的で帰納的なものであり、統計に即座に訴えるものではなく、事件が引き起こされる実際または仮定された原因についての調査を含みます。

まず想像してみてください。考慮中の事象は最初に固定した原因やそのグループによって引き起こされます。もしすべての影響する状況がこの原因の中に含まれていたら、その事象は常に同じ方法で発生することになります：誤差や偏りは何もありません。しかし、現在の仮定では、非常に小さな多数の原因が存在し、それらは互いに独立して作用し、それぞれの原因が長期的には一方への影響と反対方向への影響を等しく与えるとします。これらの仮定から一般的な観点で何が起こるかを見ることは容易です。

非常に少ない場合では、ほとんどの原因は同じ方向に働くことになります；言い換えれば、極端な偏りが生じるケースは非常に少ないです。しかし、多くの他の場合では、一部の原因は一方への影響を与え、少数の原因はそれらを補うために反対方向へ作用します。その結果、比較的多い数のやや小さな偏りが生じます。このパターンは増加し続け、最終的には中央点に近づきます。ここで非常に多くのやや小さな偏りが見られることになります：反対の影響が互いに均衡を保つケースでは、全くの誤差が生じない場合が実際に稀ですが、全体として最も頻繁なケースとなります。

もしすべての平均からの偏りが上記の方法で引き起こされる（現在十分な説明としてこれで十分です）場合、これらの偏りや誤差の頻度の配置には常に一つだけの法則があることになります。それは第5節で述べた指数関数的な法則です。

10. 事件の生産に関する知識から、これらの仮定に類似したものが実際に多くの場合で確保されていることは容易に認められる。しかし、これがあproximately（近似的に）そのような状況を占めるとしても、それは厳密な正確さに近い形で人工的に確保される可能性は非常に低いと断言することは不可能である。なぜなら、偏位の原因がほとんどまたは全く相互独立ではないからである。いくつかの原因はしばしば、複数のものが一方に向かって作用しているという仮定が、それぞれのものに単独で働きを果たす能力に影響を与える可能性がある種類のものである。例えば、標的への射撃の一般的な例では、十分に精度の高い射撃の場合、風（誤差の一因）による効果は、弾丸の正確な方向に関係なくほぼ同じであると考えられる。しかし、目標からかなり外れた場合、風は飛行線に対して直角に作用するとは限らないため、その効果は以前と全く同じではない。つまり、ここでの原因は厳密には相互独立でないことが示唆され、したがってそれぞれの結果は他のものによる影響を受けずには存在しない。おそらくこの効果は微々たるものであるが、私はこれらの総合的な要因がどのように協力して作用するかを精密に調査すれば、絶対的な独立性の仮定が非常に多くの場合で危険であり、あるいは正当化されないものであることが明らかになると考える。これらの短い議論は現在の目的には充分であり、次章を通じてこの主題についてより詳しい検討を行うことにより、その内容を深めていくことにする。

従って、現在の段階でこの主題に提供できる最善の考慮に基づいて、推論的な議論の線から直接統計への訴えを引き出す類似の結論を導くことができる。同様の一般的な結果が確立されているように思われる；すなわち、実用的な目的には十分な精度で、偏回帰の原因を調査することで得られる結果は、通常受け入れられている「誤差の法則」を与えているものであると述べることができるだろうという結論だ。したがって、指定された範囲内で両方の探求は互いに明確な相互確認を提供する。

11. (III.) なお、第三の証拠として、誤差の法則が常に一つであり、それが常に同じであることを立証するために提出される可能性がある间接的で数学的な方法についても言及しなければならない。その法則が常に同一の法則に従って集約されていることを主張することができる；すなわち、数学と天文学者によって知られている最小二乗法という方法を用いて、様々な異なるケースにおいて非常に満足できる結果を得ていることから、その方法が適用される誤差は一定の法則に従ってグループ化されているものであると仮定することができる。もしすべての「誤差の法則」が同一のタイプではなく、つまり、大きな違いと小さな違い（先ほど述べたようなもの）の相対的な頻度が一つのパターンに従っていない場合、どのようにしてその全てに対し同じ方法やルールが適用可能か？

この治療法を継続するためには、この調査について少し注意を払う必要があります。しかし、最後の議論の場合と同様に、現在の段階では主題に対する完全な議論は不可能です。一つには、数学的な概念の使用や、少なくともそのような概念を使用することにより、本稿の全体計画に適さないほど多くの数学的使用が含まれるためです。従って、私はこの問題について特別の一章を割り当てました。

ただし、ここでの議論を展開する主な理由は、そのような議論を行うと、これまで扱ってきた確率の科学の全く異なる側面が予想されることであるからです。これは特に強調すべき点であり、この点の無視は多くの混乱と一部の誤りを招く可能性があります。

これらの初期の章では、私たちは確率の物理学的な基礎を築くことに完全に没頭してきました。我々はただ一つの方法でもう一つの方法でも、自然の中で見られる特定のグループや配置が存在することを証明しようとしました。そのような現象がどのように生じるかを説明し、またそれらの主要な特性を示すことを試みました。しかしこれらは推論の基礎に過ぎません。まだこれらの基礎上に建立される論理的な過程について一言も述べていません。

したがって、確率の論理についてはまだ着手していません。

12. しかし、最小二乗法に関する方法の説明方針は、この区別を理解するのに障害となる傾向があります。著者は最小二乗法と誤差の法則を同義語として扱っているが、実際には両者は全く異なる概念であり、互いに必然的な関連性すらないことが明白です。誤差の法則は物理学的事実の表明であり、それは実践で見られるどの種類の誤差や偏差も実際に現れる頻度をより正確なほど詳細に記述します。したがってこれは科学の基礎と呼ばれるべき物理的な基盤に属します。

一方、最小二乗法は科学的な意味での法則ではありません。それは単なる規則または指示であり、我々がどのような誤差群に対しても最善の方法で対処し、その結果を導き出すための手順を教えてくれます。従ってこれは推論や論理的な部分に属します。

確かに、我々が採用する方法は、その方法を適用する事実の配置とは何らかの関連性を持つことは否定できない。しかし、二つは性質においては別々であり、この場合における関連性は必ずしも必要ではなく、至多適切さと便利さの問題に過ぎない。最小二乗法は、確実に最近まで我々が取り組んできた誤差の最も一般的で広く使用されている形、すなわち指数形式に対して適用されることが多い。しかし、誤差の他の形態も存在する可能性があり、もし存在した場合でも、その方法はそれらにも同様に適用可能である。

私は必ずしもそれが各々の場合において最良の方法であるとは主張していないが、それは可能な一つの方法である。実際、次章で示すように、それはほぼ全ての場合において良い方法であると言えるだろう。しかし、その方法の特有の優点や欠点は、我々がそれを採用する任意の場合における適用可能性を妨げない。したがって、ここでできる限りの議論からも明らかになるように、一つの方法が非常に一般的に良好な結果をもたらす頻度が高いという事実は、その方法が適用される誤差が一定の法則に従うべきであることを証明するものではない。

13. 以上は、この特定の離散法則がすべての場合において支配的であることを証明しようとする試みについての議論を述べた。次に、クエテレットがこの主題について扱う際に注意すべき誤りや混乱の原因となる点は、彼と他の人々が主張する「タイプ」の存在に関する教義である。これは、前数節で提示された一部のデータと結論から自然な結果と言えるものである。§4で既に述べた三つの物種のうちの二つを思い出す。もし私たち自身の一連の不正確な観測や試行、あるいはある一つの種またはクラスに属する自然物体の一連の集合を並べ替え、長期間的にはそれぞれの場合においてその離散法則が同一であることが本当に見出された場合、私たちは当然、「誤差法則」という表現を両者について同じように適用しようとするだろう。しかし厳密には、それは最初の例にのみ適切である。自分自身で操作を行い、何を目指しているか明確な意識を持っているとき、各々の偏差がすべて誤差であると述べることは完全に正当であろう。しかし、自然がどのような種類の物体のグループを私たちに提示したとしても、そのような状況でも「誤差法則」という表現を使うのは大胆な比喩と言えるだろう。彼女が常に何かを目指していたかのように、そしてほとんどすべての場合で目標を達成できなかったかのように話すのは、私たちと同じようにずれがあったと仮定しているからである。

[13] ここで引用した文は、原著の英語版を直訳しています。原著では、括弧内に[13]と番号が付けられており、おそらく脚注や参考文献の一環として追加情報があることを示唆しています。この部分は原文の意味を保ちつつ、自然な日本語表現に調整しました。

次のような状況を想像してみましょう。人間の正確な身長を測定するために、長い系列の試行を行ったとします。そうした場合、何らかの理由からほとんどまたは全く絶対的な正確さでその身長を測定することは稀にしかなく、また決してないかもしれません。

しかし、私たちは自分の不完全な測定が、ある特定の誤差法則に従って、実際に異なる人間の身長の系列の正確な反映となることを前提としてよいわけではありません。どのような系列の直接的な測定値においても、その誤差法則は事前に断言することは困難であり、おそらく経験によってそれを決定しようと試みる機会が十分に多くないため、それを確認することができないかもしれません。

しかし、一般的な観点から見れば、その系列の誤差法則が所謂指数法則を徴するとは決して明らかには言い難いように思われます。しかしながら、それが果たしてそうであるかどうかは別として、自然が私たち一人ひとりが同じ方法で働きかけていると話すのは、一種の言葉遣い上の許容範囲内に過ぎないかもしれません。それは（大抵は無効に終わる）特定の結果を達成しようとする試みであり、つまり一定の身長や比例関係を持つ人間を作り出すことを目指しているからです。したがって、そのような人間は典型的な人間として見做すことができるかもしれません。

14. 上記のように述べれば、人間には固定された不変の種類があり、すべての人間の個体はその種類に達するべきものとみなされるが、それからどの方向かによってずれ出ている；このずれの法則は\_a priori\_で決定可能であるという原理はほとんど狂気以外の何物でもありません。しかし、事実を少し詳しく見つめて、これらの事実の可能性のある解釈を見つけると、重要な真実に到達する道が開けます。その事実は、クエテレットの統計（その大いなる興味と価値は公然と認められるべきです）によれば、以下のようなものであります：私たちの身体のどの要素も正確な測定が可能なものを選び、例えば身長を選び、どのような比較的均質な人種の多くの異なる個体についてこの測定を行います。すると、これらの身長は既に繰り返し述べられたように平均を中心に整然とした配置を持っていることが見つかるでしょう。

「均質なクラス」とは何を意味するのか？これは適切で重要な問いですが、この条件を単純なケースに適用すると、その意味は容易に説明できます。それは測定対象の人々の国籍によって平均値が異なることを意味します。クエテレットによれば、イギリス人の場合、平均は約5フィート9インチで、ベルギー人では約5フィート7インチ、フランス人では約5フィート4インチです。成人男性の測定結果であることはもちろんのことです。

15. ここで疑問にされるべきは、英語とフランス語を別々に扱わずに、測定結果すべてを混ぜ合わせた場合の結果は何であったかということです。この質問は重要な点を突いており、私達が何を意味するのかという均質なクラスについてより明確な理解が必要となります。通常、この疑問に対する答えは大まかに正しくても、あまりにも決定的で簡潔すぎます。英語とフランス語の結果を混ぜ合わせることは、明らかに異種の要素を混ぜ合わせることであるとされ、その結果、誤差の合計法則は以前示された単純な形にはならないだろうとされます。

この答えが部分的に認められる限り、その根拠は容易に理解できます。通常の誤差法則によれば、平均値からの乖離が量によって増加するにつれて、それらが減少する数が連続的に減ります。では、英語とフランス語の高さを混ぜ合わせた場合、どのようなことが起こるでしょうか？英語の平均5フィート9インチから始めると、最初はこれらの英語条件によって決定された法則にほぼ従います。この時点では英語のデータが非常に多く、比較的フランス語のデータは少ないです。しかし、フランス語の平均値に近づくにつれて、数は英語スケールによる配置に従うべき連続的な減少を示さなくなります。ここではフランス語のデータが再び非常に多く、比較的英語のデータは少ないのです。

このような異種要素の組合せの結果は、付属する図で例示されていますが、これは非常に過大な形で示されています。

[図：2つのガウス分布の重ね合わせ。]

16. 上記の場合において、不均一性の本質と、統計を収集し整理する理由が、比較的明白に示されていました。人工的な事例から類似のケースを取り上げることで、その点はより明瞭になります。例えば、ある男が特定の場所（壁に固定されたワッフル）に向けて数千発の銃弾を撃った後、彼が射る位置が移動し、新しい位置でのワッフルに向けてさらに数千発の銃弾を撃つとします。この場合、両系列のすべての銃弾を、どちらかの中心からの距離順に収集し整理すると、実際には異なる不協調な要素が混ざり合っていることがわかります。それぞれは、他の要素から分離していれば単純かつ均一な性質を持っていました。その結果、誤差の総則が通常の形から著しく逸脱し、上記図で示したような形を取ることになります。

イギリスとフランスの高さの例は、先に述べたものに類似していますが、必要な条件が厳格に満たされている点では大きく欠けています。実際、ここには最も適切な要件が得られていないのです。すなわち、少数の固定した原因を補完する無数の小さな干渉要因からなるグループです。我々が呼ぶ「国」は、地方的または偶然の干渉要因にさらされている多くのメンバーで構成されています。フランス人の中には、ブレスト人、プロヴァンス人、アルザス人などが含まれており、これらはフランス人とイギリス人の違いよりも小さいかもしれませんが、全体として見れば非常に重要なものではありません。そして、これらの種族的な差異に加えて、気候、食事、職業による重要な干渉要因も存在します。したがって、イギリスとフランスの統計を混同する反対意見は、当然ながら程度は異なるものの、フランス人の様々な地方的および他の構成要素からなる統計を混同することにも適用されます。

17. 人間の身長に影響を与える重要な原因は多種多様であり、その中で「誤差法則」に最も近い条件を満たすと考えられるものは、我々が最も知るべくもないものである可能性が高い。食事や就労などの観察結果は、これらの要因の影響について何かを語ることができるが、親の身長が子孫の身長に与える純粋な生理学的要因については、おそらく真の知識と呼ばれるべきものはほとんどない。我々ができる最も良い仮説は、「似たもの同士は似る」という言葉に従って、純粋な生理学的要因が恒常的な要素を代表することである。つまり、最初から均質な人々で構成され、自由に結婚し、気候、食事、職業などの類似した状況にさらされる場合、基準は全体として恒常的に保たれるであろう。

このような状況において、平均身長、平均体重、平均体力などを所有する男性を、一種の意味で「型」と呼ぶことができるかもしれない。この型からの逸脱は、生理学的、物理的、社会的な無数の小さな影響によって生じるが、これらは大抵独立して作用し、通常の誤差法則に至る。このような制限と解釈を考慮に入れると、「フランス型」や「イギリス型」という言葉を使用することについて合理的な反対はないと考えられる。しかし、現在の政治国家におけるこれらの多少混成的な集団は、それぞれ異なる小さなグループに分割される可能性があり、これら各々はしばしばより明確にそのような型の特性を示すであろう。

18. この点では、既に参照した人間測定委員会の報告書は非常に示唆的である。それらはこの分類がどの程度まで進められるかを示し、クェテレットの「平均人」の発見が我々を長い旅に導くことを証明している（もし必要であれば）。彼らの結果によれば、「英語圏」の平均身長（インチ単位）は67.66である。しかし、これはスコットランド人、アイルランド人、イギリス人、ウェールズ人のそれぞれの平均身長68.71, 67.90, 67.36, 66.66から成り立っている。さらにこれらの数値もまた分割可能であり、精密な観察によれば、特定の地域（例えばノース・イースト部県）では平均英語圏の身長が他の地域よりも明確に大きいことが確認される。

また、専門職者の平均身長は労働者より著しく高く、誠実で知的な人々の平均身長は犯罪者や痴呆症患者などの人口構成要素よりも非常に高い。そして、観察が十分広範であればあるほど、これらのクラスの中で見つかる特性群（特に広範なクラス）における特徴も狭いクラスでも同様に見つかることが示唆される。

さらに、分類をどの程度まで進めても、同じ分割過程は我々が望む限り延長できる理由は存在しない。

以上の記述に加えて、進化論の教義に多少なりとも賛同する者は、上記で限定的に用いた「型」が実質的な恒久性と固定性を有しているとは考えないであろうことはほとんど言及する必要もない。もし不変の原因が変わらなければ、そして長期的には変動する原因が相互に相殺し合うならば、結果は同じ平均値を中心に集積し続けることになる。しかし、不変の原因が徐々に変化した場合、または変動する原因が相互に相殺することなく、その中の一つや複数が優位性を獲得して全体的な効果に傾斜を与えるならば、平均値は一言で言えば直ちに位置を移すことになる。そして一度移動を開始すれば、種の変異と進化という事実を我々が認識する範囲においては、その移動は継続することになるだろう。

種が変異し、進化が事実である限り、平均値は直ちに位置を移すだけでなく、その移動は一定の程度と一時的なものか、あるいは相当の程度と永久的なものかのように進行する可能性がある。我々が射撃を続けるにつれて、私たちが的めざす点が徐々に位置を変えるかのように想像できる。この位置は確実に一定の程度と一時的に移動し、おそらく相当の程度と永久的に移動する可能性がある。

20. この章を通じて、我々の例はほとんどが物理的な特性から引き出されており、人間や無機物に関することである。しかし、これらは必ずしも我々が限られた例にとどまらなければならないとは限らない。例えば、Mr Galtonは、これらの計算原則を精神的な現象にも適用することで、それらのより正確な定量化を目指すことを提案している。

このようにして達成するべき目的は、主として当該主題の推論的側面に関することであり、後ほど詳しく説明されるが、それは新たな原理を導入するものではない。科学的手法を精神領域に適用する他の試みと同様に、この提案にも反対意見があった。しかし、私にはその理由は非常に薄いと思われる。

我々の精神的な特性が正確な測定が可能であれば、通常の誤差の法則に従うことを仮定することは、大きな疑問を残すものではない。心身特性間の既知の関連性から、最も多く測定された身体的要素が大体近似的にその法則を遵守していると証明されれば、精神的な要素も同様であるという推測は高い確率を持つと考えられる。

これらの精神的な要素についての推測は、それらが正確な測定が可能であれば通常の誤差の法則に従うことを仮定することから導かれるものであり、そして、最も多く測定された身体的要素がその法則を遵守していると証明されれば、精神的な要素も同様であるという推測は、既知の心身特性間の関連性から高い確率で成り立つと考えられる。

このような測定が実際的にどの程度可能かは別の問題である。私には、それがあまり成功しないものだと考えられる。これは、私たちの精神的な質は非常に密接に結びついており、実際には互いに溶け合っているためで、比較のために孤立させることが不可能だからである[16]。この困難さは肉体的な測定にも一部存在するが、心的な測定においてはさらに深刻である。ここでは、ほとんど確かな推測を超えることは難しい。したがって、精神的な質は現在広く知られている配置の法則に従うという教義を根拠として立てるのは、私たちの評価と測定が信頼できるものである限り、おそらく十分な類似性しかない。したがって、この試み自体は不合理ではないと思われる。

つまり、人間を平均値より上と下に分けて一定数のグループに分類し、それぞれのグループが優れたまたは欠如した特定の限界に対応することを目指す自然分類システムを確立する試みは、何ら不自然なことではない。必要なのは、平均値からの離差率が広範な異なる状況下でも比較的一定であるということだけである：この場合、人間のすべての種族においても同様である。もちろん、もし分岐の法則が無機自然界における法則と同一であれば、より広範で自然な分類システムが手に入り、それ自体は評価を行う人々にとってある程度は馴染みのあるものとなるだろう。

以上の理由から、このような試みは不合理ではないと思われる。

21. おそらく、このような原理の正当な適用例として最も良い例は、ガロント博士の『遺伝的英才』という著作に見ることができるでしょう。実際、その著作内のいくつかの論理には、この章で議論してきた概念を理解している人だけが、その真髄と全貌を十分に把握できるでしょう。私たちはスペースの都合上、一つや二つの点のみを指摘するしかありませんが、少なくともその書籍のより論理的な部分を読むことで、現在議論中の教義に対する興味深い例を見ることができるでしょう。

まず一点として、不慣れな人には絶対に理解できないことが言えるでしょう。誤差法則に馴染みがない人は、優れた才能がどのようにして希少になるかという速度を、ほとんど理解することはありません。もちろん、誰しも「一千万分の一」であるという数値的な意味合いはすぐに理解できますが、その人があなたよりわずかな優位性を持つ「二千万分の一」であるとは、あまりにも少ない優越性であることを理解することはできません。彼らは、「一千万分の一」という単なる数値的差異を、何か特別に優れた存在であると混同し、実際の優越性は小さな分数的な利点でほとんど表現されるという事実を見落としてしまうでしょう。

「一億分の一」は非常に豪華な響きですが、考慮されている品質が選択の広大な範囲から無知であり、自由に数値的な偏りから解放された場合、人々は通常、そのような選択によって得られる比較的小さな優越性を驚くことでしょう。

22. ここで述べた点は、統計学的な議論において重要な点です。例えば、血縁関係によって結びついた小さなグループの人々が、ある精神的特性に関して顕著な優越性を有している場合、そのような卓越性の比較的稀少性は、これらの質素が共通に持っているかどうかを決定する際に大きな影響を与えます。このような決定は必ずしも正確なものではなく、しかし、それが行われる限りこの考慮事項は要素として入るべきです。

また、一つの民族と別の民族を比較する場合、例えばアテネ人と現代のどのヨーロッパ民族かを比較する場合、一般的な優越性を示す証拠として、一民族からソクラテス、プラトノ、そして彼らの同時代人などといったグループが生み出されることがどれほど重要であるか、一般の人々は認識していますか？この後者の例では、私たちはまた、「誤差法則」を二つの異なる方法で用いています。つまり、極端な値が大きい場合、平均値もそれに伴って大きくなると仮定します。言い換えれば、どの程度の平均からの逸脱がどの程度の相対的な頻度で起こるか（粗略に計算可能）という前提を立てています。

この真理は、広く曖昧な形で受け入れられているとしても、その証拠は、その支持理由を説明できる人々によってのみ認識されるものであることが理解されます。

しかし、同じ原理は最後の例の場合にも警告を与えます。それらは、比較の目的だけを考える場合、どのグループやクラスにおいても平均的な人間よりも優れた人物を選ぶよりも、より良い対象となることを思い出させます。彼を特定するためにはより多くの困難があるかもしれませんが、その困難を克服した後には、比較のより確実で安定した基盤を得ることができます。彼は、社会の中で最も多数派である層から選ばれます。一方、優れた人物は希薄な層から選ばれます。従って、平均値と大量数の現在の法則に従って、前者の変動は一般的には非常に少ないし、後者のものよりも小さいことがわかります。

1. \_社会物理学試論\_, 1869年。\_人間計量学\_, 1870年。

2. 同じ主題に関する後の統計については、読者はイギリス協会の人体測量委員会報告書（1879年、1880年、1881年、1883年；特に最後のもの）を参照することができます。これらの報告書は私にとってクェテレットによって得られた結果よりも著しい進歩を示しているように思われ、その統計が「範囲と数において一類のもの」であるという秘書（Mr C. ロberts）の主張を完全に正当化するものだと私は考えます。これらの報告書は軍人志願者だけでなく、ほとんどすべての階層と年齢、そして両性も含んでいます。さらに、これらは身長だけでなく、他の多くの身体的特徴にも言及しています。

3. どのような数学的な知識を持つかに関わらず、各可能な投げ方の相対的な数は（1 + 1）10 の展開の連続する項によって与えられます。すなわち、1, 10, 45, 120, 210, 252, 210, 120, 45, 10, 1。

4. つまり、それらはより密集して集まることになります。各連続する円の中心に「ターゲットの大きさ」という領域を検討した場合、その中に含まれる射撃マークの数は次第に減少します。実際には、ターゲットが他のどの円よりも広い面積を覆っていないため、最も多くの射撃マークが含まれているわけではありません。

5. 通常、指数法則と呼ばれるもので、その方程式は y = Ae{-hx{2}} の形を取ります。この曲線は y 軸に直角で交差し（平均値近傍では多くの値がほぼ等しいという事実を表現）、一定の時間後には x 軸に向かって急激に傾斜し始めます（平均値から遠ざかるほど結果が稀になるという事実を表現）。そして、x 軸は両方向で漸近線となります（どの大きさがあれほど遠く離れていても絶対不可能ではないという事実を表現；つまり、どんな過度の逸脱でも十分な長期間の経験の中で必ず遭遇する必要があるという事実）。この曲線は明らかに対称的であり、長期的には等しい逸脱（過剰と不足）が同様に頻繁に起こるとする事実を表現しています。

[図：ガウスの正規分布。]

曲線の粗略なグラフィカルな表現が上に与えられています。数学に不慣れな人々のために、いくつかの性質についての短い補足説明を加えることが適切かもしれません。

(1) 曲線のすべての例が互いに類似しているとは思ってはなりません。点線はその曲線の標本でもあります。実際、xとyのそれぞれに対する推定単位を変更することで、曲線の頂点部分を任意に鋭くまたは鈍くすることができます。この考慮事項は重要です；なぜなら、これらの任意の単位の1つを変更することによって、極端な乖離が非常に稀であるという前提に基づいて任意の対称的な誤差曲線に十分よく似た「指数曲線」を得ることができるからです。したがって、特定のケースにおける誤差法則が指数曲線でないことを単なる観察によって証明することは困難であり、統計が非常に広範であったり、実際の結果が指数形から著しく異なっている場合を除いてです。

(2) いかなるグラフィカルな表現も、時間とともに曲線がx軸に近づく速度の過度の速さを十分に示すことは不可能です。私たちのスケール上の点Rでは、曲線はx軸から15,000分の1インチ以内まで近づき、これは非常に良い顕微鏡でもしかして検出できる距離です。一方で、双曲線の場合などでは、曲線が漸近線に近づく速度は常に減少しますが、ここではその逆です；この速度は常に増加します。したがって、視覚上では、非常に短い時間後には曲線とx軸が互いに溶け合るように見えます。

6. クェテレットによって： Herschel など、他の人々も注意を喚起している \_Essays\_ の409ページ。

7. \_Proc. R. Soc.\_ 1879年10月21日。

8. ここで考慮しているのは有限な統計量のケースであることを覚えておいてください；したがって、それぞれの端には実際的な限界があります。

9. 実験はまだ（私としては）高さに関してこの非対称性を示していないという点については認めなければなりません。

上記の論理は、この調査段階においておそらく正当性が認められるでしょう。しかし厳密に言えば、ここでは自己で正当化されてもよいと思われる前提が採用されていますが、それらは多少の予測を含んでいます。これらは今後の章でより精密な検討と批判を受ける必要があります。

11. この種の平均を中心にした頻度の集中の具体的な数値例は、第4節の注に与えられていました。それは二項分布の形をしていたもので、(1 + 1){m} の展開の連続する項から成っていました。今、(Quetelet, \_Letters\_, p. 263; Liagre, \_Calcul des Probabilités\_, §34)では、このような二項分布が m が無限大に近づくにつれて、指数関数的形式へと漸近的に近づくことを示すことができます。つまり、等間隔の縦軸をそれぞれ1, m, m(m - 1)/(1·2), ... に比例させて取り、その頂点を結びつけると、得られる図形は y = Ae{-hx{2}} のような曲線を近似し、m を無限大に増加させることでそれが完全に一致するようになります。言い換えれば、有限の等しい独立した原因によって生じた誤差が存在する場合、その誤差の指数法則への近似は原因の数が増加し、その大きさが無限大に減少するにつれて完全に同一になることになります。Jevons (\_Principles of Science\_, p. 381)は、この近似の速さを示すスケール図を与えています。

ここで一つ重要な点を注意しておきます。それはしばしば見落とされる点です（Quetelet もその一人でした）。二項分布において、等しい項を持つもの（例：(1 + 1){m}）の係数は最初から対称的で、すぐに最終的な指数関数的形式に近づいてしまいます。しかし、一方で (1 + 4){m} の項を考えると（これは m 回の試行における0回、1回、2回、3回...の失敗の相対的な確率を表します）、その項は対称的ではありません。ただし、m を無限大に増加させることで、この非対称性が徐々に消失し、二つの等しい項から始めると同一の指数関数的形式へと近づくことになります。唯一異なる点は、曲線の頂点の位置が中心ではなくなり、つまり最も確率が高い項や事象は成功と失敗の1対4の比率で存在することです。

12. 「誤差の法則」という技術的な用語は、以前に §4 で述べた「平均からの乖離の法則」を指すものである。厳密には、その定義は三つの現象類型のうちの一つ、つまり第三のクラスにのみ適切であるが、便利な一般化により、他の二つのクラスにも適用される。したがって、どの場合でも平均からの乖離量を「誤差」と呼ぶ。

13. これは、クエテレットの二冊の詳細な著作、すなわち『社会物理学』と『人間測定学』において直接的にも暗示的にも述べられている内容のように思われる。

14. クエテレットはこれらの結果を正確な平均身長の尺度として提供することをほとんど主張していないし、必ずしも同じ尺度を与えるわけでもない。実際には、大規模に測定されたのは兵士だけであり、イギリス人の身長は大きなスケールでの正確なデータを提供しなかった。数ページ先で示される統計はおそらくより信頼性が高い。

15. 次の章では、この声明について説明と修正が行われる。

16. ここで言及しているのは現在広く知られている心理物理学の結果ではなく、主に感覚や他の単純な意識状態の測定に関わっているものである。

17. ガルトン氏の方法について最も簡潔な説明は、『Mind』（1880年7月号）で発表された「精神的想像に関する統計」の論文中に見つけることができるかもしれない。比較対象となっている主題は、異なる人々が過去に存在しなくなった物体を鮮明な視覚像として再現する能力の相対的な強さである。これは通常の意味での測定は不可能な主題であるように思われる。しかし、その能力を持つ人々のすべての回答を関心のある順序で並べることにより、ある程度正確に中央の人を選んで他のバッチ内の同様の人物との比較基準として使用できる。そして、中央の人以外の相対的な位置にある人々についても同じように扱うことができる。

18. 初章と終章を参照してください：本書の大部分は、その性質上、主に統計的および生傳的な詳細に占められています。

19. ガロトンの『遺伝的英才』（Hereditary Genius）を参照してください。pp. 336-350、「異なる民族の相対的な価値について」です。

\newpage

\chapter{系列の起源と因果の過程（Origin or Process of Causation of the Series）}
\label{ch:iii}

\begin{center}
\large\textit{対象と作用因・自然種と自由意志・因果性の物理的基盤}
\end{center}
\vspace{4mm}

第三章.

確率に関連する様々なグループや系列が、厳密に同じタイプであるかどうかについて議論した際には、それらが自然に生じる過程について一部の検討を行いましたが、現在はこの過程の詳細についてさらに掘り下げてみることにしましょう。すべての事象は、数多く複雑な前因後果によって生じます。実際には、その数と複雑さは私たちが全てを決定しまたは考慮に入れることが不可能なほどです。

しかし、完全には全てを決定することは決して可能ではないという点については厳密に正しいことは言えますが、推論（Induction）の場合と確率の場合では、それを行う能力に大きな違いがあります。推論の場合は十分な前因後果を見極め、合理的な確信に達するほど正確に把握することが可能です。しかし確率の場合はそうはいかないのです。

将来的な章で再び推論について扱う予定なので、ここでは詳しく説明することはありません。

我々は、哲学的な正確性を主張していないものの、事前のガイドとして役立つ分類法から始めることに便利です。それは対象と影響その対象のものという単純な分類法です。確率に関連する全ての現象（科学全般に関しても大半は同様）は、自然や人造物が特定の影響力の下で生じたものです。例えばコイン投げの場合では、対象は順に投げられたコインまたはコインであり、影響因子は投げる行為そのものと、直接的または間接的にこの行為に関与した全てのことです。これは単純かつ理解しやすい分類法で、我々が扱うすべての対象を含む意味での拡張も容易に行えます。

もし、二つ以上の場合において同じ対象が存在したり、区別不可能に類似した対象が存在し、これらの対象がすべての点で完全に類似した影響を受けるとすれば、我々はその結果も完全に類似すると期待するべきである。自然の均一性という原理の一つの応用により、先ず前の状況が完全に類似であれば、その後果たされるものも完全に類似であると確信できるだろう。同じ銭片や類似した銭片を全く同じ方法で投げれば、必ずその表が上に出ることを見ることができる。

2. 実際に見つけるものは、もちろん、それから非常に遠く離れている。対象物が人工的な構築物である場合、例えばサイコロやコイン、カードなどにおいては、これらは可能な限り区別なく作られていることが真実です。同じものを繰り返し使用するか、または厳密に同一のモデルに基づいて作られた異なるものを使うかのいずれかです。しかし自然な対象物では、そのようなことはほとんど見られません。実際、これらの対象物を調べるときには、最終的な結果（我々が説明しようとしていたもの）において見られる特徴と全く同じ特徴が再現されていることがわかります。したがって、それらの特性をさらに前の段階で検討する必要がある場合でも、ある程度はその程度で、我々が説明しようと試みた現象の特異性を再度仮定しているように思われます。

例えば、寿命を考えるとき、対象物は私たちが考慮する人間（またはその特定のクラス）であり、それらの人々の寿命です。彼らの間に存在する類似性は、主な生命器官の強さと健全さ、そしてそれが我々が「体質の良さ」と呼ぶすべての状況を含む集団的な類似性に見られます。これらの多くの状況の大半は実際的な測定に近づくことができない真実です；しかし、前章で指摘したように、そのような測定が可能な多くの状況はすでに測定され、その特性を示していることが確認されています。

したがって、我々のさまざまな器官間の既知の類似性と相関関係から、人間の体質全体を一般化して並べ替えたり、それらを構成する個々の要素を特定して並べ替えたりすると、その強さに基づいて順序付けられた場合、最終的な結果（この場合は寿命）が我々に示すような集団的規則性と中心値周りのグループ化を見つけることが可能です。

したがって、この目的を達成するためには、自然種やグループの存在が不可欠であることが観察される。当然のことながら、当たるゲームにおいては同じサイコロを何回でも投げることができ、同じデッキからカードを何度でも引くことができる。しかし、人類に起こる多くの事象は全く繰り返すことができないし、個々の個人にとって十分な統計的均一性を確保するために頻繁には繰り返されない。自然現象における我々が追認する規則性は、通常よりも遥かに多いことが認められるように、物事が自然種に分類されることによるものである。それぞれの自然種は多くの個体を含むからである。

もし動物や植物の各種が単なる一対またはそれ以上の数だけに限定されているなら、その中で統計表を集める余地はほとんど残らないだろう。あるいは、より激しい仮定を考えるとしても、各自然種内の個体数が現在よりも著しく小さかったり、その品種や亜種間の差異が著しく大きければ、それらから十分な長さの統計表を抽出する困難性は避けられないが、致命的ではないかもしれない。クラス内に多くの対象物があり、それらが相当程度類似していることにより、一つのクラスに適切に含まれることのできる条件が、確率理論をどの現象にも適用するためには重要である。

このクラス内の対象物の過度な少なさと同様の状況は、労働者の数が非常に限られているし、雇用が危険で個別の業種として分類されるような産業への特別な保険事務所の適用において見つかるだろう。例えば、火薬工場での労働者だけを対象とする保険会社が設立された場合、可能な破綻を避けるために極めて高い保険料が課せられることになる。これは適切な統計の極度の少なさを示しているのである。

4. 現在までに述べてきた対象物についての議論を終わり、我々が「要因」と呼んだものに移ります。これらの要因は、その相対的な強度の調整と、それぞれの頻度によってもたらされる効果全体が、ほぼ同一であることが確認されました。もちろん、これが異なる結果となった可能性は考えられます。第二群の条件が第一群を完全に補正し、一般的な一貫性を絶対の一貫性に変換する可能性や、逆に第二群の影響が第一群を著しく歪曲し、その効果全体の一貫性を破壊する可能性も考えられます。実際には、どちらの状況でも起こりませんでした。第二群は詳細を変えるだけであり、全体としての一貫性は以前とほぼ同じ一般的な説明に留まりました。あるいは、対象物が完全に同一である場合、例えばコイン投げの連続的な試行の場合、一貫性をもたらすことが可能です。分析により、これらの要因は無数の異なる成分から成り立っていることがわかりますが、それらの多くには我々がしばしば言及した特徴が広く存在します。これら成分の組み合わせにおける割合はほぼ同じであり、作用する強度もほぼ等しいことが確認されます。そして、より多くの例を平均化することで、ますます完璧な規則性に結びついていくことが見られます。

例えば、人間の寿命について考えてみましょう。これまで見てきたように、ランダムに選ばれた多くの人の体質は、大まかな同一性と個々の不規則性を伴うような特徴を示すことが多くなります。

これらの人が世界に出ると、様々な要因に曝け出されます。これら要因の集団的な影響が、各個人に与えられた寿命を決定します。これらの要因は当然無数であり、それらの相互作用は分析によって完全には解明することができません。それぞれの効果はまた新たな原因となり、他の無数の原因と絡み合って最終結果に反応します。気候、食事、衣服などはそのような要因の一部です、あるいはそれらを含む集団的なグループです。人の仕事の性質も重要です。一人は過労し、もう一人は不健康な職業に従事し、また一人は感染症に曝露します。

これらの相互作用の結果、我々がこれまで呼んだ対象と要因との間に存在するものは、個々の要素から独立して検討した場合でも同じような一貫性を示すような一般的な特性を持つ最終的な結果が現れます。あるいは、我々は直ちに示すように、そのような結果は大半の場合においてそうなります。

5. しかし、このような説明は実質的に何かの名に値するものではないという反論があるかもしれません。それは特定の状態がどのように引き起こされるかを説明するのではなく、他の場所でも同じ状態が存在することだけを指摘しているからです。一般的な一貫性は、通常計算のために提出される対象において発見されます。その後、その原因を探求し始めますが、結局はこれらの原因にも同じような一貫性が存在することを発見します。これは部分的には正しいですが、反論の一部は取り除くことができるものの、最終的に客観的な科学の基礎は、事物が特定の特性を持つという事実に依拠するものであることが常に残るのです。

6. この分類は、対象とそれらを影響する因子に分けたもので、主には実用的な概略的整理のためであり、読者に我々が日常的に経験する一貫性をもたらす直接的な原因の性質を指し示すのに十分なだけです。さらに後ろに一歩進むとすれば、これらは因子にまで還元されるべきものであるという主張が公正に成り立つかもしれません。対象とは我々が呼称したように、現在の状態で見られる形で最初から存在する原始的な創造物ではなく、全般的な類似性と個別の差異が完全に発達しているような一連のグループや階級は同時に生成されたものではないと考えるべきです。もし各自然種の初代親族が既存の形式から発展したものではなく特別に創られたとしても、現在見られる各群数々において特徴的な差異と類似性は我々が因子と呼称したものがもたらした結果であるという事実は変わりません。

たとえば、身長という一つの特性について考えてみましょう。その身長が特定の一群の人々の中でどのように定められているかを決定するものは何でしょうか？おそらく一部は、彼ら自身の食事、衣服、就労などの性質自体、特に彼らの生活の最初の数年間において；また一部は、おそらく彼らの親が一度やるたびに経験した類似した状況と条件によって決定されるでしょう。現在の知識水準においても、誰も残りの原因を列挙しようとする試みやそれらの正確な性質についての何らかの示唆を与えることを試みることは好ましくないと考えるのは当然でしょう。しかし同時に、多くの人はこれらの一般的な説明に該当する因子が決定的な原因であったという疑念を抱いていないと考えられます。

もしこれらの作用そのものが最終的に分析されるべき何かに分解されるのであれば、その問いに対する答えは、本稿の計画から逸脱することになるでしょう。しかし、それらの作用についての一般的な注目すべき点が存在し、ほぼすべての種類の作用に対応する場合、私たちは前章で議論した起点に引き戻されます。つまり、すべての数列が必然的に一つの固定された法則に従って形成される必要があるかどうかについて議論したときからです。

そこで私たちは、重要な原因は比較的少数であり、これらは通常の観察によって注意を向けられる全てのもので構成されていると考えることができます。そして、それら以外には、個々の作用を把握するのに十分な数の小さな原因が存在します。これらの小さな原因は、その作用が非常に多く、細かい点や不確実であるため、個別に考慮に入れることは不可能であり、またそれらを評価することすら困難です。

確かに重要な原因もまた、より小さな成分から成り立っていると考えることができます。しかし、比較的単純で明瞭な扱いが最善であると見做すべきです。なぜなら、その成分は主に集団として作用し、同時に現れたり消えたりするため、これらは一貫した統一体の特性を持っています。

7. まず概して言えると、確率の最も適した条件は次のようだと考えます：重要な原因は比較的固定され、恒常的なものであり、残りの原因は平均的に一方に作用する回数と他方に作用する回数が同じであるべきです。これらは二つの方法で達成できます。まず第一に、これらの原因について何も予測できない場合でも、それらは一方に作用する回数と他方に作用する回数が同じであるという単なる事実だけを予測できるかもしれません。その結果得られるものは、第一章で説明した単純な統計的均一性に過ぎません。しかし、これらの小原因が互いによく独立して作用することもあるし、実際的にはそれほどまでには近似してそうであることが多いです。その場合、我々は第二章で説明し、例を挙げたような一連の均一性を得ます。これらすべての原因の可能な組み合わせが一定の頻度で現れるなら、我々は単なる余剰と欠如（どの変動する量に関しても）だけでなく、その特定の量の余剰と欠如も含む一つの特異な均一性を見ることができます。したがって、この場合、我々は「平均」や「型」と呼ばれるものを得ます。それらは単なる平均とは異なります。

例えば、矢を標的へ投げるという状況を考えましょう。固定された原因はその力、矢の重さ、そして特定の点に狙う意図です。これらの要素は我々が求めるような均一性を得るために当然一定であると仮定しなければなりません。変動する小原因は最後の章で述べた無数の微細な影響すべてです。これらの小原因について、我々はただそれらが一方に作用する回数と他方に作用する回数が同じであるという事実だけを確認できるかもしれません。その結果、矢は標的より遠く飛ぶ回数と近づく回数が等しい傾向があることを発見します。しかし、これらの小原因が互いによく独立しており、影響力も多少均衡しているため、各量の余剰と欠如が同じ一般的な特性を示し、多くの投げ方において、特定の距離からの逸脱量はその逸脱量に従って一定の関数として定義されます。

8. これらの条件が欠けている場合に見られる例を参照することで、その必要性は最もよく理解できるだろう。
つまり、全体として生活の長さは比較的規則的に推移しており、また、多数の一般的な病気で、連続する年や世紀における死者数も大抵は一定であることが知られている。しかし、すべての病気にはそうではないように思われる。たとえば、発汗症、黒死病、アジア性コレラなどについて何を言えるか？
前二つは再発しないか、またはもし再発したとしても、その形態が非常に穏やかで、同じ名前を付けるに値しないほどである。実際のところ、流行性疾患よりも地元性疾患であるものは何か？これらすべてには原因があり、かつての条件が再現されることによって再び発生するだろう。しかし、一部は明らかに再発しない。その発生は非常に稀な条件に依存していたため、その発生はほとんど唯一のものだったように思われる。そして、再発する疾患の中には、その経過が頻繁に不規則で、しばしば人間の意志や意志の欠如によって大きく左右されるものがあり、これらにより、我々が述べている統計的な一貫性は完全に失われてしまう。

おそらくこの説明は、このような場合の主要な原因の一つが通常私たちが「感染」と呼ぶものであるからだ。もしそうであれば、固定されたものではなく、最も多様性を示す原因であることが明らかになる。厳格な注意はそれを破壊し、無関心はその程度を任意に増加させる可能性がある。人間の意志は政府、医師、または公衆のいずれかによって表現され、それが望むようにほぼ何でもできるかもしれないが、その侵入を完全に防ぐための予防措置は、どのようなコミュニティにもその可能性がないことを保証することはできない。

9. もし、統計的な規則性の欠如は時間が制限されていることに過ぎないと反論され、十分長い期間を考慮に入れればその不規則性は他所と同様に消滅するとすれば、我々はその回答が求める確率の条件が存在しない別のケースを提示されることを見出すでしょう。このような反論は結論的であるためには、人類の行動や思考が長期的には変化しないか、あるいは周期的な変化のみを経験するものであるという仮定が必要です。社会が長期的に向上または悪化の一途をたどるという前提においては、この議論は一発で崩れ去ります。我々が世界の歴史から学んだところによれば、これらの恐ろしい過去の脅威は私たちの歴史上孤立した出来事と見なされるか、少なくとも再現されないものとして扱われるでしょう。よって、登録官の記録簿を千年間開けても、そのような事件によって引き起こされた死亡数には連続的な一貫性を見つけることは不可能です。

この理由はおそらく、我々が総称して「進歩」または「文明」と呼ぶ無数の条件の徐々な変化に求められるべきでしょう。このような小さな事項の一つ一つも、誰かが病気を患う可能性に影響を与えます。しかし、これらの影響が一定の平均強度で均衡し合う状態から、ゆっくりと確実な流れによってその影響力が増加する傾向を示すようになるとき、統計はそのような変化を受け入れ、それに伴う安定性や永久性を持つ形でその印象を受けることになります。

次に、確率の特徴が最も顕著に見られる場所と、最も欠如している場所について簡潔に述べるとすれば、（1）それは主に自然種類の性質において prevalently 、最終的な性質もしくは派生的または偶然的な性質において prevalently 現れると言えるでしょう。すべての自然種類の特性、彼らが行うことや起こること、そのような性質によって依存する範囲で見れば、これらはしばしば規則性を示すことが少なくありません。人間について言えば、身長、体力、体重、寿命、死に至る病気など、すべてが既知の均一性を呈します。生命保険表はこれらの確率の応用がいかに重要であるかの最も親しみやすい例を提供しています。

（2）同じく、多くの自然作用力の強度と頻度においてもこの特徴は再び見られます。風や天候は、大規模な観察によってその昔から言われるような不規則性を失うことがわかります。したがって、これらの作用力によって行われた人間の作業は、最初に全く同じように作られたとしても、その後には全体的な規則性のみが現れます。例えば、私は毎年同じ量の種子を畑にまきます。その収穫量は1年目には中程度で、次の年には好天により豊富になり、また翌年には雹のために破壊されることがあります。しかし長期的には、これらの不規則性は作物の結果において等しく化されます。これは生産作用力の強度と頻度が等しくなるからです。保険業者や雹による収穫損失を補償するための事務所もこの分類に属します；既に述べたように、これらは前者の分類とは深く区別されるわけではありません。

読者にはもう一度注意していただきたいが、この固定性は一時的なものであり、すなわちここに示される数列は、振動型を持つ数列のクラスに属する。実際、自然種の固定性を信じる者は、その中から真の永久的なタイプの数列を見つける最高の機会を得るだろうが、それでも彼ら自身も時間の経過とともに特性が変化することは可能であると認めている。

主要な自然作用素においては、当然のことながら現在の平均値は現在の地質時代に限定される。私たちの平均気温や平均降水量は過去には現在とは大きく異なり、おそらく将来もまた異なるだろう。

進化論によれば、原始的な単純さと均一性から現在の多様性が生み出された過程についてのより詳しい調査は、本著の範囲からは外れている。至多、数点のヒントだけが必要である。

11. 以上は自然の対象や自然作用素の例である。この性質は、人工的なものと区別してこのようなものにのみ見られるという理由があると考えられる。これは議論と説明を必要とする主張である。読者の一部にはすぐに思い浮かぶ二つの反対例が存在する。一つはその大きな内在的重要性から、もう一つは問題の頻度から、それぞれ数分間の個別の検討が必要となるだろう。

（1）これらの中の最初は、既に述べた機械的観測の場合です。天文学者や他の機器を使用する際には、可能な限り高い精度が望まれることが多く、これは単一の観測では達成できないほど高くなることがあります。したがってこれらの状況では、多くの連続的な観測を実施し、その結果は自然に多少異なりますが、これらを利用して最も正確な真値（後の章で方法論の最小二乗法について説明します）を求めます。計算の対象となるのは、次々と得られる若干誤った要素の数です。これらの観測が人工的であるか、または意図的な結果の直接的な産物でしょうか？もちろんそうではなく、あるいは正確には何を意味するかによって答えは異なります。観測者が真実に至るまでに意図し、目指すのは当然のこととして完全な精度、つまり真の観測、またはその観測に依存する意志的な手順や前段階です。これらの観測が意志的であるかどうかは問わず、これが唯一の意図的な結果と呼び得ます。しかしながら、この結果を得ることはできません。実際に得られるのは、それから逸脱した一連の偏差であり、その結果は真実から多少離れたものがあります。これらの偏差はどのように生じるのでしょうか？これは、本章の一部の前段階で考慮してきたような諸要因によって引き起こされます。熱とその不規則な歪みの影響、空気の流れがそれに応答する効果、一部の部分での塵やそれに伴う摩擦、金属を鋳造した後も続く徐々に不均一な収縮による微小な歪みなど、これらの要因は真実から逸脱させます。このグループ以外にも、人間の意志に関連するものがありますが、厳密には意志的とは言えません。それは筋肉の不規則な作用、様々な器官や部位を我々の意図に正確に従わせる能力の欠如、あるいは神経伝導速度の異なる速さ、または刺激に対する反応の違いなどがあります。これらの一部による効果とそれに伴う調整が必要となることは、天文観測、心理物理学的観測その他の観測において「個人差」の名前で一般の人々にも広く知られるようになっています。

12. (2) それ以外の例としては、カードやサイコロを用いたものがある。この場合も前回の場合と同様に、結果は一時的に意志的な選択が現れる段階においては自由意志的である。しかし、その後の段階では、結果は多くの非意志的な要因によって生じたり影響を受けたりするので、その特性はこれらの要因に依存している。

例えば、サイコロの特定の面が出るかどうかは意志的な行動の結果であり、しかし直ちな結果とは言えない。その特定の面を選んだわけではないが、選ばれた事実は選択行為の遠因である。意志的な行動と反対に作用する多くの矛盾した要因が中間段階で混在しており、これらの要因は事前に計算することはできず、その後でも区別することは難しい。これらの要因は長期的には一定の統一性を示し、その結果にも類似の統一性を生む。

カードを一組から引くという行為も直接的には意志的なものであるが、ワーストでパートナーを決める際にカードを切るような場合とは異なり、この過程では間隔のある間にデッキをよく混ぜることによって多くの非意志的な影響が加わる。

13. 自殺や誤送信のような一時的に驚くべき結果を示す統計的パターンは、それ自体で別のクラスに属している。これらの結果につながる最終的な決断（またはその欠如）は個々の行動や不作為において重要な要素であることは確かだが、全体としての結果に対する意志的な影響がある場合でも、それは瞬時にその統一性を破壊する。

立法府の声が響くか、大きな説教者や道徳家が阻止に成功するか、または印象的な事例が影響を与えるかなど、これらの要因によって我々の道徳的統計は瞬時に変化する。この主題については将来的な章でさらに議論する予定だが、ここでは（上記のような一般的な影響を除く限り）個々のケースにおいても強力な意志的な行動が全体としては非意志的な条件、つまり健康状態や雇用上の事故などによって左右されるという点に注意baruだけである。これらの要因は実質的に人間の寿命を決定する様々な要素である。

14. これらの区別は、ある人々には余計な区別のように思われるかもしれませんが、それらの区別を無視した結果、時として深刻な誤りが生じることがあります。人間の心の直感的な産物については、その本質的な確率性という特性を持つように思えません。これらの特徴は、完全なる数学的正確性か全くの正確性の欠如、法則的な不変性か単なる随意性に似たものであると考えられます。例えば、森の中の木が一直線に生えていた場合、その木々は人間によって植えられたと断定するのは疑いようがありません。しかし、それらの木々が規則的に植えられていないことを発見したとしても、それが人間によって植えられなかったという結論を下すことはできないでしょう。部分的には、プランナーが計画を持たずに作業を行った可能性があるからです。また、その後の自然による不規則性が計画を覆い隠した可能性もあります。

実際には、心は欠陥のある道具に依存して働き、様々な矛盾する要因によって妨げられることもあり、これらの影響により作業は本来の特性を失うことがあります。例えば、ある人が不完全な観測やサイコロの目を使ってではなく、単純に思い浮かべて最初から数値結果を選択したとします。そのような系列は何らかの確率的な系列を得るよりも、心臓や脈拍を直接意志で制御しようとすることほど困難でしょうか？しかし、この提案自体は荒唐無稽です。人間は思考する際に目的を持ち、規則や公式に従って思考しなければなりません。ただし、自然の系列を模範としない限り、そのような系列を心の中で永久的に再現することは不可能です。

また別の例として、人間の努力の産物を考えます。この作業は比較的成功する可能性があります。例えば、誰かが家を建設した場合、ブロックやコンクリートの収縮、基礎の沈下などの微小な干渉要因が働きます。しかし、これらの影響が生じる効果は非常に小さく、結果として得られたものは心の直感的な実現とみなすことができます。その結果は何らかのことが起こりますか？

一列に並んだすべての家は、一人の設計者によって一度に設計された場合、隣接する家と同じ高さ、幅などです。もし変動があるとしても、それは少数で明確かつ規則的なものです。その結果は、ランダムに選ばれた多数の人々の身長や体重などとは全く類似性を持ちません。buildersはおそらく一定の計画を持っており、それを実現したのです。

15. 以下のように反論することができるかもしれません：観察を拡大し、例えば大きな都市の家々まで調査対象に含めれば、その性質は明らかになるでしょう。多くの異なる高さが集まると、それらは類似した治療を受けた多数の人間の高さと同様の分布を持つかもしれません。このようなことはおそらく見つかるかもしれませんが、類似性は決して完全ではありません。しかし、この問いを提起することは、我々が以前述べたように、個々の心の働きとその異なる目的ではなく、一つの心の働きについて話すことを意味します。多数の設計においては、単一の設計では見つけられないような変動する均一性があるかもしれません。builders（建築家）が考慮した高さは、彼らが独立して行う他の多くのことでも同様な均一性を持つものとして集まることを見つけるかもしれません。そのような結果を導く同じ二つの条件の作用を追跡することができます：多数の異なる設計における均一性と、それらの設計を無限に多様化した影響の均一性。しかし、これは一つの人間の作業がこのような結果を示すことを言うこととは全く異なります。その違いは、千人の人が互いに意識せずに歩く足取りと、彼らが連隊として訓練されたときの足取りとの違いに似ています。前者の場合には、千人分の心がどのように作用しているかが働いている一方で、後者の場合には一つだけの心が働いているのです。

この章と前章の調査は、我々が関心を持つ統計的な結果の特異な形態が実際にどのように生成されるのかについて十分に詳細に検討し、主に確率論の科学の方法に関する著作のためにその目的を果たすのに充分です。しかし、特定の統計的または社会学的研究の重要性は、将来的な章でこれらの統計の一つの具体的な応用に戻る必要があることを要求します。それは人間の行動に関連するいくつかのクラスに焦点を当てた統計研究であることに注意してください。

16. 前述の観察に対する重要な追加または修正は、ガロント氏によるものである。彼は最近、そしておそらく最初に指摘したように、ある特定の場合において因果過程の分析が可能であり、実際には観察された事実を説明するために絶対的に必要である。

たとえば、どの国の人口の身長について考えてみよう。これらの身長が平均値を中心に分布または分散する場合、上記で言及した無数の生産的要因に制限なく作用させられたとしても、我々は確かにその世代の子孫において同じような配置を見ることができるだろう。つまり、親世代における同様の秩序が前段階で再現されるべきである。

しかし、ガロント氏によれば、これらの要因が自由かつ独立的に作用した場合、同じ「種類」の配置や分布は得られるとしても、その「程度」は異なるものとなる。逆に、分散の傾向がある。つまり、「能力曲線」（p.29の図を参照）は同じクラスに属するものの、モジュールが異なる。

これについては、似たような要因が完全な反対作用なく働いた場合でもすぐに理解できる。例えば、最初に平等な財産を持つ人々が小さな金額で賭けや博奕を継続的に始める場合を考えよう。時間経過とともにその状況を観察すると、彼らの財産は同じ一般的な法則に従って分布されることが期待できる。つまり、我々が現在議論している馴染み深い法則に従うだろう。しかし同時に、最も貧しい人々はさらに貧しくなり、最も裕福な人々はさらに裕福になることも予想できる。

この点については後述する（「賭博」の章）で詳しく説明するが、ここでは確率の法則には何ら抵抗がないと仮定してもよい。つまり、極端を強化する傾向に反対するものは存在しないのである。

現在、反対に、生命現象において、例えば身長の場合や、おそらく他のいかなる方法でも特徴的な自然種の質的特性の多くも同様であるが、数多くの連続した世代を通じて、驚くべき程度の固定性があることが見出されている。高い男たちはさらに高くならないし、低い男たちはさらに低くならない。人口全体の百分率において、連続した各世代で言えることだ。ただし、当然のことながら、気候や食生活などの一般的な状況が変化していない限りである。したがってここには、極端を引き寄せることによって分散を制御しようとする何か作用している原因があるべきである。

17. これらの事実は慎重な実験によって検証された。ガロント氏の最初の論文[3]の日付において、人類の身長に関する統計データは存在しなかったため、物理的な要素（特定の種子の大きさや重さ）が精密な実験を可能にする要素から正確に評価された。これらのデータから、極端からの反復還元量、つまり極端のメンバーに対する微かな圧力によって平均値へと押し戻される量が決定され、これは理論上要求される量と比較された。それ以来、人類の身長に関する直接的な統計データが相当数得られるようになった。

我々が導き出す一般的な結論は、これらの原因がいずれも微細でないし独立しているという点である。例えば、極端が平均値ほど繁殖力が高いわけではなく、また死と病に耐える能力も等しくないという観察事実がある。したがって、彼らの数においては、それらが少しずつ薄くなる傾向がある。さらに、「反復還元」に関する明確な正の原因もある。極端の子孫は少ないだけでなく、これらの子孫は平均値に近い位置に集まり、つまり全体のグループの真の中心からわずかに移動した平均値の子孫になる傾向がある。すなわち、平均親の平均子孫に近づく傾向がある。

18. この特性についての詳細な議論と、その作用形態や相対的な効果を示す多くの巧妙な例証については、Galton氏の原著記事をご参照ください。現在の私たちの目的には、これらの特性が種の固定性を維持する傾向があるという点と、それらは「確率曲線」または「能力の法則」の一般的な形態に影響を与えないものの、その具体的な値を決定することがあるという点だけを述べれば十分です。もし射撃のような例でなぜこのような補正の必要性がないのかと尋ねられた場合、それは累積的な影響が導入される場所以外では補正の余地がないからであることが答えとなります。

当たった金額が常に増大する賭け当事者の運勢が恒久的に分岐し続け、生命現象の場合には特別な補正が必要となった理由は、各段階での新しい出発点が前段階の結果によってわずかに決定されることでした。一度一シリングを失った男は、次回はその分だけ不利な状況から始めます。補正が示唆されていない場合、生まれたときから高い男は、その子孫を優位な高さの起点から始めると言えるでしょう。射撃者の例における真の類似性としては、彼らの新しい目標点が常に最後の分岐方向にわずかに移動していると仮定することです。そうすると、射撃痕は無限に広がり続け、同様に、想像上の賭け当事者の運勢も無限に分岐し続けることになります。

1. 上記のように、これは実際には、私たちは経験によって提示される具体的な詳細において説明する必要がある同じ種類の均一性が存在することを再確認するだけです。それらは過去の段階からさらに後ろに遡るものです。

2. 実際、粗雑で未熟な観測では、誤差は主に少数の重要な原因から生じると考えられ、その結果、我々の仮説（指数法則による誤差に関するもの）は事実を正確には表していない可能性が高い。そして、誤差の頻度は、それを導き出す法則に大まかで曖昧な近似値であるだけである。しかし、天文学者がその精度に満足せず、残存する誤差源を探求したとき、三つや四つの重要な誤差源ではなく、多くの比較的小さな誤差源が重要性を伴って明らかになったことが判明し、それらはこれまで取り除かれていた重い誤差によって隠され、その影響を被っていた。これらの誤差には、刻度の誤差があり、機械の収縮による多くの他の誤差も含まれていた。また、調整に関する誤差も存在した。技術的には観測と呼ばれる誤差、温度や天候の変化による誤差、微細な不規則な運動や振動による誤差などがあった。要するに、現代の天文学者が熟知している千種類もの微小な干渉要因があった。

3. 伝統的な遺伝法則；1877年2月9日にロイヤル・インスティチューションで朗読された。また、1885年11月の人類学研究所誌を参照されたい。

\newpage

\chapter{系列の発見と証明法（How to Discover and Prove the Series）}
\label{ch:iv}

\begin{center}
\large\textit{経験的基礎と演繹的推論・理由不十分の原理の無効性・ベルヌーイの定理の限界}
\end{center}
\vspace{4mm}

第四章.

確率論の基礎となる系列やグループの特性の確認と数的比率を決定する方法について

1. 現在到達した点では、我々は特定の種類の系列やグループに所有しているものと想定される。これらの系列が科学の基礎として確立される底層にあるものとして考えられる。このような系列の特性については十分具体的に説明し、その一般的な形成過程についても示唆を行った。次に、個々のケースにおいてこれらの系列の存在をどのように確認し、その特殊な性格と性質をどのように決定するか？そして第二に、これらの系列を得た後では、どのような方法で論理的な目的のために利用すべきか？

前問に対する答えは難しくないように思われる。経験が唯一の指針である。我々が実際に存在する「事物」の系列を見つけ出すことを望むなら、その事物自体に訴えるべきであり、他處ではあまり助けを得られないからだ。一年間に生まれるか死ぬ人々の数や、一年間に焼失する家屋の数や船の数を知るためには実際に数えなければならない。このような「経験」について話すとき、我々はその用語を最も広い意味で使用することを意図している。我々は推引的または帰納的な論理が提供できるすべての援助を補完した経験を使用するという意味である。例えば、特定のクラスの人々が連続した年間に亡くなる人数の系列を見つけた場合でも、その系列を過去にも未来にも延長することは躊躇しない。このような手続きの正当性は通常の推引的原理から求められるべきである。確率と推引との関係について特別な議論が行われるため、ここではこの主題についてさらに述べる必要はない；しかし、第四章で述べられるものは、我々が使用する系列は最終的には経験のみによって得られたものであるという主張とは矛盾しないことが確認されるだろう。

2. 多くの場合、直接的な経験に頼る必要はないと断言できることが多い。たとえば、ワイスで10枚のトランプを持っている確率を知りたいとき、過去にそのようなことが何回起こったかを尋ねることはしない。もし地球上のすべての人々がワイスのパーティに分かれ、その問題を解決するために十分な調査を行うのであれば、非常に長い年月が必要となるだろう。もちろん我々は、10枚のトランプが出る可能性のある組み合わせの割合を代数的に計算し、これが私たちの問題に対する答えとなる。

また、六面のサイコロで6の目が出る確率を知りたい場合も同様である。その最善の方法は、重心から対面が見せる固形角を幾何学的に計算し、この固形角と球全体の表面積との比が求められる確率を十分に近似することである。

このような例では、実際には誰も統計に訴える必要を感じないだろう。ここでの事件が生じるための機械的な条件や他の条件は比較的少ない、確定的であり、個別に考慮できるもので、それらから構築できる組み合わせの大量性が、その結果、事件が起こる可能性のある方法の膨大な多様性をもたらすからである。したがって実際的には、事前に決定された確率の推定は容易であり、経験に基づく後方的な推測は単に煩雑で実用的ではないほど困難となる。この点と、これらの例が論理的に扱うと頻繁に単純かつ魅力的であるという点を合わせて考慮すると、これらは非常に人気があり、確率の理論を説明する適切な典型的な例として印象を与えることになった。

しかし、科学が実際には保険の対象となるような事件のみに関心を持っていたとすれば、おそらく他の見解は存在しなかったであろう。それは直接的な経験に基づく推測に依存しているという視点だけが採用されたであろう。

3. しかし、少し深く見ると、二つの例示されたクラス間のこのような明確な区別は必要ないことがわかります。サイコロやカードのような場合でも、決定条件や事件の可能性的な多様性から直接推理するかのように見えても、実際の観察結果ではなくその発生から推論するものよりも、この手続きが有効であるためには、明示的に明記されず、直接経験によってしか定義できない暗黙の前提が必要です。これは確かに非常に自然で明白な仮定であり、日々の観察から新たな重みを受けることがありますが、考慮せずにこれを認めることは適切ではありません。この点は非常に重要な問題であり、その理論に対する光を投げかけることでさらに重要になるため、私たちはこれについて詳細に検討することにします。

まず、非常に単純な例として硬貨の投げ方を考えましょう。私は二回の投げ方を連続して考える場合、唯一の可能な結果はHH, HT, TH, TTであると見ることができます。これは確定的な事実です。また、経験から、これらの結果が長期的にはほぼ等頻度で発生すると信じられています。これは当然のこととして受け入れられています。しかし、通常の見解では、この事実は観察結果ではなく、決定条件や事件の可能性性から推論するものよりも、あらかじめ知ることができると主張されます。この見解が一般的に提示される形は、上記の四つの投げ方が「等確率である」という声明を先立つことです。もし「等確率」という表現の意味を尋ねるなら、その答えには二つの可能性のみがあります。一つは観察者の心持ちから説明を求めるもので、もう一つは観察されたものの特性から説明を求めるものです。

(1) たとえば、一方では、四つの事象が等しく容易に想像できるという意味である、あるいは正確には、その発生に対する期待や信頼度が等しいということであると言えるかもしれない。しかし、これで満足するのではなく、さらに問われるのは、「なぜそれがそうなると信じられるのか？」という問いだ。私たちの期待が等しい事象の特徴とは何か？について答えることが求められるだろう。この追加の問いに対する回答を認めるなら、第二の答えに至る道筋となるが、これを認めないとしても、私たちは確率が心理学の一端であることを認めざるを得ず、その結果として事実自体に関連なく心の状態のみを考えることになる、つまり真偽に関わらず心の状態を考察することになる。すなわち、物事について推論する科学から離れてしまうことになる。この点については、他章で詳しく検討しなければならないが、論理学（そしてその一部である確率論）は人々が本当に信じていることを扱うものではなく、正しい信念を持つためには何を信じるべきかについての科学であるという事実に注意を向けすぎることは不可能ではない。

(2) 他方の答えでは、その表現の説明は心の状態ではなく、考慮される事象自体の特性から探求されることになる。私たちは次のような意味を付与することができるかもしれない、すなわち長期的には各事象が等しい頻度で発生することである。この主張の根拠はおそらく過去の経験に見出すことができ、私たちの信念の概念を完全には除外することは不可能だろう。しかし、以前のように信頼度の量的な平等を求めることと、ここでは事象自体の頻度の平等を求めることは、広範な違いがあるという点に注意しなければならない。

4. この程度まで進むと、経験への訴えは長く避けられないと容易に示すことができる。果たして、その主張（コインの投げ入れが等確率で起こること）を先験的に断言できるでしょうか？この主張を先験的に断言しようとする者は、困難に直面していることを十分に認識していないように思われます。真剣にコインが期待される通りに動くかどうかを問い始めるとき、制約条件を導入しなければならないことがわかります。

まず第一に、それは理想的なコインでなければなりません。その両面が等しく公平であることを意味します。この制約は完全に理解できるもので、立体幾何学の研究により、コインを円形または円筒形の板状体として理想化することができます。しかし、この条件だけでは不十分であり、他の制約も必要です。コインは「ランダムに」投げ入れると想定していましたが、これは意味するものとは何か、そしてこの過程をどのように理想化すればよいでしょうか？このような問いを提起することは、無駄な精巧さを提示することではありません。なぜなら、投げる直前にコインの位置を一定の側面から始めるという方法で、頭と裏が公平な確率を得るとは思われないからです。

その結果による差異は、原因こそ微細なものであっても時間とともに結果に現れる傾向があります。または、二つの側面を交互に上にして投げ入れる場合でも、同じ側面の長い繰り返しが公平な確率を得られるでしょうか？

おそらく、これらの問題について何も考えずにいると、すべては自然にうまくいくだろうと反論されるかもしれません。それはそうであり、確実にはそうなります。しかし、これは経験的な前提に依存することを意味します。そこで、我々は先述した経験的仮定に立脚していることを認識せざるを得ませんし、この仮定は避けられないものであることも明白です。なぜなら、自然の状況や私の身体的または精神的な構造が、常に同じ側が上に向かうように始まるか、あるいは実際には私自身によって任意の順序で始まるような場合を考えましょう。そうすると、反論されるでしょうが、それは公平な試験ではないだろうと。

このような問いに対する答えを求める形でこの点を追及するならば、これらの無意識的な制限は実質的に実験結果を確保する方法の一種であることがわかります。それらは「ある特定の種類の順序を確保するために一連の行動が行われるように」という別の言い方と等価です。

5. 直接経験に依存することを避けるための中間的な方法として、ある事件の確率をその事件が有利な状況である場合の数と全体で可能な状況の総数との比によって測定すると定義することが時々見られます。これは少し曖昧で不正確な表現法だと考えられます。単に有利な状況の数だけを数えるだけでは足りません、それらはまた評価されなければなりません。なぜなら、各々が等しく効果を生むわけではないからです。もちろん、その点は否定されることはありませんが、実際にはそれらを評価する方法として他に手段がないことを十分に重視されていないように思われます。つまり、それらが実際にもたらすまたはもたらす可能性のある効果を推定することによってのみ評価されるのです。

有利な状況の比ではなく、事件が実際に起こる場合の比に直接依拠することが望ましい（少なくとも科学の基礎に関しては、現在我々は実用的な計算方法について議論していないため）のは、明らかに良い方法です。

6. 上記の議論は、他の多くの確率的な例にも当然適用されます。サイコロを振る、カードや袋から玉を取り出すなどです。最後の場合においても、例えば、玉が先に戻された場合、その玉が次の回でもより高い確率で出ることになるだろうと自然に思われます。なぜなら、それらの玉の中にはその玉が再度出ることを妨げないよう、より適所にあると考えられるからです。この状況をどのように改善するかは？

意図的にその玉を他の玉の中央や底に押し込むと、過度な注意を払うことになり、また人間の意志が介入し、それが確率に対する我々の理解において根本的に敵対的な要素となる可能性があります。袋をよく振って状況を整えることを信頼しても、容易に欺かれるかもしれません。なぜなら、袋を振ることで、既に近傍にある玉同士が依然として近傍にある傾向を大幅に減らすことは難しいからです。各玉が他のすべての玉と相互作用する結果、その初めの配置は必ずしも最終的な位置には影響を与え続けるでしょう。したがって、このような状況では、言葉を吟味すれば、予想される \_à priori\_ の問題解決方法は、特定の結果を得るために手段を選択することの簡潔な表現に過ぎないことがわかります。

結論として、我々の推論は最終的にはその結果に基づくため、それが科学の基礎となるべきであると主張することが単純で哲学的な方法だと考えられます。

7. また、コインを投げての例を取り上げ、より詳細に検討し、実際にはどのような結果を得ることができるかについて何が証明できるかを見ていきましょう。コインに対して公正な扱いをするために、完璧であると仮定します。つまり、長期的には表も裏も好む傾向はありません。その場合、同じ面の繰り返しは通常の理論によって正当に得られる割合を獲得するでしょうか？簡略化のために、H を表、HH を連続して表が出るという意味で用います。この問題が残されたとき、次のような結論が導かれるべきであるとします：--H の確率が 1/2 である場合、二枚のコインを投げたときに HH が出る確率は 1/4 であることは当然でしょうか？さらに、「HT 十回」が一度に 1024 回中に出ることは必要でしょうか？数学的な大部分の学者は、この結論が基本原理から必然的に導かれるべきだと考えていますが、私にはそれが第一原理に基づくより確かな証拠よりも合理的な推引による拡張であるように思えます。

次に、二枚のコインを投げたときに得られる可能性のある結果について考えてみましょう。四つの異なる結果が生じることがあります：（1）HT, （2）HH, （3）TH, （4）TT。これらの四つが等確率であることを証明できれば、一般的に受け入れられている結論が導かれるでしょう。なぜなら、同じ議論はすべての大きな数にも適用されるからです。

8. 通常用いられる証明は「十分な理由の原理」と呼ばれるものを使います。この原理は次のように述べられます：四種類の投げ方があり得ます；それぞれの要素 H と T が等頻度で現れることを認めれば、それらの組み合わせも等頻度で現れることが必然的に導かれるでしょう。なぜなら、一つの組み合わせに対する理由は他のどの組み合わせに対しても同様に適用されるからです。

我々は、ある程度において、その原理の妥当性を認めなければなりません。上記の投げ方の場合、（1）と（3）の等頻度、および（2）と（4）の等頻度を証明することは可能であり、これらは互いに符号HとTの出現順序以外に違いがないからです。THはHTと同じであり、我々が導入した記号の順序だけが異なるのです。しかし、（1）と（3）のいずれかは必ずしも（2）と（4）のいずれかとは異なります。符号の表記を変更しても、ここに残る区別は心で認識することができます。同じものが連続して二回現れる状況は、異なる二つのものが結びつく状況から区別され、これは我々が任意に導入した表記法だけによるものではなく、この表記法の変更によって全く変わることはありません。従って、十分な理由の原理を認めるとすれば、たとえば（2）と（4）のような二種類の双子が等頻度で現れることが証明されるかもしれませんが、それらがそれぞれ四回に一度ずつ現れる必要があるとは証明されません。この方法では、これらが全く現れないことも否定することはできません。

9. したがって、公式は先験的に証明できない（結論として推測されたように）としても、経験から得られるでしょうか？ある程度においては可能です。現在の人間の銅貨やサイコロに関する経験は、理論によって割り当てられた比率に近い頻度で小さな投げ方の連続が実際に現れることが示唆されています。しかし、これらの連続がどの程度その比率に近いかについては誰も確実には言えません。小数の場合でも十分な時間を費やすまで事実を確認する感覚を得るのは非常に困難であり、大数の場合には全く受け入れられません。「ヘッドが十回」を得るための投げ方を十分に行う実験は、二人か三人によって実際に行われており、その結果はデ・モランとジェボンによって報告されています。しかし、これは平均的に「ヘッドが十回」を得る機会が一度しかないだけであるため、証拠は非常に薄く、多くのこのような実験が必要となるでしょう。

どのような法則であっても、我々が最近議論したような、長期間の投げ入れの一連の結果を描写するものについては、その有効性は経験によって限定された範囲内でしか確実には証明されません。すなわち、同じ面が繰り返される小さな回数においてのみ、その法則が完全に証明されます。より広い範囲では、我々は確信を持てます——法則が著しく誤っているとは思えません。なぜなら、それらの違いはほとんど間違いなく検出されてしまうからです。

このような観察から、我々は同様の法則が全体を通じて成り立つことを強く推測します。言い換えれば、我々は帰納法と類推によってその法則を拡張しようとします。しかし、自然現象において、提案された法則が狭い範囲で成り立つものの、広範な適用では著しく誤る例が非常に多いことを考慮に入れると、公式の真実性について我々は最高でも限定的な賛同しか示すことができることを認めなければなりません。

10. 上記の議論の目的は単に、法則が正しいかどうかを確信することはできないことを示すことです。次に、硬貨の表と裏の連続生成の原因について一時的に考えてみましょう。これにより我々はさらに疑問を感じるかもしれません。

既に述べたように、事前確率を計算する際には、制限や仮定を導入することでしかその可能性を推測することはできません。これは実質的には期待される結果を前提としているのと同等です。厳密には「特定のプロセスが行われる」という言葉を使いますが、言語分析を行い、そのプロセスが想定通りに行われない場合に感じられる無意識的な仮定を検討すると、我々の実質的な意味は「特定の種類の一連の結果が得られる」ということになります。その後の計算はこの一連の結果に対してのみ適用されます。物理的なプロセスが行われた後で、その結果に似た一連の結果が本当に得られるかどうかを知りたいのです。

現在、コインが常に同じ面を上にして投げられ、同じ回転速度と同じ高さに投げられるなど、言い換えれば、同一の条件で投げられる場合、それは常に同じ面で下りてくるはずです。実際には、個々の投げ方の結果は無限に異なるため、このようなことは起こらないことが知られています。それでも、これらの条件の平均値があるものがあり、その周りでは投げ方がより密に集まっていることが見られるでしょう。したがって、常に同じ面を上にしておけば、我々が通常期待するような系列とは異なる結果が出る可能性があると推測します。

大量の投げ込みにおいては、このような状況下で、頭か尾のどちらかが優勢を示すようになるかもしれません。もしそうであれば、各々の連続するペアの投げ方の方法にも類似した効果が見られるでしょうか？私たちは、HHやTTを常に上にしておけば、二つの頭または二つの尾の成功例の割合に影響を与える可能性があるかもしれません。

この推論を追っていけば、我々が期待するような系列を得るためには、最初から類似した系列で前件を配置することが必要であると結論付けることが可能です。投げることによって生じる変化や機会は無限に異なるかもしれませんが、長期的にはこれらの二つの系列間にも密接な相似性を見つけることができるかもしれません。

11. これは、ある程度は困難を移し変えるに過ぎません。以前、前シリーズの投げ方の比について証明する可能性があるという主張がなされたことを思い出してください；おそらくそれは後者の系列に対して再び繰り返されることでしょう。それでも、この問題は非常に狭められました。なぜなら、私たちは一連の「自主的な」行為にまで問題を縮小させたからです。人間は好きな面を上にして投げることができます。私は先ほど述べたように意識的に行動することができるか、あるいは全く何も考えずに（つまり、ランダムに）投げることができます；もし後者の場合、多くの人々が無意識的にそのような系列を生み出すことを強調するでしょう。それはそうかもしれないし、そうではないかもしれません；容易にアクセスできる決定的なデータは見つからないように思われます。私がここで心配しているのは、一般的に受け入れられている結果が実際には特定の条件が最初から満たされることによるものである可能性を示すことです。その条件は誰でも自由に満たすかしないかを選択できるものです。短い連続性は投げ方の偶然の変動によって無限の複雑さが生じることにより自ら解決するでしょうが、長い連続性は最初から意識的にまたは無意識的にその関心を考慮せずにいると問題になるかもしれません。

12. 「十分に試行すれば、いずれ可能な結果は何でも得られる」という助言は説得力がありますが、その根拠は唯推引と類推に過ぎず、数学では証明されていません。何度も述べられてきたように、この問題に対する見解には二つの異なる視点があります。一方で、一連のシンボルを「表」と「裏」呼ばれるものとして取り扱い、これらが長期的には一定の頻度で出現すると仮定することができます。それらの各々、さらにはそれぞれの組み合わせも同様に、長期的には規則的に頻度を持つと仮定します。その後、これらの様々な組み合わせを計算し、仮定されたデータから導き出せる結論を抽出することができます。これは純粋な代数的な過程であり、完全無欠で、その範囲は全く制限がありません。この視点は、私が「置換的」または「理想化した」系列と呼んだものに相当し、通常は計算の基礎として導入されるべきものです。注意すべき危険は、これを純粋な代数的概念として過度に捉えすぎることであり、これから自己意識から容易に導き出し、不適切な長さまで自由に推進することによる誤りに陥る可能性です。

あるいは、他方で、私たちは物事の性質を扱っていると考えることもできる。球体やサイコロ、出産や死亡などといったものについて推論を行うことができる。しかし、その場合、前者では許容された仮定がここでは経験によって検証される命題となる。

実際の確率論は、この二つの視点間で緊密な対応関係があることが当然前提となっている。すなわち、置き換えられた数列が実際の数列に合わせられていることを意味する。経験は十分に示している。当該例においては、大きな範囲内ではこのような対応関係が存在することを。

しかし、数学的な公式から導かれる遠く離れた推論を強制的に認める必要はありません。その試みにより、先ほど指摘した区別が目立つようになり、私たちには選択の余地が生まれます。あるいは数学に立ち向かい、それらについて議論する権利を失うか、あるいは物事側に立ち、数学を拒否するかです。

私たちは後者の形式的な正確性については疑問を呈しないが、その適用範囲外のものとして扱うべきであり、データの正確性を確信できないためです。または、これらの公式を再構築し、実際の事実に合わせることを選択する自由があるのです。

13. どんな教義に対する批評家が、その主張を説明し、反対の理由を正当化しただけでは、彼の役割はまだ半分に過ぎないと言えるのは稀である。彼にはなお、それが持つべき正当な機能や位置について、少なくとも数語で示すことが残されている。彼がそれを絶対的に無価値または意味がないとは考えない限り、それはしばしば起こらないからだ。

そこで私は言いたいのは、確率が直接的な対象への関連を断ち切られ、経験に基づいていないという点で実質的に分離された場合、それは単に一般的な組合せ論や代数学の教義に帰着されるということだ[5]。これらの考慮事項は純粋に形式的かつ必要なものであり、経験への依頼なしでも完全に論理立てることができる。私たちは純粋な数または量に関する考察から出発し、それらの考察と結びつけるのはただ算術的な計算だけである。

例えば、一便士を三回連続で表が出る確率を求めたいとする。この場合、私がすべきことはまず可能な投げ方の数を確認することだ。組合せ論によれば、二つの要素（つまり表と裏）があり、三回実行するプロセスを考えると、八通りの結果形があることがわかる。その八通りの中で一つだけが有利であるため、求める確率は八分の一と判明する。

以上の翻訳では、原著の論理構造や語調を尊重しつつ、自然な日本語に仕上げています。また、数式や変数はLaTeX形式で記述しています。

現在、確率問題の実際的な計算は必ず上記のような性質を持つことが真実であることは認めますが、数学を用いる実在科学と、数学そのものが唯一の成分となる形式的科学との間には、広範かつ重要な区別があるように思われます。経験への依頼を排除した場合、我々は調査の中介的なまたは計算的部分だけを保持するに過ぎません。サイコロや銭やカードについて話すかもしれませんが、これらは実際には私たちが用いるシンボルの名前であるだけです。私たちは扱うHとTは客観的な出来事とは関係なく、代数の残り部分で扱われるxとyと同じように機能します。確率をこのような方法で扱った場合、実際にはそれが他ならずとも適用可能な一連のシステム化された順列と組合せの体系に過ぎないのです。

このようにして、任意の純粋な推論法による治療が適用できる範囲は狭いことがすぐに理解できます。それはほとんどが確率ゲームのようなような用途に限定され、既に述べた通り、その信頼性を確保するためには様々な暗黙の制約が必要となります。この点だけで、主題に関する理論がこのような基礎に依存していることに反対する証拠となるでしょう。

一方で、実験的方法は同じ理論的な意味において普遍的に適用可能です。それは確率ゲームから提供される通常の問題だけでなく、サイコロが歪んでいて銭が完璧ではない場合なども含みます。また、統計学を様々な社会現象への無数の応用にまで広げることができます。

14. 以前に述べた確率の推論的性質に関する特定の見解は、比較的条件が少ない単純なものでなければならず、これらの条件は正確な数値的な決定に適しているため、おそらく他の例には浸透しないと考えられます。しかし、実際上はほとんど同じ理論をわずかに隠した形として広く普及している学説があり、最も純粋な経験的真理にも適用されています。この学説はラプラスの引用から最適に導入されます。

彼は最初の見かけの自然の不規則性と不確定性について述べた後で、「より深く観察すると、驚くべき規則性が見えてくる。これは設計を示唆するものであり、一部の人々にとっては神の存在の証拠だと考えられています。しかし、深く考えてみると、この規則性は単純な事象のそれぞれの確率の発展に過ぎず、より高い確率を持つ事象が頻繁に起こるべきであるという点で明らかになります」と述べています[6]。

この指摘がコイン投げで表裏の順序について述べられた場合、それは理解できるでしょう。それは単に次のような意味を持つことになります：その体質は、特定の方法で扱われたときに結果を予測することができると期待できることを示し、経験は長期的にはその予測を正当化することを証明します。

しかし、より一般的な形で自然現象の事実に適用される場合、この指摘にはほとんど意味がないように思えます。これを検討するために、個々の出生の不規則性の中から一例を見てみましょう。長期的には、男性児童が女性児童に対して約106対100という比率で存在します。もし私たちは「これはそれぞれの確率の発展に過ぎない」と被告知された場合、その声明には既に主張されている事実を多少誇張した再述しか残らないのでしょうか？

確率はただその比率そのものであり、このケースでは統計データ以外から得られる源は存在しないことが明白です。上記の指摘では、結果の順序を発生する割合の単なる数値声明から導き出す試みがなされているように見えます。

15. おそらく反論されるかもしれませんが、上記の確率とは、単なる出産数の比ではなく、その比率に依存する私たちの構造上の事実を指していると解釈すべきです。硬貨の両面が等しい関係を持ち、それが表と裏の最終的な均衡を生み出したように、私たちの構造や状況が106対100という比率に影響を与える何かがあるかもしれません。その何かが発見された場合、観察された数値は事前に決定される可能性があります。

しかし、それでも忘れてはなりません。このような原因と組み合わさると、最終的な結果を生む他の並行する条件が存在することです。硬貨の形状だけでなく、それを投げた際の「ランダム性」の性質も考慮に入れる必要があります。これらの何かは何であるか、現在では誰も確実に断言することはできません。最良の生理学者でも、この主題についてすら予想すら忌避しているようです[7]。

詳細を論じる必要はありませんが、これらの条件は父母の健康状態、状況、習慣や伝統（最も曖昧な表現を使用するなら）など、影響を与える要素から成り立つことは確実です。これらの要素がわずかな影響要因として導入されれば、結果は固定された永久的な条件に依存するだけでなく、変動する可能性があります。

これにより他の条件が加わります。これらも確率の特徴を持つと仮定されるなら、その事実を証明しなければなりません。この状況が実際に存在すること自体が、疑問視されるとともにほぼ不可能であると考えられます。これらの条件は一般的に「進歩」や「文明」と呼ばれる性質を持っていますから、平均値周辺に長期的に滞在する傾向を示すことは期待できません。

16. 確率について馴染みのある読者にとっては、当然ジェームズ・ベルヌーイの有名な定理が思い出されるでしょう。この定理は、先に示した例を一部の特殊ケースとして含むもので、一般的には次のように表現されます：長期的にはすべての事象は、その主観的な確率に比例する相対的頻度で発生することが傾向を持つことになります。この定理の数学的証明については、本作の範疇外であるため、ここでは触れませんが、果たして前節の批判に価値があるとするならば、その基礎は適切ではなく、実際には何ものにもpropriety（適切さ）をもって「主観的な確率」と呼ぶものが存在しないからです。

この定理が時々用いられる解釈や用途を見れば、それは最後のRealism（現実主義）の残滓であると見做すべきでしょう。ここでは他所で放逐されてもなお、確率という遠く離れた地域に残っているのです。これは概念を客観化する恒常的な傾向の一例であり、その概念が存在権利を持つべきでもない場合でも同様です。一定の均一性は観察されます；たとえば、ゲームの運命においては、それが使用される物体の物理的構造と密接に関連しているため、事前に推測できることがあります。しかし、ここでもその関連性は通常よりもずっと必要ではないことが、これらの物体だけでなく影響を与える諸因子との関係も考慮に入れるべきであるという事実によります。

この構造は「主観的な確率」として変換され、「均一性」を示す順列が発展するものと想定されます。最後に、この非常に疑問視される主観的確率は、ゲームの運命以外のどの場合でも同様の発展能力を持つと仮定されますが、それらがゲームの運命とどれほど類似しているかに関係なく、均一性が観察される限りです。

17. ほとんどの場合において統計的な均一性を見出すと同時に、そのような概念がいかに適切でないかは、一瞬考えれば明らかとなる。観察された現象は、これらの状況では非常に多く且つ複雑な前因の産物であることが多い。例えば、どのような社会において年間犯行数は、法の厳格さ、人々の道徳性、社会的状態、警察の警戒心など、様々な要素の関数であり、それぞれの要素はほぼ無限に複雑である。これらのすべての要因により、ある程度の均一性が生じるが、それ以上には前述したタイプの変化が明らかに現れる。平均年間数は、ある意味での制限値に漸進的に近づくという恒常的な均一性を示さないが、逆に振れ幅がある。これは原因によって異なるとしても、賭博ゲームのような永久的な均一性を特徴とするものには見られない。この事実と、観察された規則性を生み出す多くの異なる前因から僅かな数だけを選んで主観的確率として定義するという明快な任意性が組み合わさると、そのような存在の存在を仮定する必要は全くない。ただし、他の状況における誤った類推がそれを探求させた場合を除いてである。

同じ誤りの一種は、広く実践的な形として知られており、その傾向はおそらく類似の理論的源から生じる可能性がある。それは統計データを集積する傾向であるが、これも過度に行われると問題となる。

もし型が絶対的に固定されているなら、統計データを過剰に集めることが何ら害になるわけでもないだろう。我々がそれらの統計を集めることにどれだけ時間を費やしても、結果はより正確なものとなるからである。しかし、型が変化している場合、つまり、その生成に寄与する主な原因のうちいくつかが現在の強度に関して時間または空間的な制限がある場合、これらの制限を超えて集積を行っては好ましくない。早すぎる停止の危険性は容易に理解できるが、それを避けるために過度に長期間統計を集めてしまうと、徐々に或いは突然、変化した状況の影響を受けることになるという反対の誤りにも陥る可能性がある。

18. この章は、我々が関心を持つ数列の生成過程ではなく、その存在を決定するための方法の理論的基礎についての考慮事項に専念することを意図していた。しかし、二つの調査を分けることは不可能である。ここでは、古い格言が適用され、知るには原因を知らなければならないという原則があるからだ。

一度前の章で述べたことを思い出す。そこで私たちは、生成条件が（1）比較的少数のほとんど変化しない要素と（2）多数の独立かつ非常に変動する要素に分けられる一大クラスの事象があったことを確認した。もし他に要素があれば、それらは特別な考慮事項として扱われたり、調査から省かれたりするものと仮定される。

特定の場合、例えば賭博のような場合では、変化しない要素は時間や空間を通じて実質的に変化しないと見做すことができる。したがって、そのような場合において、推論的手法の扱いは最も単純で自然かつ結論的になる。さらに、事象の統計学的数据を必要に応じて収集すれば、真実に近づくことができる。

一方、社会的な確率の適用では、変化しない原因は多くの制約条件のもとでもしかして本当に変化しないものと見做すことができる。直接的にはそれらについてほとんど何も知らないし、しばしば実際には多数で不定かつ波動している。時間や場所に対する厳格な制限の下でなければ、その安定性を理論的に支持するのに十分なものとして属性付けることは危険である。

したがって、推論的手法は理論と実践において全く適用不可能になる。さらに、統計学的数据への依頼は、データを集めすぎると害を被る可能性があることを念頭に置かなければならないという注意が必要となる。

19. 最後の二章の結果は、以下のようにまとめられる：私たちは第一章で得た級数の概念を拡張した；つまり、これらの級数が主にグループとして現れることが判明した。これらのグループは、非常に広範かつ多様な経験の中で類似の型によって形成されていることが確認され、その一致の原因については詳細に議論し、説明した。しかし、系列を非常に長いものと仮定してさらに調査を進めると、これらは重要な区別によって二つのクラスに分けられることが見つかる。これらのクラスの一つ（確率ゲームの結果を含むクラス）では、生成条件が実質的に絶対的に固定されているため、そして従って統計的発生の法則も同様であるため、平均からの逸脱は有限な限界を超えることがないかのように扱うことができる。一方のクラス（通常の統計調査の大半を含むクラス）では、生成条件がより速い速度で変化するため、結果もそれに応じて変化する。また、平均からの変動はしばしば一定量を超えることが不可能であることが多い。前者を「理想的」系列と呼ぶことができる。これらのみが必要な特性に近い程度で示されるものであり、これらの不完全な他のクラスの系列を扱う際には、このような種類のものを別のものとして置き換える必要がある。しかしながら、これらの「理想的」系列は、主題に関する著者の多くによって過度に考慮されてきたため、それらに適した概念が他のクラスにも導入され、しばしばそのような概念すら適用できない場合もある。このような方法で、科学の一般的な傾向として予め定めた或いは推論的（à priori）な扱いが促進された。

1. この後者の調査は、この巻のより純粋に論理的な部分に属し、第六章の中で着手される。

2. この分野で一般的に使用されている記号を知らぬ人に対して、HHは単に二回続けてコインが表が出るという意味であることを述べておきたい。HTは最初の投げが表で、次の投げが裏出ることを示すものであり、その他の符号も同様である。

3. この十分な理由則について、受け入れる立場の人々から見ても正当性のある方法で扱おうとしているが、結果を得ようとすると矛盾が生じないかどうか非常に疑問がある。コインの両面が絶対的に同じであれば、数列の項同士にどのような違いがあるのかということが不明である。そのような場合、順序は単調で一貫したものに削減され、ただ同じことを繰り返すだけの単なる反復となり、我々が考えていた数列は消えてしまう。コインの両面が絶対的に同じでない場合、この則をどのように適用することができるのか？

4. \_形式論理学\_, 第185ページ。\_科学原理\_, 第208ページ。

5. ブライ爆出版の著書『選択と偶然』のタイトルはこれらの主題間の密接な関連性をよく示している。

6. \_哲学的試論\_, 1825年版、第74ページ。

7. 一度は一般的に信じられていた（クエテレットもその他の人々と共に引用し支持していた）見解では、親の年齢が子供の性別に関連している可能性があった。もしこれが真なら、上記の議論は完全に正当化されるだろう。実際としては、106対100という割合がすべての国やすべての時代で普遍的であるとは限らないことを注意すべきである。この主題に関する様々な統計表については、クエテレットの『社会物理学』第1巻、第166ページ、第173ページ、第238ページを参照されたい。

\newpage

\chapter{ランダム性の概念（The Conception of Randomness）}
\label{ch:v}

\begin{center}
\large\textit{一様分布の公理・幾何学的確率・円周率の乱数性・無作為抽出の数理}
\end{center}
\vspace{4mm}

第五章.

ランダム性の概念とその科学的な扱い

1. 機率に関する論文で頻繁に用いられる言葉として、既に繰り返し使用した「ランダム」という用語がある。この用語を正確に用いるための科学的な概念は、集団的な秩序と個々の不規則性（または表見上の不規則性）というものを理解することだと考えられる。前章で既に述べた通りである。この科学的な概念における必要な要素について、特に一般的な概念からの逸脱の性質や一部の理由についての簡潔な議論は、当該部分に関する主要な残存する難点を解明するために有用であろう。

「ランダム」という言葉の元来の、そして今でも広く受け入れられている意味は、科学的な用法とは大きく異なる。この言葉が注目するのは結果ではなく、ランダムな行為の起源である。また、個々の行動に焦点を当てているよりも、一連の行動に焦点を当てていることが多い。例えば、矢を「思いつきで」または「ランダムに」引く場合、我々はその動作の無秩序性だけを指摘するに過ぎない；これは特定の目標を射るための明確な意図とは対照的なものである。しかし、既に述べた通り、このような行為に対して推論の過程を適用できるのは、それらが頻繁に、あるいは延べして繰り返される可能性がある場合である。

例を挙げて説明しよう。雨粒の配置は科学的なランダム分布の典型的な例として最も適しているかもしれない。誰も特定の瞬間において雨粒がどこで落ちるかを予想することはできないが、紙を出してみると徐々に均一に点々とされることは分かっている；また、紙上の任意の2つの等面積について、これらは徐々に等しく打たれる傾向があることも理解できる。

2. I. 乱雑な配置から推論を引き出す試みは、そのような状態がどの段階で発生したかを前提にしなければならない。しかし、特定の段階での導入がしばしば困難であり、時には任意性も生じることがある。

(1) 例えば、数学家によって議論された多くの問題では、結果だけを見ており、その結果を得る過程についてはほとんど考慮していない。雨滴の配置を観察するときと同じようにである。この種の例としては、以下のものがある。直径が1インチのチェスパAWNを、各マス目の直径が1インチと4分の1のチェスボードにランダムに置き、その底面が交差線のいずれかに横たわる確率を求めよ。ここでは、パAWNを垂直にボードに雨のように降らせることができると仮定し、境界線に遭遇するものの最終的な割合を求めることになる。したがって、問題は単なる幾何学的問題となり、ボード上の特定の面積に対する全体の面積の比を決定することである。この比の決定こそが数学的に考慮される全てである。

次に、次の問題を考える。二つの場所でランダムに折られた直線的な脆い棒から三角形を作ることができる確率を求めよ。三本の直線で三角形を作るためには、それぞれの二つが第三よりも大きいという条件があるため、この問題は前例と同様の一般的な概念を含んでいるように思われる。棒が一対の場所で折れることが考えられるが、その位置は誰も正確に予測できない。しかし、これらの点が全体の長さに均等に分布する最終的な傾向があることを想像しなければならない。最後の例と同様に、数学家はこの最終結果だけを考え、それがどのように達成されるかについては考慮しない。したがって、問題は再度測定に関するものへと帰着され、より複雑な性質を持つものの一つとなる。

3. (2) ある種のケースでは、最終的な結果よりも中間過程を考慮する必要があります。しかし、その同一な概念はここで導入されるべきで、現在は前段階に適用され、従って一般的には後段階には適用されません。

例えば：最大射程（45°仰角での場合）が3000ヤードである銃から、ランダムに発射された弾丸の実際の射程が2000ヤードを超える確率は何でしょうか？ここでは最終的に一様（またはランダム）な分布は通常、銃を指す可能な方向全てが長期的には等しく表現されることとして仮定されます。つまり、水平面以上のすべての可能な方向が長期的には平等に表れることです。したがって、私たちは一様分布の表面を考慮する必要がありますが、これは地面の表面ではなく、発射者がいる中心を占める半球体の表面となります。地面で弾丸が打撃する場所での最終的な分布は一様ではありません。実際の問題は、分布密度の変動法則を見つけることにあります。

上記は、このような問題を解決するために一般的に採用される方法だと思います。しかし、そのような特定の選択には絶対的な必要性はありません。誰でも自分の発射のランダム性の概念を次のように定義することを選べると思います：人間が任意の方位点に向かって最初は等しく向き合い、その後は仰角を任意に上げることが最も確実であると仮定します。ここでの最終的に一様な分布はさらに前段階に退いています。これは直接虚想される半球体の表面に割り当てられることではなく、その表面上の高度と方位線に割り当てられます。したがって、半球体自体における分布は今では一様でなくなります—極に向かって比較的密集するようになります—そして地面での最終的な分布は以前とは異なります。

4. この種の困難は、一様分布の段階が何であるかについての不確実性から生じることが時々あります。例えば：本棚から適当に一冊の本を取り出すとします。特定の一冊を偶然に見つけられる確率は何でしょうか？この質問に対する一般的な扱い方については、私にはよくわかりません。

もし私たちがその男の位置を本棚の中間に設定し、彼が目隠しをして実際的な状況で何が起こるかを尋ねるとします。おそらく、各巻が等しく選ばれるわけではありません。むしろ、肩の高さによって示される中心付近では頻度が増加する傾向があるでしょう（左利きではない場合、開始点から正対した点よりもわずかに右側）。

もし本当にこの問題をこのような方法で計算する価値があるとすれば、我々は長い道のりを歩むことになるでしょう。私たちが射撃兵の位置（第二の仮定）を二つの独立な座標によって決定することを想像したように、私たちは彼の動きを一定数の独立成分に分析しようと試みることができます。彼から始点までの地面におけるすべての方向が等しく確率的であると仮定し、彼が本棚に到達したときに手のランダムな動きは、乱射のように制御されるものとすることができます。

以上の方法は問題を述べる一つの手段です。しかし、読者には理解していただきたいのは、私がここで主張するすべては、このような「ランダム性」の概念において、ある段階または他のどこかで、特定された量に対する最終的な一様分布というこの公理に直面することであるということです：時間、または空間（線形的、表面的、または立体的）のいずれか。しかし、その適用される段階を選択することはしばしば困難を伴い、甚至は任意性も生じ得ることです。

5. 数年前には、この主題に関する興味深い議論が『Educational Times』の数学的編集部で行われました（特にVol. VII. を参照）。数学においてしばしば起こることのように、様々な著者の間では特定のケースにおける仮定についてほぼ完全な一致があり、従って導き出される結論についても同様です。しかし、公理や定義の使用については非常に大きな違いがありました。

例えば、M.W. Crofton氏はWoolhouse氏との実質的な合意のもとで、「ランダム」という言葉には「非常に明確で確定した意味がある」と断言しました。「\_法則に従わないこと\_]という定義によってWilson氏が述べたように、この意味でのみ使用します」と。私は「法則」についての科学的な解釈によれば、「法則がない場合に推論は不可能である」と主張するべきでした。しかし、Crofton氏は無秩序性の最終的な分布の均等化を前提としており、これは彼の定義から導き出される結論です。「この無限の平行線系が法則に従わないで描かれているのであれば、[共通の]垂線沿いのどの部分でもそれらは他の部分と同じくらい密集している」と（Vol. VII. p. 85）。

Crofton氏は、どのような種類の\_不均等\_な分布も法則を意味すると主張しました。「平面内の点が任意の一部よりも別の部分でより密に集中する傾向がある場合、そこにはそれらを引きつける法則があるはずです」（同上 p. 84）。同じ見解は彼の『ローカル確率』に関する論文中でも強調されています（Phil. Trans., Vol. 158）。確かに、点が等しく密に集中する場合も、規則性や法則の一例であると言えるでしょう。

適切なランダム性を確保することが重要な実用的な結果に関わる場合、常にそのような配慮が払われ、意図的な計画や不自然さ、無意識の偏りが結果を妨げることはありません。ここでの主な例は、賭博ゲームに他なりません。サイコロを投げる際には、長期的には1から6までのすべての数字が等しく現れることが望まれますが、個々の出方は予測できないようにすることが重要です。問題の定式化において、「ランダムに思いついた数値」や「ランダムに撃った一発」というような条件を設定してもよいかもしれませんが、誰もその公正性が保たれるという前提に基づいて勝敗の機会を賭けることは冒険的です。したがって、サイコロの各面は正確に同じように作られ、このサイコロに対して投げ出される直前のどの段階でも自由な操作を行っても、結果のランダム性を妨げることはありません。

6. II. 科学的な概念が私にとって、一般的な或いは原始的な意味と異なる特徴の一つは以下の通りです。考慮に入れる分布領域は必ず有限または制限されたものでなければなりません。この制限が必要である理由は一見明白ではないかもしれませんが、いくつかの例を考えてみることでその点がどのように感じられるかが明らかになります。

例えば、無限の数列からランダムに数字を選択する場合を考えましょう。一般的な意味での「ランダム」を選ぶとは、言葉を発する前に選ぶ時間を考慮せずに数字を述べることです。これは問題ありません。しかし、一瞬考えてみると、この方法で最終的に均等な分布を確保することは不可能であることがわかります。つまり、平均的な値を常に安定して定義することはできません。同様に空間の無限についても言えます。特定の直線、面積、または体積からランダムに点を選択することは理にかなっていますが、線が無限に続く場合や無限の空間から選択する場合、私たちが「ランダムな点」の均等な分布に向かう最終的な傾向を基盤とすることができなくなります。

他の無限の量の例でも同様です。棒を水平面に投げて鉄製のフェンスの列に対して通過する確率を考えましょう。この問題はチェスの駒と類似しており、棒の平動に関する動きについては特に問題ありません（もし平動から始めると）。しかし、回転に関しては異なります。任意の線形速度に対し、ある一定の角速度以下の場合は棒が接触せずに通過でき、それ以上の場合は通過できません。前者の範囲は有限であり、後者の範囲は無限であるため、均等な分布を想像しようとすると以前に遭遇したような困難に直面します。

もちろん、この特定の難題を回避するために角速度から推定を行うと、線形速度について先ほど述べたことを適宜修正して再現する必要があります。

7. 当然、現在存在する様々な問題が、先ほど述べた内容と矛盾しているように思われるかもしれませんが、これらはすべて説明の対象となります。例えば、三本の直線をランダムに選んだり引いた場合、その長さが三角形を形成できるようなものである確率は何でしょうか？この問題に対する考え方は二つあります。一つには、三本の直線の長さが特定の長さ n よりも短いという前提から始め、n が無限に増加するにつれてその確率がどのような極限に向かうかを決定することができます。もう一つには、この問題は三つの直線の相対的な比率に関するものであると考えることができます。したがって、最も長い直線を表すために任意の大きさから始め、他の二本の直線の長さを変化させることで全ての三角形の形状が表現されることになります。どちらの場合も、無限に可能な長さの中からランダムな選択を行う必要はありませんし、明確な結果を得ることができます。

したがって、「三点問題」と呼ばれる問題では、空間内にランダムに選ばれた三つの点について、その中から鋭角三角形を形成する確率を求めます。この問題は、一球体（特にその単純さから）という閉じた体から始め、その内部で一様分布であると仮定した確率を求め、その後無限に増大するこの球体を考えることによって解決されます。このような視点からは、問題が完全に整合的かつ理解可能なものであり、なぜ空間内でのランダムな選択というよりは球体内での選択と呼ばれるべきかは私には分かりません。もちろん、異なる体から始めれば異なる結果を得ることができますし、したがって（既に述べたように、エドュケーショナル・タイムズ誌のクロフトン氏によると）無限空間はより自然かつ適切に球体の拡大によって近づかれるものと見做すべきであり、立方体や他のどの形状よりもそうであると主張されます。

再び：整数の集合をランダムに選んだとする。その中から取り出した数が奇数である確率は偶数より高いことを示せ。
この問題に対する解答では、任意の有限な数 n から始め、その範囲内で作られるすべての可能な組み合わせについて証明し、それらの中で奇数が偶数よりも多いことを示す。この事実は n の大きさに関係なく成り立つため、選択が真正な数え上げに想定される無限大からランダムに与えられるものとして考えられる。

8. これらの条件を満たせない場合、私には有限の確率値を指定する試みは失敗すると考えられる。例えば、Mr J. M. Wilsonによって提案された次の問題においては、「無限の平面上で3本の直線がランダムに引かれ、4番目の直線がその他の3つの直線で形成される三角形と交わる確率を求めよ」（\_Ed. Times\_, 再版, 第5巻, p. 82）という問題に対して、彼は次のように解答を提示している。「4本の直線中、2本は必ず三角形内にあり、残りの2本は必ず三角形外にある。すべてがランダムに引かれたため、最後に引かれた直線が他の3つの直線で形成される三角形を通る確率は従って 1/2 である。」

この解答を引用するのは、私にとってその難しさを示す例として適切だと考えられるからである。問題の文脈では、3本の直線によって三角形が指定されていると仮定している。どの程度大きいかに関わらず、その三角形の面積は外部に比べて無限大に近い領域に対して有限な比を成していない。そして、この仮定に基づいて何らかの明確な解釈を加える限りにおいても、有限な領域を通る第四のランダムな直線を引く確率はゼロと見なされるべきであると考えられる。Wilsonさんが解決した問題は私にとって全く異なるものであり、「4本の交差する直線が与えられたとき、その中からランダムに選んだ直線が他の3つの直線で形成される三角形を通る確率を求めよ」という問題である。

他の多くの場合でも、この乱雑性の概念を実際の無限に適用しようとする試みにおいて同じ困難が生じるようです。以下の問題はそのような問題の正確な類似例です：ある数がランダムに選ばれるとします。次に別の数がランダムに選ばれた場合、その数が前者よりも大きい確率を求めよ；答えは明らかに1であるべきで、つまり確実性であるべきです。なぜなら、任意の指定された数より上の範囲はそれを下の範囲よりも無限に大きいからです。

あるいは、「無限」という言葉を理解できる唯一の言語で表現すると、私が意味するのは次のようなことです。最初の数が m であり、n（n > m）までの数を選べる場合、m を超える確率は (n - m) : n となります；2n までの数を選べる場合、その確率は (2n - m) : 2n となります。つまり、いかに大きな n や m の値であっても、この表現は常に明瞭です。しかし、最初の数 m を選んだ後、n は m よりどれほど大きいものでも選べる：すなわち、確率は私たちが望むほど近づくことができます。

デモグラの優れた示唆に富んだ論文『無限』（『Camb. Phil. Trans.\_ Vol. 11.）においても似たような誤りがあると考えます。彼は「三点問題」について議論しています：すなわち、3点がランダムに選ばれるとします。それらの点が鋭角三角形を形成する確率を求めよ。彼が示しているのは、一方の辺が与えられた場合、対頂点の可能な位置を考えたときに、鋭角三角形を形成するような位置は無限に多いことを、余弦三角形を形成するような位置よりも多いという事実です。しかし、これもまた異なる問題を解いているだけです。

9. 真の無限のランダム性や、真の無限性から真正なランダム選択を行うための最も近い方法は次のようになります。円周に接線を延長し、その両方向に無限に延ばします。次に、中心点からランダムに半径を引き、つまり、接線側にある半円周が最終的に均一に交差されることを許容します。これらの半径を接線と交差するように延長し、その交差点の分布を考えます。結果として得られるのは、我々のランダムな分布の一つの特性です；すなわち、この接線上のどの部分も、いかに小さな部分であれ、いかに遠くにある部分であれ、最終的には仮想的な継続的な雨の pavement の任意の小さな部分と同じ位置を占めることになります。つまり、これらの要素的な断片は徐々に交差点で点々と密に埋まります。しかし、もう一つの重要な特性、すなわち最終的に均一な分布が欠けていることを認識する必要があります。接線との接触点を中心にした特定の形態の分布があり、これは「誤差の法則」として今後の章で議論されるべきものであり、その結果、集中は接触点に最も多くなる傾向がありますし、そこから各方向に容易に計算できる公式に基づいて薄くなることになります。このような状況が存在すること自体は、真のランダム性という概念に対する理解とは全く対立しています。

10. III. 定義やそれらの結果を除き、ランダム性に関する最も重要な問題は次のようになるかもしれません：特定の場合において、観察された配置がランダムであるかどうかをどのように判断するか？この質問については、今後の章でより詳しく議論しますが、我々はその中での困難さを克服する道筋を見出せる立場にあります。

(1) 考慮する事象や物体が無限に延長されている場合、またはその生成の方法について十分な知識があり、最終的な傾向を特定できる場合、あるいは実際的な数え上げを超えるほど多く存在している場合、大きな困難はありません。単に事実に関する問題に直面しているだけです。雨滴の場合、二つの等しい大きさの舗装面または他の表面を観察し、より一様で均等な点が徐々に濃密になるかどうかを記録します。もしそうであれば、その配置は我々が「ランダム」と呼ぶものです。たとえば、タバコパイプが本当にランダムに折れるかを知りたい場合、十分な数のタバコパイプを落として、長期的にはどの長さの破片も等しく代表されているかどうかを見るだけです。あるいは、私は材料の強度やそのような物体の分子構造についての知識から推論し、小さな破片と大きな破片がすべて同様に折れる可能性が高いかどうかを議論することができます。

11. 読者の注意は、ここでの混乱源である特定の横断分野に対する注意が必要です。現在議論しているのは事実に関する問題であり、具体的にはある最終的な配置の性質についてです。つまり、我々が議論しているのはランダムでないものと比較するためではなく、ランダムなものと非ランダムなものとの対比ではありません。数秒後に明らかになるように、設計された結果であっても、それがランダムな配置を示す特徴的な印を残す可能性があります。

星の配置について、しばしば議論されてきた問題があり、また後ほど再び取り上げることになるので、ここではその場合について考慮します。この問題は、実際に星同士の相互位置については何也不知道という事実によって多少複雑化されます。私たちが認識できるのは、それらの星の視覚的なまたは観測可能な位置だけです。これらの位置を観察することにより明らかになるように、全体としての配置はランダムなものではないことが明白です。ミルキーワイアーや他の解消可能な星雲は、雨に落ちる滴がそれよりも広い範囲に多く降った場合に、その範囲内にいくつかの不規則な斑点があるような状況と同様に明らかに反例となります。裸眼で見える星や僅かな望遠鏡力で観測できる星のみを考えるときでも、これらの配置は大部分において、大体ランダムなものであるということが確実です。これは、視覚的に見ると密度が高い集まりが、その部分がより広範囲にわたって存在しているだけであり、それ自体がより密集していないことを意味します。

星の実際の配置も同様の特性を持つ可能性があります：つまり、視覚上密度が高い集まりは、私たちが観測する領域がより広範で、より密集していないという事実から生じるものです。私たちは目前に広がっている領域がより多く存在しているわけではなく、単に広いだけである可能性があります。実際の選択肢は以下の通りです：星々が配置されている宇宙全体が比較的規則的な形状を持っている場合、その配置はランダムな秩序ではないでしょう；それに対して、その全体が非常に不規則な形状を持っている場合、その内部での配置は全体としてランダムな秩序である可能性があります。

12. (2) 問題の配置が比較的少数の事象や対象を含む場合、それがランダムなものであるかどうかを決定するのは非常に困難になります。実際、私たちは立場を変えなければならず、実際に観察されたことではなく、もっと長い期間にわたって観察を続けることができればどのような結果が得られるかを推測できると結論付ける基準で判断しなければなりません。これが「逆確率」と呼ばれる概念であり、観察された結果から原因の性質を決定する問題は、将来的な章で完全に議論されるでしょう。しかし、いくつかの導入的な解説はここで行うことが便利です。

現在の意味での確率に関する問題は、最終的な限界という概念を導入し、従って無限回の再現可能性を前提としています。私たちは有限数の発生しか見ることができない場合、その配置の性質についての直接的な証拠が失われ、生成する原因の性質に依存せざるを得ません。そして、この数が小さくなるにつれて、原因の性質を評価できる信頼度は徐々に低下します。

中間的なケースから始めましょう。小さな芝生があり、その上には白い舌花が散在しています：これはランダムな配置でしょうか？単なる視察だけで、私たちはそれがそうであると確信を持ちます；つまり、（否定的に）どのような規則的なパターンの痕跡も見つけることができず、（肯定的に）任意の中程度の面積、例えば1平方ヤードを取り、その中に含まれる植物の数が大まかに同じであることを発見できると感じています。しかし、私たちは原因の知られた性質を援用することで自己助勢できます。私たちは舌花は種子によって広がることを知っていますし、風の影響や芝生の継続的な掃除と刈り取りを考えると、カードを配る方法やサイコロを投げる方法に類似した原因が働いていることが判明します。

上記の場合、生産過程への依頼は補助的でしたが、極めて小さな系列やグループについて考えると、この依頼は非常に重要になります。たとえば、「ヘッド」および「テイル」の数が10である特定の系列について教えてもらうとしましょう。ここで範囲は決定に十分ではありませんし、取得した人物が投げているのか計画しているのかを教えてもらわなければ、その結果が「ランダム」と呼ばれるべきかどうかを判断することができません。真実を忘れてはなりません：「計画」は長期的には直接ランダムさの模倣を試みるならば必ず崩壊しますが[4]、短期間では完全にその模倣を simulare することができます。短い系列、たとえば「ヘッド」と「テイル」の系列も、投げているかあるいは意図的に選んでいるかで同様に生成される可能性があります。

13. 読者には注意していただきたいのは、ここではランダム性という問題が最終的な統計的事実として単純に扱われている点です。私は完全に認めますが、これは原始的概念でも一般的な解釈でもありません。しかし、確率の推論を行うためにはこの概念を採用する道しかないと考えられます。最終的な配置の生成元を探ると、その種類は非常に様々になります。それは（数段階前の）意識的な意図によるものであるかもしれません、たとえばカードを引いたりサイコロを投げたりするとき：あるいは多くの自然原因の極めて複雑な相互作用の結果であるかもしれません、たとえば芝生や草地に散らばった花々の配置：あるいは我々が全く何も知らない種類であるかもしれません、星々がお互いに対する配置の場合です。

数学者が数日前に述べたように、ランダムさと設計は互いに交差する結果をもたらすことを念頭においていたのはこの状況だった。前段落で言及した多くの配置においては、設計の影響があったとしても、それらが全く設計の跡が見られない他の配置とその結果を区別することは不可能である。

特にここでの最も印象的な例としては、自然数的定数（例えば π または e）や対数表における数字の配置がある。これらの数字の生成過程を見れば、ランダムさの反対概念を示すもっとも極端な例は見つからないだろう。各桁は必ずしも自由に位置づけられるものではなく、その位置は必然的に決定されている。一瞬の意図の欠如でも、計算者は完全な目的を達成することができない。

しかし、結果だけを見れば、これらの桁の行がランダム配置を示す最良の例であることは間違いない。各桁はほぼ等しく出現し、さらに進むにつれてこの傾向が強まる。また、桁同士の相互接続も同じ傾向を示す：例えば 5 という桁は 6 や 7 に続く頻度と他のどの桁にも等しい。実際、これまで計算されたすべての桁から最初の五桁を取り除き、残りを見ると、これらの桁が十面体のサイコロの結果であると考える理由は全くないだろう。

14. このようにである \_なぜ\_ かという疑問が生じます。物理的な因果関係が関わる場合、我々は通常、そのイベントが以前に指定された前件によって規則的に引き起こされることを指名することで、この質問に含まれる要求を満たすと理解されています。しかし、幾何学や算術においては前件の余地はありません。それらの場合で我々が一般的に探求するのは証明です、つまり、可能な限り公理に分解すること、または少なくともより広範な一般性を持つ承認された真実に分解することです。この特性である統計的ランダム性の存在について証明を与えることは知りませんが、以下の考察は、不確実性を転嫁し、類推の優位性がその存在を支持していることを示唆するかもしれません。

まず、有名な定数 π を考慮してみましょう。これは大量の算術的および幾何学的な関係で現れる量を表します；ここでは、円周と直径の比として最もよく知られたこの関係を検討しましょう。そのように見れば、それは単に任意選択された単位による量の測定の一例に過ぎません。したがって、π の値を数百桁まで決定するためには、我々はより高い倍率で継続的に観察することで、刻度の更なる微細な分割を見つけることができるという仮定を立てなければなりません。棒や線に対して刻度を合わせてみましょう；31インチと32インチの間にあることを見つけます。次に十分な精度で十分な精度で観察するために、刻度の次の分割、つまり10分の1インチの分割を見つけてみましょう。これらの10分の1インチの数が整数の31インチの数と関連している理由があるのでしょうか？測定する「余り」は実際には31インチを測定し終えた後、我々に手渡された新しい部分であり、同様に、更なる微細な分割につながる毎回の「余り」も同じように考えることができます。

他の多くの無理数についても同様の注記が可能である。不可約平方根のようなものも含めて考えることができる。二つの直線を直角に交差させ、その交点からそれぞれ一定のインチ数を描くとしよう；例えば、2インチと5インチを描く。そしてこれらの端点を結び、直角三角形の対角線を作り出す。

この対角線を他のいずれかの直線で測定しようとすると、実質的に平方根を求める作業をしていることになる。πの場合と同じように、無理数であり、桁の順序がランダムであることが通常であり、有理数であることは例外と考えるのが自然だろう。たとえば、二辺が3インチと4インチの場合には、対角線がそれらと有理関係にあることが見つかるかもしれないが、これらは規則的な例外に過ぎず、あるいは無限に少ない例外であり、その他の無数のケースが通常を形成している。

15. このような一列の数字の真のランダム性を説明する最良の方法は、図形的な助けを使用することである。ここでも普通の統計学と同様に、単なる数字の意味を理解するのは容易ではない。しかし、点群や線群の配置ははるかに直感的である。目は、配置に規則性の兆候があるかどうか、または一方の方向よりも他方の方向でより密集している傾向があるかどうかを非常に速く察知することができる。

では、これらの数字をどのように配置すれば、その真の特性を示すことができるだろうか？私は、ランダムな線を描くことを提案したい。そして、私たち自身の力だけに頼るのではなく、数字のテーブルを使用してこの作業を代行させることを提案する。その後、それらが任務をどのように効率的に果たすことができるかを調査することにする。

直線をランダムな形と位置で描く問題は、方向や交差の制限がある場合でも十分に知られているが、形状も位置も純粋にランダムであるような線を描くというアイデアはまだ提案されていない。単純化のために、線は平面内に限定されることとする。

このような線の定義は、特に困難な点を含んでいるように思われない。通常の言語に従って表現すると、それは各瞬間にどの方向へ移動するか他のどの方向と同様の確率で移動する点が描き出す軌跡（つまりどのような軌跡でも）と説明できるだろう。私たちはそのような線を描くことができず、物理的な手段によってもそのような線を描くことは不可能であることは確かだ。どの物体も微かながらも、各瞬間に一直線に移動する傾向があるため、瞬時に変化する方向指示に対して即座に対応することはできない。また、その最終的な理想状態の線を想像したり図示することも難しい。

しかし、そのような線のグラフィカルな近似を与えることは容易であり、この近似が目的とする理想的なものに限りなく近づける方法も容易である。

次のように進めることができる。普通の曲線のグラフ的説明のために準備された格子紙を取り用紙上から始点を選び、その交差点としてとる。そして、これらの線によって示される8つの「方位点」（北、北東、東北、東、南東、南、南西、西）と含む直角の二等分線を考慮する。これらの方向の中からランダムに選択することを提案するために、これらの数値を0から7まで番号付けし、北向きの線は0で、北東向きの線は1で表すとする。各交差点での方向を選択することは、8面ダイスによって委ねることができるが、この説明のために我々は計算されたπの値に現れる数字0から7を選ぶことにする。

このような手順をランダムに繰り返すことにより進むような軌跡は容易に想像できる；そのような軌跡は本章末尾に示されている。

この例を用いる目的は、算術や幾何学の無理数定数の一連の桁をグラフィカルに表示することであり、それについては上記が十分であると思われる。実際にいくつか試してみると、目で見てか、あるいは今後述べられる理論原則によれば、その乱雑さの説明に充分応えるように思われた。

16. しかし、この話題についてさらに進む価値があるように思われる。つまり、方向の理想的な乱雑さにどれほど近づけることができるかを調べることである。これを完全に行うためには二つの改善が必要である。まず一つは、八方向に限定せず無数の方向を許容することである。これには大きな困難は生じない；少数の桁を使用する代わりに、どの方向でも自由に停止できる円盤のようなものを用いることで十分である。しかし次に、短い有限なステップではなく、それらが無限に短くなることを仮定しなければならない。ここでは実際的な非達成性が感じられる。ニュートンが不連続の多角形から連続曲線へと移行する装置を用いることは我々には馴染みがあるが、これは理想的な曲線自体が容易に描ける（そして私は言うところだが、容易に想像できる）ためである。しかし現在の場合ではそうではない。問題となっている線は最後まで不連続、あるいは角状態で残る：その角は尖ったままになり続け、構成要素の数が増え、大きさが無限大に向かって減少するとしても、尖りは失われることはない。このような理想的なものは想像することすら難しい。これは顕微鏡下に粗い物体を見ているようなものであり、その角度が高倍率で見ても丸みを帯びる兆候がないように見えるからである。私たちの「乱雑な線」は、当然のように尖ったまま残り続けなければならないが、その尖った部分の大きさは無限大に向かって減少する。

したがって、この状況は次のようである。言葉で説明することは容易であり、ある瞬間において線がどの方向に進むかは等確率であるという概念を述べることができる。また、任意の細かい程度までそのような線を描くことも可能である。しかし、最終的な形を思い浮かべることは不可能である。実際には、「極限」が存在せず、理解できるものや想像できるもの（曲線の漸近線や継続的な曲線に対応する無尽蔵に発展する多角形）へと我々は常に近づいているというようなものはない。したがって、その最終的な到達を思い描くことは容易ではない。

17. 上記のように、このような線の近似段階を考慮する場合、私は通常の確率論のいくつかの定理が安全に適用できると考える。たとえば、そのような線が最終的に始点から無限遠方に逸れるかどうかについて尋ねられた場合、ベルヌーイの「大数法則」を援用して、その発散が相対的に大きなものである確率は非常に低いという結論を得ることができる。グラフィカルな例に戻って考えると、始点を通る縦線から右か左に大きくずれる点の結果的な偏移について最初に考える。各段階で8つの可能な動きの中で2つはこの相対的位置に影響を与えないが、残りの6つは等確率で1ステップ右または左へと進む。したがって、結果的な「流れ」は、大量のコイン投げ後の表と裏の数の差に類似している。多くのコイン投げの結果として知られているのは、最終的には事前確率（つまり表と裏の等しい数）への比例的な近似がよりよく達成されることであり、絶対的な偏移はより広範囲に現れるということである。

特定のケースに対して適用し、水平方向と垂直方向についても同様の結果が得られることを念頭において、非常に多くのこのような「ステップ」を経過した後では、私たちは私たちの出発点からの距離に対する移動全体の距離との比はおそらく小さくなるだろうと言えるでしょう。一方で実際の距離は大きくなるでしょう。また、このような線をさらに延ばすにつれてこれらの傾向がより顕著になることも言えます。このテストに関して、そして描かれた線の全体的な外観から得られる検証に関して——最後のものは上記のように比較的信頼性が高いと述べた通り——不連続なまたは無理数比を示す行の配置が一般に「ランダム」であるという点については、私は疑問を持ちません。

このようなパスの一例を見せてみたいと思われるので、πの値における0から7までの8桁の配置を示したものを付記します。これらのデータは、Shanks氏によるこの定数の計算において驚くべき業績（Proc. of R. S., XXI. p. 319）から取り出しています。8と9を取り除いた後では568が残ります；図はこれらの数字を指標として私たちのパスの進行方向に従って描かれた経路を示しています。当然、多くのステップは反対方向で2回以上取られています。この結果は私には非常に公平なランダム性のグラフィカルな表示であると感じられます。私は同様のパスが logarithmic tables から得られた数字行によって提供されたものと比較し、他の方法でもその結果を確認しました。そして、これらは大体同じです。

[図：πの桁数によるパス]

1. Prof. Skeat（語源辞典）によれば、最も古く知られている意味は、 cavalry の突撃のような狂暴な行動にあります。彼は、この語源がTeutonic語 rand（縁、溝）と関連していると考えており、満ち足りた河の狂暴で不規則な動きを暗示すると述べています。

2. 1854年1月18日のCambridge Mathematical Triposの問題論文を参照してください。

3. Mr. H. Godfrayによると、問題が一度試験で提示された際に、大多数の受験者はその問題を解くために大抵仮定していました。詳細は『Educational Times』（再版、第7巻、p.99）をご覧ください。

4. 参照してください、p.68。

5. たしかに、方向の選択肢として10つを採用し、従ってすべての桁を省略せずに取り扱う方が完璧でしょうが、実践的には8つに限定する方がはるかに手間がかかります。

6. もっとも中心部から等角螺旋の形状を想像するよりもです。

\newpage

\chapter{信念の度合いの測定（Measurement of Belief）}
\label{ch:vi}

\begin{center}
\large\textit{主観的信念の定量化批判・サンクトペテルブルクのパラドックス・極限としての客観的確率}
\end{center}
\vspace{4mm}

第6章。[1]

確率の主観的な側面。信念の測定。

1. この章は、デ・モランを追うような過度に主観的な見解に傾倒する風潮に対して、私自身で少し反発の精神を持って書かれたものでした。現在再読しても、特定の一文に対し明確な異議を唱えることができませんが、もし全く新しい形で書き直すなら、表現はより慎重に行おうとするでしょう。その点での修正を行いました。主観的な見解と私の考えに著しく異なるものの、逆方向の強調で述べられている視点を見たいのであれば、エドワース博士F.Y.の「運命の哲学」（Mind, 第9巻）を参照してください。

1. この章では、ある種の系列について明確な概念を得た後、その系列は何に役立てばよいのかを探ります。つまり、どのようにして推論を行う手段として利用すればよいかです。現在我々がとろうとする一般的な手順は、客観的な視点から主観的な視点へ移行するものであり、物自体から私たちの心でそれを捉える状態へ移行するものです。

読者には注意していただきたいのは、多くの場合において、系列を計算に適した形にするために最初の一歩として置換が行われているという事実です。前章で説明したように、この置換は一定程度の「理想化」プロセスであると言えます。実際に遭遇する系列は変動的な型を示す傾向があり、その個々の要素は時折許された不規則性を逸脱します。そのため、これらは正確に論理的に扱えるまで少し整える必要があります。自然な物体はほとんど例外なく、正確に論理的に扱うために整える必要があるように、系列が現れる形は固定型を持つ系列の形となります。この仮想的なまたは理想的な系列が我々の計算の基礎となっています。

2. このことは、確率が実在の事物について推論する学問であるという以前の主張とは全く矛盾しないものではない。このような置換によってのみ、我々は実際の事物について推論を行うことができるようになる。

自然界においては、ほとんど全ての現象が科学的な目的で求められる厳密な形から逸脱した形で現れるため、我々はそのような欠点を免れた仮想的な系列を導入する必要がある。この系列が満たすべき唯一の条件は、置換がいかなる程度でも任意的であるべきではなく、つまり真実からどれほど微細な差異であるか、ということができるだけ少ないことである。このような種類の置換は、どのような種類の自然現象も厳密な科学の対象とする際に通常無視される。

私はここでこの主題に対する不慣れが読者に誤解を与えることを懸念し、そのような推論過程における形式的な正確性への特別な逸脱を含むと誤解する可能性があるため、これを明確に注意を向けます。

また、この仮想的な系列の導入は、前章で批判された学説に対して何らの支持を与えるものではないことに留意すべきです。その学説では、我々の系列が固定され変更されない型であるとされ、ラプラスの表現を用いて「現象の確率の開発」と呼ばれていたと考えられていたのです。

この仮想的な系列は、その仮定に基づく何らかのものとは異なる点として、我々自身の実際の系列に対する必要不可欠な置換であることを認識され、可能な限り事実に近い形で保たれるべきであるという点があります。これは計算のために必要な虚構や手口であり、その目的のみのために用いられるものです。

この注意は特に必要である。選択する例は、この分野において最も好まれる例の一つであり、その置換が偶然的に不要になる場合があるからである。過去に何度も指摘された通り、物事はしばしば修飾の必要がない場合があり、実際の形で提示される状態ではすでに理想化されているからである。多くの場合、数列を適切な形にするために多くの変更が必要だが、一部の場合、特に賭博ゲームの場合においては、実用的な観点からは変更が不要であることが見られる。

3. そこで我々はこのような数列から出発し、そのような数列についてどのような推論ができるかを調査する。論理的な読者に助けるために、最初の手順は一般的に「即時」推論と呼ばれる一類のものに類似することを知らせることである。すなわち、「すべての人間は死ぬ、従って特定の一人または複数の人間は死ぬ」というタイプの推論である。この例は論理において単純かつ明白ではあるが、確率においては非常に慎重に考慮する必要がある。

明らかなことは、このような推論を取る手立てが適切であるかどうかについての見解を形成するために準備しておくべき点である。これまで我々は不規則な個体とほとんど関わりを持たないようにしてきた；それらは単に規則的な系列の一端として扱われただけである。しかし、これらを無視する考え方は長く持続することはできない。これらの出来事が総じて十分確実であるとしても、我々はそれらについての対処において、その総体だけでなく個々の状況や少数ずつでしか対処することができない。保険会社における例を挙げると、彼らは大きな数値に到達することで大部分の不確実性を排除できるが、それぞれの取引には別の関心を持つ者が関わっている。被保険者はその手続きについてどのように言及すべきか？彼にとってはこれは個々の問題となるからである。また、保険会社自体も個別に案件を受け入れるため、これらの個別なケースについて可能な限り明確な見解を持ちたいと考えている。したがって、前章で述べた確率論を扱う対象に関する議論は、このような調査を排除するもののように思われることがあるかもしれない。個々の事実に対する知識がどれほど欠如しているかを前提とし、（後述するように）一定の範囲内では因果関係を否定しても結論に影響を与えないというような状況も考えられるからである。この問いに対する答えは、現在我々がまだ考慮していない全く異なる観点へと私たちを導く必要があることを示すものとなるだろう。その最良の入門は、特定の例についての議論を通じて行われることになる。

4. 一枚のコインを非常に多く投げ上げるとき、長期間にわたって表と裏がほぼ等頻度で現れるという事実（他にも多くの事実）について確信を持つことができるでしょう。しかし、投げ上げた回数を中程度まで減らすか、さらに少ない回数に限定する場合、さらには3回や2回、あるいは甚至１回にまで減少させる場合を考えましょう。極端な例では、確信またはそれに近い状態と、絶対的な不確定性があります。これらの極端な状態の間には、信念のすべての段階があるのでしょうか？

多くの著者の中には、この主題について最も優れた権威者の一人も含まれるが、彼らは我々がこのような信念量の変化を明確に意識していると述べたり、推測したりしています。そして、その心持ちは外部の事象に対する測定とほぼ同じ精度で測定および決定することが可能であると主張しています。この見解の主要な数学者的な支持者はデ・モランであり、彼は主題に関するすべての著作の中で強力にこの見解を主張してきました。彼が現在議論している視点について最も明確な説明は『形式論理学』において見つけることができます。彼はこの視点に基づいてその体系を構築しています。彼は、我々が提出される任意の命題に対して一定量の信念を持っていると主張し、その量は性質上決定可能であるものの、特定の場合では実際にはそれを決定することが困難であると述べています。彼は確率を形式論理学の一種と見なし、次のように述べています。「前提に対する部分的な信頼が結論に及ぼす影響の研究から、前提を絶対的に真だと仮定した場合の結果の研究を分ける理由が理解できません。」つまり、ある科学がある形式論理学があり、その規則は一つの命題から別の命題を必然的に導くためのルールを調査します。それに密接に関連して、ある科学は一つの命題に対する信念量が他の関連する命題に対する信念量にどのように変化するかを調査します。

同じ見解は、すでに亡きドンキン教授も支持しており、彼は（Phil. Mag., May 1851）で述べています。「計算の対象となる量である\_信頼度\_という点については、私は一般に認められるだろうし、しばしば明示的または暗黙的に述べられてきたと理解します。」

5. この見解を批判する前に、一度注意すべき点があります。この視点が実際には誤りでないとしても、定義の目的から不十分であるという反論があることを忘れてはなりません。信頼度の変化は確率理論に関わる唯一の特性ではないのです。その科学は確かに信頼度の変化に関心を持っていますが、他の分野でも同様の性質に直面します。

推論を拡張する際や類推を使用する際、他者の証言に依存する際、過去の記憶に頼る際、相反する主張を通じて結論を導く際、あるいは数学や論理学に基づいて長い複雑な推論を行う際など、我々は前提から得られた結果に対してそれほど確信を感じることができないという経験があります。これらのすべての場合で、信頼度の量が変化していることを認識しますが、その信頼度がどのように生成され、どのように変化する法則は同じでしょうか？もしこれらを一つの調和的な体系に統一できない場合、実際にはそれぞれ異なる体系に分けられ、各々に独自の法則とルールを持つのであれば、それらを一つの科学として強制的に統合することは徒労であると言えるでしょう。

この見解は、疑問の定義を採用する著者の多くが、上記の信念の減少の例の大半を考慮から除外し、第一章で述べられた性質を持つ事象のクラスにその注意を限定していることによって強化される。

これらの著者は、実際には、ほとんどそのような信念の減少の例を考慮に入れないかのように行動しており、そのような事象は第一章で述べられた性質を持つものである、すなわち「少数者の無知と多数者の知識」である。確かに、これらの一部の例を適切な形に変形するためには、非常に恣意的な仮定を導入しなければならないことが明白であり、それらを強制的にそのような形にするために相当の力を費やす必要がある。しかしそれでも、言葉の使用を精密に検討すれば、ほとんどの場合において、個々に対する我々の知識が第一章で説明された典型的な形式から与えられた前提に基づいて仮定されていることが明らかになるだろう。

この主張は、科学の誤用について考慮する際にもより詳しく証明されるであろう。

6. それでも、上記の見解が正しければ、定義のために十分ではないと考える。しかし、それが正しいかどうかは非常に疑問である。デモグラや他の者たちによって与えられた信念側における重要性を適切に割り当てる前に、あるいは科学をその側から定義する前に、両方の立場に対して強力な反論が可能である二つの主張を確立することが必要である。

（1）我々の各命題に対する信念は、厳密には測定できるものであり、どの場合でも一定量存在し、意識の中で実感でき、標準に参照してその価値を判断することができるものである。

（2）このように把握された価値が理論によれば正しい価値であり、それは完全な信頼の正確な割合であるとされるべきである。この主張は最初に少し曖昧に見えるかもしれませんが、今後説明するであろう。

7. (I.) まず第一に、私たちの信頼度を測定する困難性についてですが、この困難性の一因は明らかに注目されるべきものです。これは、強い感情や情熱によって引き起こされる影響が量的な信頼度に与える干渉です。争点となる事項に対する深い関心がある場合、それが希望を抱くか恐怖を感じるかに関わらず、信頼度計は大きな混乱を招きます。したがって、証拠を検討する際には、心は全く情熱的でないものと仮定しなければなりません。

ラプラスや他の著者はこの点に注意し、認めていますが、これらの作家たちは私にとって、誤差の唯一の原因であると考えたり、またそれほど重要ではないと考えているように見えます。すらその唯一の原因であっても、それが重要でないとは見なせないでしょう。我々は希望や恐怖を感じる機会が非常に多く、このような影響を考慮に入れずに済ませることは、実際には私たちが信頼度全体を考慮に入れるという主張からほとんど等価であると言えるでしょう。非常に強い感情は当然例外的ですが、それでも我々はほぼすべての場面で情熱的な要素が現れることを見出すはずです。

特定の出来事に対する驚きや期待について話すことは、あまりにも稀ではありません。これらの表現、特に後者の場合、信頼というよりも何かより深いものが示唆されるように感じられます。確率論の教科書では、「期待」という言葉は一般的に「信頼」に等しいと定義されていますが、流行語として使用される際にはそのような一致性があるかどうか疑問です。

しかし、それがどうであれ、情熱的な要素はほぼすべての場面で存在し、その干渉作用は常に働いているのです。

8. 以前の要因と協力するもう一つの要因は、我々が任意の命題を信じる根拠となる極めて複雑かつ多様な証拠の中に見出される。

したがって、我々が特定の瞬間における実際の信念は最も一時的で変動性が高いものであり、その測定すら困難であることが結果として生じる。これは希望や恐怖によって心が騒ぐときだけではなく、最も落ち着いた状態でも同様に当てはまる。我々がどの程度この信条を確固たるものに持っているかという質問に対する正確な答えを得ることは容易ではない。

支持する主な論点が一つまたは二つあり、それに反対する同等の反論もまた存在するかもしれませんが、それらは信念状態が形成される全ての要因を網羅しているわけではありません。これらの理由は実際的な議論において唯一の要素であるため、我々がそのような理由のみによって確信が影響を受けていると結論付けるべきではありません。

むしろ、我々の確信は無秩序な基盤上に築かれており、無数の種類の推論や類推から成り立っています。さらに、これらの推論は当時の感情状態によって歪められ、その後の記憶の程度により薄れ、また意識の中で偶然組み合わさる方法によって時々異なる力で受容されます。

アブラハム・トッカーからの印象的な例を借りれば、我々の確信の基礎は普通の建物の堅固な地盤に比喩されるよりも、ローデムントの家屋の柱に比喩されるべきです。これらの柱は深く軟らかい泥の中にある何らかの方法で支えられています。その重力をしっかりと支えることは可能ですが、異なる部分がどのように依存しているか正確に指摘することは容易ではありません。

我々が信念量を考慮するたびに、その生成に至る論拠を思い出す必要があります。つまり、これらの論拠は選択に関係なく自ずと頭の中に押し寄せます。各々の論拠が次々と意識を通過するにつれて、信念の強度が変化します。これは混乱した人混みの中での音の一つを選ぶようなもので、ある特定の理由が十分な根拠となるかもしれませんが、詳細に検討するとその力は異なる時間では同一ではないことが明らかになります。

強い論拠によって生成された信念は瞬間的に決定的なものであるかもしれませんが、その論拠が実際に意識から消え去ったときには減少します。これは強光に眩まされるようなものです；印象はまだ残りますが、すぐに薄れ始めるのです。我々が確信の源をいかに限定しようと試みても、この状況は変わりません。

9. (II.) しかし、この変動する状態について平均的な値を算出することは可能でしょうか。その平均が理論によって指定された量と一致することを見つけることができるでしょうか？言い換えれば、長期間の実行において異なる2つの事象が起こる確率が頻度に直接比例していると考えるのは自然なことでしょうか？100枚のlottery票があり、その中には10枚が賞金となる場合、一人の人が賞金を得ることを公平に一割の確信で予想するべきでしょうか？lotteryを参照すること自体は、この主張を否定するのに十分な証拠です。lotteryは常に栄華を誇り、多くの人々がそれを支持し続けたにもかかわらず、長期間にわたって全員が損をするという最も完璧な確信を持って投資した人さえも存在します。迷信や予兆、夢などから行動する人々に対する補正は当然必要です。しかし、これらの要素を除き、数学が彼に与えるすべての助けを受け入れたとしても、個々の事象についての自然な印象が理論によって指定されるべきものと常に一致することを信じるのは難しいでしょう。本当に多くの人が、単一のlottery票を購入する望みがないと言えるでしょうか？おそらく彼らは自分自身に言っていることとは、支払った金額は無視できるほどで、大きな利益を得るチャンスがあるという事です。つまり、理論によって指定されるように確信を配分していないのです。

この見解を裏付けるのは、個々の事例でこのような行動をとる人々が、しばしばそのような行動を他の個別の事例でも行わないことを考えないという点です。言い換えれば、ここでの自然な傾向は、適切に低いべき信頼性を過大評価することであると言え、すなわち潜在的な利得に対するリスクを割り引くことである。

論議を受けた場合、彼らはその感情状態を軽蔑する言葉で表現し、それが不思議な魅力だと呼んだり、そのようなものだと呼ぶかもしれません。しかし、その存在については疑う余地がほとんどありません。ここで述べているのは、私たちの心の自然な傾向であり、教育や思考によって最終的に制御される可能性があるものの性質です。

ただし、教育を受けた人々が賭博ゲームにおいてこの傾向を大半で抑制したとしても、冒険的な投資精神はまだ商取引から完全には去っていないと言えます。検討すると、この傾向はすべての時代、階層、性格の人々に広く見られることがわかり、これを考慮に入れずに、私たちが仮定する本能を一般的に受け入れられている確率論の理論とより密接に整合させるためにそれを無視することは容認されません。

10. この問題には、しばしば考慮されてこなかった別の側面がありますが、その重要性を認めることは適切です。私たちは信頼の本能を持っていると仮定しますが、なぜそれがその後の経験によって正当化されるだけの強さである必要があるのでしょうか？本能は私たちに善のためのものとして植え付けられ、直感的に即座に作用するように設計されています。しかし、それは制御の対象となり、私たちはそれが真実と正義に沿ったものになるよう調整しなければなりません。

心理学の他の分野では、自然からの直感的な指示がそのまま残されるべきではないし、平均的には個々の変動を除いて、最終的に考え抜き、その規則を定めるときによく働くようなレベルに設定されているとは限らないと仮定しないのです。たとえば、怒りに対する反応という例で考えてみましょう。ここでは本能的な傾向があり、全体的には良い結果をもたらすものです。しかし、倫理学者はほぼ一致して、自分自身の制御努力の大半がその怒りを抑圧し、正しい方向に保つことに向けられるべきであると主張しています。怒りは一種の粗い保護として与えられたものであり、社会の中で意図的に意識的に行動するためには過度に設定されていると見なされます。

このようなことでも同じではないかと考えられます。私たちの信頼も同様に本能的なものであり、それが適切な方向性を持つよう制御されるべきであるという点で類似性があります。私はここでこの点について一言だけ触れていますが、本書で採用した確率論の理論によれば、科学に対する重要な影響はないと考えられます。しかし、科学を主観的な側面から学ぶことや、この側面に大きな注目を払うことが適切ではないという強い理由があるように思われます。

そのような理論に完全には従っていない人々が存在することは、おそらく誰の目にも明らかであり、しかしこれについては単なる例外や不規則性とみなされてきたでしょう。平均的に、そして大半の場合において彼らはそう信じているという前提が作られました。先ほど述べたように、その傾向が広く一般的であるかどうか非常に疑問で、それが通常の例外よりもむしろ規則と呼ぶべきものである可能性が高いかもしれません。そしてこれがそうであることが望ましいかもしれません：多くの良い結果は、大きな利益を追求するための諦めない態度を生み出す可能性があるからです。そのチャンスを過大評価して時折試行錯誤を続ける人間が、困難や失望の中でも重大な損傷なく続けられる場合があります。冷静で計算的な傍観者が明白に彼の「信頼度」が適切よりも高いと判断するときでも、彼は深刻な害には至らないでしょう。

また、大きな悪事の確率を過小評価する傾向もしばしば見られますが、これも良い結果を生むことがあります。多くの人々にとって死は几乎無限の大災害であり、少なくとも彼ら自身としてはそう考えるかもしれません。もし彼らがその可能性を適切な価値で常に意識していたら、実用的な生活を通じてどのようにして進めることが可能でしょうか？彼らは家に留まることになるでしょう。なぜなら外出すると殺されるか、狂犬に噛まれるかもしれないからです。

このようなことを述べているわけではありませんが、本能への回帰を推奨しているわけではないのです。実際に我々が批判的で意識的な状態に達した後では、そうすることはほとんど不可能かもしれません。しかし注意すべきは、その修正によって得られる利益は限られていることであるという点です。最も重要なのは、それが後に証明されるべき本能であり、経験の完全な範囲で正当化されていない場合でも、正確な理論を基にすることの不適切性が示唆されることです。

11. 人々が実際にはこのような正確な方法で信頼を配分していないとしても、彼らはそうすべきであると主張することができます。この指摘の意味は今後検討されるでしょう；ここでは、それが現在争点となっているすべてのことを認めることだけが必要です。なぜなら、それは私たちの信念の程度が変更可能であり、必要であれば変更されるべきであるという点を認めているからです。しかし、その信念をどのように変更すべきか？明らかに経験に基づいて変更されるべきです；それが自身だけで信頼できるものではないので、その程度は関連する事象の相対的な頻度によって決定される分数で評価されるべきです。経験が基準を提供しているなら、理論をその経験から始めることが最も合理的であることは明らかでしょう。

もし私たちがこれを行わなければ、注意すべきは、確率を私たしの外部にある事象の研究から完全に切り離して、実質的に心理学の一部にすることになることです。経験によって制約されることなく、信念が自然または本能的に組み立てられ、分配される法則に焦点を当てて私たちの注意を限定するのであれば、その後経験を引用して説明する権利はないとするのです。つまり、私たしがその正確性を保証していない限りです。例えば、ある単一の事象に対する信念が正しくても、それらの複合体（本能的な信念の法則に基づいて形成された場合）に対する信念は必ずしも正しくないかもしれません。そして、それを基に行動すると実践的な誤りを招く可能性があります。もし両者が一致しているとしても、その一致性を証明する必要があり、これは私が思惑通りではないと考えていますが、既に選択した開始点に戻ることを導きます。

12. デモグラントは、上記で批判された教義が形式論理の類推から正当化されているように見受けられる。必然的推論の法則を外部事実（例示的な観点を除いて）からの参照なしに研究できるとするならば、蓋然的推論の法則も同様ではないか？

しかし、このような類推にはあまり力がありません。形式論理は、当然のことながら現代的またはカント派の扱い方においては、思考の法則と物事の法則を区別して、これらの思考の法則が特定の対象に関連付けられていない状況で確認され研究されるという前提に基づいています。現在まで、必要的あるいは不可逆的な法則に限定されている場合、すなわち普通の形式論理の場合には、この前提は特に困難を引き起こしません。

つまり、これら主観的かつ客観的な必然性間で衝突が生じるわけではありません。一方が他方の正確な反映であるという完全な調和の中で両者は共存します。その対応関係は非常に精密であり、多くの人々は、形而上学的研究によってこのような点に注意を向けられるまで、これら二つの側面があることに気づかないでしょう。どちらかを選ぶ際には、等しく信頼して「それは当然のことである」と言ったり、「それが異なる形で同じことを言うのだから、それを思い浮かべることができない」と言ったりします。

実際には、カントの時代以降にこの精神的分析が相当程度認識され受け入れられるようになったのです。そして今でも、主観経験主義哲学の主流は、ここに二つの真正な側面があるとは認めません。一方では、主観的な必然性は固有の関連性によって形成された客観的の一貫性の結果であると主張し、もう一方では、このように呼ばれる必要性（例えば矛盾法則）は最終的には単なる文言上のものであり、同じことを別の言葉で言い換えただけだと主張します。どのような説明を採用しても、一般的な結果は以下の通りです：純粋論理の領域内で実際の思考行為としての誤謬は不可能であり、その領域での誤りは意識の一時的な逸脱（つまり注意の逸脱）によってのみ可能であるとされます。

13. しかし、主観的な均一性や法則と客観的な均一性や法則との完全な調和は、純粋論理の領域内で存在するとしても、確率の領域内では遠く及ばない。我々が信念の量を研究対象の一部として統合的に扱うと、そのような不変の相関関係は失われてしまう。前者の場合においては、誤った思考を意識的に行おうとしても不可能であり、後者の場合においては、誤った信念を持つだけでなく、常にそうしている。我々の信念の量が、それが参照する事実に正確に適合し、想像力でさえ迷い出せないほど調整されているという考えから遠ざかっていることは明らかである。むしろ、その後の判断によって承認されるべきものよりもしばしば過剰または不足していることが見られる。

我々の信頼の本能は頻繁に経験と対立しており、そのような状況では何をすべきか？我々は単にその本能を事実に合わせて修正するだけである。我々は常にこの練習を行っており、いかなる教養のある心でもこれ以外のことをすることは不可能だと感じるだろう。誰もが自分の信念をその対象から切り離すことは考えないし、かつてはそれ自身が後者に従う以上の何かをすべきだったとは思わない。

したがって、純粋論理の場合には可能であると思われる正当な防衛があるように主観的な必要性と客観性の分離、そして前者を独立した科学として扱う決定は、確率の場合には全く不適切であることが結果的に明らかになる。我々が事実に依拠せずに正しく思考する方法を見つけようと如何に関わらず、我々は常にその手順を無限にチェックし続けることで誤った信念を持つことを避けることができる。

従って、形式論理が別々の科学としての地位を主張するいかなる権利も認められるとしても、現在検討されている確率理論に対してそれを支持することは見受けられない。

14. 問題の点はしばしば次のように主張されます。ある男が二つの、そしてそれ以外の何もない選択肢に直面しているとします。そのうちの一つは彼にとって成功をもたらし、もう一つは失敗をもたらすことが確定しています。彼はそれらについてさらに何も知らない状況で行動を強いられます。このような場合、彼はこれらに対して完全に同じかつ等しい信頼感を持つことになると主張しませんか？実際の系列や仮想的な系列を参照することなく。もしそうだとすると、これは一方がもう一方と同様に確信を持っていることを意味しないでしょうか？そして、それゆえに各々に対する彼の信念は完全な自信の半分であると言えるのでしょうか？

同じように、選択肢が二つ以上ある場合も同様です：例えば、何らかの数の選択肢があり、その中でどの選択肢を選んでも結果に影響を与えない特異な点以外は区別できないとします。このような状況では、各々に対する信頼感が等しいことになると主張しませんか？そして、ここでも我々は絶対的な信念の評価を分数形式で行わないでしょうか？

このようにして純粋な確率の科学の基礎が築こうとする試みがあります。これは、論理や数学の客観的データに基づいて仮定的に信頼度の分布と組合せを決定することです：すなわち、もし事態が完全に同じである場合、我々の信念はどのように配分されるかが決定されます。当然のことながら、事態が等しいかどうかという質問は論理や数学の客観的データに基づいて決定されます。

この問題を完全に論じるには、対象的または実質的な論理の視点と、形式的または概念主義的視点との間の理由について長大な説明が必要となる。しかし、私はその主題については別の章で述べることになるので、ここでは詳細に入らないことにする。ただし、確率に関しては、一つの決定的な反対意見をすぐに提示できる。

先ほど示したような調査路線を追うことは、既に述べたように、厳格な論理的土壌から離れて心理学的な土壌に移ることであり、これらのすべての場合において、適切な問いは「人々が実際に信じているのは何であるか」ではなく、「彼らが正当に信じるべきは何であるか」である。

先ほど述べたとおり、形式的論理の場合は、意識的に誤った思考を「考えることができない」という事実から、両方の調査、すなわちそれらに対応するものについては、実質的に一つになることが認められる。しかし確率の場合には、この二つは完璧に区別され、独立している。人々が実際に信じていることや将来信じるであろうことを述べることは無意味であり、彼らが正当に信じるべきことは何かを知ることが求められている。そしてこれは、現象自体への訴えなしには決して決定できないのである。

15. しかし、上記の考慮事項以外にも、この方法は私にとって完全には説得的とは思えません。以下に例を挙げましょう。ある男がワッシュ砂地やモアカムベイで、濃霧の中、高潮が押し寄せてくる状況に置かれているとします；彼はそこで潮に引き込まれることの危険性を認識しています。遠くから教会の鐘の音を聞きますが、その鐘が彼自身と同じ側面か反対側にあるかどうかを知る手段はありません。したがって、その音を追うことで、彼は中流へと誘導されて溺れてしまうか、あるいは安全な陸地に戻られるか、どちらかの運命に直面している可能性があることを知ることができません。この状況では再現性がないため、彼にとって影響する唯一の状況が同じように見えて、その結果の死の確率は必ずしも半分であるとは限らないでしょう。このような状況における心の状態を分析することは心理学的な調査よりも論理的な調査に近いものであり、結局のところ、この問題の決定は私たちの論理的位置に関わるわけではありません。しかし、私ができる限りの内省から言えば、このような状況において本当に頭の中で回っているのは、おそらく次のようなものです：多くの曖昧な状況や立場では、我々はそれぞれの選択肢の相対的な利点と欠点を考慮して行動を決定します。つまり、一つまたは他の代替案が成功した頻度を観察したり推測したりして、その頻度で行動を決定します。これらの頻度がより均衡に近づくほど、個々のケースでの誤りの可能性は高まります；すなわち、我々はどちらか一方が正しいかどうかという点で「単なる投げ入れ」のように感じることが多くなります。現在検討している状況はその極限であり、両方の選択肢間には差異がほとんどないよう設計されており、そのためどの一方を選ぶべきかを決定する自信はありません。しかし、決定しなければならない場合、過去に類似した状況で得た precedent に基づいて行動します。動かないままでより良い情報を待つことは確実な死につながりますので、我々は知っている唯一のルールを適用し、またはそのような選択肢がほぼ均衡している状況での以前の経験に依拠しようと試みます。しかし、私は自分自身においても他の人についても、このような再現性がない状況で、過去に類似した状況で正当化され、実行されたルールを適用することによって、我々の心の状態が各々同じであることを完全には説明することができません。

16. この例は、既に述べられ、また再強調する必要がある重要な点を非常に明確に示す。つまり、確率論が扱うものは事件の統計的頻度であるか、あるいはこれらのいずれか一つのイベントについて個別に考慮すべき信頼量であり、その後の行動はその頻度や信頼性に依存し、最終的には参加者の選択によって決まる。たとえば、二名の旅行者が問題にある状況に置かれている場合を考える：一緒に行動するか、または反対方向に分かれるべきか？どちらの場合も、それぞれの安全確率は同じであり、半数であるが、明らかにその状況はどの方法がより好ましいかを決定する。夫婦であればおそらく一緒に行動することを選択し、重要な国家機密を保持している唯一の人物であれば、分かれることを選ぶかもしれない。

つまり、ここでは単一の損失の確実性と二つの損失の均等な確率という二つの選択肢間から決定しなければならない。これらの選択肢はしばしば、その確率を数学的に表現することで、互いに区別がつかないように見せられがちである。しかし明らかに、この決定は参加者の欲望、感情、そして良心に基づいてなされるべきである。確率論はこの問題には一言も述べられない。

私は他方で指摘したように、この点での混乱は、賭博における科学の応用やペトロッス問題の議論において多く見られる傾向がある。

我々は、その支持者の優れた地位から見れば、慎重に且つ詳細にこの論理を検討した。しかし、そのような検討が退屈に思えるかもしれませんが、それを無根拠な理由で否定することは傲慢であると考えた。反対意見は概ね以下の通りである：ある命題に対する我々の信度（確信度）は、それが正確に定量化できるものと仮定する場合でも（これはそうではないように思われるが）、多くの要因によって左右され、そのうちの一つが統計的頻度（確率の対象）である。すなわち、この唯一の要因だけを考慮しても、我々の自然な信度は理論が示すものとは必ずしも一致せず、経験に基づいて調整されるべきである。したがって、確率の主観的な側面は非常に興味深いものであり、検討に値するものではあるが、客観的な側面の付属品に過ぎず、推論のための科学の安全な基礎にはならない。

17. そこで、確率の科学を信念の法則の科学と定義することは、どの点でも崩れてしまうように思われる。しかし、このような単なる否定的な批判に満足するべきではない。我々が抱く命題に対する信度は正確に測定するのが難しいだけでなく、得られたとしてもしばしば誤りがあり、従って信念の対象に基づいて確認されるべきである。それでも、一般的な見解では、ある程度の精度で信念の強さを階層的に評価できるように思われる。我々が今検討すべきは、これが可能かどうかであり、それが可能であればその理由は何であるかである。

一般的には、このような尺度を作成することは容易であると信じられていることは疑いの余地がない。日常会話で頻繁に使用される表現の整備された用語集があるからだ。例えば、「私はほぼ確信している」「私は完全に確信していない」「私はこれよりもそれに対してより自信を持っている」などである。これらのいずれかのフレーズを使用する際、私たちはその表現によって明確な意味を伝えていることをほとんど疑わない。また、特定の状況下においても、どの表現を優先的に使用すべきか迷うことはあまりない。

もし私たちがそれぞれの信念の強度に基づいて物事を順序立てて並べるよう求められた場合、私たちは信念自体の意識からそれを実行できるし、証拠そのものに再度訴える必要はほとんどない。一般的な会話で使用されるより緩い命題を一時的に置き去りにして、数値データを提供する表現の中から単純な例を取り上げよう。私はこの町の全地域よりも私が住む通りよりも、この週に誰かが死亡することを確信していると感じないだろうか？そして、この町が私の住む通りの人口の100倍であることが分かっている場合、ほとんどの人は反射的に最初の命題よりも第二の命題に対して100倍の自信を持っていると主張できるだろうか？

あるいは非数値的な例を取ってみよう。私たちがしばしば、すべての証拠が心から消え去った後でも、二つの命題のうちどのものについて最も信じているのかを不躊躇に述べることができ、またその程度はどれほど信頼しているかを大まかな尺度で評価できる場合がある。それぞれの命題が支持する証拠が心から消え去った後でも、各命題を個々の優れた点のみに基づいて判断しなければならないときである。

そこで問題が提起されます。一般的な意見、つまり日常的な言語に示されるように、それが正しいと仮定すれば——当然、その重みは考慮に入れなければなりません——何かを数値的な種類の命題について述べる「部分的信頼」があると言えるでしょう。私たちが目指すことは、この信頼に対するテストや正当化を見つけることです。つまり、実質的に、その信頼が正しかったかどうかという問いに対する理解できる答えを得ることです。「それは正しいでしょうか？」という問いに対する答えを偶然見つけ出しますと、その答えは別の重要な問い——そして遥かに複雑な問い——についての多くの光を投げかけることになります。すなわち、「この部分的信頼とは何か」という問いです。

18. 私たちはまず、このような調査が普通の完全な信頼の場合にはどのように回答されるかを確認することから始めることを推奨します。そのようなステップは全く問題ありません。例えば、私たちは以下のような命題を得たとしましょう——これは帰納法によって得られたものであれ、通常の論理的推論のルールによって得られたものであれ、現在の目的には関係ありません——特定の酸素と水素の混合物は爆発性がある。ここでは推論があり、その結果信頼が生じます。今、私たちが信頼が正しかったかどうかについての問いがあった場合、どうすればよいでしょうか？最も単純な方法は、命題が真偽であるかを経験的に確認することです。私たちは物事を論理化しているので、信頼の正当性、すなわちその正確性のテストは、命題の真偽から最も容易に見つけることができます。もし私が推論によって特定の混合物が爆発すると信じた場合、その信念が正当である、つまり正しいと考慮するのは、適切な状況下で常に爆発が起こるときです。もし爆発が起こらない場合は、その信念は誤りだったということになります。

このような答えは、信頼の本質について何であるかを説明するためには、たぶんわずかなもの、あるいは全くないものに過ぎませんが、私たちの現在の目的、すなわち当否性とそのテストの意味を決定することだけを考えるためには十分です。すべての物事に関する推論において、信頼度の量を考慮に入れない限り、このような説明は十分であり、それは科学上の問題における通常の解釈となります。結論が特定的であれ普遍的であれ、この理解は完全に明白です。私たちは「人間はしばしば死ぬ」と信じるか、「すべての人間は死ぬ」と信じるかに関係なく、適切な経験の列によって信頼性をテストすることができます。

19. しかし、同様のテストを部分的な信頼に適用しようとすると、不快な困惑に直面することになります。この問題は、その主題を扱った人々が驚くほど見落としてきた困難です。単純な例が目的を達成するためには十分であり、私たちはコインの投げ入れというケースを取り上げます。私は一つのコインを投げる準備をしており、従って「表が出る」という半信半疑の状態でそれを予期しています。ここで見落としているのは、我々が過去に検討した例のように事実を参照してテストを行う場合、信頼は必ずしも間違っていないということが忘れられていることです。したがって、上記のテストは失敗します。コインは必ず何かしらの結果を出すべきであり、このケースでは表が出るか出ないかのどちらかです：つまり、コインは表が出るか出ないかのいずれかでなければなりません。第三の選択肢はありません。しかし、それがどの方法であれ現れる場合でも、半信半疑という状態が解釈できる範囲では、その信頼は間違いであることが明白です。表が出た場合、私は十分に期待していなかったため間違っています；私たちは「半分の信念」を抱いていたにもかかわらずそれが実際に起こらなかったからです。表が出てこない場合も同様で、私は過度に期待していたため間違っています；私たちは「半分の信念」を抱いていたにもかかわらずそれが全く起こらなかったからです。

同じような困難は、個別の偶然の事象に対する部分的な信頼性を正当化しようと試みるすべての場合において発生します。例として、もう一つ異なる例を考えましょう。ある男がサイコロで六目が出れば$p$1を受け取ることができますが、それ以外の目が出れば$p$1\_s\_（1\_s\_は1先令）を支払う必要があります。一般的には、彼がこの機会に対する2\_s\_ 6\_d\_（2先令6便士）を支払うべきであるし、もし彼がそう支払った場合、それは公平な金額であると認められるでしょう。この例は前回の例とは異なる点として、単なる命題に対する信頼性ではなく、数学者が「期待値の価値」と呼ぶものを取り上げています。つまり、信念だけでなく、その信念に基づく行動もより重要な位置づけを与えているのです。しかし、同じような困難はここでも再び現れます。事象に訴えると、すなわち個別の事象に対して見ると、一方の当事者が補償なく金銭を失うことが明らかになります。このような期待が公平であるという意味合いとはどのようなものなのか？

20. この問題に対する可能性のある解答は、おそらく唯一の解答として、我々が「頭が出る半信半疑」という表現の本当の意味は、平均して毎回二回に一度は必ず頭が出ることを確信しているということだと答えるでしょう。同様に、第二の例では、「公平な金額」と呼ぶことで、長期的にはどちらの当事者も利益を得ることも損失することもないことを意味するとします。この形式の解答で提起された問題点については、今後再び言及する予定がありますので、ここでの唯一の注意は、その解答が議論の中心である全般的な問題を完全に放棄していることに注意を喚起することです。それは明確には個別の事象に対してこの部分的な信頼性が厳密な意味で適用されないことを認めているからです。実際、その結果は予想外ではありません；少なくとも、本稿を通じて採用された理論に基づいて考えれば然りです。確率の使用は個々の事象に対する知識不足を前提としますので、個別の事象について他の見解や声明を正当化する方法を見つけるのは容易ではありません。それ以上の意見を述べる余地はないと認めざるを得ません。

21. それでは、私たちはまだ特定の問題に対する解決策に非常に近づいていないように思われます。私たちは特殊な例をより精密に分析するにつれて、個々の事例において何ら量的な信念のような正当化を見つけることができないという結論が明らかになることがわかります；少なくともそれは通常の科学的結論（推引的または推演的）における意味での期待されるようなものではありません。しかし、一般的な印象は、一連の共通表現によって証明され、この量的な信念の存在を支持しています。これはどのように説明できるのでしょうか？例をもう一度引用し、より精密に分析することで、我々は満足する解釈を見つける道を見つけられるかもしれません。

前回の分析（§18）では、私たちは真であるという命題の事実だけを信念の正当化として十分に停めました。しかし、その段階で止める限り、この説明は個々の事例に対する部分的な信念の正当化を提供しないことが明らかになりました。信ずる状態と不信ずる状態は、事件が起こることと起こらないことに対応する結果に非常に適切に対応しています；心理的目的のためにそれらをさらに分析する理由がない限り、私たちはそれほど深く分析する必要はありません。しかし、部分的な信念に対しては、そのような結果には対応するものがありません。なぜなら、我々が関心を持つような状況では、事件は部分的に起こることはありません。

そこで、分析を一歩進め、命題が真であるとは何を意味するのかと再度問い直すことにしましょう。これは当然ながら非常に長い且つ複雑な道を導きますが、その道の短い距離を進むことで、我々には大きな困難に遭遇することはないでしょう。前回と同じく、通常の完全な信念の場合で分析を適用し、それを例として挙げることから始めましょう。

22. 信念の本質的な性質についてどのような見解があろうとも、その状態において行動を起こす準備があることは、その心持ちに不可分な伴侶であるとおそらく認められるでしょう。信頼が変化する（特に精神健全な人の場合）ことなく、行動が変わることはあり得ませんし、逆もまた然りです。私たちの行動における何か一つでもそれが私たちの信念に関連しているものでなければなりません。まず、ペンシルに関する例を挙げて、完全な信頼を持つ状況から始めましょう。以前よりもさらに詳細に分析を行い、その後、部分的な信頼を持つ類似した例について同じような分析を試みます。

私は penny を投げる準備をしており、その penny が一方の面に落ちるか、あるいは縁に落ちないことを考えています。完全な自信で、それはそうなると確信しています。私たちの信念において何か他に含意されているものがあったとしても、少なくともこれは意味しているのです；つまり、私たちはその penny がそのように落ちることを前提として行動する準備ができているということです。すべての賭け事や他の一切の行為はこの仮定に基づいて行われています。 penny がそれ以外の方法で落ちた場合に生じるいかなるリスクも恐れずに受け入れます。これは、単一の投げ出しの場合でも長期間の連続した投げ出しの場合でも同様です。

しかし今度は penny が一方の面や他方の面ではなく、特定の面である \_head\_ に落ちることを考えてみましょう。この例は前回の例と一部類似しています。長期的には私たちの行動を賭ける準備ができています。大量の投げ出しの結果を考えるときには、 \_head\_ が毎回出ることを前提として行動します。例えば、賭け事の場合、 \_head\_ が出たときに $p$1 を支払う条件で、 \_head\_ が出なかったときに $p$1 を受け取る条件で$p$1 を支払うことに異論はありません。これは私たちの信念である \_head\_ と tail が等しく出ることを実践に翻訳したものです。

一瞬の考慮でも明らかになるように、我々の行動は多少変化する可能性がある。つまり、その結果が同じである限り、形式的には異なることができる。私たちは損をするたびに$p$1を支払うと、勝つたびに$p$1を得るという方法は、当該事例において、例外なく毎回十先令を支払うことと、頭が出たたびに$p$1を得ることと全く同じことである。これは、平均して頭が毎回出ることから明らかだ。長期的には二つの結果は一致するが、個別に考えると明確な違いがある。その違いは二つある。

まず第一に、我々は「毎回支払う」という概念から、「常に支払う」という概念へと移行する。第二に、我々が何かを得たときの支払い額は、得る場合の半分である。この差異は微細に思えるかもしれないが、我々の行動を個々の主観状態、すなわち信念に再変換すると、その効果は明確になる。私には、このような再変換により、我々が各ような命題について正確に定義できる程度の信念があるという概念が形成されたと解釈されるべきである。

我々が等しい量の損失と利益を経験し、それらを等しく頻繁に行うことで、二つの事象、つまり表と裏が出る回数が等しいという信念を得た。これは明白なことだったため、それは長期的な観点から言及されたものであった。しかし、これを常に支払うことと交換できると発見したとき、我々が偶に得る額の半分の金額を毎回支払うという状態は、各個別投げごとに半分の信念があるという状態であると自然に解釈された。

23. もちろん、そのような例は理論を確立するためには十分ではない。しかし、読者は、私たちの行動の大半が同じ結果に傾いていることに留意していただければ幸いである。これは単なる賭け事だけでなく、事件を数える必要のあるすべての行動においても、このような数値的な配分が可能であることを忘れないでいただきたい。従って、一般的な連想原則によれば、特定の頻度で起こる各個々の事件に関連する数値的状態ではなく、むしろ数値的関連性が不可分に結びつくことが非常に確実であると思われる。

たとえば、六面サイコロを投げると、一回に目が出る確率は1/6である。この事実の知識と、それがもたらすすべての実践的な結果を組み合わせることで、疑いなく、各単独の投げ方に関連する不可分の一六分の一の概念が形成される。しかし、一六分の一程度の確信を持つとは呼べないだろう。少なくとも、そのような個々のケースには、仮に存在したとしても、その分数の確信は正当化されず、説明もできない。

このような連想から、私たちは単なる偶然の事件に関しても、再現を期待しないまま確実な行動をとる。サイコロが一度だけ投げられる。私は1対5で目が出ないという賭けをするが、これは一般的に主張されるように、一六分の一程度の確信があるからではなく、むしろ長いケースの平均的な正当性に基づいてそのような行動をとるためであり、もしそれが長い系列の一部であると知っていた場合でも同じルールを適用する。これにより、両者の行動が同一になることが説明できる。さらに、連想によって、これらの状況に対する感情もおそらく非常に類似していると思われる。

24. ここで、この論文で採用した見解に基づいて、議論を打ち切ることも可能である。個々の事象が発生する確信度の「尺度」を説明することは必要であり、その説明は統計的な頻度から判断される場合に限られるため、このような一連の事象が私たちの出発点であったことから求められている。しかし、人々がまたは主張しているように、特定の確信度を持っている場合に必ずその尺度を参照する系列があるか、またそれを正当化する系列があるかについて調査することは義務付けられていない。

確信度を主観的な側面から扱い、確率を定量的信念の科学と見做す人々は、この調査を行うことが必要であるが、我々にはその義務がない。

この問題は、本題に関連して自然に生じる疑問であり、無視することはできない。私は、そのような正当化が他の場合でも相当程度適用されることを見出すことができるだろうと考えている。実際、私たちはしばしば単一の事態について決定し行動する機会に直面することはない。経験は、私たちが例示のためにこれまで仮定してきたように、一つの系列として整然と並べられた一連の出来事だけではなく、非常に多くの異なる系列からなる無数の出来事を私達に紹介するものであることを忘れてはならない。

一日の大半を過ごす人間は、ほぼ毎日のどの時間でも何かについて行動し、その結果、信念を行使しなければならない。一人の人間が自分の経験の中で再び現れないかもしれない無数の出来事の中から、それぞれ異なる種類のものであるにもかかわらず、多くの出来事を決定する機会がある。

しかし、これらが多様な種類であっても、その数が多いという事実によって、秩序は維持され、行動の道筋が正当化される。例えば、木々の植林において、ある一方的な方向で測定すると、平均的に各木が他の木から同じ距離にあるとすれば、そのような一定の秩序があることがわかるだろう。しかし、どのような他の方向でも同様な秩序が見られるかもしれない。

自然全体においても同様である；同一種類の出来事の一連の系列には規則性がある。しかし、異なる種類から順次いくつかを選んで系列を作るとする場合にも、規則性は存在する可能性がある。

このような状況において、私たちは信頼度の尺度に対する実用的な正当化を拡張することができる。例えば、ある人が生命年金を購入し、鉄道旅行での生命保険に加入し、lotteryに参加するなどを行う場合がある。これらの行為は個々の事象であり、彼がそれらを繰り返す意図がないとしても、一連の系列として整理することが可能である。

彼の行動と、したがってその信念は、それぞれの個別な事例において結果によって測定されるとすれば、正当化されるものではなく、逆に正当化されないことが示されている（§19参照）。もし彼が各種類の行為を十分に繰り返せば、それが正当化されることになるが、生活の条件によりそれが不可能である。

全く想像できるように、彼の続々と異なる種類の行動からなる新しい系列においては、規則性がない場合がある。しかし実際には、そのような規則性が見つかることが多く、これが確認される。これによって、生命年金、保険、lotteryを個別に取り扱った場合で彼が期待する利益と損失の均衡化は、これらの事象を全体として取り扱うことで達成される可能性がある。

もし彼が各々の機会においてその確率を正しい割合で評価し（それに応じて行動すれば）、彼の一生を通じて利益も損失も経験しない。また、異なる行動選択肢の中から、上記の信念の評価に従って行動する場合、つまり最も高い確率を持つ結果を選ぶと、彼は他のどの方法よりもこの方法でより多くの利益を得ることになる。

したがって、これらの「横系列」（cross-series）が存在することにより、生命年金、保険、lotteryを個別に取り扱った場合では期待できなかった行動の数が大幅に増加する。これらすべてのケースは、現在議論されている方法によって正当化される可能性があると見做すことができる。

以上のように、これらの横系列の存在により、多くの連続的な機会で利益と損失を均衡化する行動の範囲が大きく広がる。

25. 上記の説明では、私たちは個々の代理者の視点から彼の信念を正当化する根拠を提供してきたことを観察することができます。もし広範な共感が存在すると仮定すれば、その正当化はより広範に適用されます。私が宝くじで賞金を得ない確率が999対1であるという主張には、私の信念が1/1000の確信として表現される形でも非常に明瞭な意味を付与できます。正確には、非常に多くの投げ込みの中で、私は約1000回に1回だけ勝つことになることを意味します。

同様の種類の他の可能性も含めると、最後の節で説明したように、各個人は自分の人生の範囲内でこのような経験に達することができるでしょう。これは特定の一連の行動だけで達成できるわけではありませんが、他の行動と組み合わせることで達成できます。今、それぞれの人が他者の利益も自分自身の利益に比喩的に類似していると考える可能性を導入すると、上記の正当化はさらに広範に適用されます。ただし、直接競争相手の利益についてはこの考えを除くとします。

以下の文は、この見解を検証する公正な例として適切である。私は週のどの日でも死ぬことが必至であり、一周間しかない。したがって、私自身が日曜日に死ぬ確信度は七分の一となるだろう。ここでの偶然の出来事は明らかに繰り返しのないものだが、それでも、この公式によって与えられた価値を持つ信念は人全員にあるべきではないか？

前例と同様の原理がここで働いていることが見られるようだ。確かに私自身は一度きりの死を経験する機会しかないが、私は頻繁に死者となるクラスの一員であり、そのクラスからの証拠に基づいて意見を形成している。たとえば、私が生命保険で1000ポンド分を保険金として加入していた場合、事務所が私だけが日曜日に死ぬ場合のみ支払いを行うという声明に対して、7000ポンドを求めることに一定の妥当性を感じるだろう。確かに私個人としては、この配分が公正とは思わないかもしれないが、全体の人々がこのような原則に基づいて行動すれば、その合算した利益は公平に通貨化されることが可能であり、事務所にとっては全く影響がないことになる。

最後の数節の結果は、次のように要約できる：異なる事件に対する我々の信頼度は、一般的な表現によって認識されるが、その意味は明らかに何かを表している。しかし、それらの大部分の意味、そしておそらく唯一の正当性は、属する系列の対応する事件の系列の中に求められるべきである。これらの系列には、最初の見込みよりも遥かに多くの事件が関連付けられ得るということが示される可能性がある。信頼のテストと正当化は行動から見出される；全体としての系列に対するこのテストにおいては何ら特別な要素もなく、単一のイベントについて我々の信念に基づいて行動するときに行われる類似のテストとは何ら異なるものではない。

しかし、そのような形で適用されると、事態の性質上、個々の系列のメンバーに対して順次適用される。ここでは、私たちは各特定のイベントに対する行動を個別に考慮できるため、これは一定量の信頼を示すと解釈されてきた。おそらく関連づけの原理に基づくと、各単一のイベントに関して特別な精神状態が生じる。そしてこれらの関連づけは、系列から孤立させてもなお自然に保持されることが多い。

これらの単一のイベントを見つけるときには、それらは他の系列の残りと共に存在するときと同じように扱われ、それに向かって同じような感情を持つことになる。

27. いまや、我々が一連の出来事からなる系列において因果関係を必要としないことを、以前よりもはるかに明確に理解できるようになった。このような見解に対しては、個々の出来事が偶然的に起こることを認めるとすれば、その出来事の発生確率について話すことが可能なのかという疑問がよく提出されるかもしれない。人間生活の例を取り上げてみよう；30歳の男性たちの一組の平均寿命は約34年であると仮定する。したがって、彼らのどの一人に対しても、「あなたの余命は34年です」と言える。しかし、その人の生涯が因果関係による規則的な連続性を欠く可能性があることを認めるとすれば、このような話しが可能なのかという疑問が当然生じるだろうか？

この反論に対しては、因果関係の否定により個々の寿命についての期待値を特定することは、その平均寿命に関する主張と同様に、個々の場合ごとにそれが全く未知であることを考慮に入れる限り、より多くも少ないも言わないという点で、因果関係の肯定とは同じであることが指摘できる。つまり、上記のように個々について期待値を設定することにより、我々はただそのクラスの平均に関する声明をしているに過ぎない。

これらの個々の場合において因果関係があるかないかに関わらず、平均寿命については既定の前提のもとでは独立した経験に基づいており、それに対する知識は変わらない。つまり、どちらの仮説でも正当な推論は同じであり、その価値も等しい。唯一異なるのは、因果関係が否定された仮説においては、個々の「適切な」期待値について話すことが不適切であるという点で、我々に注意を強いることと、因果関係が肯定された仮説ではそのような知識が実用上得られないからといって無視されるという点だ。実際的なところでは、我々の知識量はどちらの場合でも同じであり、それは個々の平均に関する知識のみである。

28. したがって、我々が信頼度を正当化できる限界は、最初の見解よりも遥かに広いことが結論付けられる。個人が完全な自信未満の程度の信頼を感じる場合全てが確率の領域に入り得るかどうかという問題は、より広範な問いとなるかもしれない。しかし、そのような場合でも、信念が原理に基づいて正当化される可能性があるとしても、その規則は決して利用されることはできない。例えば、父親が病気の子供に特定の薬を投与するべきか否かで迷っていると仮定しよう。一方では医師は薬を与えなければ子供は死ぬと断言し、他方では誤りにより父親は自分が持っている薬が正しいものであることを完全には確信していない。数学的に何らかの数値的な確率を持つすべての事象について確信していると信じる一部の数学者は、その人の信念が何か「価値」を持っている（それが何かを特定できる場合）と主張し、例えば九分の九であると宣言するかもしれない。彼らは意味するのは、彼がそのような具体的な価値を持つ信念を持った九回のうち八回で正しかったことを示すということだ。反対側では彼の疑念も同様に考慮される。両方を合わせて考えると、賢明な父親としては彼が採らなければならない唯一の道は明らかである。

しかし、ここで重要な点は、一連の事象における特定のイベントの識別が純粋に主観的な条件によって依存している場合、例えばこの例ではその信念の鮮明さの度合いにより依存している場合、他の誰も判断できないのであり、そのためテストも証拠も見つけることができないということである。

29. それでは、しばしば試みられた主観的な基礎による科学の確立に関する努力について述べる。これらの試みは、ここでの議論のように、満足する結果を導き出すことができないと結論付けることができる。しかし、私たちは信念と行動が密接に結びついているため、上記の試みに対する防御的な説明が可能である。しかしながら、しばしば試みられるように、純粋な信念以外の感情も科学の結果によって正当化しようとすることがあるとき、混乱はさらに深刻なものとなる。以下は、アーチbishop Thomson の『思考の法則』（第2版、§122）からの引用であり、ここで考慮されている科学の応用の種類を示すものである。「確率論をその発明に関連する主題に適用するとき、即ち賭博においては、幸運と必要性に関する「道徳的な算術」という言葉でうまく表現された原理を忘れてはならない。ゲームスターが等しく有利な条件で対戦相手を見つけることは困難であり、またそのような平等性が存在する場合でも、大きな利益の優位性が深刻な損失の害と均衡することは不可能である。ブフォンは言っているように、「二人の男が全部の財産を賭けて遊びを決める約束をするとき、その結果は何になるだろうか？一方は自分の財産を倍にする一方は一切の所有を失う。損失と利益の比率は何か？それはすべてと何もないの間の比率である。一方の利益は僅かな金額であり、他方の損失は数的に無限であり、道徳的にはその大きさが彼の一生の労働でも十分に補填できないほど大きい」ということである。」

倫理的なアドバイスとしては、これはすべて真実で善いものです。しかし、この教訓が主題の科学に対する貢献と見なされるなら、それは非常に不適切であり、混乱を招くようなものに思われます。
確率論は、特定の種類の事象について数と量面での関係性に関して判断します；これは、これらの関係性についてどのような感情を持っているかという特定の人間の個人的な感想とは何の関連もありません。我々は算術の規則に補足的な推論を付けるべきではないでしょう。たとえ二倍の数が四であることは真実であっても、それから馬が所有者にとって二匹よりも二倍の喜びを与えることになるわけではありません。
二人が平等な条件で競う場合、彼らの確率は等しいです；言い換えれば、彼らがしばしばこのような方法で競い合うなら、それぞれが勝つ回数と負ける回数は同じであるでしょう。これが確率論が言えるすべてであり、状況下においてその男たちの決定や意見は彼ら自身だけが決めるべき事項です。
多くの人々はいかなる種類の平凡さにも耐えられず、財産を二倍にする見込みよりも全然失う可能性が高いというより大きなリスクの方が優先される場合があります。それらの人々のみが判断者となります。

もし楽と苦しみのバランスを導入するのであれば、個人はその計算を行う必要があります。仮定は、総合的な破産が非常に苦しいことである一方で、部分的な損失は損失の割合よりも比例的に少ない苦しみを与えるということです。従って、偶然の大損失による平均的な苦しみの方が頻繁な小損失より多いと結論付けられるのです。しかし、長期的には両方の場合で損失の金額が同じであるとしても、その推測は異なります。

もし、大きな支出を望む欲望が非常に強く、また生産量が多くて損失が容易に補填できるという状況を考えると、計算は逆になるかもしれません。このような状況では、高額賭博の方が低額賭博よりもより多くの幸福が得られる可能性があります。実際、すべての感情的な考慮事項は関係ありません；それらは至多、理論の応用に過ぎず、それぞれの個人だけが自分自身で行うことができるものです。

この主題についてさらにいくつかの注記を述べます。その内容は保険と賭博の章で扱われます。

30. このような考慮を導入することで、パルメンツァーの問題は非常に複雑化した。既にこの問題について一部説明を加えたが、ここでは非常に簡潔に触れることにする。この問題は、各セットごとの投げ方の系列に対し、私には何かを得る権利があることを示す。最低限の得物として一便士を受け取ることとする。私の得物は、それぞれの特定の種類のセットがどれほど珍しいかに比例して増加する。各種類は一定の明確な無限の順序でより珍しくなると観察されたり推測される。

問題文の表現を適切に解釈すると、私は永遠に続くと仮定される。従って、明らかに、これらのセットごとにどれだけ大きな費用を支払っても、時間と共にそれを補うことが保証される。これは数学的には、私には常に無限の金額を支払うべきであるという意味になる。この主張に対する反論は、実感的な人間はすらすらと大きな有限の金額、たとえば50ポンドも進むことはないと反論する。確かにそうだが、なぜか？

それは、彼自身や彼が支払いを受け取る者にとって、十分に価値のある投げ方が得られるまでには長生きする可能性が低いからである。更には、その手間と時間を考慮に入れるとなおさらだ。我々は、この問題が理想形で述べられたまま実際の形で現れる場合と一致することを仮定してはならない。木工師がユークリッドの二番目の公理について、彼の平面が六フィートの板で止まってしまうのであれば、それを理由に反対するのも無意味であるように。

多くの人々はこの点を見落としており、理論が数列のメンバーに対する数値的推論だけでなく、平均的な庶民が形成しがちな意見やそれに基づく行動についても正当化すべきだと誤って考えている。明らかに、このような考慮に入ること自体は私たちの適切な立場から逸脱している。これらのケースでは、我々は統計によって与えられた人間の行動のみに関心を持ち、その行動を実行する過程で経験される感情については直接関心を持つことはない。

この誤りは、政治経済学において使用価値と交換価値を混同し、パンの価値を生産コストで測定することを正当化しないという誤りと同じである。食事直後に飢餓の時に比べてパンがより価値があるからといって、その価値を生産コストで測ることは正当ではないように。

31. しかし、ある種の感情については例外を設けるべきである。これらは、異なる人物や異なる時間においても明らかにその一貫性と均質性を示すため、実際にはより優れた考慮対象としての主張がある。推理に関する科学との関連では、これらは絶対的な要素から見ればただの偶然であることは否定できないが、他の感情と比べると、これらの感情は絶対的な要素から離れない偶然のもののように思われる。

読者各位は、この章の前半部で、我々が「信念」と呼ぶ心持ちの複合的な性質について注意が向けられたことを覚えておられるでしょう。これは主に知性的な部分と、また情緒的な部分を含んでおり、経験に基づくように主張されるものの、実際には希望や恐怖その他の干渉要因を通じて経験が歪みながら作用しているのです。

信念の持ち主の心持ちについてのみ注目する限りにおいては、私はこれらの二つの部分が完全に不可分であると考えます。実際、これらを二つの部分と呼ぶことは誤解を招くかもしれません；それぞれ異なる源から生じるものの、その結果は完全に一つとなり、何らかの複合体として存在します。確率的な推論であれ、それ以外の推論であれ、このような歪みによって影響を受けやすいのです。

胆怯な男が次なる戦闘で負傷するという信じる権利があると正直に信じている場合、他の人々は同じ経験でも冷静な判断により、その確率が考慮する価値がないほど小さいことを認識しているかもしれません。しかし、そのような男の信念は、その内容を観察する限りでは健全な信念とは何ら異なるものではありません。彼の行動もまた、その信念に基づいて行われるとしても、差別化のための根拠にはなりません；彼が死に瀕した人物のように真剣に遺言書を作成することさえあります。唯一の解決策は、過去と現在の経験を引用してその信念を制御し、訂正することです[7]。これは確率が主に信念に関する法則に関連している理論に対する反対意見として提示されました。

しかし、この論文で我々が物事についての推論に関する法則に関心を持つという見解に基づくと、このような誤りや困難は消えてしまいます。私たちは物事についての推論に関する法則に注目すべきであり、信念の中に経験に依存しない部分があるとしても、それは注意の対象から消失します。

32. このような感情は、推論の科学の一環として注意を払う資格がないものと主張することができ、厳密にはそれらはその中から排除されるべきである。しかし、これらの感情が一貫性と規則性を持って生成され、かつその強度も一定であれば、これらは事実上偶然的で外的な付随要因として適正に受け入れられるべきである。これが驚きのケースである。この感情は相当程度の一貫性を示す。何かが実際に起こるとき、稀な出来事がどれほど稀であるかに関わらず、私と多くの他の人々はそのことに驚く。驚きは、興味の最も怠惰な形から「驚愕」の状態まで幅広い範囲で現れる。そして、同じ状況下において異なる時間に、また異なる人物に対しても驚きがほぼ同じように感じられるため、それはほとんどの他の予期せぬ出来事に対する精神的状態を伴うものの中でも極めて不規則なものとは異なり、一定の規則性を持つことになる。したがって、推論の科学に適切な入場権がない驚きは、その対象となるものと非常に密接に関連しているため、しばしばそのことに注意を払う必要がある。この観点に関する参照は、以下の章で時々見ることができるだろう。

次に、この主題をさらに説明するために、驚きという感情的な伴隨が、我々はそのような形で部分的に値を与えることができるという点について触れることもできる。この驚きの感情は、信頼性の通常受け入れられている分数とは二つの重要な点で異なる。

まず第一に、それは独立した存在を有する「独立性」という特性を持つ。これは自己で理解可能である。我々が示そうとしたように、信頼性は説明が必要であり、それがその後の行動によって見出される。しかし、感情に関しては如此ではない；これは自立し、自身の中で検討できるものである。したがって、第二に、それは単一の事象に対して、また一連の事象に対して同様に適用され、どのような種類の正当化にも適応することができる。

この点については、我々が以前の章で注目したように、特定の一因縁事象に対する我々の信念状態とは完全に反対である。これらの考慮は、我々についてこれらの事象に対して一定量の信頼性があるという一般受け入れられている教義を説明するのに役立つかもしれませんか？私は明らかに感情的な部分がそのような量を持っていることが明白であると認識しているので、それが無意識的に複合状態の知的部分に転移され、適用できない場所で正当化を見出すことができないという点について考えている。

33. 現在議論中の主観的な確率に対するさらなる例を挙げることができる。

一般的な言葉への訴えは常に役立つ。特定の単語を使用すること自体が、人類が何かを区別して選択し、その名前を与えているいくつかの明確な特性を見出したことを示すからである。確率が考慮に入れる種類の事象に対して、一般的な使用によって指定されたようなクラスの単語がある。これらの単語を検討すると、私はその定義が、我々が以前の章で注目した二つの視点、つまり客観的な側面と主観的な側面、事象の質と我々がそれらを考える際の心持ち、を明確に示していると考えられる。

例えば、「extraordinary」という言葉は、事象が一種の「ordo」やランクに並べられているという観察された事実を指しているように思われる。どの一つの事象もその位置を正確に予測できるほど適切な配置にはnessいなかったが、多くの事象を一括して考慮するとき、目を通すかのように線を追うと、それらは大抵順序立てて並んでいることが見えてくる。線から離れた事象のこのずれが観察され、「extraordinary」と呼ばれる一方で、その他の事象は「ordinary」または「in the line」と呼ぶことになる。同様に、「irregular」や「abnormal」も、多くの事象を検討した際に規則性や尺度による配置の観測から使われる言葉であり、原因が同じであるかどうかに関わらず、自然の中に秩序が見られるようになったときには同じ言葉を使用することがある。これは多くの事象がある場合にのみ成り立つことであり、個々の事象についてそのような配置があったと話すことはできない。

また、「law」は哲学的な意味での用語としては現在広く使用されている。この用語が導入された経緯は確実ではないが、少なくとも以下のような方法で使われるようになった可能性が高い：人間の行動に対する法の効果は通常、以前に存在しなかった規則性を生むことである。自然の中に秩序が見られるようになったとき、原因が同じであるかどうかに関わらず、同様の言葉を使用することがあった。それぞれの場合には、偶然的な逸脱があるものの、その一般性は一致していた。

一方で、「wonderful」、「unexpected」、「incredible」という言葉を観察する。これらの意味は単に精神状態を描写している；当然ながらこれらは統計的頻度における確率の意味でのみ限定されるわけではなく、ある原因から予想していなかった出来事が起こったときに驚くという意味を持つだけである。

現在、これらの二つの文種は本来、完全に異なるものであることを念頭に置いて考えてみましょう。一つは特定の性質を持つ事象を単純に示すだけであり、もう一つは同様の事象を表すが、その上で心の状態を示すという意味を持ちます。しかし、これらは同じ事象に対して普遍的に適用され、完全に同義語として使用されるため、これは確率が二つの側面から見られる方法と、それらとの間の密接な関連性が認識されていることを強調する例となります。これらの言葉は一般的には同義語として使われており、その意味を過度に解釈することは避けた方が良いでしょう。しかし、私が少し傾いているように思われるところでは、心の状態を表す言葉が他の言葉よりも広範囲に適用されることを観察することが可能であれば、これは既に述べたことの例となり、確率の領域は信念の変化が観察される領域よりも狭いことを示すものであるという観察を支持します。確率はこの変化が特定の統計的な方法によって引き起こされる場合のみを考慮します。

34. この問題の主観的側面に少し注意を向けたことは、最終的には興味深く重要なこととなります。まず第一に、単なる推測的な調査として、我々の信念量はどの命題についても注目すべき点です。これを深く研究することは心理学の領域に入ることになるかもしれませんが、それは私たちの主題と密接に関連しているため、すべてを避けることはできません。したがって、孤立した事象に対する期待や驚きの増加または減少の法則について議論し、これらの状態を個々の例で説明し、必要に応じて適切な量から逸脱した場合にそれらを修正することを考えます。

しかし、この理由よりももっと重要な理由があります。

確率の範囲内に限って私たちが議論する特定の事例においては、我々の信念を参照することなくこれらの問題を扱うことも可能です。この側面を取り入れるかどうかは、我々の好意によって決めることが可能です。たとえば、私に「A.B. が今年死ぬ確率」が「明日雨降る確率」よりも高いか否かについて尋ねられた場合、それぞれの確率を計算することができます（実際には私の信念と同じことである）、例えば前者は1/6、後者は1/7とすると、私は前者の方が期待される結果が高いと考え、より可能性が高い出来事であると判断します。この場合、信念が導入されようとされまいと、過程は全く同じです；つまり、我々の心的状態は実際には質問に影響を与えないのです。

しかし、他の場合では異なることがあります。例えば、二つのものを比較する場合があり、そのうち一つは確率的な意味での完全な異質なものであり、何らかの数値頻度を割り当てる試みも無駄であると判断できるものがあります。これらが共有している唯一の基盤は、それぞれに正当化される信念量だけかもしれません。これらの発生頻度を比較することは不可能です；一つはあまりにも稀に起こるため、その頻度で推測することができないかもしれませんし、それが一回しか起こらない可能性もあります。既に述べた通り、多くの出来事の私の信念は非常に複雑かつ広範な基礎に基づいています。私の信念は多数の対立する議論や、より或る程度遠くにある類推の産物であり、これらの証明自体は私の記憶からほとんど消えてしまっているかもしれませんが、その影響は弱いまたは強い確信として残ります。したがって、その時点では私は一定の精度で、ただし非常に僅かな精度で、私が抱いている信念量について述べることができるかもしれません。

このような異なる物事を比較することは不可能です；したがって、確率的な意味で自然かつ適切に決定される出来事との間で判断を下す場合、機会や発生頻度を引用することは不可能です。信念の尺度が唯一の共通基盤となるため、それぞれの場合においてこの量を比較する必要があります。そのテストは非常に粗雑なものとなりますが、上記の理由により全く何も提供しないよりは良いでしょう；一部のケースでは、我々が必要とする全てを提供することが可能です。

例えば、百万通の書簡の中から一つが郵便局で紛失したとします。そのような場合、ある書簡が紛失したか、私の執事がそれを盗んだ方がより確実に起こる可能性が高いのかを知りたいとします。もし後者の可能性も、前者のように数値で表すことができれば、比較は簡単になります。しかし、これは意識的に直接的に形にすることができません。

それでも、私たちが執事が不正行為を行ったという信頼度が百万分の一よりも高いと感じることができれば、それらは同質のものとなり、比較が可能となります。

35. いまや私たちは、しばしば用いざるを得なかったり、あるいは間接的に示した表現の、比較的正確な定義を提供する立場に立っています。読者も既にこの定義を求めているかもしれません。すなわち、「事件の確率」またはこれと等価なもの「特定の事象が起こる確率」についてです。

これらの用語は、系列内に存在します。無限大に近い数のクラスからなる系列の中で、ある属性や複数の属性の有無によって小さなクラスが区別されます。これは前章で完全に説明され、例示されました。この大きなクラスと小さなクラスは通常、「事件」および「その特定の方法での発生」の例として話題になります。

このような用語を採用し、適切な説明とともに、私たちは「その特定の方法で事件が起こる確率」と定義します。これは長期的な観察結果における二つの異なるクラス間の比率を表す数値分数です。例えば、ある乳児が80歳まで生きる確率について考えてみましょう。大きな系列はすべての乳児を含みます。一方、小さな系列は80歳まで生きている乳児を含みます。前者と後者の比率が9対1であると仮定します。つまり、10人の乳児に1人が80歳まで生きるとしましょう。この場合、任意の乳児が80歳まで生きる確率は数値分数1/10です。

この定義では、系列が無限大であり、固定型を持つ種類であると仮定しています。もし系列が有限であるか、または定期的にもしくは不規則に変動する場合（例えば衛生状況その他の要因により平均寿命が安定して変化している社会の場合）であれば、系列は科学の対象とはなりません。このような状況では、自然が提示する不適切な系列を、観察された事実から可能な限り離れないように適切な系列に置き換える必要があります。

これは自然の物体を厳密な科学の対象とする際に行われる作業であり、常に実行されます。

36. ここで、先ほど用いた「長期的な比例」について補足説明を加えます。この「長期」とは非常に長い期間を指し、実際には永遠に続くものと仮定します。数列の項を続けて取り上げるにつれて、その比例はまだ少し振れ幅があるかもしれませんが、その振れ幅は小さくなっていくでしょう。つまり、その比例は徐々に固定された数値に近づいていくのです。数学的にはこの固定された数値を「極限」と呼びます。この分数の値が先ほど述べたものであります。

推論によって結論を導くことが可能な場合においては、この分数は直接統計学的なデータを用いることなく、事象が発生する条件についての議論から得ることができます。第四章で説明した通りです。

ここで重要な区別が再確認されます：目の前にある実際の不規則な数列と計算のために置き換える数列との間の差異、そして「第一章 §13」で述べられた主張の意味が明らかになります。すなわち、厳密な科学的な推論は後者のみに適用されるということです。なぜなら、最終的に不規則な振れ幅を示す系列においては「極限」という概念を持つことは不可能だからです。

例えば、「特定の人々がコレラから回復する確率は50%である」という主張の意味は、その病気にかかる人々の半数が長期的には回復することを指します。しかし、十分な範囲の統計データを調査したとしても、社会の習慣や慣習が疾患の型や治療法にどのような変化を与えたかによって、我々が最初から同じ極限に近づいているとは限らないかもしれません。

系列内の二つのクラスの数の比における最終的な極限の概念は絶対的な固定性を含意します。したがって、自然がこの絶対的な固定性を提供しない場合（これは通常ゲームだけであるか、または証明されない限り）、我々の唯一の手段はそのような系列を導入すること、すなわち「しばしば述べられたように」適切な型を持つ系列を置き換えることです。

37. 以上は、概ね完璧な定義として受け入れられるべきであり、前章で同様に与えられても問題ありませんでした。しかし、この定義を現在の場所に置いたのは、この章で導入された概念と直ちに関連付けるためです；すなわち、我々自身の心の状態についてのもので、その状態は我々が何らかのイベントを観察する際に抱く確信度の量に関するものです。前章で説明したような確率を持つイベントについて考えたときの我々の信念の状態に関連しています。

信念が先ほど述べたような数値的な割り当てと同様に正確な配分を受けることを可能とするという見解には理由があったものの、それでも、我々の信念がそのような表現、「私の信念は確実性の1/10である」のようなもので説明できることが示されました。これは、そのような表現が特定の系列に対する関係性において見ることから、不明瞭または正当化されないという点を除いて、合理的な解釈を提示したものでした。

この解釈が採用される限りでは、我々の信念は上記の分数に比例すると述べることができます。これは、信念の純粋な知的側面についてのもので、それが行使される対象から考えられる形で切り離されることはありません。この知的側面には通常、さまざまな感情が伴います。これらの感情をある程度分離することは可能であり、分離した場合、他の感情と同様にその程度の正確さで測定することができます。また、個々のイベントに関しては、これらの感情が理解できるようになります。これらの感情は、事件の希少性を表す分数に多少応じて増加し、減少することが見つかります。驚きという感情も、ある程度の正確さでそのようにします。

以上の調査は、しばしば主張される「確率を表現する分数が個々のイベントに対する我々の信念の部分的な確実性に相当すること」についての真実の量を非常に簡潔な形で説明しています。この問題に関するさらなる分析は、心理学よりも確率論に関連しているように思われます。

1. この用語の通常の意味において、デ・モランが使用する言葉では、形式論理学は確率をその一分野として含み、彼のタイトル「形式論理学，または推論の計算法：必然的かつ確率的な」に示されているように、その一部門として位置づけられています。

2. \_形式論理学\_. 前書，第v頁。

3. ここでの主張を例証するものとして，最近1876年に，多くの読者が予想しなかった場所で，それが与えられました。私は，J.S.ミルの非常に興味深い神義論に関するエッセイを指しています（Many readers will readily infer, to J. S. Mill's exceedingly interesting Essays on Theism）。ここでは、その結論について批評する余地はありませんが、彼は信念が証拠によって正当化されていないものであり，おそらく正当化の可能性すら低いもの（例えば来世や神の存在に関する信念）でもありながら，それでも優れた精神状態において存続し，またそのような信念を希望として促進することも有益であると，非常に示唆的な形でその見解を表明しています。その真偽に関わらず，これらの信念は保持されることに特定の利益があると考えられています。

4. 例を挙げなければなりません。人間が行動を強いられる状況においてのみ、彼自身について何らかの信念を持っているかどうかを示すことができるとします。彼は「何も知らないし、これには興味がない」と主張し，そのような状態から信念を抽出することは不可能であると述べるかもしれません。デ・モラン（\_形式論理学\_, p. 183）が質問した場合や，ジェボンズ（第2版, 第212頁）が質問した場合でも同様です：「月の見えない側に火山があると考えますか？それらは私たちが見ている側よりも大きいでしょうか？」または「平滑多面体係数は正であると推測しますか？」これらの例では，我々が述べる問題について常に一定程度の信念を持っているという立場を示すには良い例とは思えません。理由を提示できない完全な無知状態は半信半疑に相当すると考えられます。

5. ただし，後述する第24節，第25節の原則に基づく場合を除きます。

6. この点についてより詳しく議論するためには、因果関係の章を参照してください。

7. 私が思い出す限りでは、主観的な側面と客観的な側面を区別して判断する最良の例は、列車の中で起こった出来事でした。私は恐れ多い老婦人遇し、彼女は事故に非常に懼れていました。私は通常の統計的根拠で彼女を慰めようとしました。彼女は静かに聞いており、その後こう返答しました。「そうですよ、先生。それは確かにそうですが、私があそれに関心を持たないからといって実際の危険が少しでも減るわけではありませんよ。」

8. これは、破られる可能性がある経験的法則についても依然として成り立つでしょう：現在では言葉の意味を大きく変更し、最終的な法則とみなされるものには破られないと考えられています。

\newpage

\chapter{確率における推論規則（The Rules of Inference in Probability）}
\label{ch:vii}

\begin{center}
\large\textit{加法定理・乗法定理・独立事象・逆確率（ベイズ逆推論）の仮定の恣意性批判}
\end{center}
\vspace{4mm}

第7章。

確率における推論の法則。

前章では、論理の類比から Immediate Inferences と呼ばれるものを調査しました。任意の年齢を与えられた場合に、十人中九人が四十歳まで生きるという事実が与えられたとき、特定の人々の寿命の見通しについて何を推測できるでしょうか？その推測は非常に単純なものではないことが示されましたが、それでも「比例的命題」と呼べるこれらの命題に対する信頼度の量的な意味での説明と、その量を正当化する方法が存在することも明らかになりました。現在は、より一般的に理解されるべき推論について考慮に入れる必要があります。すなわち、通常の論理学教科書で扱われるような種類の推論についてです。言い換えれば、特定の命題が包含する一般命題からどのように推測できるかを確認した後では、現在は一つの一般命題が別の一般命題から推測できる場合について検討する必要があります。

ここでいう「一般命題」は当然のことながら、確率的な観点から考慮される一般的な命題を指します。このような推論のルールは非常に少ないもので、その簡潔さゆえに長い時間を費やすことなく考察できます。現在手元にあるデータから、通常の科学の教科書で与えられる確率に関する規則を導き出すことができます。正確には、一部のこれらの規則を導くことができると述べるべきです。なぜなら、検討すると二種類非常に異なる種類であり、全く異なる基礎に立っていることが明らかになるからです。これらは形式的なものと、多少実験的であるものとに分けることができます。これはまた、最初の章で説明したような系列から、算術の単なる適用によって必然的に導かれる規則がある一方で、他の規則は特別な仮定に依存するか、経験への継続的な訴えと様々な帰納的資源による拡張が必要であることを意味します。

現在は形式的な規則の分類に主眼を置き、当面はその種類に限定します。後者の規則については、帰納法との関連性を考慮したときに完全に理解できるでしょう。

2. 構造的に厳密に定義される確率の基本的な法則、すなわち形式的な法則は、二つのクラスに分けられることができる。一方で加算や減算によって得られる法則があり、これは排他的または不互換な事象の関連性と一般的に呼ばれるものに対応する；他方で乗算や除算によって得られる法則があり、これは依存的な事象と一般的に呼ばれるものに対応する。これらを順に検討する。

(1) 簡単な加算により推論が可能である。たとえば、ある系列の各メンバーにおいて観察できる二つの異なる特性があり、これらの特性は同一個体には共存しない場合、明らかに任意の一回の抽出で一つの特性またはもう一つの特性を持つ個体数は、それぞれの特性を持つ個体数の合計に等しい。例えば、100人の新生児中36.4人が60歳以上まで生き残り、35.4人が10歳未満で死亡する場合を考える。1万人という大きな数を取ると、約3,640人が60歳以上に達し、約3,540人が10歳には到達しない。したがって、指定された範囲内で生存しない総数は約2,820人に達するだろう。もちろん、これらの割合が正確に定められていれば、その結果も正確であるはずだが、確率においては、この割合は長期的には数値が向かう極限であり、特定のケースでの正確な結果ではない。したがって、我々はただ、数値がより多くなるにつれて向かう極限に向かっていることを推測することができるだけである。

この規則は、一般の代数的形式で表すと、確率論の言葉では以下のようになります：二つの排他的または互換的ではない事象のそれぞれの確率が1/mと1/nであるとき、そのどちらか一方または両方が起こる確率は1/m + 1/n、つまり(m + n)/mnとなります。同様に、より多くのような事象がある場合も同じです。この作品で採用された原理に基づいて、このような代数的形式で表した規則は、統計的形式で表現したときと全く同じ意味を持ちます。前章の結論で示したように、例えば「ある特定の方法で与えられた事象が起こる確率は1/6である」というのは、「長期的にはその方法で一度に六回のうち一回程度は起こる」という別の言い方と等価です。

明らかに、この規則に対する補則を得ることは、加算ではなく減算によっても同様に行うことができます。一般化して述べると以下のようになります：二つの互換的ではない事象のどちらか一方または両方が起こる確率が1/mであり、そのうち一つだけが起こる確率が1/nであるとき、残りの一方向が起こる確率は1/m - 1/n、つまり(n - m)/nmとなります。

例えば、一年間に誰かが死ぬ確率が1/10であり、特定の病気で死ぬ確率が1/100である場合、他の病気で死ぬ確率は9/100となります。

ここで読者に注意していただきたいのは、この規則を「二つ以上の事象のうち一つまたはそれら全てが起こる」という言い方と、「二つ以上の方法のうち一つで同じ事象が起こる」という言い方の2つの異なる形で述べることができるということです。しかし、混乱は生じません；どちらの表現も正当であり、違いは単なる文法的なもので、最初の章（§8）で示したように、「方法」間の区別が深く多数であるかどうかによって、事象を通常的に同一視する資格があるかどうかに依存します。

ここでは、確実性を単位で表す一般的な教義の正当化について指摘することができます。特定の確率度合いはその適切な分数で表されるように、ある事象がm回に1回起こるという文句はその確率が1/mであると同等に表現できますから、この事象がm回にm回、つまり毎回無例外に起こることを述べるのと同じ意味になります。そして、それが毎回起こる事象は当然、その発生について確信を持つものであり、したがって、確実な事象の「確率」を表す分数は単位（1）になります。

また、与えられた確率が1/mである場合、その事象が起こらない確率は1 - 1/mまたは(m - 1)/mとなりますことは明らかです。

3. (2) 乗法や除法による推論も可能です。互いに排他的でない二つの事象について考えます。一方の事象が他方の発生に依存する場合、ある系列のメンバーの一部が特定の性質を持つという一定の割合を教えてもらい、さらにその一部が別の性質を持つという別の一定の割合を教えてもらったとします。両方の性質を持つ全体の割合は、上記の二つの割合を表す分数を乗じることで求められます。

ロンドンの住民について考えてみましょう。1000人中25人が一年間に死亡するとしましょう。また、そのうちの1人の死亡が発熱によるものであるということが分かっているとします。この場合、1年に1人に発熱による死亡があると推論できます。

もちろん、除法を使って逆の推論も容易です。全人口から発熱による死亡率を知り、他の死因に対する発熱の割合を知った場合、両方の分数を割ることで、すべての原因による総死亡率を求めることができます。

上のルールは確率の言葉で様々な形で表現されています。おそらく最も単純かつ適切な表現は、それが依存的イベントの規則を提供しているということです。つまり、あるイベントが1/mの確率で起こる場合、そしてそのイベントが起こったときに別のイベントも起こる確率が1/nの場合、後者の確率は1/mnとなります。この場合、前者に完全に依存していると仮定され、必ずしも同時に起こることなく後者だけが起こることはないとしますが、それは常に一緒に起こることを意味するのではなく、前者から後者が起こるという同時イベントの場合や、前者の後に続く後者の場合です。

4. 上記の推論は算術的な推論と同じような意味での必要性を持つものであり、それらを確立するために任意の仮定が必要ないことを要求しません。これらの推論において我々が仮定するのは、実際に与えられたデータによって支持され、実際には必要なものだけです。そして、これらのデータから算術を利用して推論を行います。上記の単純な例では、最も一般的な形式の算術以外は何も必要ありませんが、実践的な例はしばしばより深い方法が必要とする場合があります。そのような問題を解くためには、より高度で抽象的な算術である代数の全Resourceが必要となるかもしれません。

しかし、これらの手法への依頼の必要性がこの部分の原理に触れることはありませんから、ここではそれらについて言及する必要はありません。

5. 次に議論する公式は、以前のものとは異なる立場から成り立っており、その確固たる根拠と経験や仮説への依存を排除している点で異なります。前二つの例では、データが特定の条件のもとで与えられたと想定し、その条件下で異なる種類の事件の特性がそれぞれ独立であるか、または関連しているかということが既知でした。しかし今度は、このような条件が与えられることはありません。

例えば、10人のうち1人は黒髪であり、12人のうち1人は近視です。それから、任意の男性がこれらの二つの属性のいずれかのみを有する確率、または両方を有しない確率、あるいは両方を持つ確率について何らの結論を下すことができるでしょうか？

明らかに、これらの特性は互いに矛盾し、同じ人物には同時に存在しない可能性があります。あるいは、すべての近視者は黒髪であるかもしれません。あるいは、これらの特性は他のどの割合であっても与えられる可能性があります。

このような点について全く不知である場合、必要な確率に関する何らの推論も下すことができると考えられます。

しかし、推論は引き出され、特に多くの場合、公正な推論が行われます。問題の表面上の不定性から脱却する方法として、全男性の10%が黒髪である（これは与えられたデータの一つでした）という仮定だけでなく、短視と長視のどちらでも同じ割合で10%が黒髪であるという仌も行われます。1200人を一例として考えてみましょう。最初に与えられたデータから容易に理解できるように、このような1200人のグループの中で平均して120人が黒髪であり、したがって1080人が黒髪でないことがわかります。そして、厳密にはここで止むべきですが、少なくとも経験に再度訴えるまではここですべてを終えるべきです。しかし、ここでは止めません。仮定されたデータから推論すると、120人のうち10人（つまり120の十二分の一）が短視であり、残りの110人が長視であることがわかります。同様に、1080人のうち90人が短視であり、990人が長視であることも推論できます。したがって、全体として1200人は以下のようになります：黒髪で短視な人10人、短視で黒髪でない人90人、黒髪で短視でない人110人、短視でも黒髪でもない人990人。

この規則を最も一般的な形で表現し、確率の言葉で述べると以下のようになります：あることがpとqである確率がそれぞれ1/mと1/nである場合、そのことが両方ともpとqである確率は1/mn、pでありqでない確率は(n - 1)/mn、qでありpでない確率は(m - 1)/mn、pでもqでもない確率は((m - 1)(n - 1))/mnとなります。ここで、pとqが独立であるとします。これらの確率の合計は明らかに1となり、これは当然のこととして4つの選択肢のいずれかが必ず存在するからです。

6. 私はこの場合の推論と、前二つの場合における推論との区別を強調した。多くの読者には、その程度が過剰であるように思われるかもしれません。しかし、算術の必要な規則という非常に異なる信憑性を持つ2つの源から、確率が利用する科学においては、その区別にどれだけ注意を払っても足りないとは言えません。少数の学者は、この主題に関する著者の多くが恣意的な仮定を行ったことを否定することはほとんどありませんし、誰もが「逆確率」と呼ばれる場合における仮定が疑問視されるべきであると断じることはありません。したがって、そのような誤りや不確定性の可能性が初めて明らかになる地点で一時的な休憩と厳格な調査を行うことが最も良い道です。これらの指摘は、この主題に関するいくつかの優れた著者に適用されます。劣った著者や確率を理解せずにそれを引用する著者の場合では、さらに進むことができます。彼らが自分たちが独立性から推論した事象を仮定することを正当化できると主張していることについて、過言ではないと言えるでしょう。つまり、当該の特性に関する情報がないときには、それらが上記のように分布されていると仮定しなければならないと。そして、そのような比にわたる信頼性を割り当てるのです。これは「事象が独立であることを仮定する」と呼ばれ、この独立性から規則が確実に導かれるという仮定がなされ、逆に情報がない限り、事象が独立であると仮定することが正当化されることです。

この最後の主張の妥当性はすでに第一章で議論されています。これは経験によって提示される系列を前もって構築しようとする試みの一つであり、コイン投げにおける表と裏の等価性ほど強力な防衛ができないものではありません。しかし、「独立性」を持つ事象の意味は一瞬の考慮が必要です。

問題の状況は次の通りです。ある系列内の個体について、それぞれの質的差異（存在と非存在）によって、二つの異なるクラスの組が生成されます。論理を確立するために、その議論の下での仮定として、一つの前提が必要かつ十分であることが見つかりました。すなわち、各区分による分類は、他の区分によって生成された各クラスを同じ割合で再分割することです。独立性がこの意味で定義されると、当然その規則は成り立ちますが、特にその点に注意が向けられなければ、その用語は自然にはそう解釈されないかもしれません。

7. 以上が確率論における基本的な推論の法則であるとすれば、次に問われるのは、これらの法則が科学の主要な公式群や実際的应用においてどのような関係にあるかという問題です。答えは、これらの公式が科学に関連する限り、これらはすべて上記の基本的な法則の応用に過ぎないということです。このような応用は任意の複雑さを取ることがあり得ます。特定の例題が実際に提示される形では難解であるため、しばしば数学の最も深い定理に頼らなければなりません。しかし、これらの定理は実際には過度に煩雑な算術過程の簡略化とみなすべきであり、そのような過程は実践上、他の方法で行うことが困難になっています。

この説明は、通常与えられる規則の一部を解明するものですが、すべての規則を網羅しているわけではありません。確実な算術の法則によって証明できる規則については説明しますが、実際には帰納的推論に基づくものだけではなく、それらは完全に独立した規則として認識されるべきです。そして、多くの後者の規則が前者と結びついており、一般的な見解では、それぞれの規則が類似の起源を持ち、等しく確実であるという印象を受けることはほとんど疑いようがありません。この傾向についてはすでに何度もヒントが与えられてきましたが、その重要性は極めて高いので、個別の章（推論に関する章）でその考慮に専念することになります。

8. 上記の規則の妥当性を確立するために、本作業の一貫した計画に従って、我々は考察の基礎として、言及された事象の統計的頻度を使用しました。しかし、各公式が確立されると、その形式は等しく適切に、特定の事象が発生する「確率」を表す分数というより馴染み深い形で表現されることが示されました。したがって、現在問題提起できるのは、科学の主な対象として統計的頻度ではなく、信頼量を取る著者が（前章で説明したように）一貫してこの点を基礎とし、その規則を構築し、確率を表す分数は単なる同義表現であると考えることができるかどうかです。デ・モランは主張します。普通の論理では前提が絶対的に真であると仮定する一方で、確率の分野では「前提に対する部分的な信頼が結論に対してどのような影響を及ぼすか」を研究することがその任務であると述べています。したがって、厳密には、この見解によれば、我々は自己内省によって結論に対するこの後続的減少を直接決定することができるべきであり、統計学に依拠して信頼すべき量を教えてもらう必要がないはずです。

前章で得られた一般的な結果に共感した読者諸君は、当然ながら、私たちの意識に基づく量的信頼度に関する結論から何らかの名実一致する論理学的な科学を抽出することはできないという結論に同意することでしょう。例えば、100人の中の1人がインドへの航海中に亡くなり、その他の9人は5年の在住中に亡くなるとします。一般的には、今出発した人物が5年後にも家路に帰ることの確率は88/100だと考えられます。なぜなら、彼が到着する確率は99/100であり、そして到着した場合の生存確率は8/9であるからです。これらの分数を掛けることで結果を得ることができます。この例では、実際の論理の基礎は統計的であり、その後の過程や結果は分数に翻訳されるだけです。

しかし、信頼度の側面から見ると同様の結論が成り立つでしょうか？私は心理的な事実を認めます。つまり、我々はどの程度の信頼を持っているかは、その信頼に関連する事件の頻度に多少対応しているという事実です。上の例では、たとえば、それぞれのリスクが存在することによって、男の帰還に対する期待が減少し、両方のリスクを個別に受けるよりも少ない確率であると考えられます。つまり、これらのリスクを一定の方法で組み合わせているのです。

しかし私が認めないのは、統計データなしでこのプロセスを正確に行うことができること、そして、もし可能であればその結果が統計データによらない限り正しくなる保証があることなどです。実際には、私はこのような信頼度の過程についてあまり意味も安全性もないという見解を採用するべきだと考えます。確率は信頼度として表現されても、分数として表現される場合と同様に、正当化を求められるときには統計的な頻度に翻訳されなければなりません。そしてその後、推論を行い、その正当性も統計的な頻度に基づいて確認する必要があります。

したがって、推論の過程をそのように表現するなら、それは符号的な過程として扱われるべきであり、それが先述したような事柄に関する推論であることを示すものと考えます。したがって、私はこの後者の形式でそれを述べて説明することがより適切だと考えます。

確率論における逆確率とその必要ルールについて

9. 既に述べた通り、確率論における基本的な推論の法則は、第2節および第3節で説明した二つだけであるが、当然のこととして、派生する法則は豊富にある。これらの派生法則の性質と用途を理解するためには、主題に関する手引きを学ぶことが最も適切である。しかし、これらの派生法則の中でも一クラスは、その発生する問題に対する明確な区別が認められ、特別な検討が必要である。これは直接確率と逆確率の区別と認識されるものである。デモグラ（De Morgan）によって次のように紹介されている。

「前章では、その発生または失敗の状況を知った上で事件の確率を計算した。今度はその位置を反転し、我々は既知の事象から、同じ事が起こる可能性のある一連の状況に対する確率が何であるかを求めることになる。」従って、この区別は原因を与えられたときに効果を見つけ出すことと、効果を与えられたときに原因を見つけ出すこととの間の差異として簡潔に説明できる。

科学の原理に基づく定義についての議論と採用された前章で述べた原則によれば、このような区別の重要性は認識されないことが容易に推測できる。確率に関連するすべての対象物には共通の特徴が見られ、この特徴から推論可能な法則を直接導き出すことができる。他の全ての区別は単なる配置または管理上の違いである。

しかし、その区別が根本的なものではないとしても、逆確率問題に主に現れるような問題に対する実用的な扱いは、非常に深刻な不確定性と困惑の源となる。直接確率における任意の仮定は特に重大なものではなく、逆確率によって提出される多くの問題では、両方とも重大かつ避けられないものである。

10. この点は、特別な例の検討によって最もよく理解できる。どのようなものであれ、それが私たちの目的に適うのであれば、次のような二つの例を用いることにしよう：

(1) 9個の黒球と1個の白球が入った袋からボールを引く：その中で白球となる確率は？

(2) 10個のボールが入った袋からボールを引いて、それが白球であることが判明した場合、袋の中に白球がただ1つしかなかった確率は？

最初の例が属するクラスは既に十分に議論されている。この例の解釈は次の通りである：ボールを継続的に引き入れて再置き直すと、引いた総数に対する白球の割合は1/10に近づく傾向がある。考慮される行動は単一でありながら、その系列の一環として見なされる。つまり、我々の意見形成はその前提に基づいている。

結論として、個々には偶然のように見える事件の系列から一つ選ぶことになるが、長期的には特定の種類の事件（白球）は全体の1/10となる。この種類（白球）は予め選択される形で singled out される。その確率を1/10と述べることとは、単に物理的な事実だけでなく、それに適切に関連する他の心理的事実、感情、推論などを含む主張を意味することである。

11. 二番目の例を異なる解釈で考えるべきでしょうか？ここでも一つの事例が与えられていますが、質問の性質から判断すると、適切に参照できる唯一の数列は次のものであると考えられます：異なる袋からそれぞれ十個ずつ入っている球が引き抜き、常に白い球であることが確認される場合、最終的にどの程度の割合で、その球が一つの白い球しかない袋から抽出されたものかという比を求めることです。この問題においては時間に何らかの関与があるとは考えにくいので、この要素を考慮に入れないことにします。したがって、これらの二つの例に対する我々の見解を形成するために用いる二つの数列は次のようになります：（1）与えられた最終的な比で異なる色の球が提示される；（2）異なる内容を持つ袋が同様に提示されます。これらのデータに基づいて、（1）白い球であるという予想；（2）一つの白い球しかない袋であるという予想に対し適切な重みを付ける必要があります。このような形で述べると、問題は形式的に同一であるように思われます。

しかし、実際の作業に取り組むと、重要な違いが明らかになります。第一の例では、あまり恣意的な要素がなく、球は期待される割合にほぼ正確に出るでしょう。さらに、もし球がそうなることを証明するのが難しいという反論があったとしても、「よく混ぜられた」または「公平に分布された」という要求を提出し、あるいは事前に特定の種類の数列を得ることを確実にするための他の条件を導入することは不公正ではないと主張できます。しかし、第二の例ではそうは言えません。

12. 次の証明方法が一般的に試みられる第二の例を解くための手法は、次のようになっています。袋の中に必ず白球があるということが示されると共に、可能な前件は十種類あるとされます。すなわち、各々が十個のボールを持つ袋で、その中の白球の数はそれぞれ1から10まで異なる場合があります。このことは当然、我々の級数において見つかる項の種類を制限します。しかし、それ以上の制限が必要であり、これらの項が最終的に系列の中でどのように配置されるかの比例如何を知る必要があります。これには、袋に関する経験が必要となりますが、そのような経験は私たちに与えられるわけではありませんし、多くの類似したケースでは与えられません。したがって、問題を解決するためには仮定を立てなければなりません。次のような仮定を立ててみましょう：「上記で説明した各袋が等確率で選ばれる」として、その結果を見ていきましょう。袋が等確率で選ばれることから、それらの袋は必ずしも等しい数の白球を提供することを意味しません。むしろ、最後の例と同様に、各袋は所持する白球の数に直接比例して白球を提供します。一つの白球と九つの黒球を持つ袋からは十分の一の確率で白球が提供され、二つの白球と八つの黒球を持つ袋からは二分の一の確率で白球が提供されます。そして、この結果から容易に理解できるように、100回の選択の中で平均して55個の白球と45個の黒球が得られます。しかし、白球を提供しない選択については問題文では考慮する必要はありません。なぜなら、問題文は白球がすでに抽出されたという事実を前提としています。我々が必要とする系列は、白球を提供するものから構成されます。この系列の一部には、他の部分にはない追加の属性があります。それはどのような属性か？明らかにそれは、その袋の中にただ一つの白球しか存在しないという属性です。55回の選択の中で、そのような属性を持つものはただ1つあります。したがって、求められる確率は1/55となります。つまり、袋の中にただ一つの白球しかないものから白球が抽出される確率は55回に1回です。

13. ただし、斜体で表記された部分を除き、ここでの過程は他の例と全く同じです。この過程がやや長くなるのは、経験に即した結論をすぐに得られないためであり、結果は何であるかを推測する必要があるからです。しかしながら、その推論の妥当性自体は当然ながら最終的には経験に基づいています。

しかし、上記の部分は非常に重要なものです。この仮定が最も合理的なものであるとは限らないという点に言及する必要はありません。

異なる指定された種類の袋が等確率で選ばれるという仮定は、特定の状況下において最も合理的な仮定でしょうか？一人の人はそう考えるかもしれませんが、別の人は反対の見方をすることもあります。実際、その主題に関する優れた教科書[6]では、少なくとも一つの例において全く異なる仮定が採用されており、その場合、各数の得られる可能性が等しいという前提ではなく、各数を得る方法が等しいという前提がなされています。これにより、中間的な数の黒と白のボールを持つ袋よりも極端な数のどちらかを持つ袋の方がはるかに確率が高いことが導かれます。

この仮定に基づくと、5つの黒と5つの白を252通りで得られる一方で、1つの白と9つの黒を10通りで得られるため、後者の説明から抽出した袋の確率は最初に採用された仮定よりもはるかに低くなります。実際、その確率は今では1/512となり、元々の仮定では1/55でした。

一方の場合では、それぞれ異なる結果が等確率で考慮され、他方ではそれぞれ異なる方法で各結果を得ることが等しいとされるのです。

14. この種の不確定性は、逆確率において特に生じやすいものです。なぜなら、与えられた効果に対して特定の原因が発生する可能性を探る際に、しばしばその原因の相対的な頻度に関するヒントを提供されず、一般原則に基づいてそれらを決定しなければならないからです。実際の頻度が統計で与えられたり、その頻度がどのように形成されるかの条件が提示された場合、我々は何をすべきかを知ります。しかし、このようなデータがない場合は、予測に頼るしかありません。

上記の例では、袋が最初にどのように充填されたかについて教えてもらった場合、つまりどのようなプロセスで、またはどのような状況下で充填されたかについて教えてもらった場合、我々は何をすべきかを知ることができます。白い球と黒い球の数が等しい箱からランダムに充填された場合、ウィットワース氏の例における仮定は最も合理的です。しかし、このような情報がない場合は、完全に暗中模索で、§12において仮定されたものも他の多くの仮定と同様に信頼性と合理性がありませんし、優れた単純さを有しているだけです[7]。読者が第5章の§4, §5に戻ってみると、この特定の難題は完全に説明されています。誰もが実際的に、ある一定の特性的な配置や分布が必要であると認めていますし、その段階を選択した瞬間から、理論的困難は残りません。しかし、例の場合において、どの段階で我々が仮定を立てるのが最も合理的かを決定するときには、非常に明確なまたは理性的な指針が与えられるわけではありません。

15. しかし、上記の純粋に人工的な例と比較して、自然から提示される例を取り扱う場合、多くの不確実性が消失し、直接的確率と逆確率との区別もしばしば消えてしまうことがあります。このような状況では、原因は大抵、適切に定義されたルールによって決定されます。それらが単なる恣意的な推測の雲のようなものではなくなります。その頻度を把握するためには、テーブルを参照したり、形成される状況を調査することで推測することができます。このような場合、ほとんどどんな単純な例も、直接的確率と逆確率との区別が完全に消失し、または時間に関するものに分解されることを示すのに適しています。これは、後述する章で詳しく説明しますが、それは私たちの主題とは全く関係ありません。

もちろん、困難は上記のようなものと同じように時々ここでも侵入することはありませんという意味ではありません。既に述べた通り、それらは主題そのものに内在しているわけではなく、少なくともより単純な条件から推測される結果を決定する方法と比較してほとんど避けられないものです。つまり、意図するのは、自然の連続性や共存が比較的規則的で一貫した領域内に留まっている限り、原因の統計は効果の統計と同じくらい容易に得られるということです。一方が他方よりも多く恣意的な要素を持つことはほとんどありません。しかし、これは直ちに指摘されるように、仮説的な原因が調査に導入されると、その保証が失われるのです。

例えば、ロンドンでは年間約20人の死者が千人中発生すると知られています。また、そのうちの100人が死亡した場合に、約4人が気管支炎による死であると推定されます。したがって、未知の人々が特定の年間に気管支炎で亡くなる確率は1249対1となります。逆問題も同様な統計情報が利用可能です：ある男性が死亡した場合、彼が気管支炎で亡くなった確率は何パーセントでしょうか？ここで、男性の死が既に確定しているため、一般の平均的な死亡率を知る必要はありません。必要なのは、その疾患による相対的な死亡率だけです。確率論がこのような実際の統計に基づく推論のみに関わっていたら、直接問題と逆問題との区別はおそらく存在しなかったでしょう。

16. したがって、この主題に対する理論への貢献として、直接確率と逆確率の区別を放棄する必要があります。適切な統計情報が手元にある場合、両クラスの問題は扱い方において同一となり、それらがない場合でも、どちらの場合も解を導き出す権利はありません。しかし、この議論は別の重要な区別に再度注目を向けさせる役割があります。それは、確率計算が適切に適用される一連の例と、その他の例について思い出させます。統計データしかない場合にのみ確率計算が正当に適用されることを。それに対して、これらのルールを使用する必要がある場合でも、他により信頼性が高い情報を得る機会を放棄することになるか、あるいは人間の知性が定量的な見解を持つ権利がない事柄についての解決策を得ることになるかもしれません。

17. 私が記憶に残る限り、このような判断に対する最も近い実用的な正当化は、以下の例のようなケースによって提供される：「リストラの方法で治療された10件のケースでは7件が良好な結果を出し、3件が血液中毒を経験した。通常の包帯を使用して治療された14件のケースでは9件が良好な結果を出し、5件が血液中毒を経験した。リストラの方法の成功は偶然によるものだったのか？[8]」あるいは、別の言葉で表現すると、「短い経験から、一つの方法の方が他の方法よりも実際的な優位性があることが示された。類似の条件下での長期の経験では、この優位性が確認される確率はどの程度か？」である。

提案者はこの問題を「袋と球」の問題と見なし、以下のような例に比喩的に対応させた：一つの袋から10個の球を取り出した結果、7個が白く3個が黒かった。別の袋から14個の球を取り出した結果、9個が白く5個が黒かった。実際的な白い球と黒い球の比率は前者の方が後者よりも大きい確率があるか？この実際的な比率は当然ながら、最終的な白い球と黒い球の抽出比を真に示していると考えられる。

私はこの問題を数値解法が可能であると考慮する限り、唯一合理的な方法として、このような比喩的アプローチを採用すべきだと考える。

もちろん、異なる袋の種類が等しく存在するという前提は避けられないが、計算の正当性を主張する支持者たちの言葉を借りれば、「各種類が等しい事前確率を持つ」という前提を仮定しなければならない。しかし、この特定の例では、この前提の恣意性が通常よりも少ないと考える。これは、問題が単に非常に類似した2つのケース間での均衡を議論しているためであり、両側で不適切な前提が相互に相殺されるからである。もし一つだけの実験セットが提案され、その結果が長期的に再現される確率を評価するかのように尋ねられた場合、私は既に述べたような疑問を感じるだろう。しかし、ここではほぼ同じ条件下で行われた2つの実験セットがあり、それらの未知の条件が比較的同等に存在しているという前提はより恣意性が少ないと言える。

18. 一般的に導入される説明の例は、不適切であるように思われますが、逆確率に提出されたいくつかの問題の曖昧さを完全に理解するためには、それほど遠くはありません。自然な例も人工的な例も同様で、統計が得られない場合、調査は絶望的になり、計算のルールを定める試みはすべて放棄せざるを得ません。

例えば、いくつかの議論を引き起こした問題[9]として、ある星群が偶然分布の結果であるか否か、またはそれらの物事が（あるいは世界そのものが）偶然によって生成されたかどうかというような広範で曖昧な問題があります。

このようなケースでは、克服不可能な困難は、「偶然」や「ランダム」という言葉に含まれる意味を正確にどのように解釈するかにあるのです。第四章では、確率におけるその科学的かつ形式的な意味について説明しました。これらの用語をその関係から一般化しようとすることは、形而上学で無限や絶対性について話すのと同じような反論があります。

無限の大きさや無限の力は、ある程度理解できるか、または少なくとも何が議論されているのかを理解できますが、「\_その\_無限」は意味のない言葉だと私は思います。同様に、何かがランダムな方法で生成されたと仮定する場合も、その代理機関の性質、ランダム性の範囲などについて教えてくれなければなりません。それらのデータがない限り、何をすべきか分かりません。

このような問題に関する更なる考慮は、謙虚さを持って「論理の誤り」の章に移すことができるでしょう。したがって、私がこの主題について述べる最後の言及は、その章の中で見つかるでしょうし、因果関係と設計の章の結論においても同様です。

1. 「互いに排他的な仮説に対する不一致」や「互いに排他的な事象のクラス」という表現の方が正確かもしれません。

2. この種の例は、人間の死亡率に関するもので、カーリスルの表から引用された。これらの表は、他の表とは大きく異なることがよく知られており、デ・モランがこの点について最高の権威として、現在の英国中流階級の平均的な死亡率を最もよく代表していると述べている。

3. 私は「ほぼ」どのような割合でもよいと言ったのは、算術が仮定に課す制約があるためである。たとえば、すべての茶髪の人々が近視だと仮定することはできない。なぜなら、どの男性集団においても、前者の方が後者よりも多いからである。しかし、これらの制約の範囲は限定されており、その存在は上記の議論において重要な意味を持つものではない。

4. 「確率に関するエッセイ」第53ページ。私は彼が都市百科事典における確率に関する記事で、このような規則が新しい原理を含まないと述べていることを思い出させられた。

5. この点については、将来的な章で一般論的立場である主観的な論理システムの説明と例示を行った後に詳しく議論する予定である。

6. ホイットワース『選択と偶然』第2版、123ページ。またブール『思考の法則』370ページを参照されたい。

7. これらの仮定の防御に関する意見は分かれるが、その性質に関しては異なる。上記の仮定が疑わしいと認めつつも、最も公正な仮説であると呼ぶ著者がいる一方で、平均的な仮説だと見做す者もいる。

8. 教育時報；再版、第37巻40ページ。この問題はド・マカリスター博士によって提案され、大きな論争を引き起こした。通常、このような逆問題に関しては、どの著者間でも仮定の作成について一致しないものが多く、それゆえオッズの数値評価についても一致しない。

9. トッドハンター『歴史』333,4ページを参照されたい。

この問題に関する興味深い議論は、1850年12月号のフィロソフィカル・マガジンに掲載されたジェイムズ・D・フォーブス氏によるものである。その後の号でドンキン教授が反論を行っている。

\newpage

\chapter{継起の規則（The Rule of Succession）}
\label{ch:viii}

\begin{center}
\large\textit{ラプラス継起則批判・帰納的飛躍と確率計算の峻別・心理的信念と論理}
\end{center}
\vspace{4mm}

第8章。

\_順序の法則\_

1. ここで新しい用語を導入することについて、謝罪の言葉を述べる余地がある。他の選択肢は「帰納の法則」というタイトルを与えることだったが、そのようなタイトルを認めることはできない。理由はすぐに説明される。

1. 前章では、論理学で即時推論と間接推論と呼ばれるような確率における推論の種類について、ある程度詳しく議論した。私たちは例えば、ある男A.Bが一年以内に死ぬ確率が3分の1であるという結論から、その状況にある三人のうち一人が死ぬという一般命題が得られた場合、その意味を確認した。また、より完全な推論的性質を持つ法則の本質と証拠についても議論した。しかし、この点で止まるのは、この問題に対する不完全な見解に過ぎない。確率は実際の事物に関する真の推論の科学であるなら、それ以上の何かを導くべきである。つまり、観察された限界を超えて進出し、未観察のものについて結論を引く方法を見つけることが必要である。これにより即刻、確率と帰納との関連性に関する質問が生じる。この問題については、一定の詳細さで検討する必要があります。

確率と帰納との間には密接な関連があることは、どちらか一方について論じたほとんどの人々が気づいてきたことである。しかし、その関連性についての説明を見つけるのは容易ではなく、私には十分に満足できるものとは思われなかった。この関連性を明確に描写するためには、より狭い主題である確率に関する論文を探すべきである。しかし、ここが最も混乱が生じる場所である。

確率の領域は比較的狭く、その領域から帰納の隣接地域への侵入が頻繁に行われている。このようにして、前章で議論した算術的な法則の中には、厳密にはそれらと分類すべきではないものもしばしば導入されている。これらは全く異なる基礎に基づいているからである。

2. この混乱の起源は説明が容易であり、おそらくそれは確率の主観的な側面に過度な重きを置く習慣から生じている。これは、我々の信念量とその変動性について扱う部分に焦点を当てる傾向から来ていると考えられる。既に述べたように、この信念の変動は、確率が本当に本質的に持つべきものよりも最多で恒常的な付随現象であり、他の主題にも共通している。したがって、科学をこの側面から定義すると、他の主題もその内に入ることを要求するようになる；これらのうちいくつか、つまり帰納法は受け入れられたが、他には多少恣意的に拒否されたものもある。我々が帰納法によって得た広範な命題に対する信念は、検証以前に、推論から導かれた狭い一般化よりも弱い。この事実を考慮に入れると、ある種の確率に関する規則が与えられ、その規則によれば、この信頼度の減少が多くの場合計算できるとされる。

しかし、時間も我々の信念に変化をもたらす。ほぼ全ての事件に対する我々の信念は、その後長期間経って証拠が心から消え去ったときに、その時点よりも弱い。確率に関する論文の中に忘却に関する規則が挿入されないのはなぜか？もし彼に Tide（潮）が再び上昇する確固たる期待があると教えられた場合、彼は十年間の記憶に基づく事件に対する信念の強さについての規則も求めるべきではないか？この種の規則の違反は、我々が直ちに見るように、以前認められたものよりも数多く広範に存在する可能性が高い。実際、信念の強度を変化させる要因は無限に多いため、これら全てを一つの科学としてまとめることは不可能である。したがって、定義のために信念量を考慮に入れるべきではなく、または少なくとも付属的なものと見做す方が良いだろう。そして、主題の統計的側面から定義することで、その他の明確な境界線が描けるという事実に基づき、科学は定義されるべきである。

3. 然而，由于归纳法的重要性，值得单独讨论；一个简单的例子就能说明它対我们主题这一部分的影响。我们正在考虑某人A.B.再活一年的前景，并发现十个人中有九个同龄的人能够存活下来。然而，在形成关于他能否存活的意见时，我们会发现实际上有两个非常不同的因素在决定我们的信念强度中起作用；这些因素虽然在实践中彼此紧密相连，但我们很容易忽视其中一个，而将效果完全归因于另一个。

(I.) 有一个严格属于確率的因素；正如第六章所解释的那样，它衡量的是从一般命题推导出个别案例时我们対它的相信程度。假定十个人中有九个同龄的人能够再活一年，这并不必然意味着A.B.也会如此。我们用这样的表述来描述：我们的信念减弱的程度与9比10成比例，换句话说，这种信念由分数9/10衡量。

(II.) 但我们确定吗？十个人中有九个像他这样的人会再活一年吗？我们知道他们过去确实存活下来了，但他们是否会继续如此呢？由于A.B.仍然活着，显然这个命题在某种程度上是假設的，或者说通过归纳法得出的。然而，我们対归纳推论的信心并不如対从中得出数据那样确定。因此，在这里，有一个第二因素会减弱我们的信念；实际上这两个因素总是相伴而行，但在思想中它们可以被分开考虑。

上記で説明した二つの異なる原因は、しばしば混同されやすく、そのような例が主に抽出されるケースのクラスがこの混同を増幅させます。特定の比率において「すべての人々が死ぬ」という声明から、「彼らは引き続きその比率で死ぬだろう」という推論へと進むステップは非常に微細なものであり、そのことに気づかず、それに伴う信頼性の低下も予想されません。所謂事前（à priori）の例では、このステップはさらに微細です。力学法則の恒常性について我々がどれほど確信しているかを考えてみると、多くの人は過去に他のサイコロが等しく面が出たことを考慮に入れていない場合でも、長い目で見れば同じサイコロもそのように出るだろうと信じることを疑問視することなく受け入れます。

4. 既に指摘された（第六章において）ように、確率の定義が信頼の度合いとその変動を論じる科学として捉えられる場合、議論の核心は次のようになるべきである：上記（II.）で言及した原因が一つの単純かつ整合的な体系にまとめられ、従って任意の普遍的な同意の修正規則を得ることができるかどうか？これが可能であれば、それらを（I.）と同類に分類する有力な根拠が示されるだろう。すなわち、これらを確率の規則として考える理由がある。それでもなお、実質的に非常に異なる種類の規則となることは明白であり、必要ではなく偶然的なものである（これらの用語は哲学的体系へのコミットメントなしで使用できる）。しかし、それらを一緒に分類することで得られる簡素化という利点により、この反論は無効にされる可能性がある。この見解は、確率に関する著者たちによって、あるいは少なくとも私理解するところでは、定義と扱い方の方法から推察される形で一般に採用されている。または、これらの原因が広大な体系、ほとんど混沌（chaos）であると考えるべきか？これらは確かに分類され整理されることが可能だが、完全な普遍性を有する規則を得ることを期待することは不可能であり、常に例外となるものがある？その場合、唯一の道はそれら全てを確率から除外することである。すなわち、我々が最初に採用した起点として説明されてきた一般命題を前提とし、それが証人の証言によってその価値が低下するか、または私自身の記憶によって想起されるかに関係なく、それを与えられたものとして受け入れる道である。その命題は帰納法によって提供された多くの規則のいずれかから得られるか、あるいは推論的に導かれるか、あるいは我々自身の観察によって与えられるかである。その価値は証人の証言に依存するか、または私自身の記憶によって想起されるかなど、これらの原因やそれらの組み合わせにより低下する可能性があるが、これらは直接的に関係ない前段階の問題である。我々は統計的命題を真実とし、その獲得過程による価値の低下を無視して受け入れる。まずこの点から始めて、その後規則を適用する。実質的には、既に完成されたものとして受け入れ、その製造過程や完全性についての質問を一切しないでいることである。

5. もちろん、どの作家も上記の原因すべてを一つの計算システムに統一し、それらの包含によって生じる確証性の減少を一つの公式で表そうとする試みを行ったとは思われない。しかし一方で、彼らはこれらの要素への一切の引用を制限することには抵抗している。

科学の研究初期から、確率論の計算法を使って観察された事例から近接かつ類似の事例へと進展する試みがなされてきた。既に述べたように、実践ではこのような種類の拡張を避けることは不可能である。しかし、これらの例においては、経験の厳格な基盤からの逸脱が実際には認識されず、少なくとも論理的手続きの一環として認識されることはない。

確かに我々はそのデータによって完全に必要とされるものより少し外側へと進出し、それゆえに厳密には必要とされないより広範な結論を得た。そしてそのため、確証性が減少した結果をもたらす結論が推引で与えられたとしても、その逸脱は考慮の対象とはならないものとして省略されるべきである。未検討の事例が既に検討された事例と似ているという前提が採用され、これは十分な根拠を持つ可能性がある；計算理論はそれらが異なることを仮定し、その差異が小さいからこそ実際的な誤りは微々たるものである。

6. しかし我々が現在議論する規則、これを「連続規則」と呼ぶことができる規則は非常に異なるものである。この規則は、過去の経験を離れる事実を認識し、その逸脱の結果を計算の対象として明示的に扱うことを認める。

この規則は一般的に次のように述べられる：「既に観察された事件の再現確率を求めたい場合、そのイベントが観察された回数を1増やしたものを2増やした同じ数で割る。」

7. この則の起源を指摘することは、数学公式が適切な役割を果たすべきであることを再確認し、そして数学の名目のもとで受け入れられる驚くべき結論はどのようなものかを示すためにも有用です。第182ページの逆確率の例に戻りましょう。私たちは、ある仮定が成り立つ場合、袋の中に10個の白い球または黒い球があることが分かっているとき、一度だけ白い球が引き出された場合に、その袋の中にただ一つの白い球しかない確率を決定できることを見ました。これを行った後で、私たちは容易に「直接」この白い球が次の回にも引き出される確率を計算できます。同様に、その袋の中に2つ、3つ、そして10つまでの白い球がある確率を推定し、それぞれの仮定のもとで白い球が引き出される確率を決定します。これらの値を加えることで、次のような問いに対する答えを得ることができます：一度だけ白い球が引き出された袋の中に10個の球があるとき、次の回に白い球が引き出される確率は何ですか？

これまでは算術のみが必要でした。次のステップでは、より高度な数学が必要となり、その助けを借りて次のような問題を解決します：ある数（それが何であるかは不明）の白い球または黒い球がある袋からm回白い球が引き出されたとき、次の引き出しでも白い球が出る確率は何ですか？答えは(m + 1)/(m + 2)です。

数学の範囲まで到達したところですが、次に物理的な仮定が加わります。つまり、宇宙をこのような袋と比喩的に見なすことができ、そのようなルールが適用され、観察された事象が一定の方法でm回起こったとき、次の回にも同じ方法で起こる確率を求めることを解くために使用できるという仮定です。ラプラスは、彼の『哲学的 Essay』を執筆する時点での太陽が再び昇る（世界の年齢に関する旧な仮定のもと）確率は1,826,214対1であることを指摘しました。デ・モランは、河岸に立って10隻の船が旗を掲げて通過したのを見た男性が、次の船も旗を掲げる確率は11対1であると判断すべきであると言っています。

8. このような法則を真剣に受け取るのは難しい。証人の証言の評価における抽象的な公式の適用と比べても、その妥当性はそれほど確実なものではないように思われるからである。しかし、その有効性が議論される場合、二つの非常に異なる調査路線があるように見える。

(1) まず第一に、この法則をその主張するところのものとして、そして一般的に理解されているものとして受け取ることができる。すなわち、特定の状況における事件の再発現に対する期待度合いを測る尺度を与える法則であるとする。もちろん、本書で採用した見解によれば、平均的にこのような法則に従ってイベントが再現されるかどうかについて問い直すべきだという点に注意を向けなければならない。このテストを通じて、誰もそのような公式を正当化しようと試みることは絶対にないだろう。過去の事例が将来の再発現に対する信頼の基礎になるとは限らないだけでなく（詳細は後述する誤謬の章で述べる）、多くの場合において直接の逆説が成立している。また、このような法則は、次節で説明する重要な難題を引き起こす非常に深刻な問題にさらされている。つまり、事件が数回しか起こらなければ確実な指針がないし、一度も起こらなければどの種類の事例に該当するかについて全くの指針がない。先の旗に関する例では、なぜ単純にこの概念で止めたのか？旗のサイズや形状など具体的な特徴を指定しなかったのか？

デモグラは、このルールを制限的な意味でしか受け入れないことに注意しなければならない。彼は、これは当り確率の「純粋な帰納」に関する規則であるとし、「上記の規則によって示される確率は、単なる過去の頻度から生じる信頼性が増加する他の知識源によって拡大される最小値である」と述べている。つまり、彼は自然の均一性と因果法則の存在に関する信念が、単なる過去の頻度から生じる信頼性を補完することのみを認識している場合に限って、その規則を受け入れるだけである。しかし、これは過去の頻度が将来の再発生に対する不信根拠となるようなケースには対応していない。

9. (2) しかし、このような規則に対して別の見解を取りうる可能性もある。それは、その規則が受け入れられた理由に大きく関与しているかもしれない、非常に曖昧な認識である。

我々がこれによって表現しようとしているのは、--十分発達した理性を持つ心が分類し、意識的な推論を下すことができる場合の合理的期待の尺度ではなく、--原始的な信念の要素がどのように始まり、発展したかの法則である。もちろん、このような解釈は論理の領域から完全に撤退し、心理学への移行を行うことになるが、それは完全に正当な研究路線である。我々はフェッナーとその追随者の心理物理学測定に関する研究を発展させる試みに過ぎない。ただし、彼らのように投票箱とその内容との類推から始めず、実際の心的現象の精密な決定に基づいて調査を行うべきである。

光感覚の強度がその客観的な源の強さと共にどのように変化するかという法則は既に定められている。記憶の成長を単語や単純な音節の一連の繰り返しの数によって測定することも可能であることが理解できる。例えば、我々は実験を行い、ある発言を八回聴いた場合と七回聴いた場合で、人間がどれだけ効果的に記憶することができるかを決定することができる。

この例は、我々が想定した信頼の強度の増加を繰り返しで測定しようとする試みと非常に類似している。すなわち、もしこのような心的現象の順序での実験が可能であれば、特定の回数現象を繰り返し、その後実際に内省したり単純なテストによって、信頼がどれだけ増加しているかを確認するべきである。この視点から問題は私にとって希望のないもののように思われる。困難は十分に深刻であり、単純な感覚を測定しようとするとき、過去の訓練の影響を取り除くことや、実質的に心を開放的で無防備な状態にして刺激を受け入れる試みにおいてもである。しかし、信頼をこの方法で測定しようとすると、これらの困難は完全に克服不可能なものとなる。我々はここでも過去の訓練から自分自身を剥がすことができないほど、知恵や思考からも剥がすことはできない。私はどのようにして誰もが観察された再発現を他の経験した事象と分類しないか、そして特定の類推や帰納によって導かれず、デ・モランの「純粋な帰納」だけに頼る方法でないか回避することができると理解していない。同じような考慮は、当該ルールに対するもう一つの防御形態を否定することにも寄与する。たとえば、我々が事前に先決条件の数や性質について全く不知である場合でも、このルールに頼る方法があるという主張がなされることがある。私はその立場が私たちが実際には絶対に達することは考えられないと思っている。選択された現象があまりにも遠く離れていても、あるいは例外的であっても、我々はそれを既存の一般化と関連付けることができ、それらから結論を導くよりも純粋な先験的な考慮に基づくことなく結論を下すことができる。

10. それ以来、過去の獲得した知識を放棄したり考慮に入れることが不可能であるなら、唯一の選択肢は幼い心に作用する実験となるだろう。何らかの調査路線が不可能であると断言することは望ましくないが、その困難さは確かに莫大なものとなるであろう。そして仮に我々が事実を確保したとしても、それらは論理学において全くの無意味なものとなるであろう。問答無用で問題が解決されようと、例えば信念や感情が経験の繰り返しによって徐々にゼロ点から発展するか、あるいは単一の出来事だけで再現性のある出来事が起こったときに完全な信念を生じ、その後経験は我々の信頼を徐々に奪い去るかという証明が得られようと、成熟した心には何ら助言を与えることはない。

このような見解のようなものを受け入れている理由の一つとして、この規則が時折その下に存在していると考えるべきである。例えばラプラスは、それが実在的、すなわち客観的な推論の規則であると確信してそれを採用したにもかかわらず、彼の議論の終わりの方で生理学や心理学に関する多くの内容を考察し、彼が目指していたのは信念の自然史ではなく、その後の正当化であることを示唆している。

再び、ジェボンズが採用した奇妙な教義について述べる。その教義は、帰納法の原理は確率論に完全に依存しているというもので、これは「事実に関する知識はすべて確率的である」という主張とは非常に異なるものだ。この教義を理解するためには、何かしらの解釈が必要のように思われる。次章ではさらに詳しく議論するが、ここでは重要なのは、我々が現在抱いている反省的で理性のある段階において、確率的な推論は必ずしも帰納法に依存しているわけではないが、その一方で、どちらも他方を基礎とすることは不可能であるということだ。我々が最終的な信念を説明するためには、関連性というだけの理由では不十分であり、規則性の適用について議論する限りは、帰納法と確率性を区別しなければならないことだ。

1. ジョン・クレイグは、彼のしばしば引用される著作『Christianæ Theologiæ Principia Mathematica』（ロンドン、1699年）において、問題解決に向けた試みを行った。彼は以下のような問題を解こうとした：「移り変わりのある歴史事実について、その言及が生きた声で行われている限り、その確率が消失するときを決定する方法」（Quando evanescet probabilitas cujusvis Historiæ, cuius subjectum est transiens, vivâ tantum voce transmissae, determinare）。

2. m = 1のとき、分数は2/3になる。つまり、再出現が2対1の確率で起こると推定される。そして、この結果を受け入れる著者がいる。例えば、ジェボンズ（『Scienceの原理』第258ページ）は、「人間が初めてシャークを見たとき、そのシャークに伴侶魚がいるという事実を確認し、名付けた場合、次のシャークも同様に伴侶魚と伴侶関係にある確率は2対1である」と述べている。認識と命名を行うこと自体が、これらの魚を見たことが以前からあることを示している。その「事実」の観察特性の中でどれだけのものは再現性を考慮に入れるべきか？伴侶魚が先に現れなければならないのか、また同じ距離で現れなければならないのか？経度、海洋、季節、シャークの種類なども次の機会には再現されるべき事項として考慮に入れなければならないのか？このような疑問は尽きない。私は、どのような事件が最初に起こったとしても、定量的な目的での帰納法的問題を明確に定義することは不可能であると考えている。

3. Mind（10巻、454頁）において、ジェイコブ先生はエッビンガウス氏の『記憶に関する研究』という著作で述べた調査についての説明を読者に提供しています。

\newpage

\chapter{帰納法（Induction）}
\label{ch:ix}

\begin{center}
\large\textit{帰納的推論と確率の境界・クラス所属の困難性・生命保険と個別事象の価値}
\end{center}
\vspace{4mm}

第九章。

推論と確率との関連性

1. 前章では、未検討の例について推論を行うためのルールを検討しました。したがって、これは明らかに推引的な規則でした。しかし、その形式は通常表現される形で完全に失敗しています。そこで、確率が完全に形式的または推奨的な科学であるか、それとも未検討の例について有効な推論を行うことができるかどうかを問うことが適切です。この質問は、前二章を通じて部分的に暗示的に回答されてきました。本章では、確率と推引との関連性についてより詳細に調査し、限界を超えない範囲でその性質をできるだけ詳しく説明することを提案します。明瞭さのためには、議論を少し早い段階から開始することが適切であることがわかります。しかし、可能な限り速やかに地図を進むことで、確率と正しく呼ばれるべき領域の境界に近づくことにしましょう。

2. では、ある者が自然を最も広範な観点から調査し、既知の経験的事実を系統化することに着手する状況を考えましょう。そして、それが可能であると彼が見当たる場合、未知の他の事実を見つけることを発見します。彼は物理的で精神的な多様な現象を観察し、同時または順序立てて起こるものを含めます。これらの現象間にはどのような関係があるのでしょうか？どのような規則が存在するでしょうか？あるいくつかのことが観察された場合、それから他のものが推論できるような規則は何か？

注意深く述べると、彼は自分自身を調査することに決意し、事実自体を探求し、経験によって心が強制される法則があるかどうかという前もっての問いかけで邪魔されることなく進むことを想定します。これは、独立した必要のある精神的な法則の存在について不信を抱いているか、または心が経験によって提供される何でも観察し信じるべきであり、それ以外は何も信じないという信念から生じる可能性があります。あるいは単に後者の調査が好ましいと考えているだけかもしれません。言い換えれば、彼は形式論理を拒否し、客観的な存在の研究に取り組むことを想定します。

我々はここで、推論過程の性質、その前提となる仮定、そして適用される材料の特性をより鮮明に読者に示すために、虚構的なケースを思い描くことに過ぎないことを一時的に疑うべきではない。心理的には、誰もが知識の材料を持たず、周囲で見つかる統一性に対する偏見がない状態で自然を研究するためには、全ての精神的機能が完全に発達しているとは限らない。実際では、形式と内容、信念や関連性の法則とそれらが適用される対象は互いに作用し反応し、どちらも他方の存在を前提とする程度でしか存在しないことは当然である。しかし、このような論理体系の要件を注意深く説明するために仮定することは完全に正当であり、実際的な探求におけるほぼすべての段階で行われるべきことと同様である。

3. 彼の最初の仕事は、おそらく遅く且つ退屈なものだと考えられるかもしれない。それは個々の例を様々な区別点によって区別し、また類似性によって結びつけながら徐々に積み重ねていく過程となるだろう。これらの例はそれぞれ心の中で区別され、関連付けられ、その結果が一般化された命題としてまとめられ、その後で推論が引き出される。もちろん、これらの推論には新しい事実は含まれていない。それらは彼自身または他の者が以前に観察したことを再現するだけである。これまでに我々が行ったことは、物事を分類し、それを再度参照することだけである。したがって、我々はアリストテレスの定理に基づく普通の論理の手続きの範囲内を守っていることになる。これらの手続き（その最終的な解釈に関わらず）は当然ながら、アリストテレスの定理に従って、与えられたクラスが他のクラスの中に完全または部分的に含まれているか否か、または排除されているかを示すための方法として表現されることができる。

4. しかし、非常に短い期間の観察でも、彼の知識を広く拡張する可能性が示唆されるでしょう。経験そのものが速やかに、事件が一定の規則性で結びついていることを発見します；そして、「自然法則」があることが確認されます。事前にこれらの法則を信じる必要はなくとも、彼はすぐに実際上そうであることを発見し、それらが存在することを知りますが、その普及の範囲については確信を持つことはできません。この自然の秩序の発見は、彼の行動計画を一変させ、調査方法全体のトーンを定めることになります。彼の主な仕事は、これらの自然法則をどのように探求できるかを見つけることになるでしょう。

例を挙げて説明しましょう。例えば、私は巨大な板岩のような岩石を小さな破片に分解するために作業しているとします。最初はわずかな寸步ずつの作業で進めていきますが、存在する「裂け目」により、その過程は完全に変化します。この配置によって、数回の打撃だけで作業が完了し、私が最初に予想したよりも遥かに広範囲にわたるのです。言い換えれば、経験から示される切断プロセスそのものと他には何もなしで、物事が大幅に容易になり、また広範に及ぶ構造を発見します。このような発見は当然ながら戦術を変えるでしょう。その後の主な目標は、この裂け目の範囲と方向を確認することになるでしょう。

何か似たようなことは帰納法において見られます。自然の法則の発見は、少数の事実から推論によって一気に一連の対象全体を通じて飛躍的に結論を得ることができ、個々の事実を順次獲得するためには長い且つ疲労する調査が必要であり、実際に多くの場合では不可能であったことを可能にします。この状態が普遍的であるという証明的な証拠は存在しませんが、広範囲でそのような状況が見られることが分かったため、他の場所でも同様の試みを続ける決意を固め、失望することはありません。このような方法によって得られた命題や、それらの命題に依拠する原始的な事実から新たな推論を引き出すことができます。絶対的な確証ではなく、非常に実用的な程度の信頼性を持つ推論が可能です。私たちは信頼できる一般化を行うことができる大きな進歩を遂げました。例えば、ジョンとヘンリーは死ぬだけでなく、すべての人々は死ぬという結論に至ります。

5. 以上の簡潔な調査には、希望しますが、実質的な或いは科学的論理において帰納法の推論の性質を提示するための適切な概要が含まれています。これはミルによって引き出された区別を含んでおり、『論理学体系』を読んだ読者はこの区別に馴染みがあります。それは「公式に従って」導かれた推論と「公式から」導かれた推論との区別です。私たちは実際には経験によって提供されるデータから結論へ直接推論を行います。その一般的化を通じてこの推論を行うという方法は、便利な装いですが、私たちが自分の過去の経験や他の人々の経験を結論の根拠とする場合に特に必要不可欠であり、実際にはこの過程を二つの部分に分割するのが最も良い計画であると見つけることができます。第一段階は一般的化の確立に関連しています；第二段階（普通の論理の規則が含まれています）はこの一般化から何らかの結論を導くことができるかどうかを決定します。

6. これで、確率の役割と関連する科学との関係を確認する道が開けました。帰納論理は、上記のような一般化を見つけ出すルールを与え、それらの正確性を検証する方法も提供します。これらの一般化が普遍的な命題で表される場合、通常の論理の役割は、それらから何を推論できるかを決定することです。一方で、比例的な命題、すなわち第一章で説明したような命題で表される場合、それは確率に委ねられます。例えば、十人の乳児のうち三人が四歳未満で亡くなるという観察から、すべての乳児はその割合で亡くなるという主張を一般化するかどうかについて、帰納論理が判断します。このような命題が得られた場合、それがどの程度の価値を持つかに関わらず、それは一連の既知の系列であると認識され、すぐに確率によって取り扱われます。

この後者の場合においては、推論的要素と帰納的要素の二つの部分に分けることは、単なる便宜的な分け方よりも明らかに必要不可欠です。これは明瞭な思考と論理的な処理のために実質的に重要です。確かに、選択できる几乎每一个例子中，我们都会发现上述两个元素同时存在并结合在一起以确定我们的信念程度，但当我们仔细检查它们时，我认为这些命题的论据基础、它们産生的信念類型以及因此而産生的規則是如此完全不同以至于不可能整合成一个一致的系统。

従って、特定の帰納的公式が確率の一部であるという見解は、つまり最後の章で批判した順序則に見られるように、私はその妥当性を否定することができると思います。正確には、上記のように述べたとおり、帰納法は確率とは全く異なるものですが、ほぼすべての推論において協力していると言った方が正しくなるでしょう。例えば、私たちは帰納法を使って、10人の男性のうち4人が56歳まで生き延びるという命題が安全に一般化できるかどうか、またその程度はどの程度かを決定します。そのような命題が安全に一般化されたと仮定した場合、確率はそれを基にどのような推論が導き出せるかを述べます。

7. 以上が純粋な帰納法を確率の一部として扱う見解についてです。多くの哲学的または論理的な著者は当然ながら、全く異なる見解を持っていました。彼らはこれらの科学を非常に明確に区別し、あるいは互いに対照的に捉える傾向がありました。一方が哲学的または論理的と受け入れられ、他方が数学的として拒否されるという形でです。これについては無害な意味で「一般的な確率に対する偏見」と呼ぶことができます。

しかし、ここで別の見解を考慮に入れなければなりません。これは後者の反動によって、むしろ前者を超えるもののように見えます。そして時折、「どのような種類の帰納的推論もすべて確率的な問題である」という声明で表現されます。この見解の例を2つ挙げましょう。

古い文献を基に、バトラーの『アナロジー』の冒頭でしばしば引用される言葉から始めましょう。彼は「確率は人生の指針である」と述べています。この言葉はしばしば解釈され、実用的な結論を導く際に確率の法則や原理がすべてに及ぶものだと理解されます。

文脈から推測すると、彼の意図は確かに合理的な解釈と見ることができます。ただし、当時の時代背景において「確率の科学」が存在したとしてもという点には注意が必要です。ジェイョンス教授は『科学原理』（p.197）で、現代科学、特に物理学や数学の原理に従った類似な見解を述べています。「私は確信しています。 inductions の方法を健全な形で説明するためには、確率論の理論に基づいて説明せざるを得ないと思います。完全な知識だけが確実性を与えますが、自然の中での完全な知識は無限の知識であり、私たちの能力を超えています。したがって、私たちは不完全な知識に満足する必要があります。つまり、知識と无知が混在し、疑問を生むものである」と述べています。

8. この傾向を理解するには二つの意味で考えることができます。確率という言葉を一般的な流行語の意味で使用すると、これ以上の意図は、結論がすべて「確率的」なものであることを指摘することに過ぎません。つまり、絶対的な確実性ではなく、それらはそして絶対的に確定的ではないということです。これは当然のことながら認めなければなりません。物理的証拠その他の複雑さを理解している人なら誰もが、適用論理のどの分野においても絶対的理想な確信を達成することは不可能であると主張することはありません。抽象的な科学では仮定的な確信は可能かもしれませんが、具体的な領域における絶対的な確信はあり得ません。これは以前の章で既に述べましたが、しかし、論理の主観的見解に基づく正当化が存在しても、客観的な視点からするとそのような正当化は認められません。

9. しかし、ここでの意味は、この確信の欠如の根底が常に同じ一般的な特徴を持つということではないか。この主張が認められた場合、その力は実に強固なものとなり、したがって簡潔に述べるべきである。私たちは十分に確認したように、確率の厳密な領域における結論について「これは確定的ではない」と言及するとき、私たちが意味するのは、そのような結論が正しいのは一定の割合のケースのみであり、残りのケースでは誤っているということだ。

今、任意の帰納的な結論について、「絶対的な正当性について確信していない」と述べるとき、私たちは同じ解釈を意識しているのか？私にはそうは思えない。実際上、最終的には証明される可能性があるが、過去に私たちの調査方法を使用した際に失敗体験があったことが現在の完全な信頼性の欠如の主な原因だったという点である。しかし、これは我々が繰り返し強調してきたように、論理的な探求の領域ではなく、心理学的な探求の領域に属するものだ。

私たちは通常、過去の失敗の偶然的または反復的な例によって意識的に導かれるわけではありません。それらが影響を及ぼす範囲では、それは不確かな感覚を醸成することで働きかけるように見え、その根拠は統計学的なものではなく、少なくとも定量的には説明できない形で存在するように思われる。私たちの信頼性の一部は、最初に直接調査した人々から同情的に得られる。

また、分析によって特定の割合の過去の失敗が不信感の根源であることが検出される可能性があるが、そのような事実は表面には現れない。さらに、我々が完全な確信を持てない理由の一つは、一部のデータを記憶に依存していることだ。誰もが長い間前の出来事を記憶した場合でも絶対的な信頼性を持つことができない唯一の理由は、過去にそのような誤りがあったからではないと主張できるだろうか？

それ以外の意味においては、帰納的論理の結論に対する過度な証明力を付与することへの抗議としてのみ、その議論が確率に基づいているという言及は明白に説得力がないように思われる。

10. これで、因果律が必ずしも必要とされる範囲と、その範囲を超えた場合に必要ない理由が明確になりました。観察された事象から未観察の事象へと公式を一般化し拡張するためには、自然の順序に一定の恒常性があることが明らかに重要です。この恒常性は因果律という用語に含まれる意味の一形態です。もし人々が生か死ぬ状況が変わらなければ、10人のうち4人が50歳まで生き延びると推測する根拠がないことは明らかです。なぜなら、観察した個々の例でその割合がこれまでに見られたからという理由だけでは不十分だからです。これは因果律が主張する一貫性を認めることであり、この法則の真実性が一般化を得るためには必要です。言い換えれば、帰納的な部分におけるプロセスのために重要です。しかし、確率論に属するプロセスの一部に関しては、因果律が必要ないことも明らかです。当該の一般化が私たちに対して保証されている限り、個々のメンバーが因果律に従うかどうかは私たちにとって何の関心事でもありません。なぜなら、私たちは実際にはこれらの個々のメンバーについて推論を下していないからです。この最後の点については第六章ですでに十分に議論したため、ここでは更なる言及は必要ありません。

11. 以上の説明や、あるいは示唆は、これらの一般化を得るプロセスについて現在までに必要なものでしょう。それらが得られた場合、実際的にどのように利用するかを考えましょう。第六章と第七章で行った調査の過程において達成した点は次の通りでした：与えられた一定の種類の系列に関して、その構成要素について推論を下すことができました；つまり、その推論は特別な種類であり、その価値や意味については適切な場所で十分に議論しました。

現在、視点を少し変える必要があります。以前の章では、特定の数列が与えられていると仮定し、その数列から推論を行うことが主な焦点でした。しかし今度は個体のみが与えられ、適切な数列を見つけ出す作業が課せられる状況を想定します。この任務から始めるとどうすればよいでしょうか？

以前の場合では、このようなデータが与えられていました：十人の中の八人は五十歳で生き延びる年数は十一年です。そして、ジョン・スミスが五十歳である場合に、彼が六十一歳まで生きる確率をどのように推定できるかを確認しました。

12. しかし今度は、ジョン・スミスが現れたとします。この状況では、どのような方法で彼に対する数列を得るべきでしょうか？言い換えれば、適切な統計を集める方法は何でしょうか？実際の推論を行う際にまだ未知の事柄について真正に推測しようとする場合、問題はこのような形で現れることがほとんどです。

この問題に対する答えは、最初の見方では非常に単純なプロセスにより得られるように思えるかもしれません。すなわち、ジョン・スミスと同じ年齢の男が十一年間生きる者と生きない者の数を数えることにより得られるように見えます。しかし実際には、その過程は表面的に見るとそれほど単純ではないことが明白です。各個々のものは一つだけではなく、一つ以上の適切なクラスやグループに属するべきであることを忘れてはなりません。私たちは確かに実践上、個々のものや事件を特定の一側面から考える習慣がありますが、そのような慣習は私たちの側で行われるだけであって、義務的なものではありません。

明らかに、各個々のものは観察可能な無数の性質や属性を持つことができ、したがって無数の異なるクラスに属することができると思われます。どの一つのクラスに属している場合でも、当然そのものはより高い階級、すなわちそのクラスが種である上位の genus にも属することになります。

さらに、個々のものにはそれぞれの偶然的属性によって、他のいくつかのクラスと交差して属することができるのです。ジョン・スミスはたとえば肺結核患者であり、北半球出身者です。彼は人間であるからこそ当然に脊椎動物類、また動物類にも属しますし、より高い階級であれ何らかのものには当然に含まれます。

また、彼が肺結核患者であることにより、他のどのクラスよりも狭い別のクラスに属することになります。そして北半球出身であることにより、新しいかつ異なるクラスに属することができ、これは他のすべてのクラスとは重複しないのです。なぜなら北半球出身者の中には人間でないものも存在するからです。

13. よって、ジョン・スミスが特別なラベルなしで私達の注意を引いて現れる場合、彼がどの観点から見られるべきかという問題は、その分類に至るプロセスを相当程度恣意的とする。もし我々が彼を一般的な名前によって指摘されていれば、それは明らかにいくつかのヒントを与えることになるだろう。名前には特定の意味合いがあり、少なくとも選択すべき固定された属性群を指定する。しかし、名前と属性が結びついている場合、彼がどの名前で呼ばれるべきか、またその無数の属性の中からどのグループを選択すべきかは全く同じくらい不明瞭である。一つのことを教えてくれることは、もう一つを教えてくれることと同じ効果を持つ。

「彼がどの分類に属するかが恣意的」という言葉は、当然のことながら、論理的な決定基準がないことを意味し、選択は外的な考慮によって定められるべきである。個々の人物を見ただけでは、彼を無数の分類のいずれにも同様に参照できることが明らかになるが、それ自体で特定の分類を選択するための誘因を与えるものではない。

個々人が持つ複数の属性によって参照される可能性のあるさまざまな分類の多様性は、この章の前半部で示された推論過程に重要な影響を与え、我々が特定の主題に関連してより詳細に検討する必要があるプロセスである。

14. 現在、確率の場合において我々が直前に踏み出すステップの性質をより明瞭に理解するために、通常の論理の場合におけるその対応するステップがどのような形を取るかをまず検討することにより、それらの特異性を浮き彫りにすることができるだろう。つまり、あるジョン・スミスについて確認したいとする。彼は30歳の男性で、インドに居住しており、また癌に著しく影響を受けている。我々が確認したいのは、彼が20年間さらにその地で生き延びるかどうかである。

このように紹介される人物に関する用語は、既に示した通り、異なるクラスに属することを指す。これらのクラスに対応する命題は、以前の観察と推論によって得られたものであり、必要とするときに利用可能かつ手元にあるものとして仮定できる。これらは例えば以下のようである：「ある男性は50歳まで生きる」、「インドに居住する男性は50歳まで生きる」、「癌で苦しむ者は5年以内には生き残らない」。

これらの前提のうち、最初と第二のものは何らの結論も導き出すことができない。これらは両方とも具体的な命題であり、この場合では結論を導く指針となるものではない。第三の前提は我々の問いに明確な答えを与える。

論理的な読者にとって、ここでの過程は中間項を見つけることにより主辞と述語を結びつける方法を見つけ出すことを意味する。この場合における主辞と述語は、我々が確認したい個体とその指定された期間内での死である。これらはそれぞれ独立しており、それらが互いに関連したり無関係であれの理由から、直接的な接続手段がない。

上記で述べた様々なクラスは中間項の一組であり、それらに属する命題は対応する一長前提の一組である。これらのいずれか一つを使用することで、適切な状況下では結論の主辞と述語を結びつけることができ、つまりその男性が20年生き残るかどうかを推論できるようになる。

15. このような論理的な過程を実行する際には、読者に一時的に注意していただく二つの考慮事項があります。これらはこの場合単純ですが、確率における対応する場合では精密な説明が必要です。

まず第一に、何らかの推論が可能である限り、その推論は絶対的な確信を持って行うことができます。論理学は自らの領域において、逡巡や疑念を知りません。中間項が適切であれば、それは極端を結びつけるために、結論に対する一切の不確定性を排除します；そうでなければ、その過程はそこで終わりとなり、結論を導くことはできず、従って元の状態に戻ることになります。前提が正しいと仮定した場合、我々は結論を絶対的に知っているか、または全く何も知らないかです。

第二に、注意すべきは、上記のような大前提の可能性のある代替案のいずれもが他のどの前提とも矛盾することなく、またそれらとは全く不一致であることはないということです。これらを孤立した命題として見ると、当然のことながらそのような調和性を保証するものは何もないでしょう；それぞれ異なる述語を持ち、互いに支持や反対の関係にあるように見えます。しかし他の前提によって、各々が相互に関連付けられ、真実と一貫性の一般的な利益により、それらは矛盾することはありません。孤立した命題として考えると、すべての人々は五十歳まで生きる可能性があり、インド居住者はそうしないという前提も存在するかもしれません；しかしインドに住む一部の人々がいることを認めた場合、これらの前提が不一致であることが明らかになりますから、一方または両方が否認される必要があります。

すべての応用論理において、この自己矛盾を避ける必要性は明白で強烈であり、誰もその形式的な公理として「このような可能性のある大前提は相互に一貫している」ということを明示的に述べる必要はありません。この公理が遵守されていないと仮定することは、事実に関する二つ以上の矛盾する仮定を立てることと等しいことになります。

16. しかし、確率において類似の手順を試みるときにはその違いに注意しなければならない。普通の命題、普遍的または個別的なものに対しては、統計的な命題、特に「比例」型と呼ばれるようなものを代入する。言い換えれば、人間が20年生きることを尋ねるのではなく、1年生きることを尋ねるべきである。普遍的な命題を見つけることは不可能だが、適切な比例的命題は容易に得られるだろう。以下のようなものである：30歳の男性では、98人が1年生きる；インドの住民では、より少ない割合が生存する、たとえば90人が100人中生き残る；がん患者では、さらに少ない割合が生存する、たとえば20人が100人中生き残る。

以上に注意が向けられた両点においては、このような命題が以前考慮した命題とは明確な違いがある。まず第一に、普通の命題のように、はっきりと「はい」または「いいえ」と主張するか、あるいは口を閉じるかではなく、その中で包含される個々の対象について一種の不確定な或は逡巡した答えを与える。これはもちろん、「確率的命題」と呼ばれるものに固有の特性であるが、これにより二番目の重要な点へと導き、実際には可能にする。すなわち、これらの様々な答えは直接的にまたは形式的に矛盾しない（その性質上、比例的命題として通常は許されない）が、それらが互いに多少の対立関係にあることが、今後示さなければならない方法で明らかになる可能性がある。

したがって、このような前提から結論を導こうとする試みにおいては、困惑する状況に置かれるかもしれない。しかし、その困惑は二つの形態で現れる可能性があり、それらを順に考察しよう。

（訳者注：原著の文脈では、その後に具体的な例や議論が続くことが予想されるため、ここでの翻訳は一時的な終端とし、続きを待つことにします。）

17. 一般的な形は、個々の事例が適切に属する異なるクラスが順次包含されるときに現れます。この場合、統計データの集合は異なる結果を導き出しますが、それらはしばしば互いに非常に深刻な乖離を見せるわけではありません。その結果として、尺度の上昇につれて、より高い種類（genera）に訴えることで統計は個々の事例に対してますます適切ではなくなり、提供される情報も徐々に明確さと正確性を失っていきます。

最初に決定したい質問を思い出すと、ジョン・スミスが1年以内に死ぬかどうかです。しかし、適切な推論がないためこの事実に関する知識は得られず、第VI章で述べた説明と制限の下では、個々の知識の唯一の可用的な代替として、その状況にある男性のうち何割が死亡するかという統計的調査への回答を置き換えることに正当性を感じました。

18. しかし、すぐに疑念と曖昧さが生じ、その人物の状況とは何を意味するのかについての理解が困難になる。これらの状況自体は我々にとってよく知られているとしても、そのうちどれだけを取り入れて彼の運命を評価すべきかについては困惑することがある。例えば、最初に考えられるのは、彼が他の動物と共通している性質のみに焦点を当てる方法である。もしそうであれば、その主題に関する統計データが得られれば、おそらく次のようなものになるだろう：百頭の動物のうち九十九頭は一年以内に死ぬ。このような参照事例は不思議なことではあるが、論理的には、通常よりも広いクラスを対象としているだけである。同様に、もし好意的に見れば、脊椎動物のクラスや哺乳類のクラス（動物学者が必要な情報を提供できる場合）を選択することも可能である。

もちろん、これらの広範なクラスは拒否され、より狭いクラスが選ばれる。なぜそれらを拒否するのかと問われれば、自然な答えは「これらは非常に一般化されており、我々が直面している特定のケースと共通点が少ないため、その信頼性に大きな誤りを招く可能性が高い」ということである。正確性は保証されないとしても、少なくとも無駄な誤差の数と大きさを過大評価するのを避けることができる。

19. 以上の答えは非常に妥当である；しかし、幅広いクラスを引用するのを戒める一方で、逆方向に即ちクラスを過度に狭めることで間違えることはないという印象を与えるように思われる。そして実際には、そのような極端な状況を避けることができる。ジョン・スミスは単にイギリス人であるだけでなく、彼が生まれた地域の特定の部分や、現在居住している Presidency の特定の地域など、多くの追加的な特性も考慮に入れることができる。このような無数の追加的特性は、活力に関する問題にも一定の影響を与える可能性がある。これらのより狭いクラスを考慮に入れる理由は何故拒否するのか？確かに拒否することはするが、それは実用的な理由ではなく理論的な理由によるものである。

最初の章で説明したように、我々の数列に相当数の項が含まれていることが重要であり、それらは私たちにとって有用であるためである。そして、個々人の属性の多くは非常に稀であり、それを考慮に入れると、科学の基本的な前提と矛盾する可能性がある。すなわち、我々は正確には大量の平均値にのみ関心を持っているという前提である。より専門的で詳細な統計が良いとは言え、それらを十分に集めることができればである。しかし、多くの場合、これが不可能である。したがって、一つずつ属性を無視し、クラスの内容を拡大する余地がある。これは、我々が本当に欲しているような例と比較して、多少の多様性や適合性の欠如が増すリスクを承知しながらである。このような手順を繰り返すにつれて、我々はより多くの量を得るよりも質で失うことを避けられなくなる。

我々はそのような状況まで続け、統計が十分な観察範囲から得られるまで立場を固める。これにより、必要な安定性と一貫性の度合いを確保することができるからである。

20. たとえば、先に述べたような例では、連続するクラスの一つである人間が明確な自然種や品種である場合、この点でそれぞれ方向に完全な断絶があるため、誰もがここに立つことを自然と感じる。一方では、探求者はより高いクラス、つまり属性が少ない動物や組織された存在など、人間に参照を加えることは決して考えないだろう。他方では、地位や居住地などの観察によって生み出される下位のクラスは、通常、人間との特性的な違いが大きくないと、それらに注意を払う価値がないと判断する。

しかし、時々これらの特性は重要性を増し、そのような場合、それに対応するクラスに焦点を当てる。つまり、その特性によって区別されるクラスに注目する。たとえば、消耗性という質は、死亡に関するすべての問題において一人が他の多くの人々から大きく分離させ、肺結核患者の寿命に関する統計は一般の人々のものとは著しく異なる結果を示す。そしてその結果を見ると、肺結核患者が保険に加入できる場合でも、彼は通常よりも大幅に高い保険料で加入しなければならないことがわかる。つまり、属性は十分重要であり、新たなクラスや系列を区別するのに役立つ。

同様に、事故に対する保険についても同じことが言える。実際には、各個別の業種の会員に対して特別な保険率を設定しようとすることは試みられないが、彼らが被るリスクの違いは、その事実を一定程度考慮に入れる必要がある。したがって、業種は概ね二つまたは三つのクラスに分類される：普通のクラス、危険性のあるクラス、さらに危険なクラスなどがあり、それぞれが自らの保険料率を支払う。

21. これらのクラスのうち、一つが自然種類に対応するか、または自然種類と同等に重要な区別を含む場合、その過程は比較的容易です。通常、これらのクラスの一つが普遍的に適切なものとして認識されているため、多くの人々は選択プロセスが必要であることを全く意識していません。例外的なケースでは、体調が悪い人や危険な仕事に従事している人は除きますが、統計を集める必要があるのはその人の年齢と国の人々のクラスだけです。

しかし、これらの連続するクラスが自然によって明確に区別されず、容易に識別可能なグループとして並べ替えられていなければ、その過程はより明らかに恣意的になります。たとえば、ある男の家が火事になる確率を考慮している場合、この例でどのような属性を集めて満足するべきでしょうか？すべての種類の建物を含めるべきか、あるいは居住用の家屋だけを含めるべきか、それとも多くの木材がある家屋に限定すべきか、または暖炉のある家屋に限定すべきか？

これらの属性は全て存在し得ますが、もし存在すれば、これらは全て判断を修正する状況となります。ここでは、我々が手に入れられる適切な数の統計によって導かれなければなりません。また、保険事務所では、通常リスク、危険性のあるリスク、さらに危険性のあるリスクというような粗略な区分を実際には描き、これらのクラスに案件を分類し、そのクラスに適した統計に基づいて将来性を判断します。

22. 以上が所謂やさしい形での曖昧さですが、さらに強化された形でもこの曖昧さは現れ、最初の見方では我々に非常に大きな困惑を与えるように思われます。

上記の異なるクラスが順次包含されていない場合、統計表をどのもので使うべきか迷うかもしれません。例えば、10人のイギリス人中9人がメダガルでの滞在によって傷害を受け、一方で10人の肺結核患者中9人が同様の滞在により利益を得るという統計が考えられます。これらの統計は架空のものかもしれませんが、想像可能であり完全に矛盾していない可能性があります。

ジョン・スミスは肺結核のイギリス人です。彼の場合、メダガルへの旅行を推奨すべきでしょうか？それとも推奨すべきでないでしょうか？言い換えれば、彼の死亡確率についてどのような推論が下されるべきか？

両方の統計表は彼の状況に適用されますが、一方は我々に相当程度信じさせる一方で、他方は相当程度不信させる結論を導き出すかもしれません。これは当然、論理的な意味での矛盾ではありません。一方の命題は何か事象が必ず起こることを主張し、他方がその事象が絶対に起こらないことを否定するわけではありません。しかし、どちらかが我々に相当程度信じさせる一方で、他方は相当程度不信させるという意味での矛盾です。

これは当然のことながら個々の事象についての話であり、統計は仮定により何ら矛盾していないとされています。したがって、さらに詳しい情報がない限り、結論を下すことはできません。

23. もちろん実際には、これらのデータのみで判断を下さなければならない場合、私たちは肺の状態が人間の健康に比して出生地よりも大きな影響を与える可能性があるという考慮から選択を下すでしょう。つまり、消費性肺結核を持つイギリス人は、一般の人々と比較して、一般的な健康な同胞と比較して、消費性肺結核による寿命がより密接に関連していることから結論を下すでしょう。

しかし、これは当然ながら経験的な考慮事項を問題に取り入れることを意味します。与えられたデータは私たちに提供され、それらのみを使用する場合、その点について絶対的な不確定性が残ります。イギリスの消費性肺結核患者全員が他の気候に移された際にほとんど死亡することもあれば、ほとんど回復することもあります。

彼らが死亡した場合、これは最初の統計データと明らかに一致します。一方で、外国の消費性肺結核患者が気候の変化により利益を得た割合が通常よりも多いという事実を通じて、第二の統計データとも一致することが見つかるでしょう。

他の可能性についても同様な説明が適用されます。すなわち、イギリスの消費性肺結核患者がほとんど回復することを仮定する場合です。この問題は、ここで一般的に適用できると明らかにシンプルで簡単な一般化されたルールを援用できないため、絶対的に不確定であると言えます。

以前の場合では、十分に広範囲の統計データが存在すれば、より専門的な統計データよりも広範かつ一般的な統計データよりも優先すべきという単純明快なルールを援用することが推奨されました。しかし、ここにはどちらの統計データがより専門的であるかを知る手段はありません。

この点において困難は存在しないのであり、当分の間正当な観点で問題を捉える限りです。再び読者に思い出させます。私たちが現在どの位置にいるか；個々の知識を得られないと分かったため、類似の事例における平均的な現象についての知識に置き換えました。このステップは、問題が確率に委ねられた瞬間に取り組むことでした。しかし、事例における類似性という概念は我々を困惑させるものとします；多くの場合ではこれを回避することができますが、ここでは避けられません。

ここで類似性には二つの側面があり、これらは二つの異なる平均値を導きます。長期的には起こる現象についての二つの主張がありますし、これらの両方の主張は仮定により確認されています。その真偽については疑う必要はありません；どちらも経験から得られたものと想定されました。

24. この例が、生命保険その他の保険の基礎となる原理に対する反証（\_reductio ad absurdum\_）であると推測されるかもしれませんが、一瞬考えてみれば、これは全く誤りであり、保険の原理はこのような例にも同様に適用可能であることが明らかになります。事務所は、この仮定した状況において何ら困難を感じる必要はありません。彼らは（後で述べるように、おそらくそうしないでしょうが）、平均値に基づいて率を決定してその個人を保険することができ、彼に対して次のように言えるでしょう。「あなたはイギリス人です。私たちに来られた多くのイギリス人のうち、9割がマデイラへ行かなければ死んでしまうことが統計的に示されています。私たちはこれらの統計に基づいて率を決定し、長期的には私たちにとって全てがうまくいくことを知っています。

しかし、あなたは肺結核患者であることは真実であり、イギリス人のうち何割が肺結核患者であるか、またその中でどれくらいの割合がマデイラで死ぬかについては私たちは分かりません。しかし、これは私たちの目的には関係ありません。『9割が死ぬ』という公式は、これらの未知の比率を考慮に入れたものであり、それ自体の値は知らなくても、統計学によってその結果について全てを知ることができます。つまり、何らかの未知の要素が存在しても、それが生み出すすべての効果に対して既に考慮されているため、私たちの無知は結論に全く影響を与えることはありません。これが私たちの目的にとって十分なことです。」しかし、彼が肺結核患者であることを示したとしても、同じような言葉を彼に対して述べることができます。つまり、事務所はどちらの選択肢でも安全に手続きを進めることができるのです。

この状態は当然、物事にとって非常に不完全なものとなるだろう。唯一の合理的な計画としては、肺病を患うイギリス人の個案を隔離し、その状況に適した別途の計算を行うことである。この計算が行われれば、他のすべての表は彼らに関しては直ちに無効となり、他のどの表よりも少ない誤差で真実から逸脱することとなる。

25. 保険事務所が上記の条件で長く運営できない本当の理由は、一部の読者が想像するようなものとは全く異なるものであり、人間行動の各分野における科学構築試みにおいて多く見落とされてきた種類の考慮事項に属している。それはただ Jonah 以来の預言者たちが経験したあの不快な状況に過ぎない。すなわち、自己殺戮的な予言を発することである。つまり、一つの条件や原因も全て考慮に入れられた場合に完全に確実な結論を公表するが、その一つが予言自体の効果であり、その予言に関連する人々に対する影響であるという状況である。

上の例において、 Madeira の特定の事例に関しては、オフィスは消費性肺疾患を患うイギリス人を見つけ出し、その者がより有利な条件で保険料を支払うことにより、現状が大きく変わることを発見するまで、非常にうまくいくでしょう。しかし、もし彼らがそうした行動を取った場合、当然ながら統計は乱れてしまいます。これらの表は、イギリス人の寿命の一定割合（具体的な割合は問いませんが）が消費性肺疾患であると仮定して作成されました。その下でいう状況では、おそらくすぐにそうした仮定が真実ではなくなるでしょう。

「九人中十人が Madeira で死ぬ」と言われるとき、これはオフィスに（用語として）ランダムに或いは公平な割合で来社する人々のうち九分の十が死亡することを意味します。消費性肺疾患患者は、赤毛の人々や詩人、他の特殊なクラスの人々のように、ただそこに赴くと仮定されます。あるいは、その他のクラスよりもより多くもしくは少ない割合で赴くこともあり得ますが、常に同じ割合を維持することが条件です。これらの表はこの状態が続くことを前提として計算されています。実際の矛盾は、人々が一度これらの表を見た後でも、その状態が継続されることを仮定することにあるのです。

もしこれらの表の公表によって、それらに関連する人々の行動にどのような影響も与えなかったという仮定を立てれば、このような問題は発生しないでしょう。

26. ここで仮定されている前提は、既に述べた通り、矛盾しているわけではありませんが、説明が必要です。これらの前提を合わせて見ると、それぞれのクラスが均質であるという仮定とは不一致となります。すなわち、どちらか一方の全体に対して成り立つ統計的な比率が、その一部を取り出すと必ずしも成り立たないということです。特定の個人（つまり、肺結核を患うイギリス人）はそれぞれのクラスに属しており、もちろん両方の異なる統計データは彼ら自身について同時に真実とはなりません。これらの個々人はどちらか一方の特性と一致するかもしれませんが、両方に一致することはまずありません。おそらくほとんどの実際的なケースでは、どちらにも一致しない、すなわち中間的な特性を持つかもしれません。

今、任意のこのような不均一な体について「九分の十が死ぬ」という言及がある場合、その意味（またはむしろ含意）は、そのクラスをより均質な性質を持つ小さなサブクラスに分割できるということです。当然、一部では九分の十以上が死ぬ一方で、他の部分ではそれ未満が死ぬでしょう。これらの差異は個々のサブクラスの特性、体質、職業などによって決まります；そのような分割の数と、それらが互いにどれほど異なるかについては、おそらく非常に大きな違いがあるかもしれません。

今、それぞれのクラス全体について「九分の十が死ぬ」と述べる場合、最も自然で健全な意味は、すべてが例外なく国外に出た場合（あるいは同じことを言って、各サブクラスがその数に比例して代表された場合）の比率を指すということです。または、それらが他の比率で移動しているだけであるという意味かもしれません。しかし、どのような解釈を選んでも必要な条件は一つあります。すなわち、統計が作成されたときに各クラスからどの割合が移動したかについては、その割合を貫徹する必要があります。クラスが均質である場合、この条件は必要ありませんが、クラスが不均質である場合、この条件が確保されなければ統計は乱れてしまいます。

ここでは、二つの統計データセットが想定されています。一つはイギリス人用で、もう一つは肺結核患者用です。したがって、肺結核を持つイギリス人は一種類として二度カウントされるような状況になります。その結果、彼らの死亡率が中程度の場合、一方のクラスの死亡率を下げる同时に、他方のクラスの死亡率を上げる役割を果たします。

統計データが完全であると仮定すれば、それぞれのクラス全体を参照することにより、実際には同じ個人が各回カウントされることが推測されます。しかし、各クラスから代表者だけを選択する場合、それぞれの表に同じ個人が含まれる必要はありません。

したがって、保険契約に移行すると（私たちの仮定するマデイラのケースについての議論はまだこの点に限定されています）、一部の肺結核を持つイギリス人はイギリス人としてカウントされ、高い保険料を支払う一方で、他の肺結核患者は肺結核患者としてカウントされ、低い保険料を支払います。論理的には、すべての個人がそれぞれのクラスに登録される可能性があり、それぞれの率で保険料を支払うことも可能です。

上記で述べたことは、どのような個体もこれらのいずれかのクラスに提出されることが考えられるということです。ある他の者がその保険料を支払い、個人としてはどの率で保険契約を行うかに関わらず影響を受けない場合を考えると、クラス内の一部（またはすべて）が一つのクラスに移動し、他の部分（またはすべて）が別のクラスに移動することも制限されません。

従って、ある人物が実際には不合理でも論理的には可能であると仮定して、一部の人々が必要以上に高い保険料を支払う場合、イギリスの肺病患者（全員または一部）は各カテゴリーで保険に加入し、適切な保険料を支払うことができるとする理由はありません。当然、彼らがこの問題を考えたときには、仮定された場合においてすべての人は肺病患者として保険に加入することを好むでしょう。しかし、これを行うと、それぞれの統計が乱れます。イギリス人の死亡率は直ちにメーデイラ島における保険会社にとっては重くなることになります；彼らの最良の生命が失われることで。また、肺病に関する統計（すべてのイギリスの肺病患者が既に保険に加入していなければ）も同様に悪化します[5]。したがって、各保険料のスケールをわずかに調整する必要があり、これが先ほど述べた混乱です。ただし、明確に理解しておくべきは、本来の統計が不整合であることを証明したわけではなく、単にその統計に基づく保険システムを実行する実用的な障害があったということです。

28. 今考慮されている困難な例題は、その主題に未習得の学生には疑問を抱かせるものでしょう。これらの問題は、本稿全体を通じて主張される学問観と整合性が取れず、平均値のみに関心を持つという見解とは調和しません。おそらく、この場合における男の寿命には二つの異なる価値があること、そしてこれらは両方とも真実であることはできないだろうと主張されるかもしれません。なぜそうではないのでしょうか？彼の「価値」は単に、そのような状況にある人々が平均してどれだけ生き延びるかという年数に過ぎません。我々は彼を異なる状況で提示されています；したがって、これらの状況が異なる平均値を提供するのは不思議ではありませんか？

このような反論では、個々の結果についての未達成な予測を置き換えるために、可能な限り同じような一群の人々に対する実現可能な結果を使用しなければならないという点が忘れられています。困難と矛盾は、人々が個別のケースでその信念に正当性を見出すための説明を探そうとするときだけ発生します。

この特定のジョン・スミスについて何を結論付けることができるかと問われれば、我々は何も結論付けられません。実際には、以前の場合でも、統計データが限られている状況では個々の彼に適用できる価値としてそれを認識していました。ここでは二つの価値があり、異なる系列から得られたものであり、これらの価値が本当に異なる場合、不調和があると文句を言うことはありますが、そのような文句は個々の正当性を認めることで個々の結果に価値を与えることを許すことに権利がないのに何らかの形でそれを行うときだけ発生します。

29. それによって選択される系列や種類は完全に任意なものであるか？決してそうではない。何度も述べられてきたように、系列を選択する際には、個体の属性において可能な限り多くの相似性を持つメンバーを含む系列を選ぶべきであり、ただし十分な範囲であるという制約だけがある。つまり、上の例では二つの系列が存在する場合、これらは公平に矛盾しているとは言えず、唯一有効な批判は不完全さや目的に対する不足である。この批判は、必要以上に広い種類を含む統計データも同じ意味で適用されるべきであることを覚えておくことが重要だ。二つの異なるケースの間の唯一の違いは、一つの場合には最終的には多くの誤りを通じて真実へと導かれる道がある（ここでの真実が達成できるという意味において）という道を知っているため、その道を選んだことによるものであり、もう一つの場合には二つの異なる道があり、どちらも以前と同じように真実に導く。矛盾は各段階の正当性だけでなく最終的な傾向も説明しようとする際にのみ表面化する可能性がある。

それでもこれらの反論は、統計的頻度に基づいてただ予測されるものに対する完全性と有効性を著しく損なう要因であることは否めない。特に経験の初期段階において、少ない統計データしか集めていない場合に尤もそうなる。確率が提供できる知識は個々の場合では高次元ではなく、繰り返される誤差という特性的な欠点がある。しかし、その知識を自身の基準で測定しても、非常に低レベルの熟練度から始まることが明白である。それゆえ、これらの誤差は相対的に非常に多く、最終的にはどの程度まで減少するかによって大きく異なる。

30. ここでも他の場面でも、継続的な専門化のプロセスが進行しています。徐々に広がる経験の需要は、常に私達を過度に多様なクラスを分割する方向に導きます。しばしば唯一の苦情は、我々の目的にはあまりにも広すぎるクラスを参照せざるを得ない点であり、このクラスを便利のために分割することが必要であるということです。これは上記で示されたように、保険事務所が増加する業務により、特定の職業に従う男性たちを他の国民から分離することを可能かつ望ましいと判断した場合にも当てはまります。同様に、統計を使用する他のすべての部門でも同じです。この特定性に対する需要の増大は、実際には文明の進歩とともにどの町や国でも職業が分割される自然な方向へと導かれるように働きます。それに関連して、上記で述べた別の困惑についても同様です。推論がほとんど利用できない科学や実用技術においては、結果として我々の知識は主に経験的であるため、異なる方向を指す提案に直面することが非常に一般的です。新しい物質が発見されたり、より一般的に使用されるようになったりするたびに、それに関わる人々はこのような状況に直面します。新しい薬を使用する医師はしばしば、一方では特定の疾患に対してその薬を使うべきであるという推奨と、他方では特定の体質には試すべきではないという推奨に直面することになります。そのような体質で病気を患っている患者が現れた場合、どの推奨を採用すべきか医師は迷います。彼は即座に、より狭くそしてより専門的な統計を集める作業に着手しなければならないと感じます。これにより曖昧さから逃れることが可能となります。

31. この例と類似の多数の実践的な問題においては、当然数値データを引用し、それを援用する必要はありません。その判断が多少意識的にこれらのデータによって決定されることが十分です。例が適切であるためには、統計データが現在の知識状態において最終的な依頼となることを認めれば足ります。もちろん、経験則的法則が成分原因に分解できる場合、直接的推論に訴えることができ、その場合には統計学の使用、つまり確率理論の使用は代替されるかもしれません。

したがって、時間とともに統計的精密さの進歩は明白であり、それは人間の発展する需要に対応するためにクラスを無限に分割することによって明らかです。しかし、原始的な段階に戻り視線を転じてみると、その操作がいかに不完全な状態から始まるかすぐにわかります。この初期段階では、確率と帰納法は非常に密接に関連しており、しばしば一つになるか、または少なくともその機能が混同されることが多くあります。

32. 私たちの統計学の一般化が推論に属することが確認されたため、この一般化の過程は確率以前に存在するものと見なすこともできるし、少なくとも確率とは独立であることもできる。また、既に第六章で即刻的推論に対応するステップについて議論しており、第七章では三段論法的な推論に対応するステップについても議論した。現在の立場は、任意の数の一般化を所有しているが、個々の対象について推論を行うためにそれらを使用したいという状況にある。これは一般的な論理の一部では、結論を確立するために中間項を見つける作業に従事している場合と似ている。後者の場合では、この過程は非常に単純であり、異なる中間項の積み重ねが実質的な曖昧性や矛盾につながることはありません。しかし確率の場合には状況が異なります。ここで、個々のケースについて推論を試みる試みが頻繁に行われる場合（例えば連続法則の場合）では、この曖昧性と矛盾の全貌に直面することになります。しかし公正な立場から問題を扱うと、すべての困難は消失します。我々の推論は実際には個々の対象そのものについてではなく、それらの系列や連続について行われます。ジョン・スミスが年内に死ぬかどうか知りたいと考えていましたが、これは確定的な知識ではありません。しかしジョン・スミスは多くの属性を持つことにより、多くの異なる系列に属しています。したがって、中間項の多様性は予想される結果です。インド居住者や肺結核患者など、一年以内に死亡する男性の連続について知ることは可能であり、これらの選択肢から推論を行うことで、我々自身および他の人々（第六章で述べたように）の信念と行動が最終的に正当化されることになります。これらの系列を選択する際には、二つの対立する指導原則があります。一方では、より専門的な系列ほど良いです。最終的にはより正しくはなりませんが、常に近い正しさを得ることができます。しかし他方では、系列を過度に特殊化しようとすると、統計の不足から生じる実用的な反対意見に直面することになります。

1. 私の読者の中には、ブフォンの自然史（『人間の自然史』、§VIII）における非常に印象的な逸話に耳を傾けたことがある人もいるかもしれません。彼は最初の人間がその能力を完全に備えていると仮定し、しかし経験はすべて獲得しなければならないとします。そして、その知識の逐次的獲得について推測しています。彼の特定の結論については何を考えるかはともかく、この文書は心理学の学生にとって非常に興味深いものであり、示唆に富んでいます。

2. またダグラルド・スチュアート（ハミルトン編；第7巻、115-119ページ）も参照してください。

3. この要求は論理的推論の範囲内で確率的な限界内でのものであり、その実際的な普遍性や科学上の重要性について疑問を投げかける意図はありません。この主題については将来的な章でより詳しく議論します。

4. 特定の命題としては当然形式は同一です。前者における「一部」が後者よりも大きな割合に対応しているという事実は、純粋な論理から切り離された区別です。

5. 理由は明白です。現在マデイラで最も健康な英国民（つまり肺結核患者）は数えられず、一方で最悪の肺結核患者が増加しています（つまり英国民）。

\newpage

\chapter{偶然・因果性・設計（Chance, Causation and Design）}
\label{ch:x}

\begin{center}
\large\textit{偶然と自然法則・自由意志統計・数値定数と設計論の数理批判}
\end{center}
\vspace{4mm}

第十章。

偶然性、因果律および設計との対比

1. 前章での議論は、私たちの主題の範囲内で最も一般的な興味を惹く可能性のある問いに対する調査に道を開くことを目的としています。すなわち、ある特定の出来事は偶然性によって説明されるべきか、あるいは因果律や設計によって説明されるべきかを決定する問いです。主な困難が一般的に問題が概念的にとらえられ、述べられる際に極めて一般的なものであることに起因すると考えられるため、異なる具体的な問題点が混在していることを明確に区別することが必要となります。

I. まず最初に、我々の科学が仮定する「機会」の存在に関する非常に古い反論があります。この反論は、確率理論よりも遥か以前から存在し、現在ではほとんど無視されるべきものと見做されていますので、その考慮には長く留まることはありません。

「機会」という言葉を大文字で表記し、それが特定の創造的または管理的な代理体を代表すると主張する場合、おそらくマンイシズムの一形態に陥っていることになります。しかし、この反論の唯一の合理的な意味は、科学の原理が我々に以下のような仮定を迫ることであると考えるべきです：ある種の事象（ただし特定の事象のみ）は原因なく発生し、従って神の通常の管理から除外されます。これは以前にも繰り返し指摘された通り、全く誤りです。

確率の科学は、因果関係によるものかそれ以外でもかまわず、事象がどのように引き起こされるかについて何らの仮定も立てていません。我々が行うことは、個々の事柄について推論できない場合や推測できない場合に適用可能な一連のルールを確立し、その説明を行うことです。したがって、この反論は多少異なる形で述べるべきであり、最終的には以下のように帰結します：確率の科学が基礎とする仮定は、因果関係の存在について特定の範囲内では信じないことに矛盾しないという点です。ここで「因果関係」という用語は、単に規則的な連続性を意味することに注意してください。

この形で述べられた反論は完全に正当であるように思われます、もしくは、その根拠となる事実は認められるべきです。しかし、このような因果関係の欠如と神の超自然的管理との間にはどのような関連があるのか、私は全く理解できていません。

この神学的反論が消滅すると、物理学的な科学者たちや彼らと sympathize する人々は、同じ抗議を強調し始めました。そして、現代の著作においても時折類似の警告が見られます。
例えば、Hume の短い論文『確率論』では、「世界には確率というものは存在しないが、何かの事象の実際的な原因について我々が知る限りがなければ、理解力に同じ影響を与える」という言及から始まります。De Morgan はさらに、「確率理論の基礎は、その概念を形成した思考において存続していない」と宣言しています。「何かが過去に起こったり未来に起こるとは、特定の理由がない限り、他の何ものでもない」という主張も含めています[1]。ラプラスや他者から類似の言及も引用できます。

2. 上記で述べた論争の具体的な形態は、自然科学と物理学的な分野において最も見られるものです。確率と因果関係が不整合であるという主張は、防御的立場を占めています。これらの科目において我々が詳細について知る限りがなければ、その詳細が無因であることを前提に結論を導く根拠がないという主要な事実が強調されます。しかし、この不整合性はしばしば社会的な調査におけるより攻撃的な形で主張されることもあります。ここでは、平均値の既知かつ認められた規則性から詳細について因果関係を証明しようとすることがしばしば試みられます。数年前には、この問題に関して大きな論争が起こり、その大部分はブックルの『文明史』における必要主義者の力強い支持によって引き起こされました。

これらの例では、しばしば読者を驚かせるために選ばれた特異な事例によって、その同意を得ようとする試みが行われることがあります。選択された例は、より大きな論理的価値があるわけではありませんが、もし表現を許すなら、感情的に効果的なものであると言えるでしょう。

例えば、ブックルの歴史を読んだ人であれば、彼が年々一定の数に保たれる自殺者の数という事実について強調したことを思い出すでしょう。また、この事実が広く知られる際に多くの人々で起きた驚きや混乱も覚えておくことでしょう。

同様に、ラプラスは毎年一定の数を送る未指定宛名の手紙の数について言及し、ドーガルト・スチュアートはその特定の一貫性についてのコメントを行った点も、彼らがこの例を他の類似例よりもより注目すべきものだと考えていたことを示唆しています。

上記の例を支持する論理には一定の真実性があることは認めますが、その著者は私に見ると、確率理論と因果律の間の対立についての反対意見を過大評価しているように思われます。まず最初に述べたように、秩序や規則性というより広い概念については、それが我々が「客観的な偶然」と呼ぶものと必ずしも一致しない点があります。確率においては最終的な規則性が常に前提として提示されます。サイコロを振る際に、個々の投げ方だけでなく平均値さえも予測不能である場合、サイコロの形状や目印の数などの性質によって任意に干渉されることがあれば、当然科学はそのような事象を考慮に入れるべきではありません。

3. この点については、当然のこととして認めなければなりませんが、個々の事象の連続についても同じような推定はできるのでしょうか？不変的な順序関係（ここではより狭い範囲に焦点を当てるため）において原因と結果の関係を否定することは、科学的異端罪の疑念なくしてはできないでしょうか？少なくとも確率論の基礎に混乱をもたらすことはないでしょうか？デモグラントが見せたように、彼はこの点について強く否定しない立場を主張しています。しかし私自身は彼と同意することができません。この不一致は詳細な違いよりも、科学的な視点における違いから生じていると考えられます。

彼は個々の確率的判断が正当化されるべきであるという見解に常に傾いているように見えます。例えば、「ある事象が起こる確率比は三対二」と言及した場合、その声明を説明し、正当化することができると言います。しかし、個々の事象についての系列やクラスに必然的に参照することなく、このような声明を説明・正当化することは容易なことでしょうか？これは困難であることは明白ですが、個々の事象に関する知識が得られないだけでなく、因果関係がないため根本的に得られないとされる場合、その障害はさらに大きくなるでしょう。

この作品で採用した理論では、詳細については無知を前提としますが、この無知がどのような根拠に基づいているかは重要視されません。事象の性質から知識を得ることは不可能であり、因果関係が欠けているため根本的に得られない場合もあれば、そのような関係は存在することが分かっているものの、それを特定することができないか、またはそれを行うのが面倒である場合もあります。この無知という事実こそが確率論の理論を求める動機であり、その根拠は重要ではありません。

4. ここでの見解では、私たちは平均値に関心を持ち、あるいは平均から導き出された単一の事象が系列の一環として考えられる場合に限られます。我々は経験に基づく前提から始まります：この平均値には均一性があると仮定します。そして、この均一性が保証されている限り、個々の要素の因果関係に関する運命については完全な無関心を示しても問題ありません。その問いは以下のように形を変えます：この前提、つまり全体として平均的な規則性があるという前提と、詳細に関して無因性の乱れがあるとする信念との間に矛盾が存在するのでしょうか？私にはそれを証明するのは容易ではないように思われます。実際のところ、両方の信条は同じ人々の中で常に共存してきました。これは大きな意味を持つとは限らないかもしれませんが、それから推測できるのは、もし矛盾が存在するのであれば、それは明らかで明白なものではないということです。例えば、数百万の人々は個々の季節を総合的に見ると一般的に均一であるという信念を持ちつつ、各々の現象が季節と呼ばれるものの一環であるための必然的な連鎖関係の存在については信じていなかったり、あるいは明確に否定していたかもしれません。また、カードやサイコロの場合も同様で、ほとんどの賭博者は長期的には判断力と予測力が有用であることを認識しているかもしれませんが、機会論に関する著者は、個々の投げ方はその性質上根本的に予測可能であるという事実を証明するために多くの議論が必要だと考えているようです。

5. この論争が倫理や社会的な問題に適用される場合、主な興味はそれが自由意志の問題に対する実際的または仮定された影響にあると言えます。因为在这个领域中，因果律和宿命论或决定论被视为同一件事。

ここでは、最後の例と同様に、私たちは「均一性」と呼んだような広くはっきりとは定義されていないが、一定の規則性があるという事実を認めなければなりません。統計学は、これについて疑念を持つ人間が存在しないことを証明しました。

以前の章で既に得られたいくつかの概念から、この事実を証明するための証拠の本質と量についての理解を得ましたし、そのような証拠は専門家たちの著作からも補充的に提供されるかもしれません。したがって、自由意志が平均的な場合や個々の場合において予測不可能なような基本的な不規則性を含むと解釈されるとすれば、自由意志の否定は遠くない論理的結果ではなく、経験上の不可争議な事実から明らかに導かれる結論となります。

しかし、多くの人はこのように解釈することには至らないでしょう。彼らが心配するのは、この一般的な規則性が何かしらの形で原因関係をもたらす可能性があること、特に一定の順列的な連続性という意味での規則性だけでなく、おそらく強制との関連も持つかもしれません。

後者の関連性は完全に非哲学的であるとして拒否する私としては、前者の結果が本当に証明されたとは思えません。しかし、そのような結論はこの源から多少の確認を得ていることは事実です。

6. 自由意志に対する統計学的な議論の性質は、少なくともその形式が非常に一般的に表現される形では、私には極めて欠陥が多いと見受けられます。我々の意志決定における前件と後件は、因果関係を確立するためには明らかに非常に近接しているものであるべきであり、それゆえに直ちな連続性を仮定しなければなりません。しかし統計学的な調査では、データがしばしば広範囲にわたって離れている場合があり、実際には個別の行動群や結果群のみに関係する場合もあります。

たとえば、誤送信された郵便物の例で見ると、証明しようとしていることは、各宛名者がその状況のもとでは他の選択をすることができなかったほど「状況の犠牲者」であったということです。しかし実際には、年々その数が一定であるという統計データの積み重ねだけでは、この主張に大きな影響を与えることはほとんどありません。それに伴う忘却の傾向を他の特性と結びつけるようなデータがあったとしても同様です。

また自殺の件数についても同じことが言えます。ロンドンで毎年250人程度が自殺しているという事実は、単にその状況下において人々の行動が恒常的な秩序の一例であることを証明するためには重要なものかもしれません。しかし、哲学者や政治家にとって信憑性を増すためには、心理的要素との関連性を大幅に強化する必要があります。

実際には、ここでのAとBはしばしば符号記号で示される因果関係を説明するために一般的に使用される論理学の著作で用いられるシンボル表記において最も密接な順序関係にあるものであるべきです。しかし我々はしばしば大きな間隔を置いて、単独のAやBのみが存在している状況に遭遇することになります。

7. また、このような論理にはもう一つ欠点があるように思われ、それは統計の単なる規則性や持続性に過度な重きを置くことである。これらの統計は非常に重要な結果を導くことがあるが、自由意志に対する何か証明するためのものではないという目的からは正確には求められていない。実際、そのような結論と年間の窃盗件数や自殺件数がほぼ同じ水準で推移することとの関連性を見出すのは容易ではない。

統計的な一貫性は私にとって、少なくとも直接的には人間の行動の場合においても物理的特性の場合と同じように何を証明するものでもない。たとえば、未宛名で投函された手紙の数を数えることによって得られる直接的な証拠は、そのような一貫性が存在することである。

以前の観察結果はすでに示していたように、人間の体格や筋力の統計は平均値を中心に集約される傾向があることを我々は知っていた。そして、記憶力の迅速さと信頼性についても同様な状況が存在することを予測することができた。未宛名で投函された手紙の数を数えることによって得られる証拠は直接的で結論的なものではない。なぜなら、記憶力は単独で測定されず、無数の他の影響要因と組み合わさって測定されるからである。

それでも、統計データ全体としてはこの傾向に一定の影響を与えることは疑いようがない。これらの統計を適切に変形し拡大することで、相当な証明力を持つものにすることができるだろう。

確実な事実は、確率が自然神学に対して有利か不利かを決定するものではなく、論理や帰納法の一般原則と同様に、ただある種の性質を持つ事件のクラスについて推論を行うための一連のルールであるということだ。神信者であっても、自然の研究によって彼はその作品について見解を形成し、それから一定の範囲でその属性についても考え出すだろう。しかし、彼が予想していた方法とは異なる方法で事件の順序が生じたからといって、彼の信仰を放棄するべきだという提案は確かに不合理である。注意してobserve、それは一般的な幸福や不幸、善や悪への傾向ではなく、彼が予想していた方法とは異なる方法で生じたからであり、秩序が不規則性を示すか、直接的な作用が二重原因の機械装置を含むかに関係なくである。

8. これは個々の無知と集合体に対する知識という、何とも異なる感情を私たちの心に引き起こす可能性があることについて考えるのは、 amusement（興味深さ）と instructiveness（教訓）を兼ね備えたものである。

もしこれらの要素が観察される順序が逆であれば、個々の無知と集合体に対する知識という現象は、我々の心にどのような異なる感情を引き起こすかを考えてみるのも興味深いだろう。個々の生活や個人の行動について確実な予測を下すことは不可能であるにもかかわらず、大規模な数値に関しては自信を持って予測が可能であることが突然明らかになると、人々はしばしば驚きを感じ、時には甚至は混乱するかもしれない。

そして、この現象に対する反応として、人々は「これは神の意志を否定しているではないか！」、「完全なる運命論だ！」と叫ぶことがある。しかし、仮に我々が最初から集合体についての知識を得た場合を考える。個々の生活についての経験的な理解が存在しなければ、そのような予測は不可能であるという認識があったとしても、大規模な数値に関しては自信を持って予測が可能であることが明らかになる。

このような仮定を完全に実現するのは難しいかもしれないが、生命保険事務所で無知な書類係の例を考えると、彼は一度に一つの生活についての絶対的な不確定性を理解する突然の洞察により驚きと混乱を感じる可能性がある。そしてその叫びは、「これは神の意志を否定しているではないか！」、「完全なる混沌と運命だ！」となるかもしれない。

世界の創造者や管理主に対する信仰は、直近の事象の連続性に関する特定の仮定に限定されないが、これまで上記で述べた一つの側面からしか事件を見なれていた人々は、時々その他の側面と結びつける方法を見つけるのに時間がかかるかもしれない。

9. これまで、私たちは非常に一般的な問題について触れられてきました；すなわち、確率の基本的な公理が世界における秩序や均一性（物理的または道徳的なもの）に対する概念とどのように関連しているかという問題です。そのような困難は主に神学的、形而上学的、あるいは心理学的なものであることが多く、現在考慮すべきはより詳細な或いは論理的な問題が中心となります。これらは主に次の2つです；（1）物理現象における偶然の配置と因果関係のある配置との区別、および（2）理性ある代理行為による設計された配置との区別の問題です。

II. これらの質問の最初のものは、偶然と因果律の対立を引き起こしますが、これはすべての現象にわたる一般的な特性ではなく、特定の出来事に関して言及されます：この出来事は偶然なのかどうか？因果関係の普遍的な存在と秩序を支持する最も激しい擁護者たちは、その質問が関連性があることを認めています。したがって、彼らにはその答えを検討するための規則があると仮定しなければなりません。

例えば、ある男が食事のために訪れた家で発作を起こし、そこで亡くなった場合、そして誰かが「彼はその家で生まれたのだ」と言います。私たちは偶然な一致がただそれだけなのか疑問に思います。しかし、情報提供者がさらに「それは古い家族の家であり、死者の兄弟が住んでいた」と追加した場合、これらの事実は状況を全く異なるものに見せます。また、コルヌトが示唆するように、二つの兄弟が同じ日に戦場で殺害された場合でも、彼らが同じ戦闘で戦死していたか、または一方がフランス北部で他方が南部で亡くなったかによって、その評価は大きく変わります。後者の場合は単なる偶然と分類されるかもしれませんが、前者は説明の余地があるかもしれません。

10. この問題は、そのような観点からすると、帰納的論理の問題よりも確率の問題ではなく、なぜなら確率を計算しようとする試みはありません。しかし、ここに注意を払う価値があると思われます。もちろん、現代物理学的な概念に支配された正確な思考者であれば、一瞬も疑わないでしょう。二つの要素それぞれには、その「原因」が存在し、完全な知識があれば確実に推測できたと断言できるからです。また、私は彼はおそらく、十分詳細で広範囲な視点から見れば、過去を十分に遡って探れば、二つの事件の共存も同様に予測できた時期にたどり着くだろうと考えるでしょう。このようにして、「偶然」という言葉が原因を否定するものであるとは限らないにもかかわらず、それでも実用的な区別に対応していると言えます。「偶然」な一致を意味するとき、私は二つの要素のうちの一つに関する知識（実際的に得られるものと仮定）が他方を予期させる能力を持つことを期待させないという意味で使われると思います。両者をカバーする一般的化は知られていない、もちろん無効と想定されるようなものを除いてです。このような適用において、「偶然」という言葉は「因果的」や「意図的な」対義語ではなく、より広い秩序や規則性の概念、つまり「一貫性」という用語を適用すべき範疇に使われているように思われます。「偶然」な一致は特定の一般化の下には収まらないものであり、確定的、確率的、またはさえ説得力のあるものであると言えるでしょう。

この差異を別の表現で表すと、これらの「単なる偶然」は独立事象の例であると見なすことができる。ここで独立性とは以前の章で説明した意味での独立性である。私たちは、任意の2つの事象AとBが、他の事象が起こった場合や起こらなかった場合においても、その頻度が正確に同じであるとき、そう呼ばれるというのを見た。この状況は硬貨を投げて表と裏の順列、あるいは町で生まれる男児と女児の順列、また多くの数学的な表における数字の順列のように思われる。

したがって、街中で人を拾い、そのうちの一室で亡くなるという状況を考えると、長期的には生まれた場所によって選択される傾向は見られないものとする。そして、このような出来事において個別の事例について単なる偶然であると否定するとき、私たちが必ずしも主張することは、その特定の事例を一般的なリストから取り出し、あるような選択がある場合に属する一例へ移すということである。

11. III. 次の問題はやや複雑であり、そのため慎重に分類することが必要となる。これは偶然と設計の対比に関わる問題である。つまり、現在私たちは（前回の場合とは異なり）物理的な側面だけではなく、これらの事象がどのように生成されたかを考慮に入れ、その生成過程が偶然性に依存していたのか、それとも設計によって行われていたのかについて調査している。

読者には明確に理解していただきたいが、我々は現在、特定の配列が「偶然」であると呼ぶものかどうかという単なる事実問題について議論しているわけではない。第四章で完全に説明したように、これは十分な材料があれば単なる観察によって解決できる問題だ。我々が提案するのは、その配列が与えられた場合、この配列は設計によって生み出されたものよりも、一般的に「偶然」と呼ばれる方法によって生み出されたものの方がより可能性が高いかどうかを調査することである。この区別は、πの比を構成する数値の順列に戻って考えてみると明らかになるだろう。私は言ったように、この配列が単なる数字の順列として見れば、偶然の配列の特性を完全に満たしているように見える。もし最初の四つや五つの数字を除いて、それらは我々の多くが知っているものだと仮定すれば、それを袋から数字を引くだけの方法で得られたかどうかを示されても、誰もそれを否定することはできないだろう。彼は一生見てもそれが単なる偶然の選択の結果であることを発見しないだろう。そしてそれは当然のことであるが、単なる配列として考えれば、それは偶然の配列である：その配列はこのような配列のすべての要件を満たしている。我々が現在議論する問題は次の通りだ：与えられたような配列について、それがどのように生み出されたかをどのように決定するか。

我々には二つの異なる方法で生み出される可能性のある事象があると想定され、そのどちらか一方が偶然によって、もう一方が意図的な設計によって生み出されたと仮定されている。そして我々は、これらのいずれの仮説にも支持する確率を決定する必要がある。したがってこれは逆確率の問題であり、このようなクラスの問題に特有のすべての困難性に晒される可能性がある。

12. このような問題に対する理論的な解決には、次の二つのデータが必要である：--

(1) 二種類の原因の相対的な頻度、すなわち偶然的に作用するものと意図的に作用するものの頻度比。

(2) 各々の原因が実際的に作用している場合に、その事象を生み出す確率。

これらのデータの後者の部分は、一般的に何らの困難もなく得られるものである。確率仮説に基づく各種の事態を決定する作業は、例が適切なものであれば、算術的な困難以外は何らの問題も生じないであろう。そして設計の選択に関しては、通常そのことが実行可能であった場合、それが目的に達成することを確実にするものと想定される。

例えば、テーブル上に十 Pence が頭側が上向きで見つかった場合、偶然性か設計によるかについて問われた場合、我々は一定の点まで困難を感じないであろう。これらの Penny が投げられた場合、十回全てが表が出る確率は1024回中一度だけである；しかし、これらが意図的に配置された場合、その結果は確実に達成されるであろう。

しかし、他方の公理、すなわちこれらの二つの類別の作用の相対的な頻度に関するものについては、より深刻な困難が生じます。この問題を大体的に回答できる例は確かに存在しますが、それも稀です。そのような場合以外では、さらに問題はありません。たとえば、学校のクラスリストでブラウン、ジョーンズ、ロビンソン、スミスという四人の名前を見かけ、アルファベット順であるかどうか、あるいは功績順であるかどうかを調べる機会が生じることもあります。この場合、「アルファベット順」や「功績順」という用語の拡張的な意味では、これは偶然と計画性という選択肢に相当します。なぜなら、男児の名前の頭文字は彼の成績とは関連がないため、アルファベット順以外の他の順序でこれらの頭文字を並べると、ランダムな特徴が現れることになります。これは無理数の桁と同様です。四人の名前が無計画にアルファベット順になる確率は23対1であり、この順序が計画的に作られた場合、その確率はほぼ確実性に等しいものとなります。学校試験における二つの種類の順序の相対的な頻度については、統計があるとは限らず、必要であれば容易に入手できるものの、功績順を採用していることが多いと推測されます。仮説としてこの比が9対1まで高まるとするとしても、その場合でも現実の順序がアルファベット順であることがより確率的に高いことがわかります。

13. しかし、大半の場合においてはこのような統計が手元にないため、非常に深刻な曖昧性に直面することになります。これらの曖昧性は二つの異なるクラスに分けられ、その本質は具体的な例の議論を通じて最もよく理解されます。

まず最初に、前章で述べたとおり、統計が非常に信用できないという欠点に特にかかりやすいです。これは、新しいヒントや情報の獲得によってほぼ瞬時に割合が揺らぐという事実から来ています。

例えば、疾病中や戦闘の危機においては、死亡率に関する統計が非常に役立たない理由を考えてみましょう。今、部屋に入り、テーブルに十枚のコインが表向きに置かれているとします。そして、これらのコインが意図的に配置された確率は何であるかについて尋ねられます。これは、状況に関する特別な知識によって大きく変わります。

誰が部屋にいたのか？子供たちなのか、コインを集める人達なのか、あるいは投げ入れ遊びや賭博のためにコインを振る人がいたと予想される人達なのか？これらのコインは新しいものなのか古いものなのか？このような区別は、同じように表向きに配置された理由があるかどうかを考える際に非常に重要です。

同様の例で考えてみましょう。統計が瞬時に得られる情報によって左右されることを理解しています。

しかし、これ以外にも考慮すべき点があります。我々は一般的な外部状況、つまり代理者の性格や立場だけでなく、当時の特定の確率配置に対する伝統的な評価も影響を受けるべきです。例えば、新しいカードゲームが人気になると、新たな組み合わせが意味を持つようになります。したがって、意図的な欺瞞についての質問が可能性のある詐欺という形に変化すると、そのような組み合わせに対する現在の評価に関する知識は非常に重要となります。

これらの困難の完全な意味を理解するためには、実際の争議を引き起こした例を通じて議論することが最も適切です。数年前に、ピアージィ・スミス教授は、ギザの大ピラミッドに関する著作[4] を発表しました。その著作の主な目的は、その建築物が大きさ、比率、内容において、長さや容積などの几乎不朽の自然基準を含んでいることを示すことです。その中で、円周率πの値が、側面と高さの比によって正確に（その正確度は私には明記されていない）示されていることが決定されました。主張は、この結果が偶然ではなく、設計されたものであるという点でした。

確率仮説の評価については、サイコロやカードの場合ほど明確ではありません。実際の長さの底辺に対して、どのような高さでも等しく可能なわけではありません。ここで一定の範囲を限定する必要があります；高すぎると建築は不安定になり、低すぎると滑稽になります。また、測定がどの程度の近似度で行われたかを決定する必要があります。寸法が1インチまで保証されている場合と、100分の1インチまで保証されている場合では、一致は全く異なるレベルとなります。この点が既に決まったとして、与えられた底辺に対するピラミッドの1万個の可能な高さの中から、円の半径と周長の比を最も近いものとする高さが選ばれたことが確認されました。

残る問題は、「設計」の選択肢がどれほど頻繁に用いられたか、すなわちその「先定確率」と呼ばれるもの、つまりその比率を達成するためにbuildersがどれほど頻繁に目指したかについてである。これは、buildersがその比率を達成するためには必ずそれを達成することを前提とした、buildersの努力の相対的な頻度を指す。buildersの技術に関する我々の極めて限定的な知識からすれば、ここで数値評価を行うことは不可能であることが明白である。

もし「設計」が意識的な意図による比率の生成という意味ではなく、単なる偶然にその比率を生み出す方法を使用することだけを意味すると解釈されるなら、多くの人は大きな疑問を感じることなく同意するだろう。buildersはπの値を知らなかったと我々は十分に確信しており、それはすべての技術者がそれを粗略に理解しているような程度である。そしてbuildersがその値を知っていたとしても、それが理由となる合理的な動機を見出すことは不可能である。しかし、巧みな提案[5]に基づきbuildersが次の方法で作業を行ったと仮定する場合、状況は異なるものになる。

buildersがピラミッドの高さを決定した後、その半径として円を描き、その円周上の線に沿ってひもを引き、それを正方形に延ばし、その正方形が建築物の基礎となるようにしたと仮定する。比較的原始的な時代においては、これほど単純な方法は考えられないだろう。そして、円周の長さがひもの長さと等しいことから、基礎の周りの長さは高さに対する適切な比率を保有することになる。つまり、幾何学的な値の正確な達成はその比率を知っていることを意味するわけではないが、それを含み且つ示す方法を知っていることを意味するに過ぎない。teredoも我々と同じように円の幾何学的性質を示す穴を掘ることができるが、それには対応する知識があるとは信用しない。

先述した通り、設計の選択肢に関する数値評価は不可能である。しかし、もし精度がMr. Smythが主張したほどであれば（上記の提案を考慮に入れた場合）、多くの人々は偶然の一貫性がなかったと推測する可能性が高いだろう。

16. なお、重要な推測的な考慮がまだ残されています。上記の議論では、10,000 の可能な値の中でただ一つが有利であると仮定し、その確率的選択を計算しました。つまり、π の比が他のすべての比よりも優遇されていると仮定したのです。しかし明らかに、もしこの比が二倍になった場合でも結果は類似した意味を持ちます。なぜなら、それは円周に対する直径の比と同じになるからです。実際、スミス氏が選んだ比——底面の高さ対底辺の2倍と直径対円周の比——は明らかに多数ある比の一つに過ぎません。

また、測定結果が底面の一辺に対する高さの比が 1 : √2（すなわち正方形の一辺に対する対角線の比）または 1 : √3（すなわち立方体の一辺に対する対角線の比）であることを示した場合、このような結果も設計の証拠ではないでしょうか。この方法で進むと、我々は次々と知られた数学的比を提案し、ほとんどすべての10,000 の可能性のある値が考慮されるでしょう。

そして、我々は以下のように議論することができます：ほぼ全てのピラミッドの高さに対応する数学的な比があることを考えると、全然それらを知らなかった建設者がその中の一つに偶然遭遇することは決して稀ではないでしょう。なぜなら、ある一例に対して設計を属性付けるが、他の例にはそうしない理由は見つからないからです。

17. この反論に対する答えは既に触れられている。すべては、一つの結果と別の結果を比較したときの、\_慣用的な評価\_ にかかっている。単純化のためにコインに戻ろう。10回連続で表が出る確率は、交互に出る表と裏、または5回表を5回裏が続く場合などと全く同じである。実際には、残りの1023通りの可能性のどれかと同じである。しかし、普遍的な慣習によって、10回連続で表が出ることが特異であると選ばれた。もちろん、この慣習は非常に自然であり、そしておそらく避けられないものだが、他の場合では全く任意的である。たとえば、カードにおいて「スペードのクイーンとダイヤのジャック」は、どのような他の組み合わせよりも希少である：さらに、ベジーが導入される以前には、この組み合わせもまた他の指定された組み合わせほど興味深いものとはみなされなかった。しかし、そのゲームが非常に人気があった時期には、この組み合わせが偶然性のカテゴリーに分類され、プレイヤー間での不正がある場合、そのデザインされる確率は大幅に改善されることとなった。

それではピラミッドに戻って、運と設計を均衡させるためには、公平性のためにπやeだけでなく、sqrt{2}やsqrt{3}などいくつかのより単純で馴染みのある比も含めて計算する必要がある。しかし、実際に我々が知っているこれらの値を考慮に入れることが可能である数は無限でも、ピラミッドの建設者たちがそれらを知っていた可能性のある数は非常に少ない。設計の存在に対する確率の順序は、このような考慮によって大幅に変更されることはない。

[6] ここで、ゲームベジーが導入される以前には「スペードのクイーンとダイヤのジャック」の組み合わせも他の指定された組み合わせほど興味深いものとはみなされなかったことを強調する。しかし、そのゲームが非常に人気があった時期には、この組み合わせが偶然性のカテゴリーに分類され、プレイヤー間での不正がある場合、そのデザインされる確率は大幅に改善されることとなった。

[7] ここで、実際に我々が知っているこれらの値を考慮に入れることが可能である数は無限でも、ピラミッドの建設者たちがそれらを知っていた可能性のある数は非常に少ない。設計の存在に対する確率の順序は、このような考慮によって大幅に変更されることはない。

18．これまで述べてきたところでは、人間の意志が互いに拮抗している形でバランスを取っているのは、偶然による意志と直接的な設計による意志の二つの形式であることが観察される。つまり、人間の意志について話しているのである。この場合、我々はその位置を明確に理解できるため、ここでは我々がどこにいるかが判明している。実際、このような意志に関する問題は数値的に解決可能なものであることは稀で、ほとんどの場合全く解けないものであるが、少なくとも明確な形で述べることができる。我々は期待される答えの種類と、その答えを決定するために利用できる理由を知っている。

調査の次の段階では、通常の人間による偶然の意志と、他のすべての可能性のある原因とのバランスを取ることになる。--直接的な霊媒師の意志という言葉を使わないことに注意して、それは仮説を無駄に狭めるものではないからである。--しかし、それら全てに対するバランスを取ることになる。

心理的研究会によるいくつかの調査は、この声明が意図するものを優れた例示として提供してくれるだろう。F.Y.エドゲワース氏の最近の論文において[8]、これらの適用について完全な議論がある。しかし、彼の説明は他の計算や図表と関連しているため、私はその一部を引用するだけである。しかし、私は彼と本質的に同意している。

「1833回のカードの種類に関する予測が行われたという記録があり、その中で成功した予測が510回あった。これは1833の四分の一である458よりも著しく多く、純粋な偶然に基づく仮説においてもっとも頻繁に見られる数値よりも大きい。現在、誤差の法則により、このような過剰発生が偶然で起こる確率を大まかに決定することができる。その確率は確率曲線の尾の端点に等しい、我々が既に検討したようなものである[... ]。この尾の端点に対する全体の比は1に対して0.003である。つまり、その端点に一発当たる確率はその分数であり、純粋な偶然に基づいて成功した予測の数が510以上またはそれ以上の数が現れる確率は、約332対1のオッズである」ということになる。

19. Edgeworth氏は、私はそのように強く信じている通り、計算の目的において、厳密な意味での任意の文脈において、仮説の非偶然側面におけるデータの決定が必要であると主張しています。我々にはその相対的な頻度が何であるかを知る必要がありますし、それが目標に達する相対的な頻度も知る必要があります。また、私は彼と同じように、「その他の原因とは何か、—それはあるトリックや無意識の幻覚、または調査家によって肯定されたような思考移転のようなものであるかどうか、—これは常識と普通の論理が考慮すべきことです」という主張にも賛成です。

したがって、我々は定量的な見解を形成することはできないかもしれませんが、特定の場合には相当確固たる見解を形成できるという考えに同意します。もちろん、偶然説の支持者に最終的な決定権を与える場合でも、証明に達することは不可能であり、彼らは常に反対仮説の先ずじまいの確率を見直し、再設定する余地が残ります。我々が公正に要求できるのは、偶然説を否定する人々が他の仮説における発生確率に最低限の値を指定することであり、それに対しては実際に生じた現象の稀少性を増加させてこれを上回る試みを行うことができます。例えば、彼らが他の任意の原因よりも働く確率が332対1より低いと主張する場合、我々はその問題となる現象よりもさらに不頻繁な発生を確保しようと努力しなければなりません。思考移転の支持者が自分の信念に勇気を持つのであれば—彼らはそれを絶対的に持っています—このテストを受け入れることはありません。私は現在でも、上記のような証拠に基づいて、結果が通常の推測によって得られた確率が非常に小さいと考えています。

20. 前節で議論された問題は、少なくとも理解できるものであり、常に解決可能なものとは限らない。しかし、この章を終える前に、主題のこの部分に関するいくつかの推測について注意を払わなければならない。これらの推測は、理解可能な範囲内に完全には留まっていないように思われる。

たとえば、Arbuthnott（『Phil. Transactions』第27巻に掲載された論文において）で議論された問題を考えよう。「性別の出生における恒常的な規則性から神の PROVIDENCE に関する証拠としての論理的議論」のタイトルで扱われたものである。彼の論理が通常の目的論的なものであれば、すなわち、存在する約等比関係（男性が6%多い）は有益なものであるという主張であれば、ここには反論の余地はなかった。しかし、彼が考慮したのは、偶然と設計の交替仮説間の均衡のようなものであり、我々がここで考慮しているようなものであった。彼自身の言葉で結論を述べると、「これは芸術によるものである、運命ではない」ということである。

このような議論を正確に表現しようとすると、その ridicule すら招くようなものになってしまいます。厳密に解釈すれば、この議論は確実に二つの解釈のいずれかに限定されるでしょう。

まず一方で、運命（Chance）を人間化して考えることができます。つまり、運命が人間を生み出す能力があると見做すか、あるいは少なくとも性比のバランスを整える役割を持っていると見做すことができます。その場合、この運命と創造主との間で、どちらが現存する状況を生み出したのかという決定が必要となります。もし運命が任意またはランダムな配置を生む能力を持つ任意の代理体として定義されるなら、恐らくこの代理体が当該事例で働いていたと言えるでしょう。

男性と女性の出産の配置は、我々が見ることができる限り、最も完璧な運命の例です。個々の不規則性から最終的な均一性が生じ、各選択肢の「ラン」や順列が適切に表現されています。例えば、すでに五人の息子が生まれた家族において、次の子供が息子である確率が減少するという特定の効果は見られません。このようなほぼ完璧な独立事象の例は物理現象の中で比較的稀です。これは運命から得られる唯一の主張です。

私はもう一つの解釈を提案できる唯一の方法として、ある単一の代理体が、私たち自身の一人でも投げかけたり計画したりする可能性があると見做すことができます。そして、その代理体が当該事例でどの道を選んだのかを特定しなければならないのです。

この場合も、我々が自分の作品について適用するテストに基づいてその行動模式を判断すれば、彼はおそらく投げかけのスケジュールを採用したように見えてきます。

21. つまり、理性的な議論において、運命と設計のどちらが原因となったかを決定することは、有限な知性のケースに限定されるべきです。ここで重要な決定要素の一つとして、代理体の知識状態や特定の配置に対する伝統的な評価があることを再確認します。これらは我々自身のような存在について扱うときのみ理解できます。

例えば、星の配置についての議論に戻りますが、ミチell（議論を提起した人物）が視野に置いていたのは、偶然と設計の間の決定という点ではほとんど疑いの余地がないでしょう。彼は（\_Trans. Roy. Soc.\_ 1767）で言っています。「私が用いる主張は…、一般的な類推と現在の状況における事象が特定の原因によるものである確率の大きさから、設計的存在または一般法則を導くような種類の議論です。」

そして彼は結論します：「星が単なる偶然によって散乱された場合」には、「約50万対1の確率で、その数（1500）の星のうち6つ以上が、 Pleiades ほどの距離に近接している」ということがあり得ません。このような解釈の下では、この論争は私にとって無意味なもののように思われます。私は一時的にも、一般的な目的論的議論には一定の力があることを否定していません。これは、万物の全体的な傾向から善と知恵の兆候を追跡しようとするものであります。

しかし、創造の性質や創造主の計画について我々が何を理解しているのか、そのような知識が星を特定の形態に配置する確率を決定するのにどれほど役立つのか、または「単なる偶然」という概念が彼の推定される全般的な作用と比較してどのような意味を持つのか、これらについては我々は何を知っているのでしょうか？

22. 明瞭な言葉で表すと、この論争から二つの質問が浮かび上がります。

(I.) 星が空間に分布されている場合、当然、我々の惑星から見るといくつかはほぼ一直線上に並んでいるように見えます。彼らが比較的均一に分布していると仮定すると、そのような観測でどれくらいの数の星が近接して見えるかを計算することができます。その後の質問は、このような計算によって説明できるよりも、それらがより多く近接して存在するかどうかです。私は見解しますが、この唯一の直接的な推論は、彼らは均一に分布していないこと、つまりペアで傾向があるということです。しかし、これは理論を完全に正当化し、合理的に適用したものです。これらの星のペアが自己中心的な系として形成され、互いに回転し、実用的には他の恒星系から影響を受けないというようなさらに具体的な結論は、当然ながら天文学的考慮に基づいています。確率は、分布が均一でないだけであると単純に答えることができますが、これらの二重星系がどのように物理的な過程によって「引き起こされた」かについては、断言することはできません。

23. (II.) 第二の問題は、双星を除いた場合に恒星の分布が、人間の意図なく「乱雑に」投げ入れられたものと類似しているかどうかである。もちろん、ここでの言及は私たちが見える天界における恒星の分布についてであり、これ以外には観測できる範囲がほとんどである。しかし、それらが望遠鏡の視野を超えてどの方向にも延びている場合、これは実際の空間配置を考慮した場合とほぼ同じ議論になるだろう。

以前の章で充分に述べたように、「乱雑性」の意味について説明している。その定義を現在の状況に適用すると、問題は次のようになる。全体としては比較的均一であるが、無数の個々の偏りがあるかどうかである。つまり、大きな領域と小さな領域を比較するとき、各等面積内の恒星の数の比がほぼ等しいか否か、そしてさらに小さな領域に焦点を当てるにつれて、相対的な数がより不規則になるか否かである。 Milky Wayや他の星雲集団のような一部の例外を除いて、これは大抵の場合そうであるように思われる。少なくとも大部分の恒星に関しては、このような状況がほぼ当てはまる。

さらに問われる問題：ネブラスヒポテシスの何らかの形に対する判断、例えば支持するか否か；またはこの仮説を認めるとする場合に、特定の部分が同じ星雲から生まれたかどうかの判断などは、天文学が裁くべきである。

性比に関する附註

次の記述は259ページに適切に入れることができず、ここに追加した。

男児と女児の出産が「ランダム」であるという特徴は、一般的には場所や時間における統計に依拠してきました。しかし、もっと望ましいのは、多くの家族を比較したときのその比率を見ることです。これは明らかに自明なことなので、私はどこかでこの調査が行われていると考えざるを得ませんが、最も可能性が高い場所を探しても痕跡を見つけることができませんでした。例えば、Prof. Lexis（『Massenerscheinungen』）は、性別比の比率が自然現象の中で真の平均周りの正常分散の唯一の例であるという彼の見解を支持するために、各国の登録簿の通常の統計データを利用しています。

同じ特性が家族単位で同様に成り立つことを証明する必要があることは明らかです。特に、Sadlerの理論などは、「連続」の男児や女児の組み合わせが純粋な偶然よりも頻繁に起こると推測しています。Lexisは、この点で双子の場合が最も顕著であることを示しており、つまり、（Mを男性、Fを女性とする）M.M.とF.F.の組み合わせがM.F.よりも著しく多いということが示されています。

私は1万3千を超える男児と女児の出産データを集めました。これらは四人以上の家族で構成され、ヘラルドの訪問記録から得られました。したがって、通常、十六世紀の英国地方貴族の長男の家族を代表しています。これらの統計はまだ十分に広範ではなく、出版には至っていないものの、その傾向を示す要約を提供することで、それらの結果の概観を示します。各ケースにおける上部の数字は「観測された結果」を表し、例えば四人組の場合では、四人の男児、三人の男児と一人の女児、二人の男児と二人の女児など、それぞれの数を示しています。下部の数字は「計算された結果」を表し、つまり、観測された数値に基づいてMとFが混合されている袋から組み合わせを抽出した場合に得られる相当する数を示します。

512 家族のうち、4人の場合；

| $m^4$   $m^3 f$    $m^2 f^2$    $mf^3$    $f^4$
結果として | 81 + 148 + 161 + 98 + 24 (観測値)
男性 : 女性 | 57 + 168 + 184 + 88 + 15 (計算値)

512 家族のうち、5人の場合；

| $m^5$   $m^4 f$    $m^3 f^2$    $m^2 f^3$    $mf^4$    $f^5$
結果として | 50 + 82 + 161 + 143 + 61 + 15 (観測値)
男性 : 女性 | 25 + 103 + 172 + 143 + 59 + 10 (計算値)

512 家族のうち、6人の場合；

| $m^6$   $m^5 f$    $m^4 f^2$    $m^3 f^3$    $m^2 f^4$    $mf^5$    $f^6$
結果として | 30 + 48 + 115 + 146 + 126 + 40 + 7 (観測値)
男性 : 女性 | 10 + 56 + 133 + 159 + 108 + 41 + 5 (計算値)

より大きな家族の数はまだ少ないため、その結果を提供する価値がまだないが、それらも同じ傾向を示している。すべての場合で、観測された中央値は計算値よりも小さいことが確認される。また、観測された極端な値は計算値よりも著しく大きい。これらの結果は、家族が一袋から必要な数だけの確率的な選出と類比できないことを示唆しているように思われる。より良い類比としては、一つには男性が多い袋があり、もう一つには女性が少ない袋があると考えることである。そして一部の人々は一つの袋から、他の人々は別の袋から選出を行うというような状況を想像することができる。しかし、より詳細な統計が必要である。

男性出生数の総合的な過剰数は大きい。これは女性の過欠記入が不十分である可能性があるが、私は各系譜の最後の二つや三つの世代に限定して、より確実性を確保するために注意深く取り組んでいる。

1. \_Essay on Probabilities\_, p. 114.

2. この場合の桁数が真にランダムであるかについて懐疑的な見解が存在する（De Morgan、『パラドックス予算』、第291ページ参照）。

Jevonsはさらに、「πを多くの桁まで表すと、7という数字が他のどの数字よりもずっと頻繁に現れるべきではないでしょうか？」（『科学原理』、第529ページ参照）と尋ねた。この質問の意味は完全には理解できていない。もしサイコロ投げのような場合で不通常な確率的な差異が見られた場合に同じような質問がされたら、おそらく次のように言い換えるべきではないだろうか：事件の度合いを特定し、その統計的希少性を算出する。そして、その度合いから疑念を持つ根拠があるかどうかを判断する。

計算は単純である。7という数字が708桁中53回現れた一方で、公平な平均値は71回である。問題は、708回の試行においてこのような平均からの逸脱の確率は何であるか？通常の計算方法（Galloway『確率』参照）によれば、特定の数字に対する18の過不足の確率は約44対1である。しかし、私たちが決定したいのは、そのような離差を示す一つの数字が十個の中からどれか一つである確率である。これを（44/45）10で近似して評価し、約0.8と推定する。これは、そのような出来事が起こる確率が約4対1であり、特に驚くべきことではないことを示している。実際、Shanksによって同じ長さまで計算された他の二つの数値（tan-1 1/5およびtan-1 1/239）にもいくつかの数字は同様の離差を示している（『ロイヤル・ソシETY会報』第21巻、第319ページ参照）。

ここでは、ランダム性を決定する際に単なる視覚検討だけに頼るのには注意が必要であるという点について触れることにする。第2章でランダム性について述べた際にはこの点が指摘されるべきだった。

我々は、ランダム性を単なる目視による判断だけで決定するときには慎重であるべきだ。ここでは、πの桁と非常に長い周期を持つ循環小数の「周期」内の桁を比較することによって、その教訓がよく示される。1 ÷ 7699 の計算結果は、完全な周期（7698桁）をもたらすもので、数年前にケンブリッジ大学の2人の卒業生（ムーニンとスフィールド君）によって計算され、私的印刷された。我々がその順列の一部だけを検討する場合、ランダム性は説得力を持つように見える；すなわち、桁やそれらの様々な組み合わせはほぼ等しい数で出る。例えば、10桁ずつのグループを取り扱うと、平均値は45付近を漂っており、非常にうまく一致しているように見える。しかし、その「周期」全体を検討する場合、これらの特性が過度に強調され、ランダム性は消失する。つまり、等しさへの単なる最終的な近似ではなく、可能な限り絶対的な等しさが達成されることになる。

もちろん、この慣習的評価は、どんな順列や秩序にも根本的に異なるものではない。連続して10回の表が出る状況は、交互に出る表と裏の状況、あるいは7回の表と3回の裏の状況などから客観的には区別できない。その違いは、それが非常に一般的に注目されるという点にあるだけである。

3. 当然、この慣習的な評価は、どんな順列や秩序にも根本的に異なるものではない。連続して10回表が出る状況は、交互に出る表と裏の状況、あるいは7回表と3回裏の状況などから客観的には区別できない。その違いは、それが非常に一般的に注目されるという点にあるだけである。

4. 大ピラミッドからの遺産 第3版 1877年

5. Nature（1878年1月24日）でジェイコブ・G・ジャクソン君によって作成された。スミス氏の代替提案は、その主張を導くものであることに注意しなければならない：「大ピラミッドの垂直高さは、実際的な正方形底面の4つの直線側の長さの合計に等しい曲線周長を持つ理論上の円の半径である。」ここで、偶然と設計との選択肢について公正な評価のために、スミス氏の主張には、人間のものではなく神の設計に対する偶然の反対が認められていることに注意しなければならない。

6. コルヌ著『知識の基礎に関する試論』第1巻 71ページを参照のこと。

7. このような考慮は、当然ながら数値評価の試みはありませんが、裁判所にも導入されていることに注意を払うべきである。例えば、有名なデロス公の判決では、証拠が間接的であった範囲において、主な疑惑の根拠の一つは、デロス公がワイスでカードを得た回数が運に期待されるよりも著しく多かったこと、そしてその結果、数年連続して平均的な利益が通常より非常に大きかったことである。被告弁護士は、他のプレイヤーも疑惑なく大きな利益を確保していた例があり、そのため実際の利益に対する相当な余裕が必要であると主張した。（具体的には、これらの事例がどれだけ広範囲にわたる経験から得られたものか、また賭け金の大きさではなく成功回数によって実際の利益の大きさが説明できるかどうかについては、詳細は不明である）（『タイムズ』1837年2月11日の記事を参照）。

8. \_メトレティケ\_. この巻末には、同一著者の関連する他の出版物の一覧が見つかる。これは有用なリストである。

9. つまり、 Arbuthnott のように単に統計結果を見る場合、またはコインの投げ方を調査する場合、そうすることになる。Dr. ドゥシングの驚くべき理論（『性比の調整...』ジェナ, 1884年）が確認されれば、状況は若干異なるものとなるだろう。彼は、生理学的根拠と、男性および動物に関する統計分析を用いて、明確な補償的な過程が働いていることを示そうとしている。すなわち、どちらかの性別が優勢になると、直ちに均衡を修復するための機関が動き出すというのだ。これは、親の相対的な年齢が子孫の性別に関与しているという古い理論の修正と改善である。

クエテレ（『手紙』、第61ページ）は、黒球と白球の入ったurnを用いたある実験によって、男児と女児の出生順序に関する命題を証明しようと試みました。しかし、これは行き過ぎています。（この章末尾の注を参照してください）

10. これはダシスの花（植物とは区別して、第109ページ参照）がランダムに分布しているわけではなく、2つ以上のグループで集まるとする結論と同様です。単なる観察からこれが明らかであり、そして植物の成長に関する知識から、これらの小さなグループは同じ根から生えていることが推測できます。

11. この議論において、著者はしばしば「物理的なつながり」の確率について話し合います。この表現は誤解を招く可能性があります。通常の重力誘引の仮説によれば、すべての星は重力によって物理的につながっています。したがって、上記のように、相対的な近接性のみに関わる問題とし、不思議な近接性から何が推測できるかを天文学に任せるのが良いでしょう。

12. フォーブス教授は、すでに参照した『哲学雑誌』の論文（第7章 §18）において、ランダム分布における実際の配置を示すいくつかの図表を提示しました。彼はチェスボードにピースを散らし、各マス目の上に休むピースの数を数えました。彼の数字は、星の一般的な配置がこのような散布方法によって生じたものと非常に似ていることを示しているように見えます。

ある最近の調査により、Mr R. A. Proctor氏は、この比較的均一な分布に少なくとも二つの例外があることを示唆しているように見えます。（1）彼は、銀河帯において星が他の場所よりも明確に密集していることを確認しました。これに基づいて言えば、この論理は双星の場合と同様で、銀河帯の星々の近接性は単なる表層的な現象ではなく実際的なものであることを証明する傾向があります。（2）彼はまた、北半球と南半球にそれぞれ大きな領域が存在し、これらの領域において星が他の場所よりも明確に密集していることを確認しました。これらについては、私には確率論は何も示していないように思われます：単に均一性の否定に過ぎません。恒星がどのように配置されたかにつながる物理的な因果関係に関する推論は、天文学者が行うべき質問です。詳細については、Proctor氏の『天文学のエッセイ』第297ページをご覧ください。また、『宇宙と近接 Transit』の一連のエッセイも参照してください。

\newpage

\chapter{実質論理学と形式論理学（Material and Formal Logic）}
\label{ch:xi}

\begin{center}
\large\textit{客観的頻度観と概念論的観点の対比・仮説検証における言語の限界}
\end{center}
\vspace{4mm}

第十一章

論理の客観的扱いに関する特定の結果

1. 前回の版では、この章の大半は論理における実質主義と概念主義的視点の区別についての一般的な考察に費やされていました。ここではその大部分を省略し、また因果法則に関する詳細な議論が含まれる他の一章の大半も省略しました。私はこれらの問題についてより詳しく、そしてより適切に一般論理の原理に関する著作の中で意見を表明したいと望んでいます。

1. 論理学の学生は、実質主義と概念主義的扱いという名前が付けられた二つの方法論の間の広範な区別について馴染みがあります。この区別は、論理自体の領域内で強調され、扱われる以前には他の名称で徐々に成長していました。ミルの有名な著作によって発表された後、その区別が論理の領域内での区別として扱われるようになりました。例えば、ウィーヘルの帰納法に関する論文やハーシェルの講演を読む者は、彼らがミルが「論理」という名前で議論した同じ主題について、そして同様の方法で取り組んでいることを理解できます。しかし、これらの人々はそれを「論理」と呼ぶことに抵抗していました。

これらの著者は、常にその任務は自然の事実を客観的な形で系統化し、最も広範な扱いのもとでそれらを整理し、推論の主要な形式とこれらの形式の調査における主な実用的助具について理由から提示され経験によって正当化される主要な推論方法を説明することだと前提していました。ミルは、彼が保持できる限り古いシュオラスト学派の論理の一部をこれらの方法と密接に関連させ、それらをより詳細に展開し、哲学的かつ心理学的な基礎を提供して整理し、その結果を「論理」と名付けました。

科学の実践的な取り扱いは、その支持者自身が目指す理想とよく一致することはないし、なおさら反対派が要求するものとはずっと少ない。もしハミルトンとミールによってそれぞれ説明された論理の視点についての差異を、それぞれの著者が主張するように記述しようとすると、その定義を一つにまとめることは決して容易ではないだろう。

例えば、この二つの視点をハミルトンとミールからそれぞれ引用し、それらを比較した場合でも、実質論理は物理科学全体の中から多少恣意的な選択肢であると見なされ、概念主義的論理も心理学の領域から多少恣意的な選択肢であるとみなされるだろう。前者では思考法に関する法則の考慮が省略され、後者では結論の真偽についての考慮が省略されることになる。

もちろん、実践的にはこのような極端な状況はすぐに避けられるものであり、いかなる論争も過剰なものであっても、対立する視点を説明する者は共通の研究領域を維持することに成功している。私はここでは、この調整によって生じる制約や、いかなる一致にもかかわらずなお存続する実質的な重要な方法論、目的、検証法、および限界について詳細に調査するつもりはない。

試みることで、一つの推論に関する一般科学の一分野だけを扱う論文としては幅広すぎるような議論を開くことになるからである。

2. この調査の一分野は、実質的に私たちの注意を惹きつけるものである。論理学の領域に対する過度に物質的な説明が課される制約について考慮する必要性が生じるのは、この点が重要になるからであり、これにより我々の注意は、帰納法の範疇において僅少で一時的なものとしても重要な側面が、確率の範疇では非常に目立つものとなり、比較的永久的なものとなるという事実に引き寄せられる。過度に物質的な見解によると、帰納論理は一般的には、主として客観的事実に関すること以外は何も扱わないものとされるべきである。その役割は、事実から出発し、事実をただ提供するような方法だけを使用することである。疑わしいものは確立されたり、あるいは現在のところ放っておくことにより、達成不可能なものについては拒否し、このようにして確かな知識の巨大な構造を徐々に築いていく。

もちろん、これらすべては人間の心で行われるものであり、したがって、このような計画において概念や主観的な考え方が存在しないかどうかを調査する場合、これらの要素が何かしらの役割を持っていることは避けられない。我々の出発点となる事実は、推理を構築する前に知性を持つ生物によって捉えられるべきであり、結論として形成される「事実」はしばしば、少なくとも一時的に、人間の心のどこか以外には存在しない。しかし、ミルの著作を読む者は、彼がこの問題のこの側面についてほとんど言及していないことに気づくだろう。彼は、嘲りのような無関心で述べているように、「掘る男は当然ながら地面と鏟の概念を持っているべきである」ということだけを述べているが、これらの「概念」は、推論に限ってみれば、農夫よりも思弁的な論理学者にとってそれほど注意を払う必要がないと考えられている。

3. ここでは、推論の主観的な側面を一切参照しないという欠落は、誘導的論理学においては一定の根拠があると認めなければなりません。しかし、推論の発見者自身はその立場が全く異なるのです。彼が名にふさわしいのであれば、彼の心の中には常に事実と呼べるような概念ではなく、それらを事実にするために努力しているような考えが溢れています。一方で、彼の後を追って誘導者の発見を解明し正当化する論理学者は、この精神的な段階や不確定な段階をしばしば忘れてしまう傾向があります。彼が主に扱うのは「完全な帰納」や「根拠のある一般化」など、あるいはそれらと矛盾する誤りである「破産した誤謬」などです。これらは実際の発見者が戦場を去った後、残していった「囚人」と「死体」に過ぎません。科学全体の方法論は、事実と非事実との区別を強調し、それ以外のほとんど何も扱いません。言い換えれば、誘導的論理学に関する著作は、概念や概念自体の意味を定義する必要がなくても書けるのです。

4. しかし、より精密に観察すると、すらすらと誘導的論理学の範囲内でも、この問題に対する明確な区別が崩れることを見つけることができます。たとえば、概念や概念自体について定義する必要がないか、それを使うことがないというような場面でさえ、事実または概念に属さないような存在（最も広い用語を使用して）に遭遇することがあります。

例えば、誘導的論理学はしばしば仮説の使用を余儀なくされます。これらの仮説は上記のどちらのクラスに分類されるべきでしょうか？厳密にはどちらにも属さないようですし、また今後指摘するように、これ以外にもそのような例は数多く存在します。

確実に言えることは、推引論理の領域においては、これらの仮説がこの点で大きな問題を引き起こすことは少ないということである。哲学的な思考によって最初に現れる形はいかに不確定なものであっても、我々が論理学者としてそれらの仮説が事実となる道程を進むまでにはあまり関与しない。つまり、その形が最終的に固まり、具体的な姿になるまでは、まだ形成途中であるという意味で、我々はそれらの仮説にあまり注意を払わない。

我々はしばしば、仮説が十分に根拠があり合理的であることを推奨される。これはただ、哲学者が自分の研究室の中でいかに自由に予想しても、その予想を公に発表する前に、それらが適切な事実と呼ばれるまでには至るべきであるということの別の言い換えである。すなわち、「確かな事実」と呼べる程度には成熟しているまでにはならないと公にするべきである。これは「確かな事実」ではなく、「確かな可能性のある事実」というような限定的な言葉で呼ぶこともある。

したがって、我々は推引過程において仮説のこの中間状態にあまり注意を払わない理由は、ここではその役割は一時的であるからである。その姿勢は非常に一時的なものであり、仮説が健全であればそれはすぐに事実として認められる場所に立つだろうし、そうでなければ完全に否認されるだろう。仮説の状態は通常のものではなく、そのため我々はその特性を詳細に検討する必要がない。

ここで言及しているのは当然、推引論理学者としての視点である。既に述べたように、哲学者は彼の研究室で仮説が存在するすべての段階から起源まで十分に馴染みがあるべきである。しかし、論理学者の役割は異なるため、彼は証明に関与し、元々の調査や発見の過程とは異なる方法で作業を行う必要がある。

実際には、さらに進むこともできるかもしれません。多くの場合、仮説は読者、すなわち知識の受容者に提示されるまで、その名にふさわしくない状態で存在していることがあります。彼がそれを以前から考えたことなく、証明によってそれを確立する過程と共に初めて提案されるかもしれません。これにより、仮説は突然、未知の暗黙のうちに完全な権利を持つ事実として現れ、実質的に中間的な仮定段階をスキップして存在します。元々の研究者は長い間それを考慮し、その不確定な状態で常に頭に保ち続けていたかもしれませんが、最終的には、他の人々が共通認識とする事実のレベルまで即座に引き上げるような証拠と共にそれを世界に提示したでしょう。

この仮説が確実でない段階は、それぞれの仮説において存在するが、論理的な目的や多くの人々にとって、その存在は非常に一時的である。注意が向けば、論理学の他の場面でも同様の段階を検出できる。

たとえば名詞の参照について考えてみよう。ミルは太陽や戦争という具体的な対象と、それらに対する我々や子供たちの概念との区別を例示している。この区別の明確さは明白である。しかし、存在が疑問視されるか否か、または議論の余地のあるような対象の場合、先に述べた困難が始めて現れる。

ただしどのようにして絶滅した物体やまだ存在していない物体について考えると、問題は生じない。実際的な現在の存在は当然のことながら事実を構成するためには必要ではないからである。また、伝説上の存在についても通常の区別が維持できる場合がある。セントアーロスとグリフィンは詩人、画家、紋章学者にとって lion および leopard と同じように広く認識されている意味を持つ。したがって、ここでも「真実的」である、「事実に合致する」という主張を提出できる。つまり、これらはそのような主題に関する全般的な合意と一致しているべきであるという意味である。

必要かつ普遍的な一致がしばしば「客観性」の意味だと主張される場合もある。そして、有名な虚構物について全般的な一致が達成可能であることを考慮すれば、事実と概念との基本的な区別、および当該語彙が客観的なものに参照すべきであるという決定は、一定の歪みを伴ってもここでも正当性を持つことができる。

5. しかし、争議のある現象の場合に到達すると、困難は再び現れます。仮定された惑星や新規発見された鉱物、歴史的な疑問の事実、論理学上の争点など、これらは多くの例の中の一端に過ぎません。ある人々が熱心に主張するものに対して、他の人々は否認し、物理現象の場合には迅速な合意の希望があるかもしれませんが、経験からすれば、倫理的・歴史的な問題や神学上の問題についてはそうはいかないことがわかります。

一致している人々が自らの間でしか交渉を行わない限り、「事実」はそのような形で受け入れられますが、他の人々と通信を開始すると、即座に「事実」と「概念」の区別が失われ、一方の「事実」は反対者にとっては根拠のない「概念」に過ぎません。したがって、私はこれらの問題において、物質主義的論理学者が名詞の参照について提示する明確で一貫した説明を実施することが実際的な困難があると考えます。

名詞や用語に関連することだけではなく、特定の声明や一般的な主張を含む命題についても同じことが言えます。

6. しかし、確率の領域に足を踏み入れ、この問題を実質的な観点や現象的な観点から扱うと、我々は以前とは異なる問いに直面する。その困難性は一時的または偶然的なものとして軽視することはできない。推測と事実の間の中間的な状態は、一時的または偶然的なものではなく、ここでは通常の状態である。これは確率が特有の状態であり、その特性を表している。したがって、概念論的理論をいかに否認しても、概念論的言語の使用を完全には拒むことができない。特定の人間が来年死ぬことを証明できたり、十分近い証拠を得て実質的に確信できる場合、その死は未来の事実と呼べるかもしれない。しかし、彼の死がただ可能性として考えられるだけであれば、これは確率が特に扱う推論や推論の代用である。もし我々が好むなら、「可能性のある事実」という言葉を使うこともできるが、その意味を吟味すると、これらは不明瞭で自相矛盾していることが明らかになるかもしれない。

確かに、ここには完全な意味での事実がある。それは最終的な見解に基づく統計データであり、これについては調査したすべての者が認識し、承認している。同様に、これらの統計を帰納的に拡張して他の場所でも類似の統計が成立することを主張する言葉も使用できる。これらは普遍的な受け入れを正当化することができる。しかし、これらの声明は、前章で充分に示されたように、個別の事象に関する声明とは立場が異なる。前者の確立と議論は帰納学の領域に属し、後者のみが確率の領域である。

7. 適切な用語が欠けているため、しばしばそのようなことを表現するのに困るか、強制されるために、「事実」という同一の名前を個別の偶然性に適用することになる。例えば、男が死ぬという事実は確率的である、可能性がある、不確実である、またはどのような状況でも話すことに躊躇しないだろう。しかし、このような表現は技術的な用語が不足していることから生じる厳密には誤った使用形態だと私は考えざるを得ない。

確かに一つの選択肢が起こらなければならないことは明白だが、その確定性は私たちの考察の対象ではない；私たちは目の前に単一の選択肢に直面しており、それは著しく不確実であることが公的に知られている。これが現在注意を集めているものであり、その発生を評価しなければならない。このことは、いかなる限定条件を加えても「事実」と呼ぶ権利がない。反対の選択肢が同等の要求を持っている場合でも、あるいはそれよりも微々たる差しかない場合でもである。

このような言葉遣いは既に述べた通り、推論的論理学においては、その発生確率が非常に高いと推測されるが、その非発生を実際的な観点から考慮に入れる必要がないほど確度が高い推測であり、またこれらは一時的なものとして扱われる場合に完全に適切である。しかし確率においては、推測の確度はどのような程度でもよい；それは他の選択肢と同様に確実であるか、あるいはそれよりも遥かに低い確度である。

これらの後者の状況では、例えば男が死ぬ確率が非常に低ければ、その「事実」を高確度で不可能であるという限定付きで話すことは言葉の乱用だと考えるべきである。確率の対象となるべきものは不確実性であり、その本来かつ正確な使用形態は実際にまたは近似して確定的なものに限られるからこそ、このような言葉遣いを適切とは呼べない。

8. この状況は、確率の特徴的なものであり、その性質は恒常的である。科学の治療下に留まっている限り、我々の推測や何らかの呼び方をしたものを事実には発展させない。例えば、サイコロが目を出す確率を計算したり、人が特定の年齢を超える寿命の確率を計算したりする。このような知識への近似は、その後も何らかの方法で調査を行わなければ、それ以上のものになることは期待できない。我々は帰納法のように、何か実験的な証拠がいつでも現れてそれを確定性に引き上げるような感覚を持つわけではない。それはただ特定の強度を持つ推測であり、確率が関わっている限りその性質は変わらない。

したがって、それ以上のことが可能になるためには直接的経験や帰納的な証明に訴えるしかない。我々はしばしば述べたように、個々の事実はここで決して決定されず、ただ多くの事件の最終的な傾向と平均値のみが決定される。確率を二度目の引用によって推測の質を改善することは確かにできるが、その本質的な性質は始まりから変わらない。

したがって、この章の冒頭で示された無神論的観点における論理の説明は、大まかには正確であるが、一部再考と再定義が必要であると考える。それは通常の帰納法に関しては適切だが、確率も含む知識の獲得に関するより広い視野においては、論理学の科学が帰納法だけでなく確率も含む場合にも同様に適用されるように修正する必要がある。

9. それゆえ、科学的な思考を行う者について述べるのです。彼が発見者である原始的な研究者であれ、あるいは証明を分析し説明する論理学者であれ、彼は時間の両端と空間上のすべての方向に広がる客観的事象の世界に囲まれていると考えるべきです。これらの事象の多くは、そして常にその大部分は未知であり続けるでしょう。これらを事実として述べる場合、我々はそれらが人類の知識の潜在的な対象であることを意味し、適切な状況下では人々がそれらについて決定的かつ最終的な合意に達することができるということを意味します。科学的または物質的な論理学者には、これらの未知的事象を可能な限り既知のものに変換し、上記で述べたように経験と観察から始める特定のデータ核を中心に集約するという過程が課せられます。この作業において彼の主な資源は、前章で部分的に述べた帰納法の方法です；また別の資源として確率論があります；さらに別の資源として推論があります。

現在、このような言葉が概念論的であるとして反対されようと如何ともしがたいにせよ、これらの過程を説明する最も簡潔な方法は、「私たちはこれらの外部現象に対する概念を探求し、可能な限りその概念を事実に変換している」という表現を使うことだと私は見ています。私たちが「事実」の自然史を追跡してその起源まで遡ると、最初には単なる推測や予想として存在し、その現実性が全く否定されるか、あるいは完全に疑わしいとされ、それらは事実とは明確な差異がある段階で存在します。これらの「事実」がどのようにして生じるかについては、論理学ではなく心理学の問題です。これらがこの形態である間、論理学にはほとんど関与する要素はありません。人々は生涯を通じて様々な主題についてこのような推測を抱き続けますし、哲学者が一般の人々よりも優れている点は主にその推測の質、そしてそれらを洗練させ、検討する技術と恒久性にあります。

次の段階では、これらは概説的であるか、または何らかの形で壮大さや重要性を持つ場合、「理論」や「仮説」と名付けられます。しかし、 trivia な内容であったり、詳細に関することであれば、私たちはそれらを区別する適切な名称を持っていませんし、従って「概念」と呼ぶことに満足せざるを得ません。これらの概念は帰納論理において非常に速やかに移動します。論理学者はしばしばその証拠が強力になるまでそれらを留め置き、それが事実として世界に堂々と現れるまで待つこともあります。

したがって、既に述べた通り、この段階は論理学においてあまり詳しく扱われません。しかし確率論の全体的な業務は、これらの概念について議論し評価することです。従って、この科学に関連して、実質論学者による名詞や命題の参照に関する説明は修正されるべきです。

10. したがって、確率における私たちの立場を説明する最良の方法は次のようになります：ある事象、過去、現在、または未来についての概念を持っています。この概念は、その性質から考えられる唯一のものであり、現状ではそれが事実と呼べるほど正確なものではないとしても、もし可能であれば、確実にそれを事実にするためにその一方または他方で結論を出すことでそれを変換したいと考えています。しかし、確率のみに基づく結論を導き出す限りにおいては、これを行うことはできません。私たちができる最大限のことは、この概念を評価または評定することです。確率の役割は、そのような評価をするためのルールを提供することです。統計に参照するか、直接的推論を行うことで、私たちはどの程度この概念を信じるべきであるか、つまり全体の事例数からどれだけの割合でそうすることが正しいかを述べることができます。したがって、これらの状況においては、私たちの立場は概念を持ち、推論の過程は既に受け入れられた真実や事実に加える正当な範囲を確定することであると見做すことができます。

確率のみに基づいていた限りでは、これらの概念はそのような形で留まることになります。しかし、帰納法に移ると、それらはさらに進むべきステップがあることを示します。その最終的な段階は、事実となって不問の真実として位置づけられるまで到達するまでです。これは論理における最終的な目的地であり、私たちがその目的地に至るまで作業は完了しません。

11. このように、確率における私たちの立場を概念を持つことと、それをどの程度信じるべきかを決定することに焦点を当てるという言葉を使用しても、概念主義的であるとは思われますが、その精神は全く異なります。最終的には常に事実に参照します。私たちはデータとして事実から始まり、結果において可能であれば再び事実に戻ります。当然、確率において個々の結果ではこれを行うことはできませんが、それでも（第六章で示したように）私たちは結論を事実に基づいて正当化します、つまり長期的な経験に基づいてです。

この部分についての議論は、多少退屈に思われるかもしれません。しかし、論理学の二つの主な見解を区別する際には、その必要性が明らかに感じられたため、誤解や混乱を恐れてこれを避けることができませんでした。さらに、次の章の過程で確認されるように、この部分について十分に明確かつ満足いく状態にしておくことがなければ、いくつか重要な結論は適切に説明し正当化することができませんでした。

\newpage

\chapter{時間性と相対性の諸帰結（Consequences of the Previous Distinctions）}
\label{ch:xii}

\begin{center}
\large\textit{事象生起前後の確率・時間無関係性の原則・確率の相対性批判}
\end{center}
\vspace{4mm}

第十二章.

前章の区別による結果。

1. 今や私たちは、前章で定めた区別から直接的な結果ではないとしても、少なくともその説明後にはより容易に理解され、受け入れられる重要な結論を解説し正当化する立場にある。

まず第一に、確率においては「時間」が問題に関係しないことが明らかになる。言い換えれば、我々が議論している事件が過去、現在、または未来であるかどうかは関係ない。私たちは直面している問題は、その最も単純な形で次のようなものである：統計（推引によって拡張され、実際にはしばしば帰納的に得られる）は、ある特定の事件が一定の割合のケースでどのように起こるか、起こっているか、または起こることになるかを教えてくれる。我々はその事件について概念を作り出し、可能性としてそれを考慮するが、それ以上に我々は知りたい。つまり、我々は以前の章で量的信頼性について説明した通り、そのイベントがどれだけ期待されるべきかを知りたい。従って、我々の知識と、したがって予想の正しさに対する正当化または否認が形成された後には必ずしも来ることになるという事実から、事件には相対的な未来性があると言える。つまり、我々の知識は予想が形成された後に到着することがほぼ必至であり、しかし未来性は相対的である。問題を解決するための証拠がまだ得られていないかもしれないし、あるいは私たちはそれを手元に持っているが後でしか参照しないかもしれない。この事実から、予想が形成された後に必ずしも到着することになるという相対的な未来性があることから、私は「事件の概念」を「予想」よりも好んで使った。後者については、いくつかの点からはより理解しやすく適切なように思われたにもかかわらず、一般的な見解では絶対的な未来性と相対的な未来性の対比を示すという誤解が生じる可能性があるという欠点がある。

2. たとえば、サイコロを投げます。六回のうち一度しか目が上がらない場合、サイコロの目が上がったことを前提に推測すると、その推測は六回のうち一度だけ正しくなります。このような推測は通常の計画によれば、時間の進行方向へ進むことになります。つまり、投げられる前に結果を予想し、まだ誰もそれが何であるか知る由もない状況で意見を形成することです。あるいは、過去に投げられたサイコロについて後方から推測する方法もあります。つまり、ある過去の機会に投げられたサイコロについて意見を形成し、その後証人の証言によってその意見を修正することです。どちらの場合も、精神的な操作は同じです。統計的根拠に基づいて形成的に意見を形成し、その後特定の証拠によって修正します。

この意見は過去、現在、または未来の事件に基づくかもしれません；その後の証拠は私たち自身の目で確認するもの、他の人の証言、あるいはどのような推論でもです。証拠とは、事案を最終的に確定させるためのその後の検討や調査を指すものです。もちろん、この特定の証拠が常に得られるわけではありません。その場合、その概念は概念として残り、それが「事実」として認識される程度の信憑性を得ることはありません。これは、結果が記録されていないすべての過去のサイコロ投げについても明らかです。

確率的なゲームについて議論する際には、将来性という実質的かつ相対的な側面に限定することに明確な利点がある。過去の事象については直接的な情報が不可能であるため、すべての人々は比較的知識不足とし、その結果を形成するために知られているまたは推測される事象の発生頻度から意見を形成するしかない。一方で、事象が既に起こっている場合、通常は何かしらの証拠があり、それがいかに微細なものであっても有用である。この証拠が直接手元にあるわけではないとしても、少し時間を置くことで、少なくとも意見形成に役立つ情報を得ることができる。したがって、このような予測において実際的な将来性に限定することにより、一般的な評価では未来性と確率は密接に関連している。

3. しかし、この観点には誤りが伴い、または少なくとも常に誤解を招く可能性のある表現の不正確さがあり、これが広く受け入れられ、最高の権威者の著者たちによって承認されてきた。例えば、バトラーは『アナロジー』の中で、ミルもまた、ある形や別の形で表現して、事象が起こる前と後における不可能性を区別することに注意を向けている。彼らはこの区別を全く異なるものであると考えている。この用語が重要なものであるという点については否定することができないが、私はしばしばその表現が改善されるべきであると考える。

バトラーはこの問題についての言及は、『類推論理』において、奇跡に関する章の中で述べられています。奇跡がある強い反対の前提（彼が通常表現「事前における不確実性」に相当するものとして用いている語句）であることを認めつつも、彼はそれらを認めるために、他の同様に強い反対の前提を持つ出来事が、実際には非常に薄い証拠で受け入れられていることを示すことで努力しています。彼は次のように述べています。「一般的な推測的な真実や最も普通の事実に対する証明前の強い反対の前提は、それらが証明されたときでもほとんどどんな証明にもかかわらず克服されます。カエサルの物語や他の誰かの物語に対して百万分の一以上の反対の前提があります。例えば、ある数の普通の事実があり、そのような状況にあるもので、一つも証拠がありませんが、それが思い浮かぶと、誰しも無条件にこれらの事実は偽であると結論付けるでしょう。そしてこれは単一の普通の事実についても同様です。」

これらの言及は、かなり批判を浴びており、その内容は私には説明不明で曖昧なように思われます。文脈から推測すると、また後で同じ議論が参照される別の箇所[1]からも、バトラーは奇跡の報告と、一般的に「不確実」または「不確かな」、あるいは反対の証拠があると思われる他の報告を区別していないように見えます。そして彼は、一部の後者があまりにも薄い証拠で瞬時に受け入れられたことから、前者を同様に或る程度良い証拠で否認することは全く無理筋であると結論付けたのです。

この主題は、異常な話の信憑性に関する章において再び私たちの注意を集め、より詳しい議論が必要となるでしょう。ここでは十分に述べることとして、このような見解が一部の神学者には満足するかもしれませんが、奇跡を否定する者はこれを受け入れることは決してありません。彼は自然に反対し、バトラーの考えとは異なり、「奇跡的な要素」は「話の説明に対する小さな追加の証拠」とではなく、「全く異なる不信性のクラス」に入ることを指摘します。そしてそれは、話を受け入れる障害となる困難と反論を増やすものであり、単に追加するものではないと主張します。

ミルの言及[2]（『論理学』第3巻、第25章 §4）は異なる性質を持っています。彼は不信の根拠について議論し、「事前に起こる可能性が低い」と「事後的に起こった可能性が低い」、つまり二つの異なる特性があり、後者は常に不信の根拠である一方、前者は必ずそうではないと述べています。彼はサイコロを投げることを例に挙げます。事前にサイコロが1目が出る確率は低く、しかし後に「信用できる証人がそれを主張した場合、それを信じない理由はありません」と言います。また、「A.B.が死ぬ確率は非常に低いですが、誰かが昨日彼が死んだと話すなら、私たちはそれを信じます」なども例に挙げています。

5. これらの問題について、困難があることは認めなければなりません。多くの人々がこれらを困惑の源とし、様々な説明がその困惑を解消しようと提案されているという事実は、これに十分な証拠です[3]。前章で考慮した点については、読者にとって若干過剰に技術的であり、甚至て学術的な表現に思われたかもしれませんが、私はそれらが問題の見方についてより満足できる方法を導くことを希望します。

「不確実な出来事」という言葉を使用する際には、客観的に見て、出来事はただ単に「より少ない」または「より多い」程度の頻度であるという点を忘れないでください。最も稀な状態とは当然のこととして、その出来事が全く起こらない場合です。前章で示したように、当時の判断形成において我々の立場は、ある概念や推測（何と呼ぶかは問いません）を持ち出し、それを信用性に相当する一定の重みを付けるという点にあります。従って、バトラーとミルがそれぞれ提示する二つの例の真正な区別は、これらの概念がどのように得られたかにあるのです。一つの場合には予測のプロセスによって得られ、もう一つの場合には証人の報告に耳を傾けることによって得られるからです。

6. バターラーの例を最初に取り上げましょう。彼がいう「証明前の仮定」では、ある男がカエサルの物語について考えていることを描写しています。つまり、特定の歴史的事実について確固たる根拠なく推測している状態から、その推測の真偽に対する価値をどのように評価すべきかを考える状況です。彼は当然のように、このような推測が誤りであると結論付けられています。しかし、「証明後の仮定」は何を意味するのでしょうか？それは、無作為に採用されるものではなく、実際に証人の提案や支持によって提示された物語が真実であるというものです。後者の場合、前者は確実に否認されることでしょう。しかしながら、これがあらわすものは、またはむしろ示唆するのは、ほとんどどんな証人でもその証言は大抵の場合、単なる推測よりも遥かに優れているということです[4]。両方の状況において「事実」について話すことを好むなら、「事件」と呼ぶことができます。実際には、前章で認められたように、一般的な言葉では他の表現法は容易には適用されません。しかし、どのような表現を使用しても、実際に男の心に存在し、その確率的な価値をどのように評価すべきかが問われるものは、「事件」の概念であり、その用語が既に説明された意味でのものです。そして確かに、どちらの概念にもより重要な区別を設けることができないでしょう。一方は単なる推測に基づいていると仮定し、他方は証人の報告に基づいているという差異です。ミルが示した例についても同じ remarks が適用されます。A.B. の死以前には、私たちはその問題に対する見解は生命統計に基づく自らの推測に過ぎませんでした；彼の死後には、ある誰かによる証拠に基づいていました。おそらくその人は事実を正確に知る十分な機会があったと思われます。

7. 当前の区別が時間とは根本的に関係がないことは、一瞬考えれば明らかです。たとえば、ある人がA.B.が死ぬかどうかを知る機会があったと想像してみてください。その人が「A.B.は明日死ぬ」と教えてくれた場合、私たちは彼の言葉を信じる準備ができているはずです。また、「A.B.\_は\_すでに死んでいる」と教えてくれた場合も同じです。もし私たちがその声明についてまだ疑問を持った場合（常にその人の誠実性については十分信頼していると仮定する）、それは私自身や彼自身の場合でも同様に予想と同じレベルの推測であると考えているからです。

過去の出来事に関しても、それが過去になったという事実は何ら影響を与えません。証言者からの情報がなければ、その出来事が私たちの頭を過ぐった場合でも、あるいは未来の出来事のように疑問を感じます。

したがって、出来事前後の確率の区別は単に次のようなものであると考えられます：出来事前には、統計的な情報を用いて様々な可能性の中から推測する以外の手段がほとんどないため、予想を繰り返します。出来事が起こった後では、証言や観察などの具体的な証拠によって予想が改善されたり置き換えられたりすることが多いです。したがって、このような場合、我々の評価は通常より価値が高いです。

しかし、これらの特性が逆転していたらどうなるでしょうか？つまり、過去について推測し続け、未来について追加の証拠を見つけることができたと仮定すると、それぞれの評価の価値も逆転します。これらの価値の差は時間とは根本的に関係なく、その出来事に対する我々の予想や推測が立っている根拠に完全に依存しています。

8. 以下の仮想的な例は、上記の点を明確にするのに役立ちます。非常に短い記憶を持つ人々を想像し、かつ彼らが過去に起きた出来事を永続的に保持するためのどんな種類の記録も保存していないとします。そのような人々にとって、過去の領域は我々が未来に対して抱くように、予想や推測の領域となります；つまり、彼らは一般的な確率に関する考慮事項に基づいてそれらを判断するしかなく、直接的な具体的な証拠によって判断することはできません。しかし、また彼らの中に先知のような人々がいることを想像し、これらの先知たちは我々の様々な歴史書や伝記、そして記憶が過去について達成できるほど正確で信頼性のある未来を予言することに成功しているとします。直接的な証拠や証言、そして確率の現在の役割は単純に逆転し、その結果、事前および事後の偶然的特徴も逆転します。そうすると、後者の場合が比較して「常に不信の根拠ではない」と見なされる一方で、前者の場合には我々は常にそうであると主張されることになります。

9. これまで議論した誤りの起源を調査することは価値がある。私は、その理由は以下の通りだと考えている。しばしば、未来の出来事について推測する際には、特に賭博から得られる未来の出来事を扱う場合においては、全ての人々がその出来事に対する完全に同じ状態の無知に置かれていることがあり、またしばしばある。利用可能な情報の限界は統計に限定され、未知の出来事が様々な形のいずれかをとることを知るだけの知識に過ぎない。したがって、一人の人間による推測は他の人よりも価値がないわけではない。しかし過去に関しては状況は全く異なる。ここでは統計情報に基づいて信頼する習慣がないため、人々の推測は非常に異なる価値を持つ；多くの場合、それは我々が「肯定的な知識」と呼ぶものに等しい。これにより、一般的な評価において、過去と未来における主観的確実性との明確な区別が生じるが、これは推理の科学から見れば存在すべきではない。

この結果、過去と未来それぞれに対して、「出来事の機会」という多少曖昧な表現を適用すると、通常は異なる意味を持つ。未来に関しては、統計的根拠に基づいて推測に与えられる価値そのものである。これは、このような推測がほとんど唯一の基準となるからである。しかし過去に関しては、一部の人々には推測しかない一方で、他の人々には肯定的な知識があるため、意味が移り変わる。現在では「出来事の機会」という言葉はしばしば、統計に基づく推測に与えられる価値ではなく、他人の推測やその証拠に基づいて得られた推測に与えられる価値を指すようになっている。

10. この章の一部の声明に反するように見受けられる一連のケースがありますが、詳細に検討すると、それらはこれらの声明を明確に支持することになります。例えば、私は遠くの田舎で散歩しているとします。そのとき突然、友人が現れました。この出来事は自然と驚きを与えます。しかし、私たちは事件自体が可能か不可能かについて正確には話せないとする見解を考慮に入れると、なぜこのような驚きを感じるのでしょうか？事前に何らの推測も形成していなかったはずです。それにもかかわらず、多くの人はそのような出来事を驚くでしょう。

この質問に対する適切な回答は、私たちはそうした予想を無意識的に形成していたということです。どのような場面でも、誰しもが無意識に知っているものと知らないものを区別します。相当の期間にわたって、無数の人々が私たちと遭遇しましたが、これらはすべて私にとって未知のリストに入りました。このような区別の存在を思い出させる兆候はありませんでした。しかし、例外的な出来事が突然起こったとき、その区別が鮮明な意識に浮かび上がることになりました。

驚きの中で使われる言葉もこれを見せるかもしれません。「ここであなたを会うとは誰が予想しただろう？」つまり、「もし事前にこの潜在的な考えが意識化されたら、それをどれほど重視するだろうか？」私たちは過去形を使用し、その区別は常に心の中に存在していたことを示しています。私たちの心には、このような既成の事件のクラスがたくさんあり、何かが非常に少ない数のクラスに該当したとき、それが事実であることに気づくのが自然です。

あるいは、私は一団の兵士の一人として射撃を受け、その射撃が私だけを打ったとします。この出来事に対して驚くのはなぜでしょうか？共通の言葉で説明すると、「どうしてそれが私がいたる人間の中でただ私に当たったのか！？」と考えます。ここで我々は、自分自身と他者という非常に自然な二分法を考えています；その結果、驚きは再び、ある意味では後ろ向きなものとなり、この分類に関連しています。予想は明確には形成されなかったため、事前にその出来事を考えていなかったからです。しかし、それが起こったとき、それは適切なクラスに即座に割り当てられます。

11. この見解は以下の考慮事項によって確認されます。友人にこの話をすると、彼も驚くでしょうが、我々ほどではありません。彼の分類は、私たちはその一人である友人と、残りの人間という形になります。これは必然的な分類ではなく、しかし彼が最も容易に思いつく分類です。

さらに別の Stranger に話すと、彼の分類は異なるものになります（つまり、私たちは多数派に含まれる）。この場合、我々は彼を驚きがあることに気づかせることができません。

これらの言及では、すべての場合において驚きが存在する場合にその正当性を証明しようとする意図はありません。異なる人々は、仮定された状況で異なる反応を持つかもしれませんので、例は主に説明のために提供されています。それでも、既に議論した原理（第6章 §32）に基づいて、驚きの程度について類似の一般的な正当化を見つけることができるかもしれません。

12. これらの状況に対する一般的な回答は、驚きが生じるべきではないことを示す努力に留まり、驚きがどのように生じるかを説明することには至りません。その形式は以下の通りです。「あなたは驚く権利はありません。実際には何か特異的なことは起こっていないからです。もしこの特定の出来事が起こらなかったとしても、他の何か同様に不確かなことが起こらなければなりません。射撃が私やあなたの友人に当たらなかった場合でも、別の誰かが事前に打たれる可能性が等しく低かったはずです」

まず、この答えは、実際にはほとんどの人がこのような場合で驚くことを説明していないという事実を述べる。また、驚きの程度が先に述べたように異なる割合であることも説明していない。さらに、これは不可能なことが実際に起こったと認めることの内在的な不満足さを持ち、他のすべての選択肢も同様に不可能であると言っていることで困難を乗り越える方法を提示している。この推論から自然には、自分自身では実際には improbable なものが存在し、しかし非常に薄い証拠で確立される可能性があるという結論が導かれる。バトラーはこの推論を受け入れ、その結論に至った。ミルは、事後のイベントによって与えられる不可能性の値が異なるという考慮からこれを避ける試みを行った。この点についてさらに詳しい議論は、誤謬に関する章と、非常識な物語の信憑性に関する章で見ることができる。

13. 現在議論中の主題に関連して、しばしば確率に関する著作で強調されるが、その重要性が過度に評価されている区別について注意を払ってみることにする。それはしばしば「確率は相対的である」という主張であり、これは特定の事柄に関してそれぞれ異なる情報を持っている人々にとって異なる値を持つという意味である。例えば、AとBが4つのボールが入った袋からボールを取り出すときを考える。Aは黒と白のボールがあることを知っているが、それ以上のことは知らない；一方、Bは3つが黒で1つが白であることを知っている。この場合、Aにとって白いボールの確率は1/2であり、Bにとって1/4であるということが述べられる。

ただし、物質的な観点からこの問題を検討すると、私にはそれほど深遠な意味があるとは思えません。他のどの科学でも同様に真実である原理、すなわち仮定するデータによって導き出される推論が異なるという原理と比べて、それほど大きな違いはないと考えられます。

論理的な観点から言えば、私たちが領域の面積や山の高さが相対的であり、したがって一人にとっては一つの値を持ち、他の人には別の値を持つと言えるようにも思えます。引用された例の真の意味は次の通りです：A は長期的には白と黒が出る頻度が等しいとする系列からランダムに白を選ぶと仮定しています；B は長期的には三回の黒に対して一度の白が出るとする系列から白を選ぶと仮定しています。したがって、彼らは同一の規則を適用することで異なる結論を導き出しますが、他のどの科学でも同じ状況下において同じようなことが起こります。

たとえば、二つの人間が山の高さを測定している場合、一人が基点を1000フィートと仮定し、もう一人が1001フィートと仮定すると、当然ながら彼らは異なる意見を持つことになります。測量学はデータの真偽に関与しないものとして考えられていますが、そのデータが正確に測定されたと仮定されるのですから、これと同じことが確率論においても成り立つべきではないでしょうか？ただし、確率論が扱う特殊なデータの特性を考慮に入れる必要があります。

14. この確率の相対性という見解は、私には主観的な視点と結びついているように思われ、実際にはその特徴でもあります。この見解が弱い側面を示す例としては、このような表現に重きを置くこと自体が挙げられるでしょう。前章で充分説明した通り、概念主義の原則を追究するにつれて、私たちは物質論理の根本的な問いから逸脱し、つまり私たちの信念が実際に正しければ、その正確性はどの程度であるかという点に焦点を当てるよりも、実際にはどのような信念を持っているかという点に注目することになります。そして、その信念は当然、個人の状況や他の情報源によって変化するため、その相対性が避けられないものとして認められるべきであり、あるいは少なくともこのようなことが起こるのも不思議ではないと見做されるべきです。

したがって、私たちは確率についての見解において、「相対性」という表現は誤ったものであり、代わりに我々が持っているデータの種類や量の変化によって生じる効果を述べることの方がずっと好ましいでしょう。確かに、科学の中でこのような変化が頻繁に見られる場合があります。例えば、私は海賊船に乗客としている友人を待っています。船上には百人の乗客があり、また船員も百人います。新聞で一人が転落して溺れてしまったという記事を見ます。当然ながら、それが私の友人が溺れたと予想する確信は小さいでしょう。別の記事では、その人が乗客であり船員ではないことが明らかになります。これにより、私はわずかな不安感を倍増させました。さらに別の情報によると、それは英国人で、その航线と季節において英国人の乗客は少ないということが判明します。これにより、私は確固たる恐れを持つようになりました。そして、小さな情報が瞬時に私の期待を変えることになります。

15. 今、確率が推論と区別される特徴として、意見を形成する際に浮遊しているいかなる小さな事実にもこうした形で屈折されることの特性があることに注意しなければならない。そういった状況を我々の期待状態を「相対的な」と呼ぶことにはどのような害がないか（と主張されることがある）？

私としては、二つの反論があると考える。まず第一に、その表現が一貫性という観点から拒否されることになる。これまで採用されてきた主題に対する一般的な精神と扱いとは不一致であり、確率を推論的論理から切り離す傾向があり、それらを関連する科学として捉えることを妨げるからである。我々は、私たちの道で到達可能な真実に向けた真理を目指しており、したがって、真実そのものが相対的な意味でのみ結論を相対的と見做すことを好まないという理由があるからである。

二つ目の点では、不安定な信頼状態であるこの条件、新たな知識の獲得によってどの程度でもいつでも私たちの判断に影響を及ぼす可能性があるというこの継続的なリスクが、確率そのものに固有の特性とは思われません。むしろ、その健全で正当な応用においてはそうではなく、規則に基づいて急いで形成された判断に適しているように思われます。このような急いで形成的な判断は、通常の論理的結論の場合にも起こり得ます。実生活における緊急性から、私たちはしばしば十分確固たる信念を持って一時的に仮定を形成しますが、その証拠は当時の私たちには最終的なものとは認識していません。短い時間だけ待つと、何か他なるものが現れ、意見を変えたり、甚至は逆転させたりすることもあります。このような結論は、与えられた状況において完全に健全なものであっても、それは私たちが前もって持っていたデータに基づいて判断を下さなければならない他の誰もが同じような結論を導き出すものであり、その点では「相対的な」判断と現在議論している意味で明らかに認識されます。しかし、論理が真実の道具として扱われるべき視点を系統的に追求する人であれば、このような判断をそう呼ぶことは避けた方が良いでしょう。

16. これまで例示した確率に関する例では、結論を支える統計データが既に与えられていると仮定することが多かった。一方で、直接の具体的な証拠を得ることは不可能であり、判断は実質的に確率的なものであり、これらの統計データに基づくべきであるということが前提となっていた。言い換えれば、それ以外のより良い情報源が利用可能でないことを意味した。他方で、これらの統計データは誰でも観察できるものと仮定されており、それらよりも劣った情報を用いて判断を形成することは避けられた。つまり、ここでの状況は他の多くの科学の場合と同じであり、結論を導く人々は同じ情報源から出発していると考えられていた。

例えば、当時のゲームの確率に関わる場合においても、不明瞭さや誤解が考慮に入られることがない。また、生命統計や保険に関する統計でも同様である。判断が正確（その性質に応じて）または正当化される限り、一般的な死亡率表は誰一人として利用できる情報源以外にはない。

17. しかし、ある個人の死期についての見解においては、その人の健康状態、習慣、職業などに関する私的な知識や合理的な推測を考慮に入れることで、それぞれ異なる評価を得ることが可能である。それでも、誰もが自分の評価を正当化するためには、統計的根拠に依存したことを明示的にまたは暗黙的に引用しなければならず、これらの情報が他の人々にも事前に利用可能でなければ、それらを彼らの注意に引き出す必要がある。つまり、ここでも他所と同様に、私的な根拠に基づいて判断を形成している限り、その判断は相対的である。正当化プロセスは論理にとって不可欠であり、個々の見解を修正し、観測者全員が情報源に関して同じ基準に立つようにする傾向がある。

\#\#\# 翻訳

従って、確率の結論を絶対的かつ客観的なものとして捉えることが望ましい。これは、推引（Induction）におけるそれらと同じような意味で、ただし程度は遥かに低いとは言え、そう考えるべきである。我々の結論がしばしば新しい知識によって時間と共に極めて大きく変動することを十分に認めつつも、そのような信頼度の揺れを、特別な統計に基づく急激な判断と見做す方が好ましい。これらは単なる全観察者共通のものではなく、それらだけに依存すべきではない。このような判断が急激なものであるという指摘には、その持つ罪悪感は含まれないが、何かしらの理由により、十分な証拠（統計的またはその他）を待たずに形成されたと解釈されるべきである。この説明は事実に適合し、また本著作で主張する一般的な哲学的位置づけと一致する。

1. すべての種類の内部の不確実性が外部証拠を弱めることが自明ではないか？確かにそうだが、ここでのこの主張は何の実用的な目的にも及ばない。既に証明されているように、道徳的確信まで達する真の内部の不確実性は、最も普通の証拠によって克服されるからである。

2. 神迹（Miracles）は日常経験と類似した一般的な自然現象とは比較すべきでない。それらは異常な自然現象と比較すべきであり、その比較は神迹に対する疑念と、不普通な現象に対する疑念との間に行われるべきである。そして、そのような比較は彗星のような不普通な現象を例に挙げることもできる。

以上の翻訳では、原著の論理や語調、数式（文中にはないが）に忠実に日本語で表現しています。

3. たとえば、ジェームズ・F・スティーブンス卿は、チャンスと確率を区別することでこの問題を説明しています。彼はバトラーがこれらを混同していると言います。「非常に普通の証拠でも一千万対一という前提となる仮定を覆すための反論は、確率とチャンスの混同に基づくものです。事件の確率とはその事実を証明する能力です。一方、その事件が起こる可能性のある状況（等しく有利と想定される）の数に対する、実際に起こる可能性のある状況の数との比がその事件のチャンスであり、それに対して起こらない可能性のある状況の数も考慮に入れます」（『英刑法総論』、255ページ）。ドンキンはまた、1851年6月の『フィロソフィー・マガジン』で、「不可能性」と「信じ難さ」を用いて同じ区別を行っています。

4. 証人が自分自身がただ推測するか、または推測したと同様に信頼できない場合、両方の証言は当然同等の立場になります。このケースについては第17章でも再び触れることになります。ここでいくつかの微妙な点があり、証人の信用性に関する問題を事前に調査しないと十分には説明できません。

5. ダンテによれば、『地獄』に住む人々ではこのような状況が存在していました。彼がそこで出会う枢機卿や他の人々は、まだ地上で起こるであろう多くの出来事を情報提供することができますが、実際に起こった出来事については質問しなければなりません。

\newpage

\chapter{様相論（On Modality）}
\label{ch:xiii}

\begin{center}
\large\textit{様相命題と確率・アリストテレスからカントまでの様相論・法学における証明度}
\end{center}
\vspace{4mm}

第十三章.

模態論の概念と扱いについて

1. 仏教学問に多少精通している読者であれば、すでに我々が多かれ少なかれこの論理地図上の最も棘手で嫌悪感を覚える地域——模態論について触れていることに気づいていただいていることでしょう。しかし、主題に関する議論をより明確にまとめることも適切かと思います。そこで私は、専門的な論理学者が模態語に対して受けた主な扱いの種類について簡潔な説明を行う章としてこの章を設けます。

まず最初に述べておくべきは、模態と模態命題がこれまで異なる時期でどのように解釈されてきたかについては、固定した一貫性のある意味合いではなく、多様な見解があったということです。この多様な見解の生じ方は、論理の領域とその本質に対する見解の相応する違い、そして自然の構造と順序に関する哲学的・科学的な意見の違いから来ています。さらに、最近では、数学者が独自の視点から発展させた確率論の理論の漸進的かつ安定した発展という非常に強力な要因が模態論の扱いに変化をもたらすことを認識する必要があります。

2. しかし、これらの異なる扱いにもかかわらず、この主題についての議論には常に一定の共通点があった。量的信頼度への言及がほとんど常に存在し、デ・モラン[1]の言及を正当化するほどのものだったかもしれない。「確率は「未知の神」であり、学校人たちはこの種の命題について扱う際に、人間の判断力で決定できるかどうかという点で最も苦労したのか、それとも模態を最も苦労したのか」というような議論を行ったのである。しかし、量的信頼度への言及は直接的かつ即時的なものであることがしばしばあり、また命題の主題事項の性質から間接的に生じる場合もあった。実際、論理学における純粋に主観的な視点と純粋に客観的な視点との区別は、第十一章で私が強調しようと努めたように、初期にはほとんど認識されていなかったし、カントの時代以前でもそうだった。したがって、模態に対する見解もその混乱に自然と参加していた。この点については、次の章を通じて明確になることを願っている。この章は、論理学における命題の模態性がどのように扱われたかを簡潔に示すために作られたものである。完全な歴史的概要を提供することではなく、代表的な見解を持つ著者だけについて議論する必要があるため、時間順序に厳密に従う必要もなければ、どの著者の意見についても詳細に議論する必要もない。このような調査の結果として、私は読者（もし既にその結論を形成していた場合を除く）に確信を与えたいと思う。論理学者たちは長い間且つ公正な試練を受けたにもかかわらず、模態に関するこの主題から何か満足する成果を得ることができなかったこと、そしてそれゆえに模態は完全に論理学から排除され、確率に移されるべきであることを示すだろう。

3. 三段論法の研究が始まった当初から、その過程は自domainにおいて完璧であることが認識されましたが、そのdomainは、最も単純な扱いの場合でも非常に限定的であることが明らかになりました。純粋形式の命題、すなわち「すべて（または一部）のAは（またはは）Bである」のような命題は、実践において我々のカテゴリー的な主張の中のほんの一部を占めるに過ぎません。しかし、BとAとの関係について Various degrees of necessity or probabilityを表す他の命題に遭遇することが頻繁です；たとえば、「Aは必ずBである」、「Aは可能としてBである」などがあります。また、これらの事実が我々の判断に与える影響についても言及します；たとえば、「私は完全に確信してAはBである」というものや「私はAがBであると考える」というようなものが多く存在します。このような命題をどのように扱うべきかという問題がすぐに浮上しました。古い論理学者たち、さらには現代の一部の後継者たちは、このトピックについて考慮することを完全に拒否したようには思えません。彼らの推論体系の真実性と完備性に対する信仰は非常に強く、その体系から普遍的に意味のある言語表現として認識される形式や、同様に普遍的に有効とされる推論形式が省略されることはあり得ないと考えたことはありません。これらの形式が不便であるか複雑であるからといって、それを省略することは不可能だと信じていたのです。

4. 非常に単純な計画が思い浮かびますし、実際にその提案は繰り返し行われてきました。すなわち、このような命題の特徴的な部分を、他處で困るものを収納するのに便利な容器である述語に移すだけです[2]。別の言及されるべきモダリティがこれによって解消されたのでしょうか？そして、否定的命題の特徴的な性質である否定詞を述語に移すことで、そのような装置で否定的命題自体を廃止しようと試みたことはないでしょうか？この方法では、ある程度は何かを行うことが可能だと認められます。論理が「砒素を摂取するとおそらく死ぬでしょう；A が砒素を摂取したので、彼もおそらく死ぬでしょう」というものであれば、主項を「砒素を摂取する者は、おそらく死ぬ人々である」と述べることで、この論理を純粋型の普通三段論法に変換することができます。そして結論は同じ形で、「A はおそらく死ぬ人々である」となります。しかし、この装置は私たちをほんの少ししか進ませません。もし小前提も同様のモダリティを持つものであれば、例えば「A はおそらく砒素を摂取したでしょう」という場合、述語にモダリティを移すことは四項論理に至るまで私たちを止めることが明らかになります。

しかし、この特定の反論がないとしても、上記のような装置によって理解性や方法面で何かを獲得するとは思えません。例えば、「毒物として砒素によるすべての中毒はおそらく死因である」という命題において「おそらく死因」なるモード述語が何を意味するかを考えます。普通の純粋な命題との類推が成り立つとするなら、それは主辞全体に対する述語でなければなりません。なぜなら主辞は分布されているからです。

しかし、その場合、我々は直ちに前章（第VI章 §19-25）で議論した困難に遭遇します。述語の適用が正当化または検証される試みにおいてです。この著作で採用された見解によれば、述語「おそらく死因」を主辞全体に対して適用することとは、結局のところ述語「死因」を主辞によって表されるメンバーの一部（半分以上）に適用することを意味するのかという問いが生じます。モードを述語に移すことで、述語の解釈においてどのような意味で行われるかというような複雑な問題が生じるとき、そのステップは確かに何かを得たわけではありません。

5. 逻辑学やその著者たちに関連する、主題のすべての困難を一石二鳥というよりは簡潔な方法で片付けるためには、模態が論理とは何ら関係がないと単純に否定することで解決できることが見つかります。これはハミルトンやマンセルなどの多くの現代の著者が採用した道です。これらの著者の著作を読むと、彼らは論理学から何が欠けているかを説明することが論理学の大部分を占めているように感じられます。彼らは模態が属する対象に属し、その対象に対する考慮によって決定されるため、論理外であると拒否を正当化します。私は彼らの拒否の根拠には部分的に賛成しています。第六章で説明したように、前提や結論の模態性の度合いが正当化されるためには対象への訴えが必要であることが難しく、その点は容易ではありません。

しかし、正当化という概念について十分な意味での議論は、ハミルトンとマンセルの科学に対する見方とは異なる範囲に属します。当時の推論の真偽に関連する完全な正当化は、これらの著者があまり気に留めない点であり、事実との一致ではなく論理の一貫性のみに焦点を当てています。

ハミルトン派として話すなら、模態は形式的であるべきだと主張します。なぜなら、私たちが一つの命題に対する信頼度を正当化するためには対象への訴えが必要ですが、それほど正確な評価は対象への訴えなしでも十分に可能であるからです。これだけのことは形式的なものとして扱うのに十分ではないでしょうか。

ハミルトンが模態を否定する際の説明は、決して明瞭かつ一貫したものとは言えません。彼は既に述べた通り、形式的模態と実質的模態（つまり、真偽）を区別することも欠けており、形式的模態について議論する際にアリストテレス派とカント派の視点から見た模態の多様な側面を区別することも忘れています。この差異の量と重要性については後ほど述べますが、ここで指摘しておくべきは、ハミルトンが（第1巻、257ページ）形式的模態を否定する理由として、「必然的、偶然的、可能的、不可能的な区別は、命題の事実に完全に依拠しなければならない」と述べている点です。彼はこの点については正しかったと思いますが、その後少し進んで（260ページ）、「純粋な命題と模態的命題に関連する三つの用語の意味を説明」する際に、カント派の広範で異なる三分類、すなわち推論的、主張的、疑問的な分類を与えています。彼はこれらの二つの論理観点から生じるこの両方の模態に関する説明について聴衆に注意を喚起するための十分な警戒心を持っていないことが問題です。

6. ある種の模態論理の三段论は、形式的ではないからそれを排除する理由が不十分であるように思われ、一言触れることができる。前提「任意のAはおそらくBである」は、「大部分のAはBである」と等価である。明らかに、この二つの前提「大部分のAはBである」「大部分のAはCである」から、「一部のCはBである」という結論を導くことができる。これは、A、B、Cの性質に関わらず常に成り立つため、通常の用語では形式的三段论とされる。

しかし、マンセルはこのような三段论が形式的論理に属すると認めない。彼の理由は、デ・モランの論理作品に対する詳細な評価および批判[5]の中で述べられている。デ・モランによる「数値的に明確な三段论」の導入が主な原因である。

マンセルは上記の特定の例を議論せず、より広範な代数的形式について話している。彼は特別な数値的な例[6]でこれを検討する：21個のYのうち18個がXであり、21個のYのうち15個がZであるという前提から、「12個のZがXである」という結論を形式的論理からは拒否し、その理由は数値的な判断が分析的ではなく合成的であり、量に関する直感に基づいているからである。

ここではこれらの理由について検討することはできないが、彼の批判はカントの算術的判断に関する教義を受け入れることを要求している。したがって、三段论が形式的でないという理由ではなく、分析的ではないという理由に基づいて拒否すべきであると考える方が良いだろう。

7. ここでは、モード推論のplexperitiesを解消するもう一つの実用的な方法について触れなければなりません。これは専門的な論理学者ではなく、一般的な思考者にとっての資源であり、この章で既に指摘した最初の回避法と同様に、非常に部分的な適用範囲を持っています。その方法は、推論過程において前提を純粋なものとして扱い、その後結論にモードを再導入することで構成されます。これは完全な確信度の一種の資格付であると考えることができます。

各前提がほぼ確実であるか、または何らかの理由からその確実性からの逸脱量に関心がない場合、この粗雑な方法は十分に機能します。私は理解するところでは、このような例（§4）で意識的に推論を行うほとんどの人々の頭の中で行われる過程がこれです。前提「砒素を摂取すると死ぬ」と「その男はそれを摂取した」が完全に真実であるかのように進めると仮定し、おそらく真実であると仮定して結論「彼は死ぬだろう」と導くでしょう。しかし、前提が確実でないことを念頭に置いていれば、結論は限定的な承認の下でしか保有されないと記憶します。その承認の程度を考慮に入れない限り、単なる性質から言えば、「彼は Likely 死ぬだろう」と述べることは完全に正しいです。この場合、前提からモードを一時的に排除し、結論に再導入することにより、モードはその役割を果たします。

このような過程は明らかに非常に粗雑で不完全なものであることがわかります。実際、これは正確な考慮から正しく考慮されない点が一つあります。それは、完全な信服度からの逸脱の程度を示す異なる表現の濃淡を考慮しないことです。これは当然、モードが占めるべき重要な事項です。したがって、最も適切かつ科学的な方法で質問を決定する必要がある私たちに求められる粗雑な手段であると言えるでしょう。

8. 任意の応用論理の使用においては、当然ながら先ほど述べたようなプロセスを経る必要があります。純粋数学以外では、我々が推論する前提が絶対的な信頼性を持つことはほとんどないでしょう。したがって、結論には絶対的な信頼性からの逸脱が必要となります。しかし、前提が真であるという仮定のもとで推論を行い、我々が意識的に感じている何らかの確実性からわずかな逸脱は、結論に対する精神的な補正として保有されます。

しかしながら、このような考慮事項は主に一般的な応用論理に関連するものであり、三段論法や帰納法の規則が実践で適用される際の注意点やヒントに過ぎません。しかし、モダリティを扱う場合、完全な確実性からの逸脱の範囲は十分に大きく、言語を通じてそれを示し評価することが必要となります。そのような逸脱を測定し表現するための科学的な議論が必要です。これは前提と結論においても同様です。

したがって、モダリティを扱うためのこのような計画は、我々が公然とこれを考慮に入れないことに決意した場合と同じく、実際の論理的調査の領域からその要素を排除するものと言えます。

9. 逻辑学に模態を導入する人々が模態性の異なる扱い方について議論する前に、この部分における混乱の可能性を明確にするために何か述べておくべきである。これまで主に取り組んできたケースや、本著作の前半部においては、模態性（ここでは逻辑学的な用語として適切に使用する）が一般的には個別命題や特定命題の中に見出される特徴を有することが判明していた。これは個別の事例について平均値に基づいて判断する際に現れ、個々の命題が取り扱われたクラス全体の一部のみを指摘しているという事実を表現したものである。

例えば、「五十五歳の人々のうち三分の三是二十年内に死ぬ」という平均値に基づく与えられた状況においては、「AB（その年齢）は二十年内に死ぬ可能性が高い」または「指定された期間内にはABは三分の二で死ぬだろう」といった曖昧な形式やより正確な形式の命題と対面した。ここでは模態的な命題は自然に個別命題や特定命題の形で現れる。

このように、模態性が個別命題や特定命題の形で現れることが一般的であるという点について説明しておく必要がある。

しかし、普通の論理に移ると、普遍的命題がモダル性として述べられることが見受けられるかもしれません。これは主に必要または不可能と呼ばれる命題について多く当てはまるでしょうが、確率的なものについても同様です。我々は「すべてのXはおそらくYである」という形を目にすることがあります。類似の状況でこれまで採用されてきた解釈と同じように、このような命題のモダル性とは、その普遍的命題が正しく作られる回数に対する比的な数値を意味すると言わなければなりません。これには特に困難はなく、唯一の違いは前者、すなわち個別または個人的なモダル性の正当化がその包含する普遍的命題の範囲内で得られることに対して、普遍的命題のモダル性の正当化は他の一連または順序のある他の命題群を探さなければならない点にあります。この命題を似たような一連のものに参照し、その真となる場合の割合を考えなければなりません。しかし、この区別は根本的なものではありません。

実際には、普遍的命題が個別の命題よりも、このような引用による正当化を受けやすい確率は低いことは事実です。しかし、既に頻繁に指摘されたように、我々が関心を持つのは、私たちの命題が実際にどのように得られるか、または人々がそれらを最も自然にテストする方法とは無関係で、最終的な正当化に訴えるものであり、それが統計的性質であることが十分に示されている点です。したがって、「すべてのXはおそらくYである」という命題を述べるとき、我々が意味するのは、我々がその判断をする際に、半数以上の場合において正しく、半数未満の場合において誤っているということです。

11. この段階で、可能性の不確実性が遭遇される。All X is Yという状況における確率について話すとき、その補対的な状況である「All X is not Y」も含めて考えたり、暗に含意したりする。しかし、この後者の選択肢自体にも不確実性がある。例えば、失敗のケースでは、No X is Yという状況が起こる可能性もあるし、あるいは一部のみのXがYでないという状況も起こる可能性がある。これらの仮定は元の命題と矛盾するため、そのいずれもが元の命題の不成立の一例となる。実際的な面では、これらの中から選択するために様々な既知の規則や推論を用いることが一般的であるが、当然ながらそれらへの依頼は厳密には論理外のものである。しかし、このような不確実性の存在そのものが、そしてそれが主題に訴えることでしか明瞭化されないという事実自体は、模態的概念を普遍的命題だけでなく個別命題にも適用する際の真の障害ではない。

12. 模態論を論理に導入する方法が回避または拒否されたいくつかの例について述べた後では、模態論を受け入れるためには多少なりとも意図的にその主張を認める人々による模態論の扱いについて簡潔な説明に入らなければならない。

まず最初に検討すべきは、模態の性質に関する諸意見がどのようなものであったかである。すなわち、模態は主張の素材に対して最初から主要な影響を与えるものなのか、それともそうではないのか、そしてもしそうでないなら、それは何を意味するのか。この問いに対する考察においては、既に述べた通り、早期の論理学者たちには現在では一般的に維持されているような明確かつ整合的な主観と客観の区別が見いだされないことが明らかである[8]。その結果として、模態に関する彼らの扱いを見ることができる。常に主張の主観的側面に一定の関連性があったことは明白であり、この場合においては信頼の性質や量に関連していた。しかし、最初からその特性を直接評価し、それ自体で見なされたというよりは、それが何であるかを決定する素材から考慮した結果として初めて評価されたことが明らかである。

例えば、一般的に受け入れられている学術的またはアリストテレス的な分類は、必然的、偶然的、可能的、不可能的などに分けられる。これは主観性のほとんどが基盤となっている分類であり、純粋な心理面では必然的と不可能的は全く同じであることが明らかである。一方が命題の真実性に対する完全な信頼を含み、他方がその矛盾する命題の真実性に対する信頼を含むからである。同様に、心理面からは偶然的と可能的な区別も容易ではない。

したがって、この見解によれば、主張の模態は素材の性質への参照によって決定された。ある主張では素材の性質が述語が主題と必然的に結びついていることを決定し、他では不可能であることを決定するなど、そのような区別がなされたのである。

13. この四分割は、現代の心にその人工性が明らかであるため、批判する必要もありません。自然の研究をわずかに学び、帰納的証拠に対する適切な評価さえすれば、現象間のすべての関連性が、必然的と呼ばれるものであれ偶然的であれ、依存する一貫性は極めて深い調査が必要であり、明確に区別されたグループに単純に分割されるわけではなく、したがって、純粋な論理学者はそれらに関わるべきではないということが明らかになります。

以下はGroteの『アリストテレス』（第1巻、第192ページ）からの抜粋で、この四分法の起源を示し、その当時の科学と一致していること、そして現代の科学とは全く異なることを説明するのに役立ちます。「問題的命題と必然的命題との区別は、アリストテレスの心に於いて、既にこの章で触れられた彼の本体論と物理学における重要な特徴的な教義に対応しています。彼は我々が見たように、宇宙の大円域から外側の星間球まで太陽球までの天体物質が必然的かつ永久不変な存在であり、その回転と影響も永遠に一様であると考えていました。そして、この天体物質の善き影響により、月球球と地球中心（四元素とそれらの化合物を含む）の内面空間においては、規則性を促進する傾向が存在し、これを「自然」と呼びます。しかしながら、独立した不規則な力である「自己発生」や「偶然」によって部分的に制御され、これらは根本的に未知であり予測不可能なものでした。

アリストテレスによって名付けられたこれらの不規則な系列は、論理学における問題的命題の客観的な対応物です。これらの下月球系列において、未来について「あるいは起こるかもしれない、または起こらないかもしれない」ことが最高の知識でも達成できる限界でした。肯定または否定の確証はあり得ませんでした。一方で、天体物質の必然的かつ一様なエネルギーは論理学における必然的命題の客観的な対応物であり、この物質は単なる存在ではなく、必然的かつ変化しない存在でした...彼は論理学における問題的命題を、話者の無知を表現するだけの主観的なものではなく、何かより深いものと考えており、それはすべての人々にとって根本的に未知なものと対応しているものです。

14. この哲学が衰退した後においても、その基礎となった模態分野の区分は、デ・モーガンの指摘[10]によれば、中世時代において重要な役割を果たすことが可能だった。彼は言いますように、「人々はしばしば単一の証拠に基づいて一方か他方の判断を行う必要があり、十分な具体的知識がなければ、その主張を検証することができなかった」。古い論理学者は「週間内のどの日でも、カタイやタルタリーから肉と血で四輪を走る人や、地上に植えられたようなポリドロスのように、多くの驚異的な物語が証明される可能性がある」と考えていた。したがって、より良い推論がない場合においても、事実のうち哪些は受け入れるべきであり、哪些は受け入れるべきでないか（必然的であるもの、不可能であるものなど）を粗略に分類するためには便利だったかもしれない。そしてこれらの仮説（これらは現在ではただの一時的な推論として見なされるべきもの）を使って、彼が得た物語を検証することが可能だった。

不可能のクラスに属する命題は、ほとんどどの程度の証拠でも考慮に入れる価値がないほど大きな反対の前提を持つとみなされ、その他の推論も同様に扱われたかもしれない。

15. しかし、この古い四分法がモーダル性の区分は、すべての哲学的正当化が失われた後も、論理学者たちによって受け入れられ続け、継続されました。私が確認できる限りでは、カント以前の名実共に優れた論理学者はほとんど誰も、それらに対する重要な変更を恐れずに敢行しなかったのです。

ポート・ロワールの論理学も、カルナップ主義に基づいて設立されても、伝統的な主張を繰り返し、極めて簡潔に述べています。この古い形式への従順は、多くの場合において相当な矛盾や混乱を招くこととなりました。これらの形式は、我々が見たように、論理の領域についての客観的見解に基づいて構築されましたが、その見解は必ずしも多くの場合で厳格に遵守されていませんでした。

実際には、カントの影響が感じられる以前に、この見解は程度と方向によって異なる範囲で放棄されるようになっていました。例えば、多くの者は四分法に真と偽を加えることでリストを拡張しました；また時々、可能、仮定、確定も追加されました。これは、客観的見解から主観的見解への移行の傾向があることを示唆しているか、または少なくともそれら間で揺れ動きがあることを示すものであるように思われます。

16. カントの模態に関する見解は、多くの人々にとってほとんど馴染みがある。彼はこの観点から判断を、アポデイティック、主張的、問題的と分類している。これらの分類についての数と相互関係については後ほど述べるが、現在はそれらが示す一般的な見解に焦点を当てる。この点では、すぐに明らかになるように、大きな立場の変化が達成されたことがわかる。

「必然的」や「不可能的」は、命題の事実性を認定するために証明が必要とされていたが、一方で「アポデイティック」と「主張的」はほとんどどのような事実でも適用可能であり、これらは私たちの意識に訴えることで区別されるだけである。さらに、「主張的」から「問題的」への区別は完全に主観的かつ個人的なものであり、一人の人間であっても時間によって異なる場合がある。ある人が真実だと知っていることが、別の人は疑問視する可能性がある。

アポデイティックな判断は、私たちが受け入れるだけでなく、思慮の上で逆転させることができない判断である。主張的判断は単に受け入れられるものであり、問題的判断は疑問を感じさせるものである。

このような見解は概念論的またはカント派の論理観の必然的な結果である。多くの論理学者がこれを受け入れており、カントの追随者だけでなく、彼の影響を受けてきたほとんどすべての人々も同様である。例えば、ウーベルフはカントとの根本的な点で意見が異なるにもかかわらず、この見解を採用している。

次の問題は、模態の分類は何種類を認めるべきかということです。アリストテレス派や学術論理学者たちは、我々が見たように四つの分類を採用しました。これらのうちいくつかの関係は特に明白ではなく、特に可能と偶然の間の関係は解釈に非常に不明確な点であり、解説家たちもその点については意見が一致していません。しかし、これら四つの用語を使用して模態的な信憑度の段階的なスケールを示そうとしたつもりではなかったという点については、相当明確であると見ることができます。したがって、それらとの現代の確率論との対応は僅かなものであり、その議論は特定的または歴史的な興味に留まるべきものとなります。

デ・モラン（『形式論理』、第232ページ）は、学術論理学者たちは偶然を「おそらく」よりも「不確実」と定義し、可能を「おそらく」よりも「不確実」であると理解していたと言っています。この主張の根拠については私には分かりませんが、それにより彼らが現代の確率観念に近い考え方を持っていたという信用を与え、それは多くの後継者たちよりもはるかに近いものであることが明らかになります。そして、その多くは[12]。

私は相当な調査を行った結果、この問題に対する主観的な結論を導き出すことができました。それは二つの主要な見解が存在したということです。一部の学者たちは（例えば、ブルーガルディック、第1巻、第32章）根本的には模態は二種類しかないと認め、偶然と可能は等価であり、必要と不可能も同様であると主張しました。これは、主題の性質に焦点を当てた人々が自然に導かれる見解です。

一方で、表現形式の形態に焦点を当てた人々は、四つの命題形の類似性と、それぞれが明確な相互対立関係にあることを必要とする点から、四つの模態間の区別を強調すべきであるという結論に至りました。したがって、彼らは（例えば、クラカンスプー、第3巻、第11章）偶然は「現在存在するが将来存在しないもの」であり、可能は「現在存在しないが将来存在するもの」と主張しました。一部の者は物理的可能と論理的可能との区別に対応させる見解を採用したようです。

18. カントの分類に到達したとき、私たちははるかに明確な地盤に立っている。これらの用語の意味は非常に明瞭であり、また疑いの余地なく、それらが信念の段階的な尺度を割り当てる傾向があることも明らかである。これらが形而上学的な観点からではなく論理的な観点から見れば、その利点は多く存在する。もし我々が内省に頼り、判断自体のみの研究に限定し、それらが基づく根拠を除外すれば、確かに思慮深く明確な区別があるように見える。それは考えることすらできないほど信頼性が高い判断；その逆は可能だが、我々はそれを真実と受け入れる判断；そしてただ可能性として持つ判断の間である。

ただし、論理的な分類として見れば、カントの配置は私にはあまり役立つものとは思えません。現在考慮しているような論理的目的においては、実際には二つの分類に分解されるだけであるように思われます。

断定的と主張的な区別については、その区別が信念量ではなく、これらの判断を覆そうとする試みによって強制的に精神に与える影響の暴力性に関わる差異である点は、メタ物理学的または心理学的な理論に基づいてそれを受け入れる人々でも認められるでしょう。それぞれは完全な信頼を伴っていますが、一方だけは反駁可能であり、他方は反駁不可能です。

主張的な判断に対する信念は、疑問の余地がないほど確固たるものであるため、それ自体が問題的であるものとは異なる点から、断定的と比べて量的な信頼度において差異はありません。つまり、主張的な判断は尺度の上部にすでに指針が動かされている状態であり、これを断定的に変換するには、その位置を固定するだけであると考えることができます。残る論理的な差異は、問題的と主張的な区別に限られます。前者は私たちがどの程度疑問を持っているかの判断全般を含み、後者は確信のないもののみを含むのです。

したがって、可能性が適切に扱うべき範囲全体は、未分断で一つの compartiment に投げ込まれます。一クラスが境界線まで包含し、もう一方はその境界線だけを制限するという意味での「分類」ではほとんど話せません。したがって、この見解に基づいて実際的には、確率の数学的な研究者としては期待されるように、モダリティは完全に消えてしまいます。

19. 不十分且不系统之思想家，以及那些以巧思为显著特征的学者，
接受或提议了多种其他排列方式。当然，在可嘉地希望将我们的逻辑形式与普通思维和语言更好地协调一致的愿望下，这种尝试是有一定合理性的。
実践中，如前所述章节指出，每个人几乎无限制地承认各种模態形式，例如“可能”、“非常可能”、“几乎确定无疑”，等等。人们或许认为，通过忽略対某些这些形式使用技术等价词，我们可能会失去対某些可能推理類型的逻辑控制，从而在精确性上甚至落后于普通思维。

关于此類额外形式，我认为，所有由未受確率论影响的作者引入的形式，
除了制造更多的混乱之外，并没有做其他事情。因为这些作者并未尝试有序地排列他们的术语或确定它们之间的相互关系。
当然，我们已经提到了真实的和虚假的；可能的、假定的以及确定无疑的。
这些细分似乎在奥卡姆（普兰特，III. 380）早期阶段达到了顶峰，
他认为一个命题可能会因被“真实地知道、被知晓、虚假的、未知的、书写过的、延展的、概念化的、相信的、意见中的、怀疑的”而模態化。

20. 機率の科学が発展した以来、論理学者たちはより良い機会を得て、彼らが目指すべき目標を知るようになった。確かにその試みは完全には成功していないと言われても、それらは少なくとも前人との比較において明確な進歩であると言える。トムソン博士[14]は九つの分類に分けた。彼は、模態の度合いを昇順に並べ替えると、判断が可能性的、疑惑的、確率的、思考者自身に対する道徳的な確定性、特定のクラスや学派に対する道徳的な確定性、全般的な道徳的な確定性、物理的な確定性（限界あり）、物理的な確定性（無制限）、そして数学的な確定性であると述べている。他の多くの分類も同様に挙げられるかもしれないが、どの数学者でも認めるように、このような分類の統一性を確保しようとする試みは完全には不可能である。ここでは、数学的確率論の有益な影響を感謝すべきである。この理論が学習されるにつれて、明確な信念の段階を区別しようとすることは、連続的な量に亀裂を作ろうとするものであることが認識されたことだろう。微弱な推定から強い推定へと進む過程、そしてその推定から道徳的確定性への移行においては、我々は徐々に上昇し、自然の停泊点は存在しない。この連続性を証明する必要はここではないが、その材料は本作のほぼすべての章から得られるだろう。読者は、数と測定が可能な場合における我々の信念の基礎が、このような形で明らかにされていることを思い出すだけでよい。したがって、信念自体がその根拠から切り離されることなく、これらの特性が自身にも適用されるべきであることは明白である。

従って、古い意味での模態性を追求しようとした試みは、結局のところ放棄せざるを得ない。その試みが達成しようとしていた全てのことは、現在では確率論という理論によって比類なきほどよく達成できるようになった。確率論には無限に分割可能な数値尺度があり、これにより古い分類システムはすべて真の自然性を欠いていると見做されるべきである。少数の分類のみを採用するシステムは、全体的な可能性の範囲を一貫したクラスとして残すことが多く、多数の分類を採用するシステムは避けられない曖昧さと不確実性に陥る。

21. 純粋な命題と模態的命題との区別に対応して、より複雑で満足いくものではないのが、純粋な三段論法と模態的な三段論法の区別である。後者の場合では当然、前提が一方または両方とも模態性を有するとき、結論に対する影響が議論されたのである。このような考慮に達したときのみでないと確率論の適切な土台に立つことはできないが、このような粗雑な材料でほとんど何も成し遂げることはできず、全体的な模態的システムの固有の笨重さと複雑性はここから明らかになる。おそらくこの複雑性に関連して、学問の人々や他の人々によって記録された厳しい言葉[15]が発せられたのかもしれない。

アリストテレスはこの主題について複雑な調査を行い、その追随者は当然類似の道を歩むことになった。この章で僅かな概要に当たる私の目的からは、これらの探求を説明しようとすることは全く適切ではない。アリストテレスの模態論と現代の確率に関する見解との間の関連は確かに存在するが、それは非常に薄いものであることが指摘されているからだ。四つの模態を考慮に入れるだけでも複雑だったものが、「可能」や「真実」、「偽り」もリストに加えなければならないとすれば、さらに耐え難くまで煩雑になることは明らかだ。以下に示す例[16]は、論理学者がどのような議論を行ったかを示している。「一つの前提が確定的で、もう一つが可能である場合、確定的な結論を推論することはできるだろうか」、「不可能なことが与えられたとき、必然性を推論することはできるだろうか」、「一つの前提が必然的であり、もう一つは『欠如』である場合、結論が必然的であることはできるだろうか」といった議論が無限に続く。

22. カントの模態観に基づく議論では、このような種類の三段論法の議論が一層単純化され（ここでは三つのモードのみが認識されるため）、厳格な確率とより密接に結びつく（モードは信服度の近似度合いに近い性質を持つため）。しかし、一方で、これらの導入にはそれほど正当化の余地がない。論理学者たちは、彼らが拙速な装置によって達成しようとしていることが、確率が最高の精度で実現できるものであることを知っているべきだという点から、そのような導入は少々正当化されないかもしれない。前者の方法は、後者の方法と比べて粗雑で正確性に欠けるものであり、古代バビロンの夜番の最も粗糙な測定法が現代の天文学者の精密な計算に比されるようにである。実際には、カント論理の一般的な支持者の中でもこのような考慮を進めるのは一部だけであり、ハミルトンやマンセルのような者はそれらを完全に拒否していることを我々は知っている。そのような主題について扱う人々は、以下の結論を下す。つまり、両前提が断定的であるとき、結論も同様であり、両前提が主張的または問題的であるときも同じである。一方で、一方の前提が断定的で他方が主張的である場合、後者が「弱い」ものとして結論に認められる。そのような規則は、ウーベルの『論理学』[17]において英語読者にも説明されている。

23. しかし、これらのモーダル性は正確な思考の道具として見なされるときには、確率の精密な算術表現によって置き換えられても、私たちの一般的なモーダル性表現が改善され適応され、より正確な使用に適応できないかという問題は残ります。確かに、それらが認識する量の大きさについて、根本的な区別を設ける試みは無意味であることは真実です。その量の大きさは我々が述べたように連続的であり、しかし正確な用語の重要性と可能な限り数値的な区別の認識の重要性を考えると、モーダル性表現の使用について合意が得られれば実質的な進歩となるでしょう。

私たちは既に第2章 §16で（参照）ガロント氏による提案を注意しました。それは、多くの種類の量が平均値を中心に規則的に集まりやすいという事実に基づく自然な分類システムの可能性についてです。例えば、「良い」という用語は最も優れたものから下に測定した最初の四分位に適用されることが提案され得ます。「中立的」は中央半分、そして「悪い」は最後の四分位に適用されます。

似たような理由で、改善された用語がいつか普通の生活における一般的なモーダル言語に導入されることも無理ではありません。例えば、「非常に不確実」という表現は可能な限り99対1以上のオッズを持つ出来事に限定されることが合意され得ます。その他の類似した表現についても同様です。

確かに、このような分類の適用には困難が伴います。なぜなら、すべての分類の適用において我々は規則作成の能力と規則に基づく分類の能力という二つの障害を乗り越えなければなりません。つまり、クラスを決定したとしても、数や測定しにくいものに対して偶発的な事象がどの程度のオッズであるか判断するためには多くの疑問と議論が生じるでしょう。

確実なことは、事件に対する勝率や敗率を知っている場合、それらを明示的に述べることは常に可能であり、最初に語彙の使用について合意する必要はありません。しかし、そのような方法で述べるとしばしば迂回と冗長性が生じ、また模態的な用語が実際的な使用中である限り、それらが一般的に理解されるべき確率の程度を定めるために合意を達成すること自体は害がないだけでなく、非常に有益であると考えられます。ベントhamは、この問題についての日常生活の言葉から正確なものを得ることは絶望的だと感じ、数値基準への直接的な訴えを支持していました。彼は、裁判所での証人の証言に対する彼の信念の程度を示すために0から10までの刻みのある装置を使用することを提案しました。同様に、判事も自分の結論の強度を表現することができました。しかし、このような数値スケールの使用は任意であり、ベントhamは多くの人々がそのような方法で信念の程度を測定することに迷う可能性があると認めています。

この章を通じて、模態語を古い論理体系によって提供された現代的な確率の最も近い類似物として考慮してきました。それらをそう考える理由は、それらが我々の信頼度の程度における若干の段階を認識し、測定する試みを一部表しており、そのような考慮事項が論理的推論にどのように影響するかを調査することを模倣したからです。

しかし、確率の方法がその種子を含むのは模態論の中でも、当科学の真の対象、すなわち最も適切に扱われるべき対象のクラスは、学術的ロジックの別の部分によってよりよく表されます。これは、古い意味での「辩证法」が通称されていた分野でした。アリストテレスによれば、辩证法は修辞学と一種の姉妹芸術であるように思われます。その主な対象は、論理三段論法であり、その前提が確定的ではなく確率的なものであるだけ異なりました。このような前提をトピックスと呼び、それに関連する三段論法をエンティメムと呼んでいました。これらは「兆候と可能性」から始まると言われていました。

25. このような推論を一般的に説明する用語は、現在の意味での確率に適用されるものに非常に似ています。確率性や確率的な三段論法という言葉を聞くと、我々はまずその結論が類似した性格であると考えるかもしれません。しかし、これは誤りです。最初に、この辩证法の領域はあまりにも広く、科学的または材料的帰納法の全範囲もカバーしていました。辩证法の前提の特徴は確定性の欠如であり、確率（私はしばしば強調したように）が考慮するのはそのうちの一類だけです。一方で、帰納法は別の一類に関心を持っています。また、前提の不確実性についての最小限の議論も行われていませんでした。この問題について探求を試みるときのみ、我々は確率の領域に入ったと見なされることが可能であり、模態がこの問題に取り組んだという理由から、彼らが私たちの専門的な対象に関与していたと考えられています。

26. 私が多少の知識を持つ古い論理学の中で、クラカンサープの論理学はこの主題について最も詳細な議論を展開している。彼は論理学の対象となる三段論法の扱いを二つの部分に分けている：「証明的」三段論法と「蓋然性」三段論法である。後者には一冊の本が完全に捧げられている。この中で、十三種類の異なる「位置」（loci）[20] の性質と結果は調査されているものの、これら全てが「蓋然的」として扱えるという意味での解釈は非常に明確ではない。

旧来の論理学者が現代の確率に関心を持つような前提を所有していたら、かつ自身の方法でそれらを扱っていたら、彼らはその前提をこれらの「位置」の一つとして配置し、したがってそれらの考慮を自己の辩证法の一部とする必要があったに違いない。しかし、彼らがここでの唯一の試みが前提と結論の蓋然性を認識するだけであり、それ以上の努力を払わなかったという点においては、その部分の論理学は現在の確率概念に対する関連性は非常にわずかであると考えざるを得ない。私には、これよりもっとも近いものとしては、この章の冒頭で述べた方法の一つに過ぎないと考えられる。

27. 論理学が信念と証拠に関する質問を直接かつ明確に扱う唯一の科学ではないように、そのような困難は他の場面でも知られてきた。模態に関して言えば、これは法曹界において明らかにそのような状況であったようである。したがって、この主題に対するこの応用について、専門的なトピックを触れる立場にある Layman にとって適切な簡潔さの中で、いくつかの注記をここで述べることは便利であると思われる。

一時的にモダリティの本質を思い返してみよう。
これらは、数値記号を使わずに異なる信頼度や信念の程度を区別し、その影響を受けた前提が結論に与える結果を決定することであると理解している。

さらに、何か種類の尺度を作成または保持するためには、測定の基準となるものを使うことが必要であり、そのためモダリティにおける確証度の基準を達成し認識することは、同じ調査の中でほぼ不可分に含まれている。この意味での用語においては、法の領域でも模態的な困難が明らかに現れたことは事実である。

類似の試みが法の領域でも遭遇した類似の困難により、私たちの承認度をどの程度の尺度で配置するかについての明確な合意に達することを目指して努力してきた。当然のことながら、専門家に限定される特定の科学分野では、他の外野者がほぼ同様の権利を持つ科目と比較して、このような合意を得るためにはより実現可能である。

前者の場合での適用範囲は狭く、専門家の習慣や伝統により基準が相当程度保たれる。しかし、法曹界においても、この方向で非常に目立った成功が達成されたとは断定されないようだ。

28. 科学的、厳密に法的形態においては、その難しさは主として、証拠を信憑性の程度に基づいてクラス分けすることへの試みを通じて現れました。これは、私が理解するところでは、ローマ法においても、そしてローマ法から発展した大陸諸国の法的体系においても同様でした。「多くの証人の直接証拠は『全証』でした。次に『半分以上証』が来ました、その後『半分以上大証』と『半分以上小証』がありました；一定の数の半証を加えることにより、例えば商売人帳簿の提出とその補足誓約を通じて全証が成立させることができました。この原則に基づいて拷問は告白を得るために使用されました。告白は、その使用が正当化されると思われる状況に対する補完的証拠でした」［21］

ベントハムによれば、イギリスの学派における対応する尺度は以下の通りでした：直接証拠、強力な推定、確かな推定、軽微または冒険的な推定。ブラックストーンを含む他の者たちによって認められることは認められていますが、私はこれらの分類が現在では全く広く受け入れられているとは理解していません。

29. 以上では、模態的三段論法に思い起こされる点があります。単純な模態命題の根底にある困難が現れる主な実用的な形は、司法的確信を得るための基準を特定することの試みです。「確信」という言葉はここでは、当然のことながら、メタ物理学的に「アポデイティック」であるものとは意味が異なります。実務においてそうしたものは希少または決して確保できないからです。ここで言う「確信」は、この上ないほどまでの conviction ではなく、合理的な人間が全ての欲求を満たすのに十分だと感じる程度のものです。

再度考えれば、正確な思考家たちには、これほどの探求が絶望的なものであるように思われるでしょう。連続する量の中で自然な断層を探す試みに匹敵するものだと言えるでしょう。確かに、限られた知性の鋭い且つ熟練したクラスにおいては、確信の基準を含むあらゆる他の基準も同様に正確に保持し、伝達することが可能であることは疑いようがありません。これは好みや意見の問題で、個人的な判断の特異性が不一致と混乱を引き起こす可能性が高い分野において特に可能です。

しかし、そのような合意はほとんどが技術的かつ慣習的なものであり、現在の場合には比較的未熟な人々が参照できる基準が必要です。したがって、この基準は単に記憶によって保持されるだけでなく、特性から一般に認識されなければなりません。このような基準が確保された場合、その重要性を過大評価することはまずありません。

しかし、弁護士の演説、判事の指摘、陪審団の判決から推察すると、実質的に名にふさわしいものは決して達成されていないように思われます。

30. 譲渡の可能性としては、裁判所が被告の罪を疑念しない限り起訴する義務がないという頻繁な保証が最も近いものかもしれません。全疑念を排除し、何らかの反対仮説を提示することも不可能であると想定される絶対的な確信は、人間の事務において達成不可能です。しかししばしば言われますように、彼らは「合理的な疑念」がない場合に起訴すべきであり、つまり日常生活における重要な事項に行動する上で障害となるような疑念がなければ起訴すべきであるとされています。このような不達成確実性を求める警戒は、非常に有用かもしれません。しかし、私たちが日常生活の重要な事項において行動する際に基づく確信は固定された基準ではなく、その内容によって極めて異なることがほとんどです。報酬が大きいほどリスクを冒す準備ができることも多く、逆もまた然りです。どの程度の確信があろうとも、適切な状況下ではそれに基づいて行動することに準備ができていない度合いはほとんどありません。

実際には、多くの著者は法的手続きにおいて一般的な意味での基準が存在しないことを完全に否定します。例えば、Sir J. F. Stephenは、彼の『英国民事刑法』の中で、最後に指摘されたような基準を注意し否認した後、「私たちの法律によって認められた唯一の基準は何ですか？それは法により証拠として提出されることが許可され、十二人の選ばれた人物（裁判員法に基づく）がそれを聞いてその事実を確信していると述べるためのものです。彼らは偏見を持ち、恐れ多く、冒険的であり、知識不足であるかもしれません。しかし、誓約、人数、および財産資格は可能な限りこれらの欠点を中和し、価値や長さの基準に適応する条件に対応しています」と結論付けます。（p. 263）

これを認めることは、その基準を達成不可能であると放棄することとほぼ同じことである。陪審員が有罪判決を下すための証拠は、私や他の人々にとって「合理的に確信すべき」ものとして考慮する基準となるかもしれないが、もちろんこれは陪審員自身に対する基準とは提案されない。もし彼らの判断が一貫性と統一性を持つのであれば、外部からの指示が必要である。人間が「どのように」確信を感じるべきかを尋ねるとき、上記のような答えを与えることは、彼はその程度に確信を持つべきであることをただ教えることである。もしも陪審員が密接な専門職で構成されているのであれば、先述のように伝統的な基準を保持することができるかもしれない。しかし、彼らは普通の一般市民から選ばれた人々の選抜であり、過去の決定を将来の行動の指針として提出することはほとんど不可能である。

31. したがって、英法では、論理的模態学説とその法的な類似物が基礎とする主な仮定、すなわち異なる確信の度合いが十分に正確に区別できることを前提として、明確に拒否していることが明らかである。そして、確信の基準に関する付随的な問題については、私たちが「合理的な確信」以外は何も望まないが、それ以上は必要ないという判断や推測を下すだけであると表明するだろう。

この声明は、論理学者や科学者にとって標準的なものとみなすことができ、彼らはその正確性に全くの理由で異議を唱えることができる。しかし、厳格な理解性さえ考慮すれば、その正確性には全くの主張もないと主張することができる。

もし人間が現在の確信度が合理的であるかどうか知りたいのであれば、彼は何処に訴えられるべきか？彼は自分の精神状態を「重要な人生の事柄」で経験されるものと比較することはほとんど不可能であろう。既に述べたように、これらは固定された価値を示すものではない。同時に、このような表現が全く意味を持たないとは思えない。人々は特に重要かつ興味深い問題について言語を使用し続けることは、何か意味を持つものでなければならず、それが無意味であるとすれば、その使用は継続されないであろう。

したがって、我々は「合理的な確信」が粗略ながら伝統的な基準を反映していることを結論付けることになる。既に述べたように、これはどのクラスの専門家にとっても完全に実現可能であり、従ってしばしば密接に関連する人々にとっては不可能ではない。日常的な生活のためにこのような表現を使用したとしても、そのような慣行がどれだけ近い程度で一般的な基準を達成しているかは信じ難く、しかし陪審員が直接裁判官の影響下に作用して「完全に確信している」と断言する場合、それは初めから考えられるよりもより固定された基準に近づいている可能性がある。

それでは、この基準は何らかの方法で確認できるのでしょうか？つまり、「合理的な確信」が陪審員によって確保されるべき実際的な価値は何かを決定するための手段は存在するのでしょうか？公式な声明がない場合、論理学と確率論の学生は自然に二つの方法に訴えるでしょう。これらについて一瞬だけ触れることで、この調査を終えましょう。

$$32.$$ それでは、このような基準を確認するための手段は何らか存在するのでしょうか？つまり、「合理的な確信」が陪審員によって確保されるべき実際的な価値は何かを決定する方法は存在するのでしょうか？公式な声明がない場合、論理学と確率論の学生は自然に二つの方法に訴えるでしょう。これらについて一瞬だけ触れることで、この調査を終えましょう。

$$32.$$ それでは、このような基準を確認するための手段は何らか存在するのでしょうか？つまり、「合理的な確信」が陪審員によって確保されるべき実際的な価値は何かを決定する方法は存在するのでしょうか？公式な声明がない場合、論理学と確率論の学生は自然に二つの方法に訴えるでしょう。これらについて一瞬だけ触れることで、この調査を終えましょう。

これらの最初の試みは、裁判所の確信度の基準を間接的に決定することを目指す。これは単に、例えば陪審団の決定が正しかった頻度を決定することで達成される。この作業が不可能であるように思われるかもしれませんが、実際には、証言や判決の確率に関連する多くの問題に対する多数の恣意的な仮定により困難化されているのであり、理論上の要素に欠けているためではない。私たちは考慮している裁判所を再審する無欠陥の上級裁判所を持つ必要はなく、私たちがテストしたい裁判所と類似した材料で構成される大きな数の通常の代表的な判例を提出すれば十分です。

ただし、常に人間による判断について深刻に影響を与える事項に対する判断が「独立」であるという怪物的な仮定を立てることを前提として、一致と不一致の統計学が私たちの目的を達成する。たとえば、特定の人数の陪審団が特定の多数決で決定した場合、その判決は10回中9回正しかったと言えるかもしれません。この種の結論についてフランスの裁判所に関してポissonは、彼の大作『判決の確率』の終了部分で試みたが、私は彼がその結果に大きな数値的正確性を付与したとは思わない。

次に、わずかに希望が持てる手段として、法的「推定」のいくつかのケースを参照することを考えるのもありです。「結論的推定」とは以下のように定義されます：「結論的、または他では強制的あるいは絶対的な法的推定とは、事実が異なることを証明するいかなる証拠によっても覆せない、特定の主張を支持するために必要な証拠量を決定する規則です。」[27] これらの推定は多くの教科書で説明されており、しかし大部分は非常に曖昧な事項に過ぎず、統計的な頻度で把握することは難しいようです。実際、一部の権威者は、これらの推定が確率の程度を指定すること自体ではなく、すべてが疑問視されると問題となる遅延を排除するためのものであると主張しています。

それでも、これらが確率の程度を指定したとしても、それは立法者や判事が合理的だと考えたことを示すよりも、陪審団自身がそう考えるようになったことを示すものではないでしょう。

ある人間が消息を絶ってから、彼が死んだと推定されるまでの時間を想定する仮定は確かに存在します（再登場によって否認される可能性もあります）。もしこの時間が関連する人の年齢に依存していたら、我々が望むような基準が得られることでしょう。生命表を参照することで彼の生存期間の確率が決定され、その確率が法に適合する尺度を间接的に定義します。しかし、実際にはそうではありません；選択された時間は年齢とは関係ありません。最も近い例（それもあまり意味はありません）として、女性が子供を産むことができなくなったと推定される年齢があります。これは53歳でした。

この場合の確率的な見込みについては、Quetelet の『社会物理学』（第1巻、第184ページ、注記）に近いアプローチを見ることができます。彼が引用する当局によれば、約5500回の出産のうちで一度だけ、母親の年齢が50歳以上であることが示唆されます。これは、50歳以上の妊娠の確率を完全には定義しないかもしれませんが、早期に家族を持つことによって将来の子供を持つ確率が多少低下するという事実を考慮に入れると、その程度の確率的見込みを示しています。

これらのヒントは、司法的な確定度の尺度を測るためのものであり、実質的に何らの価値もありません。これらは法が完全に一貫した科学的証拠の体系である場合における問題の論理的な解決策の可能性として、またはそのような解決策を求めることの方法について理論的にはどのような形で探求されるべきかを示す指標としてのみ注意に値するものです。もし英法がこれらの事項に関して完全に一貫した科学的証拠の体系であれば、このような状況下では、その手続きから正確な確定度の尺度を導き出すための努力を行うことになりますが、それを公然と表明することを拒否しています。

1. \_形式論理学\_, 第232頁。

2. この文脈は、ウィーティリーの『論理学』（第2巻、第4章 §1）に非常に混乱した部分にあるいくつかの声明の意図と見なされる。"モード付きの命題を純粋な命題として述べることは、モードを一つの語彙に付加することで可能であり、その命題は前記のすべての規則に従うことになる；...「全ての知識が有用である可能性が高い」、「おそらく有用」というのはここでの述語である。"彼は明らかに、次の注釈で述べられている真偽のモードの区別を引き立てるような区別を描いていないように見える。驚くべきことに、ハミルトンすら真偽のモードを同じ扱いにするというような区別を置かない。彼は「これらについては［前者について］、状況は全く同じである；モードは述語の一部に過ぎない」と述べている。（『論理学』第1巻 257ページ）

3. 当然、私が指しているのは、「AがBを不正に殺した」のような例で、AがBを殺したことについて、時として単純な断定ではなく、修正された断定と見做された場合があるというような例である。（ハミルトンの『論理学』第1巻 256ページ）明らかに、このような状況における修正は本来的には述語の一部に過ぎず、「AはBを殺した」および「Aは不正にBを殺した」という文法的な区別が存在しない。実際、一般的な名詞であるモードを拒否するあまり保守的だった論理学者たちは、偽の種類を「材料的モード」と呼び、真の種類を「形式的モード」と呼ぶという共通の方法を採用した。前者は、修正が本来的には述語または主語に属すべき場合すべてを含む；後者は、修正が述語と主語との実際的な結合に影響する場合のみを残すためのものである。（ケックェルマンの『論理学体系』第2巻 第3章）私は信じるところでは、これは一般的な神学的区別であった。

否定接続詞が関係項に属するか述語に属するかについての議論に関する詳細については、ハミルトンの『論理学』第1巻 253ページを参照されたい。（訳者注：原文は英語で「For some account of the dispute as to whether the negative particle was to be considered to belong to the copula or to the predicate, see Hamilton's \_Logic\_, I. 253.」）

4. 彼はこの主題について短い議論を他處でも行っている（『討論集』第2版、702頁）。その中では、やや異なる見解が述べられている。実際、モードはここでは論理に認められるものの、それは genus と species の関係によって分割される範囲内でのみ認められることであり、もちろんこれはそれらを完全に拒否することと等価である。その場合、これらは他の genus と species の関係の例とは根本的に異なる点で差異がない。

5. 『手紙・講義・評論集』45頁において彼は、私にはやや異なる判断を与えたように思われる。

6. この命題は以前の章で扱った比例的命題の一つではないことに注意しなければならない：ここでは21個のYが存在し、そのうち正確に18個がXであると意味されるべきであり、平均的に21個のうち18個がそうであるというわけではなく。

7. さらに先の320頁で述べた通り、私は古い論理学における確率的三段論法や言説的三段論法の扱いは、これとほぼ同じものになったと考えるものの、それらは異なる視点から問題を捉え、非常に異なる言葉で表現したに過ぎない。

8. しかしこの区別は決して完全には無視されていない。例えば、Smigleciusは確実性と必然性のモード的影響について議論する際に、「確実性は認識に対して関連している：しかし必然性は事象において存在する」（『討論集』第13回討論、第12問）と言っている。

9. またWhately（『論理学』第2巻、第2章 §2）は必要性、不可能性、偶然性の素材について話しているが、それらが完全に古い見解から属するものであるという疑念を示していないように思われる。

10. 『形式論理学』233頁。

11. この主題はしばしば全く省かれた。例えばWolfはそのような議論を全く行わなかったが、確率的命題や三段論法については多く述べており、Leibnitz同様に「確率論の論理」を何か新しいものとして望んでいた。私は彼が確率に関する著者の影響を受けたと考える。その中で重要なものはほぼ全て一文（『理性哲学または論理学』§593）の中で言及されている。

ラمبرトは独自の立場を占めています。数学的概念や記号が関与する多くの他の分野でも、彼の論理的な態度は完全に伝統的とは異なります。たとえば、彼の『新しいオーガノン』における「Von dem Wahrscheinlichen」の章をご覧ください。

12. プランツによる精密な論理史において、この主張には最小限の権威が見つかりません。

13. "Hi quatuor modi magnam censeri solent analogiam habere cum quadruplici propositionum in quantitate et qualitate varietate"（ウォリスの『論理入門』、第2巻、第8章）

14. \_思考法則\_、§118

15. "教養が豊かな学問家たちがそのような適切な材料からどれほど複雑化させたかを、アカデミックな場所としてここに汗を流した。模態論のasinusは味わわないだろう。" ケッケルマン、『論理体系』、第2巻、第3章

16. \_スミゲリウスの論文集\_、インゴルスタット、1618年

17. プランツによる『論理史』（オカムとブリジダン下）で、これらの学問家たちがそのような適切な材料からどれほど複雑化させたかの記述をご覧ください。

18. "εἰκòςとσημεῖον自体は命題であり、前者は一般的な確率を述べる、後者はよりまたはそれほど確かな指標であることが知られている事実を述べる。前者はほぼ普遍的な一般命題で、「大多数の忌妬する人は憎悪を持つ」などです。後者は単独の事実や一般的な信念に関する命題であり、しかしそれは符号とみなされるのは相対的に、他の命題から推論できるものであると仮定されるのです。（マンセルのアルディッチ、付録Fでアリストテレスのエンティメムと言語的三段論法についての説明があります。もちろん、グロートのアリストテレス『トピックス』などでも参照できます）"

19. 「この証明はまた、確実な証明と推論の間で異なる点がある。確実な証明では結論と前提が必然的になり、一方、確率的な推論では結論が確率的であり、前提も確率的になる」（Crackanthorpe, 第5巻、第1章）；これはデ・モrganが既に引用した文脈での言葉で、論理と確率の区別を述べているものである（参照：第6章 §3）。

おそらくこのような見解の発展がレービンツが期待していたものだったかもしれない。「私は一度以上述べたことがある。新しい種類の論理が必要になるだろう。それはアリストテレスの『トポイクス』ではそれほどしなかったことだが、確率の度合いを扱うものである」（『新試験論』、第4巻、第16章）。彼が現在の数学的確率理論を考えていたのかもしれないが、科学の幼少期においては、特定の主題が明確に考慮されているかどうかを断定するのは当然難しい。レービンツ（トドハンターによる歴史によれば）はまだ議論されたことのある機会問題について最大限の関心を持っていた。

20. 「loci」とは一般的な前提のクラスを指した。実際、これは主前提に対してカテゴリーまたは予備状態に相当する関係があった。クラカンスプは次のように述べている。「しかし確率的な推論から推論される可能性のある場所から導くことは、真に推論的確率の特徴である。そしてこの点で証明とは異なるのは、証明者はただ4つの場所とそれらの必然的なものだけを使用するからだ...しかしこれら以外にも、確率的に推論できるための多くの他の場所がある；それは種類から、個体から、付加項から、対義語から、その類似物などである」（『論理学』、第5巻、第2章）。

21. スティーブン『英刑法総説』、第241頁。

22. 『裁判上の証拠の原理』、第1巻、第6章。

23. カント派の論理学者によって主張されているとしても、--真実犯人であると疑われた者に対しては、彼自身が罪を認めた後でしか法的罰則が執行されてはならない。なぜなら、証拠や他の告発が彼に対するものであっても、その反対の可能性は常に残されるからである"（クルグ『論理学』§131）。

24. ジェイムズ・ジョン・モンロー氏は次のように述べています。「ある男が娘を殺害したと疑われているとします。普通の陪審団で有罪判決を得られないほどの証拠でも、大陪審団では真実の起訴状を提出するかもしれません。そのような証拠でもないものでも、 coroner（裁判官）の陪審団は彼に反対の判決を下すかもしれません。さらにそのような証拠でもないものでも、衡平法裁判官はその男が法院の被保護者となる監護人に任命しないかもしれません。そして、そのような証拠でも影響を与えないものでもあるとき、父親は死床でその男を自分の娘の監護人に選ぶ前に二度と心配するかもしれません」。

25. この作品の証明の性質について一般に、および司法上の証明について特に論じている部分は、すべての論理学者にとって読み物として非常に価値があります。しかし私には、著者は証明を主観的なものと考える方向にあまりにも進んでいるように思われます。すなわち、人々が満足するものではなく、彼らが満足すべきものとしての証明を考えているからです。彼は立法的確実性の基準と価値の基準を比較しています；後者は一定量の金に相当し、その金属の希少性や一般的さに関係なく一定であると定義されています。同様に、人々がより信用性を持つようになったとしても、基準の内在的な価値は変化するでしょう。

26. この問題は将来的な章で詳しく議論される予定ですが、ここでは若干の説明を加えるために一言述べておきます。証拠を最も単純な要素にまで還元し、二つの平均的な判断精度を持つ裁判所や判事のみを考えます。これを x と表記します。その一致する確率は x2 + (1 - x)2 となります；なぜなら、彼らが両方とも正しいかまたは両方とも誤っているときには一致するからです。この一致の統計的頻度が分かれば、すなわち最初の判断が第二の判断によって確認される頻度が分かれば、x を決定することができます。

27. テイラー关于证据的部分：引用の後半は非常に不明瞭です。

\newpage

\chapter{確率論における諸誤謬（Fallacies）}
\label{ch:xiv}

\begin{center}
\large\textit{事後判断の誤謬・無作為と作為の混同・倍額賭け（マーチンゲール）の破綻・大数の錯覚}
\end{center}
\vspace{4mm}

第十四章.

論理の誤謬

1. 論理学の著作では、一般的に誤謬についての一節が設けられ、つまり論理の規則を違反する様々な方法について説明と分類を行う。確率と論理の類似性は十分に高いので、ここでも同じ計画を採用することが望ましい。著者が自らの見解を述べる際には、当然ながら他の誤った見解も説明し批判する必要がある。しかし、広く流布している最も一般的な誤謬は、主として曖昧な一般世論の誤解にとどまるものが多く、具体的な支持者はほとんど存在しない。したがって、これらの頻繁に発生する誤謬をいくつか集め、可能な限りその源まで追跡することにより最善の配置を見つけることが最も適切であると考えられる。しかし、規則的な分類と整理を行うことは大変手間がかかり、あまり価値がないと思われる。

したがって、この作品の専門性に従って、論理的関心に関連する問題や、確率を倫理学や社会科学への適用に関連する問題に限定して主に取り組むことに決定する。ただし、原則的な誤りが含まれている場合を除いて、機会と技術の問題について孤立した議論は避ける。

2. (1) これには、既に数回言及され、かつ前章で部分的に検討された最も肥沃な誤謬の源がある。これは、事件が起こった後にその属するクラスを選択し、従ってその事件の希少性と予測の難しさを判断し、その後その経験的印象を事前に想像した場合に移すというプロセスである。この過程自体は完全に正当化されている（必要でないとしても）ので、時間は厳密には確率論における問題に入らないからだ。したがって、このようなクラスを選択する際の注意を払うことで誤謬は生じないが、その注意はしばしば無視される。

ここで例を挙げて説明すると、ある男が扉に描かれた小さな目標の中心に穿たれた穴を見せて、「百ヤード離れた場所から古い猟銃で射止めた」と述べ、その声明は正確でしたが、重要な事実を隠していたという驚きを与えた。彼の主張は正しかったが、発砲当時の具体的な事実を抑えていた。実際には、男は barn のドアに一般的な方向性で目を向け、命中したのであり、その後でターゲットはその位置周囲に描かれた。このように虚偽の説明が他にもしばしば無意識的に行われていると考えられる。我々は類似の原理に基づいて事象を判断し、合理的ではないほど驚きを示し、事象の発生を決定する際に、実際にはその後に追加された私的な根拠に基づいている。

バトラーが十二章で既に議論した「キセウスの物語」についての言及も同様である。彼は歴史から一連の事象を選択し、その事象を予測できた人物を想像する。私は既に指摘した通り、特定の証拠がない状況で適切な予測をする確率が低いことは一つのことだが、一部または完全に証拠に基づいた物語の真実性を評価することは全く異なることである。しかし、これらの視点のいずれかに対する心象を他の方程式に転移するのは大きな誤りである。これはターゲットを描き終わってから、射撃が中心に収まっていることに驚くようなものだ。

3. この論理的な誤謬の一つは既に議論されたが、しばしばこの問題に関する困難な点を明確にするために、より一般的な形で再検討することで役立つ。

当論議の範囲における例では、一般的に個々のものが特定のクラスに明確には指定されていないが、適切なクラスを選択することには大きな困難はありません。現在、六面骰子四度連続で六目が出るというような出来事について考えているとしましょう。このような投げ方は非常に稀な出来事として扱われ、その起こる確率は36 × 36 × 36 × 36 - 1対1、すなわち1679615対1となります。これらの表現の意味は、頻繁に指摘されている通り、単純にいうとその出来事は非常に稀に起こるということであり、数値的に言えば1679616回中1回しか起こらないということです。

4. しかし、現在、その投げ方が実際に起こったと仮定しましょう。六目四度連続が出たということが既知または報告されている状況で、その投げ方を観察することを考えましょう。同じ表現、「非常に稀な出来事」という言葉はしばしば使用されるが、この表現は異なる意味を持つことがあります。

（1）まず最初に、最も正確な意味があります。確かにその出来事が起こり、それが何であるかを知っているため、確率の法則を利用する必要はありません。しかし、私たちはそれを知らない人として自分自身を想像し、確率だけを頼る立場に立つことができます。このような投げ方が1679616回中1回しか起こらないという事実から、私たちが予想したとしても、その正しさは同じ回数であることが示されます。ただし、これは統計的な根拠に基づく公正な予想であると仮定する必要があります。

5. (2) しかし、別の非常に異なる概念が時として導入され、特にその事件が経験の証拠によってではなく、証人の報告によって知られている場合に特にそうなります。この場合、「事件に対する確率」（第12章で指摘したように）とは、我々がそのイベントを予想する際に正しく予想できる回数の比ではなく、証人がそのイベントを報告する際に正しく報告できる回数の比を意味することがあります。当時の推論の基礎はここで新たな土壌に移されます。以前の場合では統計データは投げられた回数とそれらの頻度でしたが、現在の場合では証人の声明とそれらの真実性が基盤となります。

6. (3) さらに、人々が特定のイベント（質問の投げ方）に対する確率について話す場合にも別の意味が時として意図されます。彼らは「ここに一事件がある、私はそれを予想する際にどのくらい頻繁に正しく予想すべきか？」という問いではなく、「ここに報告がある、それはどのくらい頻繁に正確であるか？」という問いでもなく、何か異なることを意味します。すなわち、「ここに一事件がある、そのイベントはどのくらい頻繁に特定の原因によって生じるか？」ということです。

例えば、ある人がサイコロが連続して同じ目を出しているのを見て、非常に不思議だと述べる場合、彼はしばしば単に自分の予想が正しくない確率や報告が真実である確率だけでなく、その投げ方が公正なサイコロによるものかどうかという問題も考えていることを発見します。もちろん、このような質問を確率の観点から扱う理由は全くありませんが、適切な統計データがあれば、そのような観点から議論することは可能です。

これらの統計データは、特定のサイコロの投げ方や特定の証人の報告ではなく、問題の一つである投げ方が公正に行われた回数と、公正に行われなかった回数を含むものとなります。この見解に反対する主な理由は、このような統計データが得られないこと、そしてもし得られていたとしても、個別の状況に基づいて意見を形成することよりも、平均値に基づく議論を好むべきであるという点です。

7. 読者には、この区別を説明したところの例を容易に思いつくことができるでしょう。しかし、私たちは一方だけを簡略に検討します。私は大きな図書館の一冊の本の中に金券を隠し、部屋を去ります。そのとき、ある人が私に伝えます：私が出てから、他人が入って、その特定の本に向かい、それを持ち去ったということです。この報告を聞いて多くの人々は、非常に不確実であると答えます。しかし、「極めて不確実である」という表現を分析すると、おそらくこの叫びには、上記の意味の2つ、もしくは3つが含まれているかもしれません。

（1）その意味はこれで、つまり、報告が真実であり、 strangerが無罪であるという仮定のもとで、稀な出来事が起こったということです。多くの本がそういった形で持ち去られても、特定の本だけが選ばれたことは珍しいことです。そのため私はその出来事を予想していなかったので、それが実際に起きたときには驚きます。しかし、人間は驚く権利がありますが、（当面のところこの視点に限定する場合）、その驚きを事実を否定する根拠とするのは論理的な正当性がないということです。これはこの章の冒頭で説明した誤りと同じです。私が予想していなかったこと自体が、それを知らされたときにそれが間違っている理由ではないのです。

（2）あるいは、報告の事実を疑わずに、その報告自体に確率論の法則を適用することができます。そうすると、「私目前のこのような話は一般的には偽である」という意味になります。そして、そのため私はそれを信用することができないでしょう。

（3）あるいは、報告の真実性を受け入れつつ、 strangerが本を無作為に選んだという事実は疑うことができます。その場合、「説明されたように本を選ぶ人々の中には、本当に無作為に選んだ者が少ない」という意味になります。大多数は、何かしらの方法で本の中身を知っていたり、それを疑っていたりして、そういった本を選んでいるのです。

以上の三つの意味のいずれも、適切で正当な解釈であり得る。唯一必要なことは、どの意味が心に浮かんでいるのかを慎重に確認し、それに応じて適切な統計を選択することである。第一義は確率自体において最も正当な使用法を提供するが、その欠点は、当該時点において他の情報によって事象が既知であれば、確率の機能が代替されることである。

したがって、第二義または第三義の方が心に浮かんでいる可能性が高い。これらの状況では、もし実際的に確率を利用できる場合でも、当該時点においてそれは重要な推論を導く手段となるからである。しかし、その欠点は適切な統計を見つける困難さと、特殊なケースの状況によって極めて混乱される可能性があるという点である。

8. (II.) 以上に述べた傾向とは密接に関連しているのが、真の確率的な選択を有意義な選択と混同する傾向である。具体的な対象物について扱っているとき、特に希少性が優れた品質に対応している場合、ほとんどの人はその区別を認識している。例えば、他国で見かけた優秀な軍隊の群れがあるとき、誰もが平均的な連隊から来たのか、それとも選抜された連隊から来たのかを尋ねるだろう。

しかし、この区別の対象が不慣れな物質である場合、特に比較的希少性だけが関与している場合、誤謬はより巧妙で誘導力のある形を取ることがあり、数々の例を通じてその注意が必要であると考えられる。

しばしば結果は、実質的な論理の誤りではなく、考慮すべき事象の確率順序に対する微細な誤解に過ぎない。たとえばピラミッド問題においては、πを単独で考慮するか、あるいは他の類似定数と共にこの定数を分類するかによって結果が異なることがわかった。しかし、7という数字の評価におけるπの代表から実際的な欠如が何か特異なことであるかどうか（第248頁参照）を判断する上で、このような考慮は全体として決定的な役割を果たす。唯一の問題は、この平均値からの逸脱が何か特異なものであるかどうかであり、その答えは我々が予め定めた桁を指しているのか、あるいは十進法で最も高いまたは低い平均値を持つ桁を指しているのかによって決まる。

また、コルヌ (\_Exposition de la Théorie des Chances\_, §102, 114) の問題では、特定の偏差が平均値からフランスの単一地方統計における男性と女性の出生比の報告において種類的な違いを示すものであり、これはランダムに選ばれた単一の地方に関するものである場合にのみ意味を持つ。しかし、それは個々の地方ではなく複数の地方のそれぞれの報告の中で最大値である場合、その偏差は偶然なものである可能性がある。この答えは、我々がこのような考慮を念頭においているか否かによって異なる。

以上のように、これらの問題は確率的な観察とそれに対する解釈の微妙な違いに依存しており、それが数学的論理や事実に基づくものではなく、主観的な判断によるものであることが示唆される。

9. このように誤解される可能性は、ある現象の原因を統計のみによって決定しようとするとき、その原因が期待される方向について全く不知である場合によくあります。このような状況では、巧妙な人間でも広範囲にわたる調査を行っても、これまで観察された事実と一致する説明を見つけることは決して困難ではありません。ピアジー・スミス氏が大ピラミッドの測定に費やしたわずかな労力よりも少ない努力で、私は現在書いている机の寸法からいくつか重要な常数や基準の説明を探し出すことができると思います。

この種の結論の最も奇妙な例は、オックスフォードとケンブリッジのボート競技の年間レースにおけるそれぞれの成功度と太陽黒点の頻度との間に驚くべき関連性があるという発見を行った著者の研究に見ることができます。

このような状況では、私たちの通常の実用的な対処法は、問題に関連する既存知識を援用することです。これにより、限定された適切な事例に基づいて説明が十分に説得力を持つように見せられても、そのような仮説の多くは否認されます。しかし、統計を援用した理論に対しては、統計を裁判所として認めなければなりません。ボート競技の理論も、この理由だけで確立される可能性があります。つまり、長期的な経験が太陽黒点の頻度と一方または他方の成功の優越性との関連性を確認できれば、二つの事象は実際には独立して起こるものではないという事実を受け入れる必要があります。どちらか一方が他の方に原因を与えるか影響を与えるか、あるいは両方が共通した状況によって原因または影響されるものであると疑わなければなりません。

10. （III.）この章の冒頭で説明した論理的な誤りは、確率の規則によって観察されたまたは報告された事象を判断しようとすることから生じたものでした。しかし、その過程において誤った統計を使用して判断を行いました。これと密接に関連するもう一つの論理的な誤りは、そのような事象の一部のみを選んで特徴づけ、確率への訴えをこれらの特性に限定することから生じます。

たとえば、十二便士を投げて十一枚が表になった場合を考えましょう。このような出来事が見られたとき、多くの人々は「全員表が出るところまで近かった」というような感覚を持ち、何か非常に重要なことが賭けられていたら、その出来事に大いに興奮します。しかし、十二枚全てが表になるという意味での私たちの位置づけとはどのような関係があるのでしょうか？

私は、この誤りは一般的な誤解であると考えています。それは十一枚の表を無意識的に既に確保されたものとして見ることで、それらを保持し、残りの一枚を確率によって獲得するチャンスとみなすというものです。十一枚は心の中で片付けられ、確定した（既に起こっているため）、そして十二枚目の出来事だけに確率の概念を導入します。

しかし、十二枚目の結果も同様に既に起こっていますから、それにも確率の概念を適用する必要があります。言い換えれば、一枚が裏になった場合に不満を感じてその結果を取り消し、運試しを再び行いたいと思うなら、十二枚全てを公平に投げ直す必要があります。

この場合、もちろん次の投げでも彼の状況はより良いものとは限らないということです。彼が拒否した同じような結果をもう一度出すかもしれません。彼が行っていることは、連続的な投げ方と混同しているという点で誤っています。十一枚が表になった場合、十二枚目も当然ながら表になる確率がありますが、提案された例では状況は全く異なります。

11. 上記の例では誤りが明らかである。
しかし、サイコロやコインのような単純な事柄からより複雑な問題へ判断を下す際、特に「僅かな脱出」などと呼ばれるような状況においては、類似の誤りが稀ではないと考えられる。例えば、僅かな脱出体験を経た人物はしばしば驚きや不安に満ち、恐れに近い感情でいっぱいになることがある。事件が過去に終わったとしても、その時の感情は厳密には不適切である。もしもこれらの感情が単なる本能的反応や関連性の結果であれば、論理学の対象外となる。しかし、より可能性が高いように思われる場合、これらの感情は過去の瞬間において事件がまさに起こる直前に自分自身を移入し、その状況で我々が予想するであろう感情を想像によって生じさせたものであるとすれば、この過程は推論に類似しており、正しくも誤りもある。つまり、警戒心は仮定的な予測において公平に考慮されるべき危険性に対して比例しているか否かであり、仮定的な予測において公平に考慮されるべき危険性が完全に移入された場合に限り、誤謬は存在しない。しかし、しばしばこの過程は不完全に行われ、事件の一部は現在知っている通りの形で決定されたり「調整」されたりし、残りの部分だけが将来の可能性として公平に考慮されることが多い。

例えば、ある男が友人と一緒に外に出ており、その友人のライフルが偶然発砲し、弾丸は彼の帽子を通過しました。彼は自分が何が起こったか心配になり、おそらく「どうにも近いところだったな！」と述べるでしょう。

私たちは彼が単に発砲が非常に近くで起きたことを意味している以上の何かを意味していると考えることができます。彼はおそらく、「そのときの自分の死亡の確率は非常に高かった」というような漠然とした考えを持っているかもしれません。驚きや恐怖は、単純に発砲が非常に近くで起きたことから生じる身体的・本能的なものである可能性がありますが、彼の精神状態を分析すると、その根底には上記と同様の論理の誤りがあることが見つかるかもしれません。事故に関連して確率を述べたり考えたりすることは、特定の事象を類似した特徴を持つ一連の事象に帰属させ、その後で期待される結果がどれほど頻繁に生じるかを考慮することです。

この場合、最も自然に選ばれる系列は彼と友人が一緒に射撃旅行をするというものです。その時点までには計画的な行動が行われており、それ以降の出来事は偶然の事故でした。このような機会で1000回に一度程度に発砲が偶然起こり、そのうちの1000分の1しか彼の友人の頭の近くを向いているわけではありません。もし我々がこの事故を確率的なものと見なすのであれば、正直にはこのような方法で（最初の例の言葉を借りれば）「コインを投げて再決定」するべきです。しかし、そうはしていません。彼は発砲が頭の1インチ以内に限界を設け、その概念から確率の概念を取り除き、その後でその救済的な1インチから逸脱する可能性があるという概念を再導入します。

12. (IV.) 今度は賭博に関連する論理の誤りについて注意してみましょう。この分野における多くの一般的な誤解は、科学の基礎に対する完全な無知と混乱を示しており、ここではそれらを議論することも余計です。しかし、以下はより説得力のあるものであり、相当な鋭さを持つ人々にとって困惑の原因となりました。

最単純な形で述べれば、次のようなケースがあります。
ある人物AがBと対戦しているとします。Bは個々の人物か、賭博銀行などのグループであるとします。最初にコインを投げて決めるところから始めますが、Aは自身が装置を持っていることを主張し、この装置により彼は勝利を保証されていると述べています。
もし最初の試行でAが勝利した場合、彼は明らかにその点を確立します。しかし、もしAが失敗した場合、次回からは1シャリングではなく2シャリングを賭けます。もちろん、この結果、Aが第2回目に勝利した場合、彼の前の損失を補填し、さらに1シャリングの利益を得ることができます。
その後、彼は毎回の失敗ごとに賭け金を倍に増やして続けます。つまり、最初の成功時に彼はすべての前回の損失を補填し、さらに1シャリングの利益が残ります。
しかし、これは確率の前提に基づいてゲームが構築されている仮定により、いずれかの時点で必ず起こる必要があります。したがって、彼は最終的に勝者となることを保証できます。
また、Aは任意の金額を獲得することも可能です。まず、最初の賭け金をどれだけでも大きくできるという明確な理由からです。たとえば、1シャリングではなく100ポンドなどに増やせます。そして、彼が一度行ったことは再び行えるという点からもです。
彼は自分の1シャリングを積み上げ、その後、長くても短くても、その場面によりますが、同じ結果となる2回目のプレイを行うことができます。したがって、単なる継続性によって、彼は任意の金額を集めることが可能です。これは運に反するかのように見えていても、実際にはそうではありません。

13. この見解は、現在がその議論に最も適した機会であるため、誤謬の列挙の中に分類しました。しかし厳密には、これはパラドックスと呼ぶべきもので、結論自体は完璧に正当です。唯一の誤謬は、そのような結果を得る方法が神秘的であるという認識にある点です。実際には、与えられた条件のもとでの最終的成功を保証する方法は非常に簡単です。この問題について探求する価値があるのは、その原理からであり、一般的な回答が困難に対応していないためです。

たとえば、しばしば主張されるように、銀行はまたは現在では、投資家が勝負の額を選べるわけではないとし、一定の額を超えるものは賭けを承認しないという点が挙げられます。これは確かに真実ですが、我々が仮定する仮想的な問いに対する答えには当然のものではありません。その問いは、このような状況が許されるものであることを前提としています。

また、この可能性は信用が提供されることに完全に依存していると主張されます。なぜなら、投資家が運が回復するまで運勢が持続しない場合、彼の財産は尽きてしまうからです。つまり、彼が十分長く続けば、最終的には彼の運勢が持続しないため、すべての利益が失われてしまうのです。これは確かに真実ですが、成功の条件ではなく原因ではありません。

両者が最初に合意して、ゲーム終了まで支払いは行われず、Aがゲームを終了と判断する権利を持つという状況を仮定することができます。それでもなお、通常の賭博では、固定額または偶然の額で勝つことを保証することは不可能である一方で、この提案したようなスキームを採用すれば可能になることが真実です。そしてこれが説明が必要な状況です。

ここでの困惑は、ある不確定性から確実性が神秘的に生み出されたという仮定にある。詳細な出来事が全く予測不能であり、平均値によって推測されるところではどちらの側にも偏りがないゲームにおいて、一方の参加者がどうしてか確実に運を自分の都合に引き寄せているように見えるからである。これは、ある硬貨が公平なものと見做されながらも、特定の装置を使って永久的に投げて出る表の数が半分以上になるような人物の例と並行しているように見える。

しかし、これは大きな誤りである。実際のところは、Aが利益を賭けているのは一切ないということだ。彼が賭けるのはただ、利益を得るまでの待ち時間の回数だけである。このような状況を考えよう。私は、ある金額を指定してもらうことに同意し、その金額をすぐに私に戻す代わりに一ポンド多く返すという条件で、私が提供する任意の金額を与えることを約束させることができる。鋭さがなくても明らかだが、彼が一ポンド、十ポンド、百ポンドと選ぶかどうかは私にとって関係がない。

今、これらの金額の中からどちらを選ぶかを彼に任せることなく、両者が運命に任せることに同意する場合を考えよう。たとえば、彼らは毎回袋から数字を引いて、それがAがBに対して指定された条件で与える金額となることにする。状況は変わらない。Aは毎回一ポンドを得るため、運の要素がこの点には一切触れていません。唯一影響を与えるのは、この一ポンドが不確定な差し引きの結果であるという事実であり、たとえば10から9を引いたり、50から49を引いたりする場合があるということだけです。彼が運に任せるのはこれらの数値そのものであり、それ自体は全く問題にならない。

個々の個人や企業に対して、このような条件で引き続きゲームを続けることを同意してもらうと提案するのは、あまりにも露骨な要求である。しかし、その問題は原則的に同じであり、形式的にもほとんど同じである。好きな金額を提示され、それを返す代わりにそれ以上の1つを加えてくれることを約束し、1, 2, 4, 8, 16, などという数列の任意の数多くの項を与えられること（あなたがその数列の次の項を持っている場合）、これらは等価である。唯一の違いは後者の場合において、結果を得るためには多少の算術的な形式美が必要であるということである。

同様に等価なプロセスは、私たちは選択ではなく運によって決めるという方法を選んだ場合にも見られる。これは実際に行われていることだ。勝ったときに賭け金を倍にする条件で引き続きゲームを続ける人間は、その数列の幾何級数の項が許可される特定の数だけ続くまで運に任せる以外は何も運に任せている。

15. 特に、議論の対象となっている特定の数列には特別な価値があるわけではありません。その数列は各回の賭け金を前回よりも2倍にするものであるが、重要なのは最終項が全ての前の項を上回ることだけです。この幾何級数より速く増加する他の任意の数列でも、目的を達成できるか甚至に優れています。また、ゲームが公平なものである必要もありません。運の影響は、数列が各回でどの程度まで延伸するかという項の数だけであり、最終的な結果は常に算術的に固定されています。コインを投げた場合でも、系列が2項以上に延伸する確率は半分以下であり、彼の固定した利益は比較的規則的に得られます。ただし、彼が限定された時間内にゲームを行っている限りは、系列が200項まで延伸しても彼には影響を与えません。彼は賭け金を倍増またはそれ以上にすることで、より長い時間を要するだけです。例えば、千枚の当選券を持つ抽選で単一の賞金を得るための公平な賭けを続けることができ、無制限の信用が提供される限り、そのような賭けは安全であるとされています。

16. この観点から見れば、問題は十分に単純ですが、注意すべき2つのポイントがあります。

まず第一に、これは確率的な出来事（この場合、コインの投げ方）とそれらに基づく私たちの行動との区別を強調する。前者は実際には確率的な対象であり、後者は必ずしもそうではない可能性がある。後者の行動は、多くの面で絶対的な確定性を持つように設計されることが十分に考えられる。専門的な賭博家にとっては、この点が馴染み深いことだろう。

なぜコイン投げの賭け方を通常公平に行うのか？それは「公正な」系列は「公正に対応」されるからである。表と裏はランダムに現れ、平均的には等しく現れるし、賭け金も固定されているか、あるいはランダムに設定されている。もし一人が常に同じ額で表を押さえれば、長期的には当然ながら勝ち負けしない。また、賭け金を毎回変える場合でも、その額が過去の勝敗によって依存していない限り、同様である。

しかし、彼自身が希望し、相手も同意すれば、彼は自分の賭け方（例題の場合）を調整して、運命が「公正な」チャンスを得ない状況を作り出すことができる。ここで、人間の選択と設計の要素が、自己の特性としてみればただ物理的な確率性を示す出来事に作用し、これらの要素が消え、算術的に確定した結果を得る。

類似の例は他にも挙げられるが、我々目前の例はその巧妙な隠蔽性において優れている。

17. 今述べた言葉の意味は、以下のケースと比較することでよりよく理解できる。男性出産と女性出産の比率が等しいという仮定のもとで、男性子供を望む一般的な傾向により多くの結婚がこの事件まで持続し、それ以降は続くものとする説明が試みられている[5]。これは、その結果として、一つの息子だけを持つ家族や二つの息子と一つの娘を持つ家族など、一連の家族構成にわずかな偏りがあるように推測される。

この論理は非常に誤りに満ちており、ラプラスが『Essai』において指摘していた通りです。そして、実際の偶然の一連の出来事における運を変更することができるか否かを示す最良の例としてこれほど適切なものはありません。実際の出生数がそのような方法で影響を受け得るとは、ある賭客が頭が出た後、コインを隣に渡し、頭が出た時点で止めるという行動によって、頭と裏の割合を半分以上に増加させることができるか否かを仮定することと同じです。つまり、彼らは頭が出た時点で止まるまでだけ行う必要がないという荒唐無稽な主張で、この段階では説明を省くことにします。「マルティネージ」の重要な点は、賭けが置かれている出来事自体が影響を受けない一方で、賭け金そのものが調整され、運が一方に傾き得ることです。

18. 二つ目の点として、この例は繰り返し強調されてきたことですが、確率の系列は厳密には無限であると仮定されるべきであることを前に押し出します。したがって、どちらかの一方向に結果を出すため、我々がそのどちらかの側から停止を要求する機会を与えられると、理論の正しさに反する結果を得る可能性があります。現在の場合においては、Aが自分の都合で停止できるようにという約束が必要です。つまり、彼が有利な投げ方のポイントで停止できるようにという約束が必要です。この約束がないと、Aは無限に損を被る可能性があります。

一方の側が任意に停止を要求し、その前にそれを拒否することができると仮定すると、ほとんどどんな公平な遊びでも非常に偏ったものに変換することができます。実際の場合においても、Aは毎回賭け金を倍々と増やす必要はありません。固定した賭け金で遊んでも、ゲームが公正であるとし、彼が信用で遊べる場合、特定の一方が任意の額を獲得することを保証することができます。

19. (V.) 非常に不確実なことが全く起こらないと仮定することは一般的な誤りです。この誤りは、言葉で述べると明らかに犯すものではなく、あまりにも明白であるからです。不確実なことは稀に起こる出来事の意味であり、その出来事が時々起こることを仮定しないなら、それは不可能と呼ばれるべきです。この誤りは、ほとんど例外なく曖昧な一般理解の中でしか発生しないもので、確率に関する著作では豊富に反証されており、ここでは簡潔に触れれば足ります。もちろん、私たちは確率の定義から、非常に稀な出来事の組み合わせを「絶対に起こらない」という言葉で表現することは誤りであることが分かります。このような表現は、緩い一般的な過言として流布するかもしれませんが、厳密には矛盾を含んでいます。そのような稀な出来事に関する真実は、デ・モランによる引用文によって以下のように最良に説明されるでしょう：

「極めて improbable なケース、たとえば、袋の中に20個のボールがあるとき、同じボールを5回連続で引くようなケースが誰も到達することはないと言われています。特定の個人が単一のサイコロで12回連続で目を出す確率は、十二分に高いでしょう；実際、そのような事件が12回の試行で起こる確率は非常に低い（約20億分の1未満）ので、特定の国が特定の世紀においてその経験を持つことは稀です。しかし、一時的に停めましょう。この議論が適用される対象は何でしょうか。

稀にサイコロを扱う人は、三回連続で目が出ることなどほとんどあり得ないと信じるでしょう；頻繁に扱う人にはそのような事実があることがしばしばあります。これらの道具を使用する非常に熟練した使用者は、さらに珍しいシーケンスを見たことがあります。

今 imagine しましょう。ある社会がサイコロを頻繁に投げ、六回連続で目が出るというシーケンスを観察し記録したとします。そのような経験から十二回連続の目が出ることを否定する議論は依然として使用されます。そして別の社会が十分な練習を積み、十二回連続で目が出ることを見たとしても、二十-four 回連続の目が出ることを疑う方法を使用するかもしれません。そして同様に、無限大まで繰り返されるでしょう。

想像力は長いシーケンスを思い描く能力が、それを示し並べる能力よりも遥かに優れているため、地球上の全ての砂粒に対応する電報を作成することも容易です。この論理上の誤謬は、経験を超える事件数があると仮定することで生じます；これらのイベントはすべて、私たちが経験したシーケンスから構成されています。過去はその成分の質に関して完全に包含していると判断します。

しかし、非常に大きな数のシーケンスに対しては、我々の経験は一握りの麦粒よりも少ないでしょう；それすらも一束の麦粒よりも少ないかもしれません。」

20. この根深い誤謬の起源は、説明が難しいものではなく、私たちの実践的な生活の要請から生じるものです。確実に言えるのは、我々は全ての状況を思い出すことができないということです。したがって、世界で何かを行うためには、多くの稀な状況を完全に考慮に入れることができません。また、この必要性は私たちの短い人生によって強化されます。数学的には、十分長い寿命を持つ者は狂犬に噛まれる確実な出来事であると述べることができます。その出来事は不可能ではなく、ただ improbable であり、時間とともに必ず起こることになるのです。しかし、これは多くの他の不快な状況と同じく、実践上ではほとんどの場合、完全に無視されるべきもので、それらは必然的に私たちの思考や期待から等しく消えてしまいます。そして、その出来事が重要なものではなく、稀なサイコロの投げ方のようなものであるとき、数学的な計算に慣れていない人々にとっては、それを可能であると認識するためには大きな想像力が必要となることがあります。

時々、この実践的な「ゼロ点」に対する不確実性の極限をどのように定義すべきか試みがなされています。論理学者や確率を数学的に評価することに抵抗感を持つ人々によって行われたこのような試みは、前章で議論したモード上の困難の一形態であると見なされるべきであり、ここでは再検討する必要はありません。しかし、数学者の視点からこの問題を考えた人々の意見について少し追加の言及が必要かもしれません。

これらの原理のうち、おそらくブフォンがその代表格である。彼は（『道徳的算術』§VIII.）で、この実用的な零点を1/10,000という確率に相当するとして評価している。この分数を選んだ理由は、彼が手元にあった死亡率の表から見つけることができる。すなわち、56歳の男性が次の日間に死ぬ確率を表すものである。しかし、通常の状況下では誰もこの確率を考慮に入れる余地がないため、これは実用的に価値がないと評価される。

明らかに、この結果はほとんど完全に任意性に基づいており、実際には彼の理由は、選んだ便利で単純な分数を採用するための経験的な微かな正当化に過ぎないものと考えるべきである。しかし、他のどの分数も選択された範囲内で広い限界内であれば、その正当化は明らかに同じく可能であった。

21. この誤差の一種には、時折その重要性が強調されるため、より特別な検討が必要かもしれません。上記のように、どんなに不確実な出来事でも、十分に状況を変化させたり、言い換えれば試行を長く続けることで、最終的にそのような出来事が起こるという疑いはありません。たとえば、サイコロを数回投げると双子の目が出ます；3つのサイコロを続けて投げると三重目の目が出ます。どんなに不思議な出来事でも、連続して六面体が千回出ることもありますが、十分な忍耐と活力があればいつかはそのような出来事が起こります。

この結果をアルファベットの文字に適用します。すべての文字が入った袋から一つずつ文字を取り出し、その後すぐに元に戻すという状況を考えます。これらの文字が順番に出されると仮定すると、一般的には無意味な文字列になるだろうと予想されるでしょうし、人間が知る言語の単語に並び替えられるとは考えにくいです。しかし、それほど頻繁ではないかもしれませんが、理論によれば、そのような結果を前提とするのは一般的で、読者もこの仮定を受け入れることが容易であると想定されます。つまり、試行を十分長く続けることで意味のある単語が現れ、さらに言えば、我々が選ぶどんな本でも、例としてミルトンの『パラダイス・ロスト』やシェイクスピアの戯曲も最終的にこの方法で生成されるでしょう。

このような文字を袋から一つずつ抽出するプロセスを数日分表記しようとすると、この巻には十分なスペースがありません。しかし、最終的には望む結果は得られるでしょう。

多くの人々は自然に、天才の作品が偶然によっても生成されたと提案することを恥ずかしいと考えるかもしれませんが、他方でそのような主張に対して議論する人々もいます。もしこれが正しいとするなら、世界自体もこのような方法で偶然によって生成された可能性があるのでしょうか？

22. まず、比較的単純で確定的な、そして理解しやすい問題から始めましょう。すなわち、シャクスピアの作品のような偉大な人間の天才の作る可能性が偶発的に生み出されるか否かという問題です。この可能性について、シャクスピアの作品を生み出す力は、天才的な才能と単なる偶然に限られず、時間さえ与えれば誰でも行うことができる第三の選択肢として存在します。つまり、純粋な機械的手続きが可能です。誰でも、ほとんど痴人にもかかわらず、十分な時間を費やせばこれを達成することができます。

例えば、必要な数の文字を偶然ではなく、並び替えるためには意図的に、そして置換理論によって提示されたルールに従って準備し、配置する場合を考えましょう。アルファベットの文字とその数は有限であり、これらの文字が発生できるすべての順序は適切な順番で訪れますから、言語で表現できる全てのものがいつか達成されます。

このような結果に誰かが驚くべきことは何もない。その可能性は以下の理由から生じる。

当々手元にある文字の数、そして言葉の数は限られている；したがって、我々が言語で表現したいいかなる内容も、それに応じた制約を受けることになる。考えられる思想の変化は実質的に無限であり、口調やその他の要素によって話す言葉の変化も同様である。しかし、言葉に至っては制約が生じる；この制約の性質は心で明確に理解できるものの、実践上は認識するほど大きな制約とはならない。これは、生成可能な組み合わせの数が想像力を超えるほど巨大であるためだ。

したがって、答えは明白であり、多くの他のケースにも適用される。つまり、何かを達成する方法の数に有限な限界を設けることは、その人が十分長く試行すれば成功することを決定づけることになる。偉大な天才がこのような状況でそれを実行する場合でも、彼は偶然の人間や愚者と自分の主張が争われる可能性を受け入れる必要がある。

たとえば、シャイレットが八つ或いは九つの指定された言葉だけを使用する制限があった場合、後者の代理者が彼と同等の地位を主張できる時間は非常に短いだろう。しかし、英語全般の範囲を自由に使用できたため、彼の名声はこのような方法によって攻撃される危険性には晒されていない。

23. 世界が偶然によって生み出される可能性というケースは、議論の全く異なる領域に私たちを導きます。
ここでは、想像力がその完全な意味の一端も理解することができないほど崩れてしまうとしても、図形の性質と用途について理解力が十分な範囲内で扱うものではありません。
理解力自体は適切な領域から逸脱しており、問題の条件を定義することは不可能です。袋から文字を取り出すときには何をしているのかよく分かっていますが、偶然によって世界を生み出すとは何か？以前の例に比喩的に類推すると、ある種の代理者が前提となることが予想されます。--したがって、以下のような仮定は他のどの仮定よりも不自然ではないかもしれません。ある存在を想像してみてください、創造者ではなく、半神人であるようなもので、彼には一定量の材料が与えられ、それらを配置し、行動の法則を与えることを盲目かつ偶然的に行います：このような世界が実際に経験しているように生み出される確率はどのくらいでしょうか？

もし真剣にこの質問に答えることに価値があるとすれば、そのアルファベットの文字数やそれを使って書く文章の長さについて誰かが私たちを助けることができれば、結果を示す方法を考えることができます。しかし、少なくとも次のように断言できるでしょう：--この小さな本の長さで逆確率が与えられるためには十分な図形を見つけるよりも遥かに不十分であることは明らかであり、我々がそれらを書き出すまでに世界が生産したすべての紙も使い切ってしまうでしょう。

24. 後に起こる非常に稀な出来事に関する困難が最も誘惑的な形で現れるのは、おそらくこのようないくつかの例でしょう。「あなたは（ある人々は言うかもしれません）、そのような出来事は時として起こる acepta; いや、実際に無限の時間の中で時々起こることもある acepta。では、なぜ私たちは今この機会がこれらの可能性のある出来事の一つではないと知ることができるのか？」

これに対する答えはただ一つでなければなりません。統計や確率に関わる場面ではいつも同じ答えを出す必要があります。「現在のところ、このような出来事が起こっている可能性がありますが、それは非常に不思議なほど低い確率です。あなたやあなたと同じ立場の人々がこの主張を数兆回も繰り返す中で、そのうち一つだけが正当化されるかもしれませんし、おそらくあなたはその一人ではないでしょうし、現在の機会がその出来事の一つでもないでしょう」

25. これは、確率の評価において一つの大きな数と別の大きな数を区別する実践的な困難のもう一つの形態であり、遺伝に関する議論におけるその重要性から一時的に注意すべきです。人々はしばしば、質的特性（精神的・肉体的）が親から子へと伝えられるという教義に対して反対意見を唱えるために、逆に多くの例があることを挙げます。実際には、そのような例の大部分であると言います。この反対意見を提起することは、特定の人間が知性や他の優れた能力で目立つ可能性のある非常に大きな確率を全く理解していないことを示しています。遺伝に関する議論は、そのような確率が存在することを含みますし、実際に存在すると主張します。これは、定義上の問題である部分があり、稀少性の度合いを優れた能力と見なすときに考慮すべき程度に依存しますが、合理的な意味での用語を使用して考えると、失敗の非常に大きな割合であっても、遺伝教説を支持する証拠には依然として莫大な優越性があることが明らかになります。たとえば、4000人中1人がそのような区別を示す程度の優れた能力が含まれるという度合いを考えましょう。これは大きな差異であり、イギリス合衆国の全体的な成年男性人口約2000人の中にはこのような人物が見つかるため、それは目立った天才を意味するわけではありません。今 imagine そのような人々の子供たちのケースを調査したとします。そして、200人のうち199人が同じ区別に達しなかったことを発見しました。多くの人は、これは遺伝的な伝播がほとんどないという結論的な証拠であると考えるでしょう。有利な例が一つだけでも、199の反対例に対して提出できるとは、そのような理論全体が崩壊したと彼らは考えるでしょう。もちろん、この誤りは明白であり、このような数字がある状況ではほとんど誰もそれを避けることができないでしょう。しかし、一般的な会話や他の情報源から判断すると、実際には、4000人中1人がそのような区別を示すという例が有利な例の一つだけであるとしても、理論に約20対1の確率があることを十分に認識することが非常に難しいことがわかります。遺伝的な伝播がなければ、4000人の息子のうち1人は父親と競争するでしょうが、実際には（もちろんここでは仮定的な割合を使用しています）、200人中1人がそのような区別を示すことを私たちは見つけます。したがって、統計が十分に大きい場合、偶然性の比率20対1で単なる一致が改善された影響があったことがわかります。実際には、非常に大きな数の意味を理解する能力が非常に低いので、大規模な数字に関連する比を頭の中で保持することが難しいため、算術的な考慮事項によって自己チェックを行わないと、一定の限界を超えたすべてのものは等しく巨大で曖昧であると扱い、評価しがちです。

26. (VI.) 機会と推論との関係について議論する際に、繰り返し観察された事象が再び現れる確率を推定するために一般的に与えられる規則の主張を検討しました。私は、そのような規則を得る試みやそれを証明しようとするすべての試みは必然的に無意味であることを示そうとしたのです。これらの理由が決定的なものであれば、そのような規則を使用すること自体は当然誤りと見なされるべきです。ここにいくつかの例を追加することで、より適切な分類ではなく単なる一連の規則ではなく、可能性のある形態を三つのクラスに分ける方が良いということが示されるかもしれません。

(1) ある場合において、何かが数回観察されたとき、その何かが再び起こる確率は結果的に増加します。これは一般的な規則と一致し、最も頻繁に遭遇する可能性が高い状況です。我々が「何かが起こること」について話す際の表現の必要性から、この一般形での規則を容認することは不可能です。しかし、それを少し特殊化することで、より馴染みのある形になります。例えば、個体の一連の中で二つ以上の特性が頻繁に共存していることを観察した場合、その特性が将来も同様の関係にあると結論付けることは力強いでしょう。しかし、我々の信憑性は単に観察された一致回数だけではなく、より複雑な考慮事項にも依存します；これらの議論については、推論的証拠に関する定規的な著作を読者に参照していただく必要があります。また、例えば、二つの事件のうち一つが他方を数回引き起こした場合、その最初の事件が再び起こったときには多くの場合、次々項の発生に対する期待感が生じます。しかし、その期待度は単純な公式で定義することは不可能であり、その程度は主に二つの事象間の関連性の程度にも依存します。この主題に関する明確な規則を立てる試みは原因法則についての議論へと導き、それらが存在する可能性を推測できる状況についても議論することになるため、これ以上の検討はここでは放棄せざるを得ません。

27. (2) あるいは、第二に、過去の再発生が未来について推論する正当な根拠を提供しないこともあります。これは確率論において最も適切に扱われるケースです[10]。そのような場合が確率論においてどのように扱われるかは、科学の基本的な概念を念頭に置いて考えればすぐに理解できます。そこで我々は、推論のために無限に延長された項の系列に導入されます。これらの項に関する詳細な情報は特定の点以外では提供されず、我々の知識は、例えば10個中1つが属性Xを持つという統計的事実に限定されています。

今、これらの中から5つの連続した項がXであると仮定します。これにより第六項もXであるというヒントは何でしょうか？明らかに何もありません；この過去の事実は何を教えてくれません。我々の推論のために使用する公式は以前と同じです：10個中1つがXであれば、次の項がXである確率は9分の1です。そして、連続して何らかの種類の項がいくつ存在しても、この公式は変わりません。

[10] これは確率論において最も適切に扱われるケースです。

28. この公式によって与えられた答えが正当化される方法は、しばしば誤解されます。そこで、数学的な概念に不慣れな人々のために、説明を加えることが有益です。

Xが連続して12回起こったとしましょう。これは明らかに異常な状況です。このような状況が継続的に起こるとは、統計によれば長期的にはXが10回中1回しか起こらないという主張に反するからです。しかし、この異常はどのように克服されるのでしょうか？つまり、我々が現在までに経験したような長い連続のXの後で、Xが十分の一よりも頻繁に起こる結論を避ける方法は何でしょうか？多くの人々は、過去の優位性を補填するためにその後Xの割合が減少する必要があると考えているようです。しかし、これは誤りです；将来のXの出現割合は常に同じでなければなりません；統計的な公式に従うならその割合は変更されません。

実際、異常な状況を修正するためには、数列の中で新しい項が継続的に流入することによってのみそれが達成されます。これらの新項は長期的には混乱を中和しますが、特別な適応があるわけではなく、単にその圧倒的な数の重みによってです。したがって、連続の中でどの段階でも、前の項の数や性質に関係なく、10回中の1回はXであるという事実は依然として成り立ちます。

有限の項しか考慮しなければならない場合、たとえその数がどれほど多くても、我々が述べたような干渉は、その後の比に特別な修正が必要となることは事実です。しかし、無限の項を扱う場合、これは異なります；その系列の「極限」、つまり私たちが扱う対象は、これらの一時的な干渉によって影響を受けません。系列の継続的な進行においては、実際にこのような干渉が増えることが多く、それらはより例外的なものになります。しかし、どの時点にいても、系列の無限延長に向けて過去または将来を見渡すと、その項が配置される最終的な極限比が変わらないことが確認されます。そして、これは上記のように確率論における厳格な規則に関連しています。

最も一般的な例としては、コインを投げることでしょう。四回連続で表が出たとしても、人々は今や「五回目も表が出る」という確率が前回と同様であることを十分に理解しているでしょう。そして、その後も同じ確率が維持されることになります。先の段落では、これらの平均的な干渉が長期的には中和される方法が説明されています。

29. (3) 他にもまれなケースがありますが、それらは全く未知ではありません。そのような推論（順序則から得られるもの）が真実と直接反対になる場合があります。あることが頻繁に起こるほど、再び起こる確率は低くなるかもしれません。これは有限の源から物を取り出す場合（袋からボールを取り出し直さない場合）、または繰り返し自体が連続を防ぐ場合（誤報を与える場合）など、特定の状況で発生します。

私は、最後の二つのクラスで述べた結果について、あるような一般的な帰納則や、あるいはその原理に基づいて信じていることに全く問題ありません。しかし、これは各特別な場合における規則の有効性を認めることと混同してはなりません。私たちは規則が個別の事例や、それらのクラスに適用される根拠について話しているのです。これは以前の章で指摘したように、規則そのものに対する基礎となる理由とは全く異なることです。

もし誰かが普遍的な規則として「すべての人々の証言は信じるべきである」と定めるとしましょう。そして私たちはある人物が嘘をついた例を提示するとします。その場合、彼が私たちが他の人々の言葉によってそれが嘘であると信じていることを示して、規則を救うとは考えられません。

しかし、一般的な規則としてではなく、普遍的なものではない規則として「人々の証言は信用できる」と述べることは完全に矛盾しません。そして、信頼できない第二のクラスの人々を別けて、さらに誰もが第二の規則の根拠を求めた場合、それを最初の規則に基づいて説明することができます。

もし我々が必要律について話しているなら、このような対立は古い論理的な「クレタ人」パズルのように無意味なものとなります。しかし、帰納的または類推的な拡張の場合には全く問題ありません。

30. 以下は、上記で述べた三つの異なる結論を明確にするために使用される、よく知られた例です。連続して10日間雨が降りました。AとBはそれぞれ第11日の雨の有無について結論を下し、Cは過去の雨が何ら意見形成のためのデータでもないと主張します。誰が正しいのでしょうか？事前に決定することはできません。直接的な観察に訴える必要がありますか、または独立した根拠に基づいて該当する分類を見つける方法を見つけなければなりません。

例えば、ある国は必ずその国内で雨を生成すると知られている場合、私たちは第三のルールを選択します。これは有限な供給から推論を行う場合に相当します。また、私たちの国の雨が地球のどこでも生成される可能性があると信じられる理由があれば、過去の降雨が連続した湿った天候の継続について何らの光を投げかけないと結論付けるでしょう。そしてその結果、第二のルールを選択します。

あるいは最後に、雨は長期的な時期やスペルで来ると知られている場合（熱帯では）、10日間連続して雨が降ったことから、私たちは一つのこれらの時期に入ったと信じ、したがって次の日も前10日の天候と似ていると考えます。

このような推引規則のすべてのこれらの形態は可能であり、事前に一つをもう一つより好む根拠がない場合でも、普遍的な予測の公式を確立しようとすることは不合理であるように思われます。私たちができる唯一のことは、これらのルールのいずれが実際的に適用されるかについてどのような状況があるかを確認し、そのような状況でそれを使用することです。

31. (VII.) （V）で議論されたケースでは、私たちは実用的な観点からはほとんど無視してもよいと判断した極めて小さな確率について言及しましたが、そのような確率が個々に小さいとしても、その数が多い分、有限または相当な大きさになる可能性がある場合、それらを考慮に入れないという誤りはしばしば見受けられます。数学的な観点からは、これらの確率は時として「積分」されて有限または相当な大きさになる可能性があると表現されるべきです。

例えば、私たちは一つの困難に直面し、その脱出方法が一つだけ明らかに可能な形でしか見えない場合があります。この困難を克服するための試みは、その見かけ上極めて低い確率にもかかわらず、他のすべてよりも遥かに高い確率であるからこそ受け入れるべきだと主張されます。私は実用的な観点からは、このような道を選択することがしばしば合理的であることは認めます。なぜなら、私たちは一つの仮定に基づいて行動しなければならないからです。そして、大量の不可能性の中から最も不可能ではないものを選ぶのは自然で正当なことです。

しかし、行動を強制されない場合、そのような決定的な優先順位は求められません。このような説明に同意しないことは完全に合理的であり、またそれが真実であると確信していないことを明確に述べることも正当です。そして、現在のところ他の説明を受け入れるわけではありません。私たちは実際には判断を留保した状態にあり、行動が強制されない限りは、そのような状態が完全に正当で合理的であると認識しています。

一つの選択肢が個々の他の選択肢よりも決定的に確率が高いかもしれませんが、全体としての他のすべての選択肢と比較すると決定的に確率が低いかもしれません。これは自明なことですが、人々がしばしば見落としているのは、この事実と感じ方を述べるとは必ずしも矛盾しないということです。そして、行動が強制されると、その選択肢が本当に正しいものであるかのように積極的に行動することも可能であるという点です。

32. 特定の例を挙げると、この考え方についてはしばしば「スピリチュアル」現象に関する議論において私自身に浮かんできました。我々は、他の可能な解決策が提案されないからと、承認を求められます。確かに、個々の代替案に対する明白な困難性があるかもしれませんが、そのような解決策がどれだけ存在する可能性があることを十分に認識しているでしょうか？各々がほとんど信じ難いとしても、それらすべてを一括して集めて提示された解決策に対して天秤に乗せることは公平です。我々は他の個別の代替案に対する同様の手法でも同じように対応します。我々は論理的な観点から問題を見つめているのであり、スピリチュアルな解釈やそれ以外のすべての解釈を平等に扱うことに完全に同意しています。

各代替案の支持者は、議決権を持つ会議委員会のメンバーと似たような立場にあります。他の誰も他のメンバーの提案を支持しないように努めます。一人ひとりは他のどの提案も否决するのに効果的に貢献できますが、自分の提案を通すことは不可能です。緊迫した必要性によって彼らは実質的な不作為から脱却し、最小限の抵抗点で反対勢力に突破される可能性がありますが、そのような圧力がない限り、彼らは希望を失う状況に陥ります。

33. 精神性の解釈が科学的な扱いを受けることができ、そのような解釈は時として証拠面前で導き出される結論であると仮定する場合、私はこのように考えるべきではないか。個々の説明に対しては、「信じられない、受け入れられないと」述べる必要があるかもしれない。そして、急いで決定を下すような状況が（おそらくそうならないだろうが）現れるまでには至らない限り、我々はその問題を一時的に保留するべきである。

おそらく一方の説明（すべての可能性が考慮された場合）が真実であると主張されるであろう。これは容易に承認できる。また、バトラーによく引用される原理に基づいて、「最も説得力のある説明」を即座に受け入れるべきだと主張されることもおそらくあるだろう。その説明の絶対的な価値に関わらず、他の説明と比較して最も説得力があるものであるからだ。これは明らかに誤りである。もし状況が我々の決定を先取りさせる必要性を生じさせた場合に、「この説明を受け入れるべきだ」と述べることは、その説明を現在は否認していることを意味する可能性がある。唯一の理にかなった立場は、真実が様々な選択肢の中にあると認めつつ、一つの説明よりも他の方が優れているというような偏見を持っていないことを明確に認めるべきであるということだ。それ以外のどの説明も信じていないと言えるだけのことである。

34. (VIII.) 「結果に基づいて判断する」という一般的に指摘される誤謬は、確率論の領域に属することから、ここでは一言触れることにする。この誤謬は、その本質が科学の一般原理を理解した者にとっては非常に明白であるため、わずかな注意を払うだけで十分である。ある意味で、すべての命題は結果に基づいて判断されるべきであり、これは単に経験による検証という別の表現に過ぎない。しかし、普遍的命題に対する検証と、比例的または統計的な命題に対する検証とは最も異なる。前者は単一の反例によって覆われるが、後者は反例を容認し、それすらも含意している。しかしながら、他人の助言が不幸にして失敗したからといってその助言を非難するか、成功したからといって自分自身の助言に信用を与えることは非常に一般的である。

もちろん、結論が明示的に確率的なものであった場合、反対の結果を受け入れる準備をするべきであり、甚至は連続して反対の結果が現れたとしても同様である。しかし、その系列が長期にわたって続く傾向があるときだけ、疑惑と疑問が生じるべきであり、それらは所与の確率度と経験によって示されるその離差の程度との比較によってのみ生じるべきである。一つの失敗に対しては、「助言は妥当で、長期的には再び同じ状況下で提供する」と答えるべきである（ただし、長期的な正当性が確認できる場合に限る）。

35. 以上での区別は慎重に考慮すべきであり、確率論と通常の論理で扱われる命題の種類の大きな違いとそれに伴う証明の本質的な差異により、同じ一般的な外見をもつ議論が前者では有効である一方で後者では誤謬となる可能性があるし、逆もまた然りであることが明白である。

例えば、全般的に肯定的な命題を単純に逆転させるという有名な論理の誤謬を考えましょう。すなわち、「すべてのAはBである」という命題から「すべてのBはAである」という命題へと移行することです。一般的な論理において、結論が前提と同じくらい確実であるべきであれば、このようなステップには言及する余地はありません。しかし、推引や確率の目線で見れば、このようなステップが時として正当化される場合があることが明らかになります。「一部のBはAである」という事実だけでも、任意に選んだBがAであるという予想を立てることができます。少なくともそのような信念には理由があり、統計的な頻度に関する情報がない限り、その信念の価値を定量的に評価することはできません。

私は、ある人間が「すべてのAはBである」という命題から特定のBがAであると結論づけた場合に、彼は必ずしも必要条件ではなく、ただ確率的な推論だったかもしれないという主張を正当化できると考えています。同様の議論は、「中間項未分布」の論理的誤謬についても適用されます。

逆方向の例を考えましょう。つまり、確率が機能しない一方で、通常の推論が立つべき状況に類似している場合です。私は百万通の書簡の中の一通が郵便局の責任で失われることがあると知っています。友人に手紙を送りますが返事はありません。この状況から、彼のせいであるという理由があるのでしょうか？

この事件（すなわち、手紙の紛失）は確かに起こりましたし、その原因として考えられる二つの要因のうちの一つの「価値」つまり統計的な頻度が正確に定められていると仮定すれば、もう一つの原因が作用した可能性について何か推論できるのでしょうか？このような状況下で何も知らないという主張は、未知の項を持つ方程式が理論的に解けないものであることを主張しているようなものです。

このような例については、確率の逆則に関する章で十分に議論されており、ここでは読者にその結論が全く導き出せないことを思い出させるだけでよい。つまり、私の友人が書簡を返さなかった原因がポストオフィスにあるよりはるかに可能性が高いという推測を立てるための証拠は得られない。彼が受け取った書簡に対して返事を省略する頻度について我々が知っていることや、そのような仮定があることなくこの問題は未解決のままである。

上記で示された類似性が一方程式に一つだけ未知数を持つ方程式と一致しない理由は、確率の観点から見れば実際には二つの未知数があるためである。我々が扱うのは比例的または統計的な命題であり、現在の事例において書簡が失われたことが示されたが、そのような割合で失われるという情報が与えられていなければ、私からの相手が返事を省略する頻度を決定するための一未知数しか存在しない。もしこの後者の情報が与えられていたら、私たちはこの問題を最も容易に解決できたことだろう。

1. これはF.Y. Edgeworth氏によって1887年4月号の\_Phil. Mag.\_で議論されている。

2. \_Journal of the Statistical Soc.\_ (Vol. XLII, p. 328) この推測は冗談であるか？

3. ゲーミングの人々の中では長い間、\_Martingale\_という名前で知られており、Babbage氏の論文（1823年の\_Royal Soc. of Edinburgh\_の\_Trans.\_）はその一部について議論しているが、その後に続くセクションに関する話題にはほとんど触れていなかった。

4. 保険とギャンブルの章ではさらにこの区別に注意を向けられる。

5. Prevost氏によって\_Geneva Bibliothèque Universelle\_（1829年10月号）で議論されている。Quetelet氏は(\_Physique Sociale\_, I, 171)でその説明が注目され、おそらく受け入れられている。

6. \_Essay on Probabilities\_, p. 126.

7. この理論的なまたは絶対的な、非常に稀な事象の無視は、天文学や他の観測家によって時として推奨される実用的な無視とは混同してはなりません。このような拒否の限界を示す基準として知られるチャーヴェンヌの基準については、メリマン氏の『最小二乗法』（p. 166）で説明が見つかります。しかし、これは長期的には誤差のバランスが小さくなるという理解に基づいています。非常に稀な事象は意図的に拒否され、無視されるわけではありません。

8. 計算過程は容易に示すことができます。例えば、対象となる著作には約35万の文字があります。アルファベットの26文字のいずれかが各回選ばれるため、可能な組み合わせ数は26{350,000}となります；この数字は、対数表から推測すると、ほぼ50万桁が必要となることが分かります。これらの組み合わせのうち、スペルの変異を排除した場合に有利なのは1つだけです。したがって、必要な35万文字を選択するたびにランダムに選んだ結果を得る確率は、この数で1を割ったものとなります。

もしこの確率が小さすぎると考えられ、ある人が上記のランダムな選択を何回繰り返さなければ、成功に対する2対1のオッズを得られるかと尋ねた場合も、これについては容易に示すことができます。各回での事象の確率が1/nであるとき、n試行で少なくとも一度その事象が起こる確率は1 - ((n - 1)/n){n}となります；これは連続して失敗しない限り実現します。対象となる例ではnが非常に大きい場合、代数的にこれを約2対1のオッズと同等に示すことができます。つまり、必要な量の文字を上記で表現された極めて大きな数回選んだ後でも、成功を得る確率はまだ2対1であり、それからそれが何であるかを認識する必要があります。

9. どの言語にも可能な限り長い寿命が付与されても、このような時間の期間を考える算術的な観点に比べれば、その長さは完全な無意味化してしまいます。

10. 当然，ここでは、平均値が最終的に収束する極限が既知であると仮定しています。これは原因の知識から得られるか、あるいは以前の広範な経験から得られるかもしれません。例えば、サイコロは公平なものであることが分かっているとします；この事実が分かっていない場合でも、観察された結果から偏りを推測する必要がある場合は、そのようなケースは最初のカテゴリーに属することになります。

11. ただし、ギャンブラー以外の場合です。占い師のホディンによれば、彼自身がスパで知遇したある賭博家によると、「特定の組み合わせが頻繁に現れたほど、次のコップ（一回）で再現されない確信度は高い：これはすべての確率理論の基礎であり、運命された機会と呼ばれる」（『カードシャーピング暴露』、第85ページ）。

\newpage

\chapter{保険と賭博（Insurance and Gambling）}
\label{ch:xv}

\begin{center}
\large\textit{不確実性の低減（保険）と増大（賭博）・数学的期待値と道徳的期待値の限界}
\end{center}
\vspace{4mm}

第十五章.

保険と賭博。

1. 読者各位が最初の数章で確立され、説明され、その後豊富に示された確率論の基本的公理を思い出すと、個々の不規則性と平均的な規則性という相反する特性は、異なる人々によって異なる評価されることが自然であることを容易に認識できるでしょう。ある人々にとっては、不確定要素自体やその結果がいかに苦痛であるか、あるいはそのような要素を避けるための手段を採用することによりそれらを軽減することが望ましいかもしれません。他方で、生活における最終的な規則性は、少なくとも特定の分野においては、彼らにとっては単調さが如何に疲労をもたらすかから、その特性を一部変更し改善するための方法を模索することを希望するかもしれません。これらの精神的傾向と最も単純で明白な満足法について簡単に議論します。

ある人々には、我々が述べたように、世界は変化と不規則性に満ちており、その結果としての不確定性もまた多々存在します。文明は特定の方向でこれらの特性を減少させるのに一定の成功を収めましたが、他の方向ではそれらを増大させることにもなっています。そして、現在においてその影響が全体的にどの程度効果的であったかを確実に断言するのは容易ではないかもしれません。

変化と不規則性の減少は、我々が必要とする食料や衣服を供給する主な製品である場合など、様々な形で例示されています。これらの製品に関しては、飢饉や不足が比較的未知のものとなりました。特に、文明度の高いコミュニティにおいては、これにより物価の変動は以前よりも著しく減少しました。しかし他の方面では逆のことが起こっています。例えば、ファッションは大規模なコミュニティ内で同時に多くの人々を同じものを望むように誘導し、その結果として大きな価値の変動が生じることがあります。

さらに、信用や通貨に関連する事項においても、以前には存在しなかったような大きな頻繁な変動が生じています。これは、政治経済学の領域に踏み込むことなく、一連の要因によって引き起こされるものです。破産はその性質上、現代の時代から生まれたものと言えるでしょう。

これらの対照的な結果を完全には評価しようと試みることは困難ですが、重要な事項において実際の不確定性が全体的に減少した一方で、財布に影響を与えるような事項ではそれらが増大した可能性があることを指摘することはできます。また、実際に不確定性がより大きくなることがない場合でも、その結果がしばしばより深刻なものになるか、または我々がこれらの結果を評価する見方を高めるなどして、実用的な観点から言えるほど大きな影響を与えた可能性もあります。

2. しかし、上記の問題、すなわち最終的な利益と損失の均衡はどのように決定されるかに関わらず、多くの人々が生活の一部における現在の不確実性が自分の好みに適していないという事実は疑いようがない。彼らは何らかの方法でこの不確実性を軽減することができるだろうか？個々のケースに関しては、慎重さと予見性によって何かが行われる余地はある。私たちの家屋は火災による被害を受けやすいように設計されないよう計画されるべきであり、また消防設備が近隣に用意されていることも考慮に入れられるべきである。疾病や事故から死を防ぐためには、慎重な生活を選択する人々によって何かが行われることができる。しかし、上記のような予防策は確率の問題を導入しない。これらの後者の考慮事項は、私たちは詳細の不規則性から救われるための平均的な規則性を呼び込むときだけに現れるものである。確かに、その不確実性自体を除去することは不可能だが、私たちが行動することで、その不確実性の結果が自分たちにとって小さくなるようにすることができる。または、自分が関心を持っている人々にとって小さくなるようにすることができる。

生命保険の例を取り上げてみよう。専門家は自分の収入に依存しており、それが一瞬で消えてしまう可能性があることを知っている。これは彼が阻止できない状況である。もし彼がそのような状況にある唯一の人間であり、または同僚と組むことができない場合や、不愿意的话，那么除了尽可能节省开支并留出儲蓄余地外，别无他法。

3. 彼はしかし、容易な脱出方法がある。
彼がすべきことは、自身と同じ状況に置かれた他の多くの人々と合意し、共通の基金を作り出すことである。彼らは、そのグループ内の寿命が平均よりも長く生き残った者は、寿命が短くなった者の家族を支援することに決めるべきだ。もし少数だけがこの決断に同意したとしても、それほど大きな利益を得ることはできないだろう。なぜなら、それでも彼等は当初の不確実性からわずかな段階でも離れないからである。重要なのは、相当数の人々がこのような方法で協力し、平均的な比較的一定性を享受できる利点を得ることであることを知っているように、多くの人々がそのような合意に達することだ。

4. 上記の単純な考察は、保険の本質を完全に含んでいる。賠償に関する合意が特定の金額までに限定されているという事実、そして各メンバーの死亡時に一時的な総体的寄付を求めることではなく、固定された年間保険料を支払うことで代わる支払いを行うこと、といった点は、単なる利便性と配慮の結果である。保険とは、生命や雇用の不確定性の結果に恐れを感じる人々との相互契約であり、その集団的な規則性を個々の不規則性をできるだけ中和することにより、不確定性を軽減するためのものだ。彼らは、このような契約によって利益を得た一人に対し、別の者が同じ量損失を被るか、あるいは多くの人がわずかな損失を被ることもあることを知っている。また、彼らはほとんど誰もが、管理に必要な費用を差し引いた上で保険料を支払った金額と同等の利益を得られるという「公正な」契約を見つけることができないと理解していることも知っている。しかし、それでもこの状況を選択する。

これは、比較的固定された財産を残す方が非常に不確定な財産よりも断然良いと考える人々からなる集団である。彼らは完全に自覚しており、このような保険システムが採用されない場合と比べて、管理のための避けられない費用により、比較的固定された金額は平均的な財産よりもわずかに少ないことを理解している。

この著作は通常の保険に関するものではないため、このような団体の正確な性質についてさらに述べる必要はありません。それらは様々な種類があることに注意するだけです。時として、これらは我々が比較したように、特定のグループの人々間での相互約束によって、一定の範囲で損失を補うものである実際的な協会です。このカテゴリーには、相互保険協会、利益団体、一部の機能に関しては職業組合、そして様々な名称を持つ無数の他の団体が含まれます。

時として、所有者や株主による利益のために運営される企業のようなもので、これらは他の産業企業と同様です。私はこのように考えますが、一般的な生命保険協会の大半はそのようなものです。また時として、それは国家が管理を引き受ける場合があり、これは当国の郵便局の保険業務の場合と同じです。

5. この原則の適用範囲には必然的な制限はありません。一部の国民の大多数が安全への強い好意から、ほぼ生活全般をカバーする大きな保険協会を設立することも十分考えられます。彼らは確かに不確定性を完全に取り除くことはできませんが、それでもその最悪の結果を大幅に軽減することは容易です。彼らは生命、家屋、船舶、他の突然かつ全額の損失が可能な物などだけでなく、事業自体を保険する決定をするかもしれません。これは破産だけでなく、偶然の不利益な年も避けることを含む、同一の交換原則に基づくものです。

このような目標がこの言葉で紹介された場合にいかに不慣れであるかは明らかですが、保守的な階層に対する疑惑の名目にもかかわらず、それには多くの注目が向けられています。実際、これは根本的には共産主義と同義であり、すべての事業分野を考慮に入れる普遍的かつ強制的な保険協会で定義されることが可能です。少なくともいくつかの形態では、社会や家庭生活の領域にも侵入する可能性があります。

虽然如此全面的計划不大可能在任何大規模范围内实际推行，但这一原則本身却正在各个细节领域稳步扩展。例如，工会常被指责的一个问题是它们过于频繁地试图通过工资均等化来确保工人的利益，从而在一定程度上防止了无能现象的发生，正如它们正确且明智地防范疾病和失业一样。再如，商人互助保护协会为因不得不進行法律诉讼而造成的偶尔損失提供保险。许多城镇还設有针対小偷小摸行为的社团，旨在避免同样的不确定甚至可能严重的損失。作为为人民而非由人民提供的保险实例之一，当然有英国济贫法这一巨大范例，在此制度下，贫困者不仅因疾病和年老，而且因任何原因（除故意懒惰外）而失去工作时都可获得保障。虽然他们确实没有全额支付保费，但因为他们大多承担了部分市政及郡级税收的负担，因此无疑应被视为支付了一部分保费。此外，在公共和私人服务的一些领域中也采用了从工资或薪水中扣除一定百分比的方法，以实现半强制性的死亡、疾病或退休保险制度。

6. 保険と関連し、確率の応用として位置づけられる一方で、当然対照的に位置するのが賭博である。厳密には、これらいずれの行為も他方の逆数とは呼べないとしても、それらは反対の精神傾向から生じるという点では正しいと言える。

ある人々は、生活が自分たちにとって単調すぎると感じたり、あるいは予測可能なものがあまりにも多く且つ優越していると感じる。このような人々は、変化を求めるため二つの方法を選ぶことができる。一つには、比較的予測不可能な出来事の結果を強化し、濃縮することである。これらの出来事自体が重要でないために、強い感情を引き起こすほどのものではないが、賭け事によってそのような結果を強調することができる。もう一つは、完全に人間の予測から隔絶した個々の可能性に基づいて金銭的な結果を決定することである。これが本格的な賭博であり、主にカードやサイコロ、ルーレットなどの手段を通じて行われる。

我々が述べたように、賭博心は不確定性と多様性による興奮を求めている。したがって人々が長期間賭博を続ける場合、特に賭け金が自分の財産に対して比較的低いときには、その不確定性は必然的に減少する。完全に経験豊富な賭博者であれば（ゲームの技術が重要な役割を果たす場合）、通常よりも高い技術を持つと見做されるべきであり、これはリスクと興奮を伴う職業であると言える。更にはその卑劣でしばしば不正なものという性質を考えればなおさらだ。

一方で、その技術が平均以下に低下しているか、あるいは技術が重要でないゲームで賭け金の割合が僅かに不利な状況にある場合（例えばカジノ）、彼は趣味の一環として賭博を続けることができる。この趣味が彼にとって毎年一定の出費を伴うものであるとしても、それは他の多くの趣味と何ら変わりないと言える。

以上のように、賭博は予測不可能性や多様性による興奮を求めることから生まれる。その結果、長期間賭博を続ける人々は、特に賭け金が財産に対して比較的低い場合、その不確定性は自然に減少する。完全に経験豊富な賭博者であれば（ゲームの技術が重要な役割を果たす場合）、通常よりも高い技術を持つと見做されるべきであり、これはリスクと興奮を伴う職業であると言える。更にはその卑劣でしばしば不正なものという性質を考えればなおさらだ。

一方で、その技術が平均以下に低下しているか、あるいは技術が重要でないゲームで賭け金の割合が僅かに不利な状況にある場合（例えばカジノ）、彼は趣味の一環として賭博を続けることができる。この趣味が彼にとって毎年一定の出費を伴うものであるとしても、それは他の多くの趣味と何ら変わりないと言える。

7. 当然、すべての場合において、上記のような合理的な説明を提供することはできません。多くの初心者と狂信者は、必ずしも勝者となることを確信しながら着実に損失を被りますが、そのような無知や精神的な歪みが、彼らが例外ではなく通常の存在であるほど広く普及しているとは思えません。確かに何か非常に一般的な動機があるはずです。それらはこのような資源によって満たされるもので、それがただ一つの愛着であるとすれば、完全な予測が不可能な場所でのみ見つかるような多様性とそれに伴う興奮を提供するものです。

ここでは、人類における様々な動機の相対的な頻度についての一般的化は非常に困難です。しかし、静かな普通のウィスト（カードゲーム）プレイヤーが「愛」で行われる試合と小さな賭け金がある試合との間の違いを考えると、損失や利益の興奮を強調したいという願望の影響は高い価値を持つと考えざるを得ません。

私は一時的にギャンブルの実践に伴う実用的な危険性を軽視するつもりはありません。賭け事を続けると長時間経つにつれて、賭け金を増やしたいという几乎不可抗力のある衝動が生じることが指摘されています。そして、この傾向により不確定要素が再導入されることが主な危険と害となるのです。

リスクと不確実性は、文明化された生活においても依然として標準的な特性であり、そのような傾向を新しい分野に拡張すること自体は非常に深刻な見通しを提供しないものです。しかし、経験が示すように、それがあまり軽視できないほど危険であるとされる fascination によって人々が過度の長さまで追求する恐れがある限り、その懸念は正当化されます。

8. 上記の賭博と保険に関する一般論は、私にとって唯一の理にかなったかつ健全な分割原則であると考えます；すなわち、その原理に基づいて、様々な名称で知られる賭博や保険の異なる実践が、その結果をもたらす精神に従って整理されていることからです。もし現在使われている名称に基づいて判断しようと試みるならば、何らかの系統的な分類に至ることはありません；実際、これらはすべて、支払いと受領が必ずしも確定的でないという共通特性を有しており、その結果、名称はしばしば交換されることがあります。

たとえば、抽選会と事故に対する通常の保険組合について、過程だけを考えるならば、実質的に全く同じです。それぞれにおいて、小さな確定的な支払いと大きな不確定な受領があります。多くの人々が小さな保険料を支払いますが、その中で少数のみが賞金を得て、残りは出費に対して何の返答も受けません。また、両者とも総合的に見れば、収入と損失が事業を営む者の利益を考慮に入れた上で相互に相殺されることを目的としています。しかし、当然のことながら、保険者が賞金を得る機会を考えると、抽選会での偶然の受領と、事故による骨折や体調不良に対する一部の補償としての受領とは全く異なることが明らかになります。

また、賭博の言葉はほぼどのような種類の保険にも適用できるように容易に拡張することができます。実際、デ・モランは生命保険を個人が会社と結ぶ賭け事として説明しています。彼が特定の年齢を超えて生き残らないことを賭けています。若くして死ぬと金銭的に利益を得ますが、老いて死ぬと損失を被ります[4]。

ここでも、技術的には完全に正確な表現であることに注意してください（なぜなら、無生産的な冒険に対して意図的な資金リスクを負うことは、賭け事の定義に該当するからです）。しかし、金銭リスクを単なる冒険の視点だけで見ることと、リスクと損失をできるだけ減らす目的で賭けることとは、広範な差異があります。実際、スポーツ生活の言葉が導入されるべきであれば、私は保険者を「自己防衛」（hedging）として表現すべきだと思います。

9. また、トントインでは、しばしば保険と呼ばれるシステムがありますが、その一部は確かに適切ですが、実質的には単なるギャンブル取引に過ぎません。私は信頼できる情報によれば、それらは長らく廃止されており、かつては特にフランスで人気がありました。

この計画では、国家やトントインを管理するいかなる社会も、最後のメンバーが死ぬまで利益を得ません。生存者の数が減少するにつれて、同じ総額の年金がまだ生きている者たちに分配され続けます。そのため、各々は徐々に増加する金額を受け取ることになります。

したがって、通常の計画では最初に死亡する人々が最も多くの利益を得るのに比べて、早期に死亡した者は最も少ない金額を受けることになります。なぜなら、各メンバーが死ぬとその年金は彼や彼の関係者にとって絶対的に終了するからです。したがって、この全体としての事態は実質的には巨大なギャンブルシステムであり、誰が最も長生きできるかを見るためのものです。国家は共通の保証人となり、監督に対する重い手数料を得ます。その手数料は自然にこのような取引を奨励する国家の唯一の動機となります。

フランスのトントインの一つについて記録がありますが、97歳の寡婦が最後の生存者として73,500ルイの年金を受け取ることになりました。彼女は当初、単純な300ルイを投資しただけでした。このようなシステムは、時々保険と呼ばれても、真の保険の精神とは根本的に対立しており、現存する財産の不平等を増幅させようとします。ここで保険者は長生きすることを賭けますが、通常の保険では若く死ぬことを賭けます。

最後の例を挙げてみよう。一般的な見解では、ルージュ・アンド・ノワールのテーブルを運営する銀行やその利用者を賭博とみなすことが多い。技術的にはそうかもしれないが、実用的な観点から見れば、そのような銀行は普通の保険会社と同じくらい確実で安全なビジネスであり、おそらく製造業や商業都市における多くの通常の商売よりも収入面ではより安定している可能性が高い。

銀行はその資本に比例して多くの小さな取引を行い、顧客もその所得に比例して大きな取引を行うことが多い。当然のことながら、前者が年々利益を得るためには、テーブルが明らかにわずかに有利であるという事実に依存する。しかし、個々の損失と比較してこれらの利益の安定性は単純に、つまり実質的に、確率論における根本的な一つの大切な対比を示すものである。

これまで保険と賭博をそれぞれ自然な衝動の産物として見なし、経験から人間の判断がその両者にとって大きな支持を持っていると考えてきた。しかし、一般的な道徳的判断は前者を称賛し、後者を非難する。この判断は主に、保険の正当性と賭博の不正義性に対する一種の前提に基づいているが、これは確かに多くの真実を含んでいるものの、しばしば一方的な解釈により誤りが生じる。つまり、保険の義務は過度に強調され、賭博の実践は実質的な自己利益の犠牲であると広く認識され、結局のところそれは一貫した誤謬であると見なされる。

この前提は保険の認知された利点から抽出され、その後賭博の実践を非難するために引用される。しかし、そのようなやり方では、人類が賭博の支持に投票してきた規模と歴史を考慮に入れれば、賭博の実践もまた勝ち得る資格があることを十分認識していないように思われる。

$$
\text{保険と賭博はそれぞれ自然な衝動の産物であり、経験から見れば人間の判断に大きな支持を持っている。}
$$

しかし、一般的な道徳的判断は前者を称賛し、後者を非難する。この判断は主に、保険の正当性と賭博の不正義性に対する一種の前提に基づいているが、これは確かに多くの真実を含んでいるものの、しばしば一方的な解釈により誤りが生じる。

$$
\text{つまり、保険の義務は過度に強調され、賭博の実践は実質的な自己利益の犠牲であると広く認識され、結局のところそれは一貫した誤謬であると見なされる。}
$$

この前提は保険の認知された利点から抽出され、その後賭博の実践を非難するために引用される。しかし、そのようなやり方では、人類が賭博の支持に投票してきた規模と歴史を考慮に入れれば、賭博の実践もまた勝ち得る資格があることを十分認識していないように思われる。

$$
\text{つまり、もし人類の無言の投票の規模と歴史だけを考えれば、賭博の実践も確かに勝ち得る資格を持っていると言える。}
$$

もちろん、私たちが扱うすべての場合においては、総合的な財富は変化しないことが明白です。お金はただ一人から別の人に移動するに過ぎません。一方の損失は他方の利益と正確に等価です。

少なくとも、我々が便利と感じて始める近似真実とは、このような見解です。もし財富によって与えられる幸福が常にその量と直接比例していたら、保険を推奨する理由もまたギャンブルを非難する理由もないでしょう。後者の場合については明らかです。私は50ポンドを失ったかもしれませんが、他の人（一人か複数人かもしれません）はそれを得ました。それによって彼らの幸福が私の減少分と完全に均衡するでしょう。

保険の場合にはやや複雑な点があります。これは、保険金支払いがランダムではなく、可能な限り損失を補うために特定の条件が設定されているからです（既に指摘した通り、それがあらゆるロトであるなら単なる抽選となります）。私は鉄道旅行で保険に加入し、事故で足を骨折します。3便士の切符代を支払い、保険会社から200ポンドの補償を受け取ったとしましょう。しかし、幸福が常に財富と直接比例していたら、この200ポンドによって得た幸福は16,000人が3便士を支払って返り咲かない損失の総合的な減少分と均衡するでしょう。

11. 保険の実践は、私は考えると理論的な興味を惹くような多くの問題を生じさせることは稀であり、従って長時間にわたる議論をする必要はありません。事実は、まだ十分な期間と広範囲にわたって適用されていないし、統計的に正確な扱いが可能な事柄に対しては特にそういった状況ではありません。ただし、一つの分野では例外です。もちろんそれが生命保険です；しかし、その主題は詳細な統計的観察を必要とするものであり、適切にルーチンに基づいて実施されています（正確には）。この点を例示するためには、デ・モランの著作を参照すれば十分です。彼はプロフェッショナルな保険精算士でもあり確率論の理論家でもあり、この種の考慮事項から得られる例題で推測的な確率論の扱いに助ける機会がほとんどなかったと述べています。

12. 一方、賭博に関しては異なる状況です。様々な興味深い単独の問題が議論されており（ペトロブスキーの問題が最も有名です）、重要な推測的な問題もいくつか提起されています。その一つは、賭博の実践の欠点に関するものです。多くの著者が、過度な興奮やそれに伴う greed 、長期的には成功を保証するためだけに必要な安定したビジネス習慣からの離脱、財富が不正な手に流れるなど、明らかな道徳的および間接的な悪影響だけでなく、賭博の実践が必然的に損失を引き起こすことを証明しようと試みた人物もいました。

二人が対戦する場合、いずれかの一方は最終的に破産することになる。これは公平なゲームであっても、十分長く続けばそのような結果が避けられない。収入が少ない方が当然ながら幸運な生存者となる確率は低い。一方に有限の収入があり、他方に無限の収入がある場合、彼らが十分長く対戦すれば明らかに前者が最終的な破産者となる。この主張によれば、これは実質的に個々のギャンブラーの立場を反映している。「一人の相手だけと賭け事をするわけではない。投資家は一般の人々との取引を行い、その資源は実質的に無限大である世界と対戦している。彼自身が何も賭けるものがなくなったことで自分の運が悪くならない可能性があるが、世界的な資源が枯渇することで利益を得る可能性は存在しない。誰もが不平等な戦いに身を置いている：有限の資金で無限の敵と戦っている」ということである。

上記の議論では、彼が対戦する相手がギャンブル場における銀行のように共有財布を持つ一団ではないという点が見落としている。これらの各々の相手は彼自身と同じ立場にあり、彼らそれぞれが最終的に破産することを示す類似の証明が可能である。これは当然ながら論理的な矛盾につながる。ギャンブルはただ一方から他方に金銭を移動するものであり、したがって実際に失われるものは存在しない。

13. 彼らが極端な状況で金銭の運命は、見通し通りです。まず、有限のプレイヤーを仮定します。十分長く続けば、最終的に金銭は誰かのポケットに入るでしょう。元々の財産が等しい場合、幸運な生存者になるのはどのプレイヤーも同じ確率であります。その場合、何らかの数値的な根拠から、彼らのいずれか一人に対する賭け事の見通しは不利であるとは主張できません。元々の財産が異なる場合、最初に最も多くの金銭を持っている者は、特定の法則に基づいて最終的に勝者となる確率が高いことが示されます。その場合、彼自身には有利であり、彼の資産が少ない競争相手には不利であると言えます。

プレイヤーが有限ではなく無限の場合を考え、破産した者はテーブルから退き、新たに参加者が入るという状況を仮定します。この結果は複雑ですが、その一般的な傾向は容易に識別できます。もし誰もが破産したときにのみ退場するという不合理な仮定を立てれば、古いプレイヤーたちは徐々に裕福になる状態となります。この場合、新規参加者にはますます悪い見通しが続き、彼らが富んだ相手と勝負する確率は非常に小さくなります。しかし、これは合理的ではないため、破産した者は退場し、一定の財産を獲得した者はも同様に退場するという仮定の方が適切です。これにより、賭博集団全体の総資産と平均的な財産が比較的安定します。

特定のプレイヤーが破産する確率や、破産する前にどれだけ持続できるかは、プレイヤーの初期収入、ゲームのルール、賭け金額、他の考慮事項によって当然に異なります。しかし、明らかに一人が失うものは他者が等しい量を獲得することになりますから、特定の賭博者は以前述べたような仮定に基づいてその行動が推奨されないことを警告されるべきです。

14. さらに、この見解の正当性を補強するために、読者は最後の例で示された状況が、多少異なる言葉で表現され、若干の事情の変更を加えた形で、私たち周囲のあらゆる場所で無休無歇に行われていることに注意していただきたいと主張することができるでしょう。これを商取引国の商人や取引者の競争と呼ぶことにしましょう。その一般的な結果は十分に知られています。技術的な要素が問題となる限り、彼らは適切にはギャンブラーとは呼ばれないかもしれませんが、運とリスクの側面から見れば、公平にはギャンブラーと呼ばれるべきでしょう。また、多くの事業分野においてこれは必ずしも例外ではなく、相当程度の範囲でそのような状況が存在します。ビジネスが放漫に進めばこそ、一般的な観点からはこの比較は適切です。それぞれの場合において、破産する者と成功を収める者が見受けられます。また、他の条件が等しい場合（coeteris paribus）、大きな財産を持つ者はその確率的に有利であることがわかります。各々は、ある意味では集団的な商取引世界と競争していますが、それぞれの競争相手も同じことを行っているため、私たちは誰か一人（特に貧しい者を除く）に注意を促すことができません。その努力が最終的には不利になるという警告を。

15. この結果を最も決定的な形で見せたいのであれば、全く異なる事件の例に類似性を見出すことができます。すなわち、\_姓の運命\_ にあります。私たちはこの点ではすべてが賭け事であり、そのゲームはニューカマーかモンテカルロでのものとは比較にならないほど厳格に行われます。完全な扱いを考慮に入れると、この主題は非常に複雑なものですが、単純な例でも一般的な傾向を示すことができます。

仮に1000組の異なる姓を持つ夫婦からなる集団があり、それぞれが4人の子供を持ち、成長して結婚するという状況を考えましょう。約16分の1の家族は全員が娘であるため、通常の慣習では約62の姓が次の世代で永遠に消滅します。また、そのうちの4分の1の家族はただ一人の息子を持つことになり、それぞれは当然その父親と同じ状況に置かれます。すなわち、その名前は単独の代表者となります。したがって、次の世代では約16分の1の姓が消滅し、このプロセスは続けられます。各々の世代で消滅する数は次第に減少しますが、生存者の安定性が増えることに伴い、その数はより多くの人数を持つことによります。しかし、この過程には何らの制約もありません。

16. この類推は非常に密接であり、名前が消えていくことに対応するのは破産した賭博者たちで、名前の数が増加することに対応するのは運の良き勝者たちです。それぞれの場合において最終的な目標は同じで、もちろんそれは極めて遠いものですが、他のすべてを犠牲にして唯一生存するというものです。その一つの姓もまた次々と消えていくこと自体は、誰が家系の歴史について調査した人でも気づいたことでしょう。そして、確率論に疎い人々によって様々な想像的な説明がなされてきました。

しばしば見落としてしまうのは、各世代における生存への「運の転換」の極めて微細さです。上記の数値例では、人口が毎世代で倍増することを前提とし、非常に有利な仮定を行いました。古い国や人口が厚い国においては、数字が非常にゆっくりと増えている場合、平均的な効果的な家族（つまり、結婚に至る子供の平均数）がただ二人であるという仮定の方が近いでしょう。この場合、どの時点でも単一男性で表される家族は四分の三の確率で絶滅しないことを維持し、もちろん薄まることの過程はより速くなるでしょう。

17. 最も興味深い賭博の欠点を証明しようとする試みのクラスは、「技術的には道徳的な期待」と呼ばれるものから「数学的な期待」に異なる点を引き立てるものである。後者の数学的期待は単純に、問題となる冒険の平均的な金銭価値と定義される；すなわち、獲得（または損失）する可能性のある金額とその確率との積である。たとえば、私は一枚のサイコロで1対4という賭けで王幣を賭けて、一回の投げでエースが出るか出ないかについて賭ける場合、多くの投げの中で平均すると、5ポンドの利益に対して4ポンドの損失があり、つまり各6つの出現ごとに3先9便（3シリングと4ペンス）の平均的な利益がある。これが真のまたは数学的期待である。

一方、「道徳的な期待」とは、この数学的期待が主観的な価値として評価されることを指す。つまり、通常通り金銭的な財産を数える代わりに、その財産が感じられる価値で数える試みが行われる。計算の要素は、ポンドやシリングではなく、実際または将来の楽しみの合計となる。

したがって、将来的な利益の現在の価値を計算するときには、獲得確率ではなく、主観的な価値を対象とするべきである。

この道徳的財産と物理的な関係についての正確な関係は、さまざまな程度の恣意的な仮定がなされている。ある著者（ブフォン）は、特定の金額の道徳的価値が獲得者の総資産に対して逆比例すると考えている。また、異なる仮定から出発する別の著者（ダニエル・ベルヌーイ）は、その実際の量の対数に比例すると主張している。さらに、第三の著者（クラマー）は、その量の平方根に比例すると述べている。

18. 歴史上、これらの提案はパルメンツィ問題の結論と実用的な常識との調和を望む気持ちから生まれたものであり、物理的評価に道徳的な評価を代入することで期待値の総合価値を有限な合計にまで低減することが可能となった。したがって、この理由に基づいてこれらの提案は大きな興味を持たないため、我々が見たように、適切に理解すれば問題には深刻な困難は存在しない。

しかし、同じような提案は賭博が必ずしも不利であることを証明するために使用されることもある。そしてこれは両当事者にとってである。例えばベルヌーイの仮定を考えてみよう。2人の人がそれぞれ50ポンドから始める場合、10ポンドを賭け金として等しい賭博を行うと、幸せ度が損失するということが容易に示される。つまり、60ポンドを持つ人からの喜びは40ポンドで止む人に比べて劣ると考えられる。

19. このような議論の形式は一般的であるが、その形態では結論が十分に説得力があるとは思えない。確かに反論できるのは、証明されたものは不均衡が悪いことだけであり、50ポンドずつの2つの財産よりも60ポンドと40ポンドの1つずつの財産の方が良いという理由からだ。例えば、元々の財産がそれぞれ60ポンドと40ポンドだった場合、結果として幸せ度は増加する可能性がある。これは、裕者の方が負けて両者の財産を均衡させることが起こる場合に特に真実である。これが正しいのであれば、我々は次のように示さなければならない、それは非常に容易なことである：もし他の選択肢が実現し、2つの財産がさらに不均衡になる（つまり65ポンドと35ポンドの場合）と、そのような幸せの損失は等分化によって得られる幸せよりも大きく補填されるだろう。そしてこれらの2つの仮定が等確率であるため、長期的には損失があることになる。

上記の議論は、幸福と財産との関係に関する一般的な結論を述べる方法が、むしろ誤解を招くものであることを示唆しているように思われます。そして、仮定される幸福と財産との関係についての前提に基づいて考えると、賭博が不利である要素は、財産の不均等性の増大傾向にあるという実態が明らかになるのです。

20. この事実は二つの問いを提起します。一つは私たちの主題に関連する推測的な興味を持つ問いであり、もう一つは人生における最高の重要性を持つ問いです。まず第一に、特定の道徳的状況や幸福の変動に関する仮定を除いても、賭博が既存の財産の不均等性を増大させる傾向があるという事実は存在するのでしょうか？理論的には、これは間違いありません。最も単純なケースとして、二人が一ポンドを争う場合を考えましょう。最初に彼らの運勢が等しいと仮定すると、結果的に不均等が生じます。一方で、運勢が異なる場合、不均等が増大するか減少する可能性は半々ですが、増大の割合の方が減少の割合よりも大きいので、最終的には運勢が等しいときと同じような結果になります[11]。より一般的な視点から見ると、これは我々が大量の数値について述べる結論の根底にあるものであり、つまり、絶対的な分散が絶対的に増大（相対的には減少する場合もある）するという結論です。そしてもちろん、財産の分配について話しているときには、これらの絶対的な分散を指しています。

21. このことは理論上の結論である。実際の賭博行為に反対する補強因子がどの程度作用しているかは、その専門家たちの判断に委ねるべきである。外野者が新聞や他の公的情報源から得られる情報を基に推測すれば、これらの補強因子は相当大きいものであり、従って賭博を批判する論拠としては不安定なものであるように思われる。多くの大金は公然と競馬場や賭博場で浪費され、そしておそらくその大部分は比較的貧しい人々の収入を増やすために使われたに違いない。しかし、この問題の解決は確実には、確率論に関する著者が期待できる程度の情報源より優れた機会を持つ者たちに委ねるべきである。

22. 一般に結論として導き出すべきことは、この「道徳的運命」の原理を賭博反対の議論として引用した人々は、彼らが認識していたよりも遥かに複雑で広範な問題を提起していたということである。彼らが取り組んでいたのは、「財富の分配は最大限の幸福を確保する傾向を持つものを何とするべきか？」という問いだった。これは平等性によって还是不平等性によって達成されるべきか？彼らが賭博を非難した根拠となった教義に従って真剣に追求すれば、社会主義者の理想を選択すべきであることが明らかになったはずだ。そして、彼らは賭博に対して表明した批判と同じ理由で、現代文明の傾向全てを非難するべきだった。例えば、激しい競争、投机性操作、信用取引の拡大、機械的な発明、事業規模の大幅な拡大など、これらは既存の不平等を増加させる影響を持ち、社会主義者によって批判されてきた。これらの傾向全てをこのテストで評価し、非難すべきかどうかを確定するべきであるか？読者が「道徳的運命」の仮定が導き出すべき問題の種類を見たいと思うのであれば、F.Y.エドゲワースの『数学的心理学』を読むことを推奨したい。私が知る限り唯一、これらの問いに取り組んでいる著作である。

1. 小児麻痺に対する接種について、前一世紀の確率論者の間で議論が交わされた問題は、同一原理が非常に異なる事例に適用される典型的なケースである。保険の対象となった損失は小児麻痺による死亡であり、支払われる保険料は病気やその他の不快感、そして接種から生じる極めて小さな死亡リスクであった。そのような議論が発生した際には、一定の軽微な損失と潜在的な大きな損失とのバランスについての困難性が浮き彫りになった。17世紀には地中海への旅に出る前に、モロッコ海賊によって捕らえられるリスクを保険でカバーし、ransom（人質解放金）が支払われるよう確保するという慣行があったことがあった。（詳細については、Walfordの『Insurance Guide and Handbook』をご覧ください。科学的な精神とは言えませんが、保険に関連するあらゆる事項に関する情報が含まれています。）

2. 以上の比較は、共産主義が目指す理想が保険と似ていることを意味するもので、その実現過程や強制手段を考慮に入れると、当然ながら異なる観点から捉えられるでしょう。次節に述べる労働組合の行動についても同様です。

3. 最近見つけたこれらの主題についての最良の議論の一つは、能力と知識が十分に備わった著者が書いたもので、Mr Proctor の『機会と運』の中にあります。しかし、彼はギャンブルに対する愚かさではなく、不正さを全面的に非難する極端な態度に陥っているように私には思われます。また、賞品の総額が支払金額よりも少ないという理由だけですべての抽選会を「不公平」としてすら「詐欺」呼ばわりするのは言葉遣いとしては強引だと考えられます。多くの抽選会が支持された際における故意に隠匿した情報や、誤った報告を意図的に広めた差異は、正当な意味での詐欺と同名を冠すべきバブル企業とは異なるべきであるという点から、両者に同じ名称を適用することは適切ではありません。

4. 「火災保険は事務所と当事者の間の単純な賭けであり、生命保険は一連の賭けの集まりです。未来の出来事の結果が利益や損失をもたらすすべての取引においては、賭けの原理が存在します。」『ペンニー百科』、賭けに関する項目。

5. 『方法的百科全書』、トントネに関する項目。

6. 当然、政治経済学の考慮事項を導入する場合、修正が必要となります。まず第一に、デ・モランが繰り返し強調しているように、保険事務所は、単なる保険事務所だけでなく、貯蓄銀行でもあります。事務所は保険料を投資しますので、それにより通常よりもより大きな金額を支払うことが可能となります。また、ギャンブルの場合においては、個々の資本の大規模な損失はほぼ必ず実質的な財富の破壊を引き起こすことになります；更に言うと、実際にはギャンブルは生産的用途から非生産的用途への恒常的な財富の移転をもたらすことが多くあります。

7. 『選択と運』（第2版）、第208ページ。

8. 私が知る限りでは、最初にこの問題を真剣な形で扱ったのはMr Galtonであると考えられます。（詳しくは、1875年の第4巻の『人類学研究所誌』をご覧ください。完全な数理的解決法はMr H. W. Watsonによって示されています。）

9. ベルヌリ自身は、その結論を実際の経験に基づいて形成したわけではなく、しかし彼が最初に仮定した主観的な価値に関する前提が、所有している金額（x）に対する微量の増分（dx）の主観的価値が逆比例するという仮説は、その上で述べられた対数法則へと即座に導き出すものであり、これは現在では心理学者にとって馴染み深い前提と同一であることに注目すべき点がある。これは通常フェンシャーの法則と呼ばれるもので、彼は多くの精密な実験を基にして、我々の単純な感覚における数々のケースにおいて証明している。しかし私は、そのような法則が財産による幸福度の遥かに複雑な依存関係においても通用することを主張したとは信じていない。

10. この道徳的な幸福を表現する公式は c log x/a である。ここで x は当時の物理的財産を表し、a は幸福が消失されると仮定される小さな値を表すものであり、c は任意の定数である。二人の者がそれぞれ x の財産を持っており、y を均等な賭けで冒険する場合、幸福度における均衡は次のようになる：
   c log x/a と 1/2c log (x + y)/a + 1/2c log (x - y)/a
   すなわち log x2 と log(x + y)(x - y)
   すなわち x2 と x2 - y2
   前者の値が常に大きい。

11. この点は次のように明確に理解できる。二組の賭博者を考える。それぞれの一組には $p$50 を、もう一組には $p$30 の持ち金があるとする。もし富者が一方で貧者が他方で勝つ場合を想定すると、これら二つの結果は長期的には平均的なものとして適切に表現されるであろう。なぜなら、この状況では二つの選択肢しかないし、長期的にはこれらの選択が等しく頻繁になるからである。明らかに $p$50 の財産と $p$30 の財産がそれぞれ二つずつあり、それらが $p$20 と $p$40 と $p$60 の財産へと変換されることを意味する。これは明確な不平等性の増大を示している。

\newpage

\chapter{証言への確率論の適用（Application of Probability to Testimony）}
\label{ch:xvi}

\begin{center}
\large\textit{証言信憑性の統計的評価・前提条件の不成立・証言連鎖の数学的分析}
\end{center}
\vspace{4mm}

第十六章.

証言の信頼性に確率を適用する

1. この著作で採用された原理に基づいて考えると、いくつかの問題の分類が、表面上は既存の権利を持つように見えていたとしても、確率の科学から除外されるべきかどうか疑問が生じる。おそらく最も重要なのは、証言の信頼性についてのもので、これは直接的または間接的に私たちに届く、陪審員の洞察力と信用性を通じて伝わるものを指す。科学に関する多くの論文は、さまざまな程度の信用性を持つ証人の報告の組み合わせに対する信用をどのように付与するか、あるいはより多くまたは少ない人数からなる陪審団の判決に対する信用を評価する基準について議論している。ポissonという偉大な現代数学者は、この主題に関する詳細な論文を書いている一方で、LaplaceやDe Morganなど他の多くの著者の著作の大部分は、類似した調査を検討するために費やされている。

これらの権威者から異論を唱えるのは慎重であるべきだが、私はこれらの問題が確率の観点から考慮されるべきではないという理由があるため、そのような見解に同意する。以下はその主な理由である。

2. 以前の導入章において、推論を行うための適切な命題が提示されずに個体が現れ、我々に必要なグループや系列を選択し、その中に該当するものを参照することを課すという過程について詳細に調査したことを思い出すだろう。つまり、事象が発生する頻度から確率を計算するための確定的な条件が直接経験によって得られたり、論理的に推測されたりしている場合とは異なり、我々はより適切な多数の条件の中から必要な頻度の条件を選択しなければならない。問題がこのような段階で提示された場合、当然ながら統計を拡張して比例命題に至る前段階の過程が既に行われていると仮定することができる。したがって、我々はそのような命題の大量を持っていると想定し、主な疑問はそれらの中からどのものを使用すべきかにある。

この選択は一定の程度において任意性を持つことが示された。個体は多くの異なる特性を持ち、そのため彼は異なる系列やグループのメンバーとなることになり、それぞれに異なる平均値を提示する可能性があるからである。現在では、このような問題が生じる困難が重要でなくなる実用的な条件について、以前よりも詳しく検討しなければならない。

3. この種の条件のうち、一つは非常に単純で明白なものである。
我々が提示される異なる統計データが実際には著しく異なる結果をもたらさないことを要求するのである。たとえば、40歳の男性が年内に死亡する確率について調査している場合、彼が赤毛である事実や特定の県または町で生まれたという事実を考慮に入れることが可能である。これらのそれぞれの状況は個体を特殊化し、従ってそのケースに対する適用可能な統計の範囲を制限する。しかし、このような性質を考慮することにより、平均値が全く同じままであるか、またはそれほど重要な変更が生じるかぎり、誰もそれを考慮に入れる価値がないと考えられる。

確かに、どのような特性についても絶対に結果に影響を与えないとは断言できないが、このような特性についてはその影響は完全に無視できると確信することができる。もちろん、確率の計算に最も適した極端な場合においては、即ち純粋な機会ゲームの場合には、これらの付随的な特性は全く関係ない。実際的に公正な一組のカードについて、すべてのプレイヤーが知っている情報以外の詳細な情報は結果に関する何らの情報をも提供しないのである。

また、異なる統計データセットが上記のようにほぼ同じ結果を示さなくても、それらは一つのグループに属するすべての統計がその結果において実質的に一致するように分類されるという非常に似たような結果をもたらす可能性があります。例えば、肺結核患者が生命保険に加入したい場合、その健康状態を考慮に入れるか否かによって統計結果は著しく異なることになります。ここでは二つの統計データセットがあり、これらは明確に区別され、自然種類の差異とほぼ同等の地位を占めることもあります。したがって、それらは実質的に異なる結果を示します。さらに詳細に分類する場合、最後の段落で述べたような微細な特性を考慮に入れるかもしれません。これにより、より個々の人と似ている人々向けの制限された統計データセットを得ることができます。しかし、これらの統計が十分に異なる結果を示さない限り、そのような作業を行う価値はありません。つまり、特定のケースに対して適用可能な異なる命題は、実質的に必要な少数のグループに分類され、それらのみが考慮されるべきであり、これにより実際的な選択肢の範囲は大幅に狭まります。

以上の原著英文を4. これらの条件の理由は、難しくはありません。これらの条件が存在する場合、個々人を参照すべき系列やクラスを選択するプロセスは非常に単純で、選択は特定の推論目的にとって最終的なものとなります。保険の場合など、例えば、我々が決定しなければならない質問は非常に単純なものであることが多く、A.B. は特定の年齢の男性でしょうか？そのような状況にある50人中1人が年内に死亡するでしょう。さらに他の質問が必要な場合でも、それは以下のようになります。A.B. は健康な男性でしょうか？彼は危険な仕事に従事していますか？しかし、これらの質問に関しても、問題となるクラスは限られている上、その境界線は相当正確であるため、個々人をそれらのいずれかに参照することは容易です。そして一度クラスを選択した後も、さらなる考慮事項から解放されます。なぜなら、他の統計が重要な平均値を著しく変えるものではないと仮定されているので、既に述べた性質以外のものは考慮する必要がないからです。

確率論におけるゲームの例は、これらの条件がほぼ理想状態で満足していることを示す好例を提供します。公平な賭博のために彼らが適切であるような同じ状況は、確率の例にもなります。サイコロを投げる場合、すべての人々は結果に対する知識と無知について平等に立っています。唯一のデータは、各面が平均して6回中1回出ることです。

5. 今度は証言の信憑性という問題に取り組む場合において、上記の条件がどの程度満たされているかを検討してみましょう。これらの問題は基本的な調査から成り立っているため、その一つとして以下のような例を考えることができます：ある証人が十分の一の確率で嘘をついているとします。この声明について私はどのような結論を下すべきでしょうか？このような方法で人間を虚偽度の段階に配置する可能性に対する反対意見は公平であるかもしれません。しかし、それについては置いておき、その証人は顔に適切なクラスが刻印されていると仮定し、従って、一般的な理由からその信用度に応じた信頼性を主張することができます。

しかしながら、この種の問題に対する反対意見は他にも強固なものがあります。

6. 前節で説明した「個体」として適切な統計のクラスや系列に割り当てる必要があるもの、つまり本場合においては \_声明\_ である。特定の事例では、この声明は既に証人の信頼性の度合いによって与えられた種類の声明のクラスに割り当てられているが、その判断を形成する際にこのクラスに限定するかどうかは私どもが任意であり、実際的に意見を形成される機会がある場合には特にそうである。しかし、この声明の場合と同様に、当該保険者の寿命に関するものと同じく、声明には与えられたクラスを特定する性質以外にも多くの他の特性が観察できる。後者では（その年齢以外に）出生地、職業の種類などがあったように、前者でも（特定の種類の証人の声明であるという事実以外に）発言された時間や場所、そして特定の状況がある。声明が作成されるときにはこれらのすべての質的または属性が私どもに存在しており、もちろんこれらの中からどれだけを考慮に入れる権利があります。現在私たちが直面している問題は、このような要素が証人の声明の真実性についての判断結果にいかほど重要でないか、生命保険における相応する要素と比較してどの程度重要であるかという点に過ぎません。このような質問に対する答えは疑いなく明確です。通常の状況では、我々が一人の死亡の可能性を判断するための条件について知ることができるすべてをすぐに理解しますので、一般的な死亡統計に満足することになりますが、証人が話す場面で誰もその信頼性の数字だけを引用することなど考えません。声明が作成された状況はむしろ重要であり、その重要性は圧倒的です。証人の外見、声のトーン、利益を得る動機があること、さらに我々がこの問題について深く考えれば次第に明らかになる無数の他の要素など、これらすべてはどんな感覚的な人間でも割り当てられた平均値から除外することを考慮に入れるでしょう。彼は実際にはこのような方法で判断するよりも、戦闘中にある兵士の寿命がどの程度であるかをカーリルスやノースハンプトンの死亡率表を使って決定することなど考えません。

7. これらの状況下で、証人の信頼性を判断するためにまだクラスに参照し、より限定的な平均値によってその真実性を評価することができるという反論は成り立たない。例えば、このような状況下で彼のような外見と行動を持つ人々のうち、九割が嘘をついていることが多く見られるという推測から判断するとするが、全くクラスへの参照は行われていないし、何らかの統計情報に即した判断も直接的に行われていない。少なくともその時刻において意識している判断や、その後で正当化のために参照すべき何かではない。この決定は、観察者の感覚の鋭さと一般的な理解力に依存しているように思われる。

証人の信憑性に関する統計は、実際には他のすべての統計と同様に永久的に適切でないものである。腿切断後の回復率についての正確なパーセンテージを知ることはできるが、そのようなテーブルだけをもとに判断をする外科医は存在しないだろうか？もちろん、当然のことながら、彼が形成する患者の回復の見込みに関する意見が過去に目撃した死亡と回復の割合に基づく可能性があることは認めないでもよい。しかし、これが実際の状況である場合、これらのデータは明確には背景に隠されているものだと言える。直接的にそして即座に統計を引用して判断形成をするわけではなく、重要な中間過程として彼が案じるべき要素と非要素を見極め、評価するプロセスがある。鋭い感覚、記憶、判断力、実用的な洞察力が必要となり、従って統計への直接的かつ単独の依拠は以前とは異なるものとなる。

外科医は脳裏に二つや三つの例を思い浮かべるかもしれない。手術自体が同等であったが患者の体質が異なる場合がある。その要素は適切に考慮に入れなければならない。また、体質が類似していたが手術がより厳重だった他の例もある。そしてそういった具合である。

したがって、最終的には統計への依拠はあり得るが、それほど直接的ではない。統計の価値は、判断と記憶というやや曖昧な媒介を通じて評価され、その観点は非常に異なるものとなる。

8. 何かウィスト（Whist）というカードゲームを知っている者は、ここでの主張された区別を適切な例として思い出すことができる。ゲームが進行するにつれて我々の推論の性質に変化があることを思い出すとよい。ゲームが始まった段階では、我々は正当に確率理論に依拠している。各プレイヤーが行動するルール、そしてそれらに基づいて他のプレイヤーがどのように行動すると推測するかのルールもまた、特定のカードの組み合わせが観察される頻度（あるいは計算によって保証される頻度）に依存している。

例えば、四枚以上のトリム（Trump）を持っている場合、なぜすぐにそれらを出し切る必要があるのか？それは確率論的な理由から、純粋な統計学的形式で述べることができる理由から、我々は大半のケースで相手のトリムを引き出すと確信しているためである。そして、その結果、自分たちがコマンドを得ることになる。

同様に、ゲームの初期段階で認識される他のすべてのルールについても言える。

しかし、ゲームが進行するにつれて、これは全て変化し、最終的には我々自身や他人が統計的な頻度に基づいて作成したルールを信頼することはほとんどない。多くの他の考慮事項が加わり、我々はパートナーと相手の性格や行動パターンに関する知識に部分的に影響を受け始める。また、実際に展開されるゲームの過程から得た多数の判断を素早く組み合わせる必要がある。その根拠は当時ではほとんど理解できなかったり説明できなかったりし、その後忘れてしまった可能性もある。

つまり、現在目の前にある特定のカードの組み合わせは、全ての情報が揃っている人にとっては、適切に分類される明確なクラスには容易に収まらない。

9. ミル（『論理学』、第3巻、第18章、§3）は、確率理論が証人の信頼性に適用できるかどうかについて類似の批判を述べています。しかし、彼は他の理由も追加しており、これらは私には同等の正当性とは思えません。彼は言います、「一般的な常識は、人類やそのどのクラスの人々の証言の正直性と真実性に資格を持つ他の要素の平均を計算することは不可能だと指摘します。そしてもし可能であったとしても、そのような平均値はほとんどすべての証人の信頼性が平均値よりも高いか低いという事実から導き出される指南针となることはないでしょう」。

この後者の批判は、寿命の推定において死亡率表から生命の長さを評価する場合にも同様に適用されます。異なる証人の信頼性には、異なる生命の長さほど広い範囲の変動が見られるためです。もし信頼性に関する統計データが得られ、そのようなデータが得られた際に適宜参照できるなら、長期的には他の同種の統計データと同等に正確な推論を提供することが可能でしょう。しかし実際には、このような統計は自然かつ正当に無視されるべきです。なぜなら、個々の声明においては、新たなクラスに属する可能性のある状況が存在し、そのようなクラスは異なる統計に基づいており、より高い確率で当該特定事例での正解を導き出すことができますからです。

このような種類の多くの場合においては、我々の意見形成過程自体が変化します。既に述べたように、このプロセスは正確な意味での統計への依存から脱却することになります。

10. 証言に関連する確率問題は、これまでに述べた例に限定されないものである。私には見えますが、特定の証人の信用性を評価しようとする試みや、その信頼性に関して彼を特定のクラスに分類しようとする試みをすべて拒否しなければならない。したがって、このようなデータから始まる多くの問題、「一度に間違う確率が十分の一である証人」などから、適切でないものとして放棄せざるを得ないとしても、それは証言が科学によって多少異なる方法で扱われるべきではないという結論を導くわけではない。特定の証人の正直性を正確な方法で評価することは不可能でも、ある種の意見を形成することができるかもしれない。例えば、ヨーロッパ人は東洋人に比べてこの点で優れていると述べることができるかもしれない。また、10人の人間（例えば陪審員）が一致した判断よりも、5人が一致した判断の方が優れている理由と範囲を説明しようとすることもできる。さらに、驚くべき話が通常のものよりも信頼性が低いことを決定する基準を設定するために何かを行うことも可能である。ただし、その点について正確で確定的な結論に至ることはできないとしても、この最後の問題は次の章の中で更に議論される。

11. その種の問題に対する科学的考慮から完全に除外する必要がなくなる視点の変化は、他の方向でも現れます。例えば、既に指摘されているように、患者の病気の個々の特性は、大多数の場合において手術医が統計的な直接的な引用によって彼の回復を判断することなく、そのような考慮事項が彼の判断にいかなる間接的な影響を与えようと、十分重要であることが示されています。
しかし、多くのケースについて意見を形成する必要がある場合、例えば大きな病院の場合などでは、統計学は再び重要な役割を果たし、適切に認識されるべきものとして正当化されます。一つの病院や治療法と別のものを比較できるようになるでしょう。この違いの基礎は明白です。個々人の特性が、私たちはそれらを個別に判断する際に平均値から離れる準備ができていたとしても、グループのケースについて判断する際には同じような混乱を引き起こさないからです。その結果、平均値は意見を形成する最も安全で利用可能な基礎となります。そして、確率が再度適用されることになります。

しかし、このような場合において確率を援用する方法は認められるかもしれませんが、それらが理論の手法と資源を説明するのに特に適切な例とは思われません。実際、ここでの数学的計算の精密さが、データから始める際に得られる相応の正確さや精度を得ることは不可能であることを考慮に入れると、その長さは全く正当化されないほどまで推し進められているという結論に至るかもしれません。

たとえば、一つの例を挙げましょう。一貫性への愛着が普通の感覚を圧倒するような事例を見つけるのは難しいでしょう。特定の主張に対する賛成票数と反対票数は関係なく、実際の過半数が同じであればそれが重要ではないという結論が認められているのです。つまり、十二人の陪審員の中から十人二人が一致する判決と、八人の陪審員全員が一致する判決は同等に扱われるべきです。しかし、この結論はポissonによって認められています。その前提条件は次の章で説明します。

また、証人や陪審員間の完全な独立性はほぼ必須の仮定となります。しかしこれをどのように確保するのでしょうか？直接的な共謀を除いても、人類は大抵の場合、同情の影響に大きく支配されています。この影響は政治的偏見、階級的偏見、地域的な感情などといった様々な名目で存在し、十分な程度まで意見形成に影響を与えます。これは個々の修正が必要となる状況を我々が特定の証人の信頼性を評価する際に考慮すべき点と同様であり、それが十分な自信を失わせるほど統計や平均値に基づいた判断を下すことを可能にするのです。

本稿は主に確率論の一般原理を説明し、確立することに焦点を当てているため、この主題については次の章で与えられる簡潔な議論までとさせていただくのが適切でしょう。

また、証言の評価において多くの場合で確率論の原理に頼るのを避けるべき理由があることも述べておきたい。しばしば、おそらくほとんどすべての場合において、我々は絶対的に決定に達しなければならない状況には遭遇しない。少なくとも有罪者を無罪と判断することを免れるという観点から言えば、決断を下さない方が良い場合もあるかもしれない。個々の事例における真実性を獲得することがより重要である場合があり、我々は必ずしも結論を形成しなければならないわけではなく、その結論が長期的には正しい方向に導く傾向があると言えるだけなら形成しない方がよい場合もあると考えられる。

\newpage

\chapter{異常な言説・奇跡の信憑性（Credibility of Extraordinary Stories）}
\label{ch:xvii}

\begin{center}
\large\textit{並外れた主張の事前確率・証人の一般的信頼度と個別信頼度・奇跡と異常の峻別}
\end{center}
\vspace{4mm}

第十七章.

不確かな物語の証言の信頼性について

1. 今や、以前に一時的に触れられたが、詳細な調査が必要な問いに戻ってみよう。すなわち、不可能に思える物語を立証するために証言が有効であるかどうかという問いだ。前章（第十二章）で、バーティルの主張について検討したことがあることを思い出す。その主張は、ミルも多少支持していたとされるが、証明以前に大きな不可能性であったものが証明後には非常に小さな不可能性になるかもしれないという主張だった。これに対し、我々が形成する信頼性の異なる評価は、事実が過去か未来であるかどうかとは関係なく、全く異なる原因によるものだということが指摘された。つまり、我々がその瞬間に抱く事件の概念（それが実際に存在している現象を正確に表現するためには判断しなければならない点）は二つの非常に異なる方法で私達に提示される。一つの場合、それは統計から得られる確率に基づいた私たち自身の推測だった；もう一つの場合、証人の主張であり、従ってこの場合の訴えは事件自体の頻度ではなく、証人の信頼性に移るべきだった。したがって、その概念や「事象」は統計的根拠に基づく推測から、特定の証拠に基づいた主張へと移行し、従って全く異なる視点で見られるべきだ。

2. しかし、この原理が導き出す結論は驚くべきもののように思われるかもしれない。つまり、事象が起こる頻度から証人の信頼性に訴えるという転換によって、主張の確率や不可能性は証人の正直さだけに関わるのか？もしそうだとすれば、嘘をつかない人が主張した物語は何でも信じるべきなのか？

この問題を解決するためには、証言が通常私たちに提示される状況についてより詳しく検討する必要があります。明瞭性の観点から数値的な例を用いるのは当然必要ですが、ある証人が全体として長期的に見れば十回に九回は正しいとされる声明を行ったとするましょう。このようにして私たちは平均が与えられ、その証人の過去または将来に行うすべての特定の声明が含まれる平均信頼性を得ます。

3. 以前の章（第9章 §14-32）で充分に示されたように、特定の平均値が指定されているという単なる事実は、統計や平均に基づく判断の最も自然かつ適切な根拠が見つかる場合でも、それを常に従う理由にはなりません。一般的な平均はしばしば個別の状況に正確に対応するためには修正されるべきです。例えば、死亡率に関する統計では、特定の年齢の人々によって提供される広範囲の表ではなく、職業、住居、または生活様式が異なる人々によって提供される狭い表を好むことがあります。

ここで読者に便利なことに、このような方法で真実性を確保できるとは限らないことを思い出させます。十分に広範かつ多様な経験を通じて一般的な平均値を一貫して従えば、長期的には真実へと導かれる傾向が生じるでしょう。より特殊な平均値が私たちに提供できる利点は、少ない且つ軽微な逸脱のない真実への同一の傾向を与えることです。

4. それでは証人の話を再考する。彼があらゆる種類の主題について非常に多くの声明を行っている場合、私たちは九つのうちの十個で彼が真実を語ることに自信を持つことができるだろうし、残りの一回だけは誤った情報を提供することになる。これは定義や一貫性の問題であることは明白だ。

しかし、このような一般的な平均値に依拠するよりも、特定の状況に特化することはできないだろうか？彼が述べる声明の特性的な状況を考慮することで、その生存率を彼の体質や生活様式によって特化させることのように、同様に彼の証言の可能性を特化させることはできないだろうか？

もちろん、これを行うことは可能だ。そして、このような考慮に基づいて生活上の実用的 contingency で行動する場合、私たちはそのような計画を採用しないと非常に愚かな行動を取ることになる。

この平均値を修正する二つの方法が提案される：一つは実際の機敏さが最も Likely に採用するであろう方法であり、もう一つは確率に関する著者がほとんど採用する方法である。前者は次の考慮に基づいて修正を行う試みを行なう：証人の一般的な平均値に依拠せず、その状況に対応する種類の推測的な修正を付与し、これは私たちの経験や観察に基づくものである。つまり、経験によって特定の種類の物語が一般の物語よりも真実である可能性が高いか低いことを証明しようとする。

後者は異なる方法でより体系的に行われる。ここでは、問題となる各ケースの誤りの性質とその数を分析することで、特定の種類の物語が正確に報告される確率が一定の数値比で高くなるか低くなるかを示そうとする。

5. これらの方法についての議論に進む前に、著者たちがしばしば十分に注意していない区別を指摘しておく必要があります。あるいは、読者の注意を十分に引いていない点についてです[2]。証言が与えられる方法は大まかには二つあります。

まず第一に、それは互いに排他的な質問に対する回答の形で与えられることがあります。つまり、その質問は「はい」または「いいえ」と答えるために作られています。「Aが起こったか、はいまたはいいえ？」というような形です。もちろん、このような可能な答えは相互に矛盾しており、一つが正しくない場合、もう一つは正しいはずです。例えば、黒球と白球の袋から球を引く例を考えましょう。証人が二つの色しかないことを知っているか、あるいは少なくともその中の一つか他のものを名乗るだけに制限されている場合です。黒球が引かれても彼が誤った場合は、「白」と答えることになります。

このような種類の証言を数値的に例示する一般的な方法は、賞品と穴のないlottery（抽選）や二色のみのボールの袋を考えることです。証人がただ「有罪」または「無罪」と言うだけであるという陪審員の役割に対する例が重要である理由はおそらく、陪審員が他の仕事を持っていないからでしょう。

一方、証言はより原始的な声明や情報の形を取る可能性もあります。Aが起こったかどうかという質問ではなく、「何が起こったか？」と尋ねることもあります。この場合、証人が真実を述べたとするなら、以前と同じように彼には一つだけ正しい方法があると仮定しなければなりません；特定の事件が発生することを当然想定していましたからです。しかし、誤った場合、彼は多くの間違いの方法があり、可能性として無限に多いかもしれません。通常、そのような誤りの声明は数が限定されていると仮定されますが、実際的にはそれほどでもない結果が出ることもしばしばです。この状況を数値的に表現するには、袋の中のボールがただ二色だけではなく、すべて異なるものであることを想定します。例えば、各ボールに連続した番号が付けられているとしましょう。

もちろん、世界で述べられる多くの声明は、いずれかの選択肢に対する答えでもオープンな質問に対する答えでもない場合があるかもしれません。たとえば、ある人がただ「スパードの7を引いた」と主張するだけであり、彼が尋ねられたかどうか、「そのカードは引かれたのか？」や「どのカードが引かれたのか？」など具体的な質問がなかったかも分かりません。さらに、これは普通の歴史的な声明の場合でも同様に当てはまるかもしれません。

この反論は非常に適切で、追加の困難を構成するものとして認識しなければなりません。私たちができることは、状況から証人が持っていた選択肢の数をできるだけ正確に把握しようとすることです。彼が争点となっていることであり、詳細に入らない形で証言した場合、実質的には「はい」または「いいえ」と答える二つの選択肢しかなかったと公平に考えることができます。一方で、より原始的な物語を述べたか、あるいは（同様の意味を持つ）詳細を述べた場合、彼には広範な比較的多くの選択肢が存在したと見做さなければなりません。

これらの二つの例は、黒球と白球の例で、その中では一つだけの形態の誤りしか起こらないものと、番号付けされた球の例で、その中では多くの形態の誤りが起こる可能性があるものであり、これらのみを考慮する必要がある。実際においては、質問の方法によってこれらの二つが徐々に混同される可能性があり、私たちがどのように質問を構成するかによります。見ましたか、A が B を打ったかどうか？と、何を見たか？という基本的な質問だけでなく、中間的な導入質問として、A は何をしたか？彼は B に対して何をしたか？など、任意の数だけの追加質問を導入することができます。このようにして、誤りの可能性のある開口部を徐々に狭め、直接的かつ互換的な質問に近づくことができます。しかし明らかに、これらすべての状況は、区別される球の数が仮定された減少によって数値的に表現できます。

6. §4で述べた二つの計画のうち、後者のほうから始めましょう。それは唯一の方法的な科学的手段であり、提案されたものの中で最も優れたものです。袋の中に1000個のボールがあり、その中の一つだけが白いボールで、残りはすべて黒いボールです。ランダムにボールを引きますが、証人（その信憑性は9/10）が白いボールが引かれたと報告します。この特定の件について彼の声明を多数取ってみましょう、例えば1万回。つまり、この袋からまたは全く同じ種類の袋から1万個のボールを順次引いて、それぞれの場合に彼の報告を行います。彼の1万回の声明が彼の一般的な平均値の公平なサンプルとされる場合、仮定によれば、そのうち9回中8回は真実であり残りの1回だけが偽であることがわかります。

これらの偽の声明の性質は何でしょうか？状況を考えると、彼に間違える方法は一つしかありませんので、答えは簡単です。1万回の引く中で白いボールは10回出てきますから、そのうち9回は正しく報告され、残りの一回だけが誤った報告となります。黒いボールが引かれた9990回の場合、彼は8991回は正しいと報告し、999回は間違えて白いボールと報告します。したがって、全体として彼が「白いボールが引かれた」という声明をした1008回のうち、999回は誤りで9回だけが正しかったことになります。

つまり、このようなケースにおける彼の特別な信憑性は、9/10から9/1008に低下しました。通常表現すると、後者の分数は彼のこの特定の声明が真実である確率を表します。

7. 現在、証人が単に一つの誤り方法を持っているのではなく、多くの誤り方法を持つ場合について考えます。仮定として、ボールは1から1,000までの番号が付けられ、証人はこの事実を知っています。ボールが引かれたとします。そして彼は私にそれを25番だと教えてくれたとします。彼が正しかった確率はどのくらいでしょうか？以前と同じように計算すると、1万回の抽選でこのボールは10回取得され、そのうち9回は正確な番号として名付けられます。しかし、このボールが引かれていない残りの9,990回では、証人は多くの誤りの機会に直面します。

ただし、一般的には、彼の誤りがすべて間違った数字を報告する形であると仮定されることが一般的です。そして、何故特定の数字を選んだのか明らかな理由がないため、彼はおそらく全ての間違った数字を等しく報告すると考えられます。したがって、彼の9,990の誤りは、単一の正しい場所に戻るのではなく、多くの異なる誤り方法に均等に分散されます。これらの状況のうちの一つだけでは25番と名付けられることはありません。

したがって、彼が25と名付けた10回中9回は正しかったことになります。つまり、この場合でも彼の平均的なまたは一般的な信頼性は特定の場合にも同様に適用されます。

8. この結論の真偽については、証人の誤りや虚偽が限定された範囲内で行われること、そしてその範囲内で彼の答えが完全に公正に行われているという仮定が認められるなら、論理は完璧に堅固であることは当然である。しかし、これらの仮定は非常に恣意的ではないか？つまり、当選券やボールの袋のような当り前の抽選法を、普通の生活における条件が曖昧で不確実な状況で完全に正確にするというような方法は、実際的に利用可能になるでしょうか？

例えば、私は長い間知っている人物から情報を得たとします。この人は私の信頼性や判断力についてある程度の知識を得ていると仮定しましょう。そして、その人が特定の困難を冒したり費用をかけたりすることで私に有利な何かがあると告げます。

証人の一般的な信頼性については問題ありません；これは我々が事前に決定します。しかし、彼の話自体について、難しさや曖昧さはどこでも現れます。この袋の中には何個のボールがあるのでしょうか？実際的に、どのような種類と内容の袋から抽選が行われたと考えているのでしょうか？

ここで合理的な判断の材料が存在していないように思われます。しかし、上記の例で得られるような修正された信頼性を得るためには、これらの質問は一定の正確さで答えなければなりません。なぜなら、この方法の主なポイントは、事件が実際に起こらない回数をどのように決定するか、そしてその結果として証人が誤ってそれが起こったと主張する回数を推測することにあるからです。

当然、このような考慮は私たちの判断に多少の影響を与えることは認めますが、その影響が通常の状況下で数値的な決定可能性を持つものであるとは限らないということだけです。たとえば、我々が何か種類の運が向いたように告げられることがあります。しかし、実際にはそのような幸運は訪れていない場合でも、これは一度ではなく繰り返し起こります。しかし、誰がそのような状況で、空白が実際に抽出されたにもかかわらず賞金が誤って公表される回数や性質を最もわずかな程度でも示すことができるでしょうか？

したがって、私にはこの手続き方法から何らかの実用的な価値を持つ数値結果を探すことはまれであり、おそらく決してないと思われます。

9. 当然、当選 lottery の場合、または黒球と白球の袋の場合についての結論は、証人の証言がそのような状況下で非常に大きな信頼性を失うという点から異議が出されることがあります。しかし、経験上、そのような状況はほとんど存在しないと考えるべきです。しかし、これらの状況が仮定されている場合でも、与えられた結論は十分に正しいと思われます。人間が単純に「はい」または「いいえ」と答えることを要求する場合、また、「いいえ」を正しく言うべき回数が「はい」よりも遥かに多くなるような一連のケースを仮定します。そして、証人がほとんどどの程度でも信頼性があると仮定すると、否定すべき正当な機会がより多くなるため、彼の肯定的な回答の大部分が実際には誤った「いいえ」よりも正しい「はい」になる可能性があります。これは当然、彼の「はい」の平均的な価値を低下させ、それらを比較的信用できないものにします。

以下の例を考えてみましょう。私は、日常的な事実について信頼している庭師がいます。ある朝、彼が私の犬が逃げ出したと伝えてきたら、私は十分に信じることでしょう。しかし、彼が犬が狂っていると言った場合、私はその声明を受け入れる際に大きな疑念を抱くはずです。これは、犬が本当に狂っているときよりも彼が間違える確率が高いわけではありません；しかし、経験から、他の症状（例えば発作）の方が狂気よりも遥かに一般的であり、多くの狂気の主張はこれらの症状についての誤った主張であることが分かっています。これは、白いボールが引かれたという主張の多くが実際には黒いボールが引かれたことを指摘する誤った主張である場合と似ています。実際には、ほとんどの場合において彼を他の点では完全に信頼しているにもかかわらず、その特定の証人の主張を過度に軽視することについての難題を感じる人はいないでしょう。

10. この問題に対するこの解釈が困難とされた特定の一例がありますが、私にとってはそれが決定的な確認となるべきものであると考えます。ただし、これは我々が予想される限界内でのことです。その例は、証言の正確性が半分である人物についてです；つまり、単なる「はい」または「いいえ」が要求された場合、彼は正解と誤りが等しくある人物を指します。他のどの程度の正確性も与えられた場合でも、実践的な判断と経験から近づいた確認を得ることは非常に困難です。証人の評価方法について我々は慣れておらず、その数値的な真実度をどのように理解すべきかほとんど把握していません。しかし、半分の正確性を持つ人物については（Mr C. J. Monroが巧みに指摘したように）、彼のような証人について我々は十分な経験がありますが、異なる名前で知られています；これは、彼が全く知識がない問題に対する質問に対して自信を持って答える人物と実際には同じことです。そのため彼はただの推測を提供するだけです。

lotteryに一つの賞しかない場合、証人の証言が半分であるという者が私たちは賞を獲得したと報告したとき、計算するとその証言は全く価値がないことがわかります。私たちは賞を獲得した確率は、彼が口を開かなかったとしても同じで、1/1000です。明らかにこれが正しい結果であり、関与のない証人が何も変えることなく状況を保持しているだけであるからです。厳密には、彼は実質的にほとんど証人とは言えません；証人の自然な役割は事案を調査し、判断と正直性に基づいて多少なりとも確認を加えることです。しかし、至らぬ者よりは必ず改善があるはずです。しかし、彼の単なる推測に耳を傾けるなら、私たちは自分自身で推測することと同じであり、その場合、私たちが正しい確率は事件の頻度によって測定されます。

同じ規則を別のケース、つまり番号付きの玉の場合にも同様に適用することは困難です。そこで、間違える方法が多くても半分の声明では正しかったという者が実質的に優れた証人であるかもしれません。彼が多くの間違いの方法を持っていても、半数の声明で正しいなら、明らかに彼は何か程度の注意を払っており、単なる推測ではないことがわかります。はいまたはいいえの場合、誰でも半分の回で正しくなることができるかもしれませんが、真実が一つであり誤りが複数ある場合とは異なります。1000個の玉があり、特定の玉が引かれた状況で単純に価値のない証人を表すために、その証言の確率を1/1000と仮定する必要があります。これが行われれば当然ながら同じ結果を得ることになります。

従って、証言の真実性の図形，すなわち一般の証人の証言の真実度を表す分数は、一種の相対的なものであり、絶対的なものではないことに注意すべきである。これは、彼が仮定される答弁の一般的な性格によります。等価の内在的真実性と価値を持つ二つの証人がおり、一つは単に「はい」「いいえ」と答える一方で、もう一つはより原始的な主張を敢えてする場合、それぞれ異なる分数で表される。前者は非常に容易な任務を自ら課したのに対し、後者はその任務が遥かに困難である。

したがって、証人の実際の真実性と慎重さは、彼の答弁の性格の一要素に過ぎない。その他の要素は、先ほど述べた通り、彼の主張の本質である。従って、このような問題において、平均的な真実性を証人に割り当て、それぞれの種類の答弁の場合でも同じようにこれを受け入れる一般的な手法は便利ではあるが、ほとんど適切ではないと考えられる。しかし、この考慮事項は、議論の大半が「はい」または「いいえ」という種類の答弁のみに関心を向けているため、その重要性はそれほど大きくない。

したがって、証言の真実性の図形は、単に証人の真実性だけでなく、彼の主張の性格にも依存することを理解することが重要である。

12. 以上の方法でこのような問題を扱うべきであるという説明はこれで終わりとする。確率論について学んだことのない人々がこの問題に取り組む方法は非常に異なるものであると考えられる。彼らはおそらく次のように考えることだろう。彼らは自分たちの中で言葉にしてこう考えるかもしれない。「ここには証人が語った物語があり、その証人は十回中九回は真実を話すとしよう。しかし、この種の物語は経験から見れば非常に一般的に誤って語られていることがあり、十回中九十九回はそうであるとしよう。」そしてこれらの相反する信憑性の要素を持ち合わせた彼らは、何か方法でこれらを均衡させる試みをするだろう。厳格な規則に基づく決定法を提示することは不可能であり、その困難から脱出するためにはどうすべきかという具体的なルールは与えられないであろう。おそらく、恣意的な判断を避けるために、彼らはさらに経験に依拠したり、以前の教訓による無意識的な記憶に基づいて、どちらの相反する推論が彼らの目の前にある特定のケースにおいて優勢を占めるべきかを決定するために導き出される新たな推論が必要であると指導されることだろう。読者は容易に理解できるであろうが、このような形で問題に対する一般的な解決策は不可能である。これは単に第九章（§14-32）で既に述べた、二つの異なるクラスや、甚至は特性において対立するクラスのいずれにも等しく属する個体の現在の熟知されたケースである。その個体の特性がどちらのクラスに最も自然かつ正しく関連付けられるべきかを一発で決定することはできない。新たな推論が必要である。

13. さまざまな著者は、困難から脱却するためのルールを提案してきました。奇跡に関する論争はおそらく、この種の提案が一方か他方で最も豊富に生み出された機会でした。おそらくこの論争に関連しているフレーズは、「対立する不確実性の競争」という表現が頻繁に類似した議題での議論において使用されるようになったものです。このような表現の意味は明らかに、ある主張の真偽について判断を形成する際に、その主張が二つの非常に異なるクラスに含まれている場合、それぞれのクラスに属するか否かによって我々の意見が大きく異なるということです。

つまり、ある主張は当然ながら特定の証人の声明や種類の一つのクラスに属し、また特定の不確かな物語の一つであるという理由で別のクラスに属します。前者のクラスにはその真実性に対する有利な要素があり、後者のクラスには不利な要素があります。

「対立する不確実性の競争」という表現について話すとき、これらの異なる証拠源が共存している場合、それぞれは何かしらの形で証明力を持ち続け、競争を引き起こし、最終的には一方または他方が勝利することになると考えられています。例えば、ヒュームは「一つの確率から別の確率を引く」と述べ、信頼性を残りの部分に割り当てるべきだと主張しています。トムソンは『思考法則』の中で、一つの確率が完全に他の確率を凌駕すると述べています。

14. このようなルールには、いかなる哲学的な価値も見受けられないと私は思われます。
これらのルールは個々のケースでは確かに有効かもしれませんが、それらは普遍性を主張する資格はありません。甚至に競争という概念自体が、この問題において必要不可欠な要素であるという考え方も捨て去るべきです。

実際、困惑が生じる方法をもう一度考えてみましょう。その結果として容易に理解できるように、既に解決不可能とされた一般の先験的な解法によって表現される困難の一形態に過ぎないことが明らかになります。私たちが持っているのは声明だけです。この場合、証人が長期的には十分な数回の証言の中で10回中1回は嘘をついていると言いましょう。つまり、彼は主に真実を語ります。しかし一方で、経験から得られる様々な証人による様々な機会からの証言によれば、この声明は大抵偽であることが確認されています。数値的に言えば、私たちはその声明が100回のうち99回は偽であると仮定します。

このような状況では、前篇で説明したように完全な行き詰まりに直面します。私たちの科学には、意見を形成するか、または一方倒しの決定を行うための原理が提供されていません。なぜなら、先に述べた通り、結論は無数にあるすべてが等しく可能であるからです。例えば、証人の全ての異常な主張が真実である可能性もあれば、全てが偽である可能性もあります。あるいは、真と偽で何らかの割合で分かれている可能性もあります。

統計的な情報にまで到達し、証人が一般的には正しいと言葉を発する一方で、その言葉は大抵偽であるという事実を受け入れたとしても、そこで止まるべきではありません。経験へのさらなる訴えが必要であり、彼のような特徴を持つ証人の特定の内容の主張について、その状況における経験がどのように立つかを確認する必要があります。あるいは、より広範な視点から、彼のような一般的な性格、地位、共感を持つ証人が提出したような類似の主張について、その状況における経験がどのように立つかを確認すべきです。

15. とても不思議な物語が、多くの場合、おそらく大半の場合は他の物語よりも信頼性が低いという点は、完全に認めなければなりません。つまり、そのような物語を二つの異なるクラスに分けると、証人が同じであったり、類似した証人が、前者を主張するときよりも後者を主張するときに比例してより多く間違っていることが見られるでしょう。しかし、これによって必ずそうなるとは思われません。例えば、ある人々にとっては物語が非常に不通常であるという単なる事実が、彼らの主張に慎重になることを引き起こし、その結果、その正直性を増す可能性があると考えることができます。これが真なら、我々はその不思議な物語を受け入れるよりも普通の物語を受け入れるのにさらに容易さを感じるかもしれません。

このような仮定は私には全く強引なものとは思えません。例えば、次のような状況を考えましょう：二人の人が、一人はただの普通の正直性と知性を持つ人で、もう一人は科学的な自然学者です。彼らがいくつかの一般的な出来事について述べる声明をします。我々は両方を信じます。それからそれぞれが、自分たちが見たと主張する不思議な自然現象や怪物について報告します。おそらく多くの人々は、自然学者の報告が後者の場合でも前と同じくらい素早く受け入れられるでしょう：しかし、同じ物語が無学な観察者から来ると、それはかなり慎重に受け取られます。その違いは何から生じるのでしょうか？それは自然学者がこのような事柄についても最も普通のものと同様に十分注意を払い、正確であるという信念からです。一方の人物に対しては、全く同じ信頼性を感じません。誰かがこの程度まで行かないとしても、おそらく両方の人物からそれぞれの種類の物語を受け入れる信用度の違いについて、自然学者から来た場合よりも他の人から来た場合に比べて、その差はかなり小さいだろうと認めることでしょう。

16. この部分について論じている最中に、用語「不思議な物語」の使用に関しては、相当の曖昧さとそれに伴う混乱があることを指摘しなければなりません。純粋確率の領域においては、「不思議な物語」という言葉は単に不通常の出来事である主張を含む物語を意味すべきです。少なくとも、この作品全体を通じて一貫して採用され維持されている見解は、そのように定義されるべきであるとされています。

我々がこの意味で使用する限り、この用語の意味は明確になります。これは純粋に客観的な参照であり、実際または将来予想される相対的な統計的頻度が非常に低いことを示す単なる言外の意味を含みます。多くの出来事の中から特定の種類を選択し、それがあまりにも稀に起こることを仮定します。このことは不通常または不思議な出来事と呼ばれます。

これは前章で豊富に示されたように、その希少性により、我々はその発生を個々の場合において非常に期待しない傾向があります。このような予想下での推測も稀に正当化されるものではなく、したがってそれが実際に起こったときにはしばしば驚くことになります。このように理解する「不思議な」という用語は確率の観点からは唯一の正当な意味であると私は考えます。

しかし、この単語には別の、そして非常に異なる用法があります。これは帰納法に関連し、あるいは証拠の科学全般に関連するもので、その限られた部分である確率論よりも多く関連しています。この意味での「不思議な」、さらに言えば「 improbable（不可能な）」出来事は、単に極めて統計的な稀少性を持つものではなく、正確に予測することは期待できなくても、中程度の証拠でも比較的容易に受け入れられるものではありません。むしろ、その出来事が実際には証拠に対する抵抗力を持っていると見做されるべきです。

これは、証人の信用性が根本から影響を受けるような何かに関連している場合に当てはまります。つまり、そのような主題についての彼の主張は他の主題についてのものよりも本質的に劣ると考えられるのです。例えば、彼の偏見や情熱、または迷信を刺激するようなことがあれば、このような状況では、証人の誤りが私たちを誤導することの頻度（第6節で述べたように）を計算して、その信用性を評価するのは不合理であると考えられます。代わりに、私たちは証人の平均的な正直さを一括りにして修正する傾向があるでしょう。

17. 真の確率論において、如に先ほど述べた通り、各事象にはその発生頻度が明確かつ認識可能な程度がある。それは極めて稀少であっても、どの程度まで稀少であるかを我々が仮定するいかなる限界でも稀少であっても、それでも全ての者がその発生に依拠するデータを理解し、認めている限りは、その事象が何らかの経験範囲で現れることが期待できる範囲を把握することは可能である。個々のケースにおいて期待することなく、また短い期間内においても期待することなく、しかし十分に広範な経験範囲では確実に期待することはある。したがって、争点の一方で実際にどこにも何時にも発生しないことを認めない議論者に対して、類似の神異の物語をどのように扱うことができるだろうか？袋や球などの機械的または数値的な事象の配置は、神異の現れを公正に説明するための例として認められるべきであるか、あるいは多くの場合、そのようなものの中でも「非常に不思議な」または「極めて稀少な」という名目で分類される多くの事象に対して認められるべきであるか？

特定の神異が報告されたある叙述者によって起こったと主張されている場合に反論する人々は、その神異が確実にいつか発生することを否定しているわけではない。つまり、何時かには必ず起こることを示すものとして何かが起こったという証拠があるかどうかの問題ではなく、それゆえこのような例としては適切ではない。

極めて稀少な事象であっても、その希少性がいかに大きくても、神異の例としての不適切さは単純な考慮から明らかである。私は誰もが、非常に不思議な組み合わせの偶発的な出現を認めているなら、自分が実際にそれを見たと仮定した場合でも、その事実について大きな疑念を持つことはないだろうと推測する。一方で、真に科学的な傾向を持つほとんど誰もが、目の前で神異が起こったとしても、それを受け入れることは容易ではないだろう。その時驚くかもしれないが、その後すぐにそれを信用しないか、あるいは意味を空洞化して説明することで、その全てが神異の意味から排除されることになるだろう。

18. これらの証言に関する問題について、全体として見れば、一般的に意図される方法でそれらを扱うのにはほとんど意味がなく、特定の状況下での証言の評価のために具体的な規則を得ることにもほとんど役立たないと思われます。証人の信頼性が測定可能であると仮定しても、実際の困難は、その事件の失敗が考慮されるべき範囲を決定する不可能性と、証人が彼に向けられた質問に対する明確な回答の程度という特性にあります。これらの特性については、数値的な評価が必要であり、それを持たないと必要なデータを持っているとは言えません。

したがって、ほとんどの人々にとって実用的な資源は、その声明が作成された状況の性質を考慮して当然の推測に基づく直接的で即時的な修正を行うことですが、これは根本的に科学的ではなく規則化することは不可能です。私には、証人の一般的な信頼性に一律に依拠する単純な計画が何かあることを主張する余地があるように思われます。つまり、彼の話が普通であれ異常であれ、その出来事が時折起こるか起こらないかを十分認識した上で、信用性を同じレベルに置くことに決めるべきです。真実であるかどうかは別として、この方法では平均的に誤りを犯すことが恒常的であり、また大きな程度まで誤りを犯す可能性があります。もし彼の話の性質に基づいてその信頼性を専門化する合理的で正確な方法が存在すれば、より確固たる立場に立つことができることでしょう。しかし少なくとも、平均的に何を期待できるかは分かることです。十分な数と多様な声明から彼の証言を受け取り、常に同じ信頼性のレベルで扱うことで、長期的には大きな誤りを避けることができるようになるかもしれません。

証言の組合せについて、いくつかの補足的な言及を加えます。ここでは新しい原理は導入されませんが、それに伴う複雑さは自然に増大します。仮に2人の証人があり、それぞれの信憑性が9/10とします。この場合、100回の証言について、前回採用した手法により、両者が共に正しい答えを出したのは81回で、共に誤った答えを出したのは1回、残りの18回では一方だけが正しく他方は誤っています。

しかし、彼らは同じ内容を報告するとしていますので、比較すべきは彼らが共に正しい答えを出した回数と、共に答えた回数全体（正誤問わず）です。前者に対する後者の比が、その場合の証言組合せの信憑性を表す分数となります。

この点を決定する際に唯一の難しさは、彼らが共に誤ったときにどれほど頻繁に一致するかをどのように定義するかにあるため、明らかに両者が正しいときには必ず一致します。この調査は当然ながら、彼らが誤る方法の数に依存します。

まず、単純な「はい」または「いいえ」（§6参照）の場合を考えます。同じ例として、1000個の球の中から1つだけが白い袋を想定します。前回と同じように計算すると、10万回の抽出（すべての可能性とそれらについての報告全体の完全なサイクルを得るためには必要な回数）において、両者が正しい白さの報告に一致したのは81回で、誤った報告に一致したのは999回です。

したがって、その証言が確認されたときの物語の確率は81/1080となります。つまり、等しく信憑性のある2人の証人が一致したという事実は、単独の証人によるものよりも約9倍に報告の可能性を高めます。

20. しかし証人たちは多くの方法で誤りを犯す可能性がある場合、彼らの一致した報告はその報告が真実である可能性が高いことを示します。例えば、1000番号付けされた球体の例では、両者が数字を間違えるとしても（密約なしに）同じ誤りを犯すことの確率は非常に低いことがわかります。

一方で、最後の場合では、すべての組み合わせによる誤りが彼らを必ずしも事件が起こったという主張へと導くことになります。したがって、特定の数字（たとえば25）が選ばれたと主張するためには、999 × 999回中1回だけであることがわかります。実際の出来事に対する確率は80919/80920となり、単一の証人がいる場合よりも著しく高いことが分かります。

したがって、二つ以上の証人（もちろんそれ以上も含む）が、彼らが多くの異なる誤りを犯す可能性のある事柄について一致して主張する場合、その有利な証言の組み合わせは出来事が起こった確率に著しく寄与します。ただし、密約の可能性がない限りです。そして、誤りの機会が実際には無限に多い極端なケースでは、その組み合わせはほぼ完全な確信を生むでしょう。

しかし、この条件すなわち密約のない状態は、ほとんど達成できません。実質的には、私たちは誤りや疑念の主な源が、ある種の密約の可能性にあることを認識しています。この危険を完全に排除することは困難であり、またその存在する場合でも数値計算には提出されません。したがって、私は詳細な報告に関して証言の組み合わせは単一の証言よりもさらに確率論において考慮すべきではないと考えます。

21. 事象の信憑性を確率論の法則で十分または適切に考慮することは不可能であるという点は、既に§17で指摘されています。しかし、この不可能性の根拠がしばしば不充分に理解されているため、数ページ追加することでこの点を強調する機会を得ることができます。もしこれは偏りと見なされるのであれば、主題の重要性と、当該科学の法則に基づいて一度たりとも試みられてきた Various Writers の一貫した努力が、その正当化となるでしょう。

必要不可欠な前段階として、神異（奇跡）を定義する必要があります。おそらく多くの人々は、「因果律の中断」という言葉で神異を呼ぶことを認めるでしょう。これは完全な定義とはならないかもしれませんが、少なくともその説明としては正確であると受け入れられるでしょう。それは真実であり、しかし全貌ではないかもしれません。神異が何であれ、これがその物理的な側面です：これが科学の対象となる点です。因果律の中断が考慮されなければ、その事象は通常のものとして扱われ、特別な意味が付け加えられたと見做されるでしょう。つまり、型やシンボルではなく、神異とは異なる存在となります。この点こそが現在主に攻撃を受ける側面であり、そのための防御も一般にその支持のために展開されています。

現在、この定義や説明は、事実に関するいくつかの仮定を含んでおり、理論的な要素も含まれていることが明白です。その仮定は、因果法則が自然全体で普遍的またはほぼ普遍的に支配しているという点にありますから、それらの法則違反は特異かつ目立つものであるということです。この仮定はなされており、しかし、これ以上の何かが必要であるとは必ずしも言えません；つまり、因果法則の性質について主観と客観の二つの主要学派を支持する必要があるわけではありません。この定義は、因果法則を単なる前因後果の一貫性として理解するだけであっても、あるいは事件間には単なる順序を超えた深いそしてより神秘的なつながりがあると主張するだけでも、同様に適切です。

「因果」の用語をこの最小限の意味で使用することは、両学派共通のものであり、一方はその不十分さを考慮に入れることもできます。したがって、「因果法則の中断」という言葉を使用しても、どちらの立場にも拘束されることはありません。

22. この定義によって示唆される問題の側面は、私たちが避けることができないものであることに注意すべきです。確かに、自然界における普遍的な法則と秩序の存在についての仮定を避けるためには、 Miracle（奇跡）の定義から異なる視点で定義しようとすることが時々試みられてきました。例えば、Miracleは「創造力の直接的な働きかけ」や「啓示の兆候」と呼ばれることがあります。さらに曖昧にすればするほど、「異常な出来事」と呼ばれます。しかし、これらの定義を採用することによって何ら得られるものはありません。これらの定義が神学者にとって満足させることは確かにできますが、物理学の研究者にとっては一瞬もそのような定義には満足しません。彼はすぐに自身と他の科学者の信念である普遍的な法則の存在について主張し、その定義とこの教説との関連性を問い始めます。

「あなたが述べた奇跡は、これらの法則のいずれかに違反するものでしょうか？それともそうではありませんか？」という質問に対する答えが必要となります。そして、その答えを提供しなければなりません。もしそのような回答を拒否し、普遍的な法則の存在についての信念を否定すれば、当然ながら彼との議論から外れることになります。

この事実を認識する必要性は、物理科学の研究が進むにつれて日に日に強まっています。そして、この問題の側面はいつか必ず対処しなければならないので、それを最初に扱うことが適切です。その科学的な形での困難は当然ながら現代のものであり、その起源となる教義も現代のものです。しかし、これは哲学的概念の成長の一例に過ぎません。多くの他の例が挙げられるように、ある哲学的概念の成長は徐々に議論の性質を変えることになり、最終的には議論の場所を移すことがあります。最初には全く関連がないと思われたような議論の中でさえもです。

23. これまでの道は明るい。異なる学派の弟子たちはこの段階までで共に進むことができる。なぜなら、上記の教義を定める際には、それが完全な真理を含んでいることを暗示したり認めたりすることを避けてきたからである。しかし、ここから先には二つの道が分岐する。その性質、起源、方向は想像できるほど異なる。この調査は単なる逸話に過ぎないため、それぞれの特徴について簡潔に述べることにとどめ、支持される主張を示す試みは控えることにしよう。

(I.) 一方では、因果律の原理が最終的な原理であるとすることができる。そのような用語を使うことで我々は、幾何学の公理から意識的に調査を開始することとは異なる意味で、自然の研究によって私たちが引き寄せられる最終的な結果であると理解していることを示すものではないことを注意しなければならない。つまり、我々が調査を通じて自然に導かれる結論であるという意味である。

調査の範囲において、事件は常に一定の規則性で続くことを見出し、さらに（一部の人々は加えるかもしれないが）この経験が私たちの性質に支持されたり甚至は要求される傾向や法則があると我々は考えることによって、すべての調査が基盤となる一つの基本原理としてこれを認識することができる。

(II.) 他方では、異なる種類の原理を認めることができる。経験の範囲においてこの一定性があることを認めつつも、超自然的な支配力や計画、秩序と呼ぶべきものがあり、その観点から設計が述べられることが可能であると認めることである。適切な区別を用いて言えば、これは事件を覆い隠すものではなく、むしろそれらの下に存在するものと理解されるべきである。

24. 明らかに、特定の奇跡や異常な物語について、一方または他方の前偏見のもとで議論に入ると、その信憑性に関する問題は非常に異なる観点を呈する。しばしば忘れがちだが、事実に対する相違は真摯な論争の条件の一つであり、最終的な原理に至る相違はすべての論争を破壊するものであることを忘れてはならない。現在の対立の可能性は、将来の調和の可能性と同様に、このような状況において完全に排除される。奇跡に関する多くの一般向けの文献では、この欠点が存在しているように見える。信憑性についての論争における主張が述べられ、検討されても、反対派と支持派はそれぞれの立場がいかに根本的に異なるかを十分理解していないことがしばしば見受けられる。

25. 以下は、我々が述べている不整合性の一例を示すためにある程度の程度で用いることができる説明である。ある海軍士官は、太平洋の遠くにある珊瑚島に一人で上陸し、そこには海岸に並べられた石が「十」の形で配置されているのを見つけたと報告している。もし我々がその報告の真実性について議論を起こす場合を考え、証言の有効性に関する考慮事項のみによって事実関係を決定しようと試みるなら、無住民であるか否か、そして彼らの習慣や慣習の性質という問題を導入しない限り、我々は多くの現在の奇跡に関する議論が不十分な性質を持っていることを理解することができるだろう。この主題に対するすべての例は不完全であるが、人間の行為の痕跡として推測されるものが他の非知能的な自然原因の秩序に干渉しているという点で特定の行動を検出するような事案は、どの例よりも適切な例と言えるだろう。ここで議論から省略されているものは明らかに重要な要素である。住民が存在しないと仮定すれば、我々はおそらくその報告を信じないか、あるいは著しく過大評価だと考えるであろう。住民が存在すると仮定すれば、問題は証人の信頼性だけではなく、彼らの性格やそのような行動の目的という要素も考慮に入れる必要があるだろう。

26. このような考慮は、しばしば議論に取り入れられるが、私にはその程度が不十分であるように思われる。パライーの言葉を引用してしばしば主張されるのは、「一度神を信じれば、奇跡も信じられなくなる」というものだ。しかし、この主張は確かに修正と拡大が必要である。むしろこのような声明の方が適切であろう。「神の働きが、道徳的世界と物理的世界全体にわたって追跡できるなら、奇跡も信じられるだろう」。

実際的には、これは次のような意味を持つことになる。「何か目的があると考えるのを認めよ。それがどのように進行するかは不明瞭でも、我々は事件間の関連性だけでなく、それらの結果を追跡するだけでなく、何か計画的存在の概念を形成できる」ということだ。パライーの助言は、「精霊がいることを認めるなら、なぜカップが割れるのか説明できる」と聞こえるように聞こえてしまう。しかし、そのような主張は適切にはなっていない。

我々が議論している困難を抱える人に対しては、神を常に自然と関連付けることなく信じるという態度は全く不可能である。神を信じる唯一の形は、無限で秩序あるシステムの制御霊としての存在であり、その信仰が彼の心に導入されるべきだ。

実際的には、私には奇妙なようだが、反説的とも言えるが、完全に帰納的な科学精神に浸透している無神論者が、広大な体系の一環として奇跡を信じるよりも、そのような無神論者が神信者となって孤立した奇跡を受け入れる方が容易である可能性がある。

27. ゆえに、ヒュームや彼の後継者たちは、クリスチャン啓示が単なる奇跡の一連であるとみなしてその信憑性について論じることで実際的に慎重を期した。これは、生体をその構成する各部位の集まりと同様に、啓示を奇跡の一連として適切に捉えることができるという見方と同じようなものだ。多くの一般的な議論において問題となるのは、攻撃者と防御者がしばしば本当に立っている異なる基盤について認識が欠けていることである。証拠や例証は無数に提示されるが、これらは一方の議論者の心の中で、その質問における存在しない因果律を確立しようとする目的とは逆に、その原理自体の一部として理解されているため、単なる要件を達成していない。物理的因果律の存在を特定の事例において疑わせる試みは、その議論者がそれらの例証が示す原理の追加の例であると感じるとき、河を堰くために雪を shovel 込むようなものに似ている。このような例証は論争の時期には豊富にあるが、実際には反対する疑心者の中での最初の接触でその形を変え、予想とは逆に洪水を増幅させるだけである。

1. 前章では、これらの例に対する確率の規則を厳密に適用することは適切ではない理由が述べられた。しかし、数値的な正確さへの近づきは達成不可能でありながらも、日常生活中では証人の信憑性についての区別を認識していることは確かである。この章の目的には十分な粗略で実用的な区別が十分である。便利のために、理論を説明するために数値的形式で例を述べるが、それによって実際的にそのような正確さが達成可能であると推測することは意図していない。

2. この著作の初版において、私はこの指摘に対して自ら罪を認めるべきである。その結果は、この部分についての扱いが曖昧で不完全なものとなり、一部では誤りも含まれていた。

3. これらの結果の一般化された代数的形式は次のようである。ある事象の事前確率を p とし、証人の信憑性を x とする。その場合、彼が事象が起こったと主張したとき、実際には起こった確率は px / (px + (1 - p)(1 - x)) であり、彼が事象が起こらなかったと主張したとき、実際には起こった確率は p(1 - x) / (p(1 - x) + (1 - p)x) である。

今後に述べるいくつかの注釈を例示するために、読者はこれらの式のいずれかが p に等しい場合、それぞれで x = 1/2 を得ることに注意する。すなわち、信憑性が 1/2 の証人は、この種の事象の事前確率に影響を与えない。

一方で、これらの式をそれぞれ x と 1 - x に等しくすると、それぞれで p = 1/2 を得る。すなわち、この種の事象が起こり易さと同様に起こり難いとき、通常の信憑性は影響を受けない。

4. Todhunter の歴史書、第 400 ページ。\_Philosophical Magazine\_, 1864 年 7 月号。

5. 「したがってこれらの二種類の経験が相反する場合、それらを引くだけである。一方または他方の意見に確信を持って立つべきだ。」（奇跡に関する論文）

6. このような考慮は、主題に対する数学的な扱いにおいて導入されています。第5節で最初に述べた最も単純なケースから始まる問題の一般的な代数的解法は、もちろん次のようになります。事件が起こる確率を \( p \) とし、証人の正直さを測定する指標を \( t \) とするとき、彼の声明が真実である確率は \( \frac{pt}{pt + (1 - p)(1 - t)} \) となります。これは事件が起こるか起こらないかに関わらず、同じ程度に嘘をついていることを仮定しています。しかし、本文で述べられている状況に対応するために、事件が起こらない場合の彼の正直さを測定する指標 \( t' \) を仮定し、その結果上の式は \( \frac{pt}{pt + (1 - p)(1 - t')} \) となります。ここで \( t' \) と \( t \) は通常の事件と不思議な事件に対する彼の信頼性をそれぞれ測定します。この形式的な解法は、第14節と第15節で述べられた主張に対しても確かに合致します。しかし、\( t' \) の決定には、私は以前指摘したように、継続的に経験への依頼が必要となります。いかなる計画も上記のものを利用することによって採用される実際的な方法は、数学的者の方法とは精神と計画において非常に異なると思われます。

7. ラプラス（『 Essays 』1825年版、第149頁）は、たとえば100個のサイコロ（当然公正なものであることが知られている）がすべて同じ面を出すという事象について述べています。そのような状況を見ると、当時は混乱し、他者の確認が必要となるでしょうが、十分な調査後には誰も幻覚を仮定する必要を感じないでしょう。しかし、この事象の確率は分子が1で分母に77桁の数が含まれる分数であり、想像力では全く評価できないほど微小です。ただし、このような例は多少仮説的であることを認めなければなりません。実際には、そのサイコロが公正であるという信頼度は、そのような驚異的な確率にさえ近づくものではないでしょう。つまり、ここでは第14章の第28節から第33節で議論された問題の異なる側面を区別することが困難です。

8. 第一版では、この点について「決定的に不十分な形で」述べられていたと今となっては思われる。しかし、（第7節で示されたように）この計画が実際には特定のケースにおいて採用できる最良の理論的な方法であることを忘れてはならない。

9. この原理に基づいて、405ページで述べられた驚くべき結論が成り立つ。ある確率 p である事象について考え、同一の証言の信頼性を持つ m 番目の証人がその事象が起こったと主張し、 n 番目の証人が否定する場合を考える。上記で与えられた算術的推論を一般化すると、事象が主張される確率は

$py^{m} (1 - y)^{n} + (1 - p) y^{n} (1 - y)^{m}$;

（つまり、事象が起こった場合と m 番目の証人が正しい n 番目の証人が間違っている確率の積；\_加えて\_、事象が起こらなかった場合と n 番目の証人が正しい m 番目の証人が間違っている確率の積）に比例することがわかる。したがって、事象が正しく主張される確率は $py^{m} (1 - y)^{n}$ である。従って、私前にある主張が真実である確率は

$\frac{py^{m} (1 - y)^{n}}{py^{m} (1 - y)^{n} + (1 - p) y^{n} (1 - y)^{m}}$

であり、これは

$\frac{py^{m-n}}{py^{m-n} + (1 - p)(1 - y)^{m-n}}$

と等しい。この最後の式は m - n 番目の証人が一致して主張する主張が真実である確率を表している。

10. バトラーがこの計画概念に強調した点は、私としては『類推』という著作における大きな長所の一つだと考えられる。

\newpage

\chapter{平均の性質と種類（Nature and Use of an Average, and Different Kinds of Average）}
\label{ch:xviii}

\begin{center}
\large\textit{算術平均・幾何平均・最頻値・中央値・平均二乗誤差とゴルトンの測定法}
\end{center}
\vspace{4mm}

第十八章.

平均数の性質と用途、および異なる種類の平均数について

1. 平均数の主要な種類の性質と特性について、普通のイギリス読者にとって理解しやすい説明が必要である。代数学の一般的教科書は、平均数には主に3つの種類があることを示唆している、すなわち算術平均、幾何平均、調和平均であると述べている：これにより、統計学者がほとんど関与しない2つを含みつつ、彼にとって重要な2つ以上の平均数を除外している。私が読者に提供できる最良の3つの参照文献は以下の通りである。（1）医科学辞典の「Moyenne」条目、ドクター・ベルティヨンによる。この文章は主にクエテレットの視点から書かれている。（2）1878年の数学物理学部会誌「Abhandlungen d. Math. phys. Classe d. Kön. Sächs. Gesellschaft d. Wiss.」に掲載されたフェッナーの論文。この論文中には、統計学者にとって特に興味深い多くの異なる種類の平均数についての議論が含まれており、彼の中央値に関する説明は驚くほど詳細で有用である。数学的な知識はほとんど必要ない。（3）1885年のケンブリッジ哲学会誌「Camb. Phil. Trans.」に掲載されたエドゲワース氏の論文「Observations and Statistics」。この論文では、一般的な調査が通常行うようにして誤差法則と平均数の一種だけを扱うことではなく、広範囲の研究領域をカバーしている。

1. 私たちはしばしば平均数や個々の対象や事件との比較で示される一貫性について言及する機会があり、それらがどのような異なる種類の利用可能な平均数があるか、そしてそれぞれが具体的にどのような機能を果たすかについてより詳細に議論することが便利である。

現在の意味での平均という漠然とした概念は、私にはただいくつかの対象が互いに類似している点において何かを「中間的」なものであるというだけのものだと考えています。当然、これらの対象は互いに類似している点がある場合のみ、それらを一緒に分類する考えが生まれますし、また異なる点がある場合のみ、それらを区別することができるのです。

この原始的な形態において平均が私たちに対して果たす役割は、グループ全体を便利に保持することです。つまり、そのグループの事物の量的側面の代表的な値を提供し、特定の目的のためにグループ内の任意の単一のメンバーを置き換えるために使用されるのです。

科学は、平均や中央値という名前で正確化する最初の概念の熹微が現れると、一般的な名称の一部の機能を果たしているように思われます。一般的な名称の主な用途は何でしょうか？それは多数の対象を一つにまとめること、共通の性質を持つ多くのものをグループ化することです。

通常の一般的な名称は、多くの属性に基づいており、大半が質的なものですが、平均はそのような目的を果たすために、単一の量的属性に基づいています。これは特定の大きさや程度に注意を向けます。羊飼いが自分の羊について「全体して50シリングで売れるだろう」と言ったり、農民が「10フィート近い長さの棒がたくさんある」と言ったりするとき、厳密には一般的な名詞や集合的な名称の等価物として使用しているわけではありません。しかし、特定の大きさのタイプや標本を選出し、全体をその類似性によってグループ化することで、非常に近いような使用方法に近づいています。

羊飼いは自分の羊について考えていますが、単なる一般的な意味ではなく、特定の金銭価値を持つ羊として考えています。一部は多め、一部は少なめですが、指定された価値に十分近いため、名詞のように一緒に分類することができます。多くの粗略な量的記述も同様であり、「8日間の時計」や「12石の人」などと話すときもそうです（ただし、これらの例では最大または最小の値を指定することもある場合があります）。実際には、このような量的な側面についての一般的な概念を素早く伝える方法は、上記のようにタイプを選択するか、物事が想定される範囲内にある特定の限界を設定することでしかないように思われます。

2. これまで、比較という概念、つまり一つの集合を別の集合と比較するという概念は、そのような平均によって必ずしも導入されていない。比較を考え出すと、平均について何を意味しているのかをより正確に説明しなければならない。可能な種類の平均が非常に多く、実際には無限にあることが容易に理解できる。原始的な量をどのように処理するかにより得られる中間値という一般概念から、我々は任意の数の分類を導き出すことができる。しかし、私たちの目的にとって考慮すべき平均は、たった三つや四つである。

（1）まず最初に、算術的平均または平均がある。この平均を得るためのルールは非常に単純で：すべての量を足し合わせ、その合計を個数で割る。これは科学的な思考とは無縁の人々にとって唯一馴染みのある種類の平均であるが、これの単純さと馴染み深さに惑わされず、この平均が特定の状況において適切な理由があることを忘れてはならない。なぜなら、実際的な異なる値の中間値として安全かつ正確なものであるのは、日常生活における多くの一般的な目的、例えば取引や販売などにおいて、個々の違いを考慮するか、すべての対象を平均に等しいと仮定するかで同じ結果を得られるからだ。牧畜者がその羊の平均価格が50シリングであると言った場合、正確な意味は、その群れについて（そして彼自身について）に関わらず、個々が異なる価格で売却されるか、すべてが「平均」価格で売却されるかでは同じ結果になるということだ。したがって、彼が一年の販売と別の年との比較を行うとき、去年の羊は48シリングで今年の50シリングに対して平均を用いることで、大きな単純化が行われ、目的に完全に正確であると言える。

3. (2) 次のケースを考えてみよう。ある特定の人口が100年で倍増した場合、私たちは平均的に年率1%の増加があったと言えるだろうか？前回の場合とは状況は少し異なるが、その類似性は私たちの目的には十分である。答えは明確にノーだ。

例えば、100個の何かが$p$100で売られた場合、私たちは平均的な価格を$p$1と定義する。これは全体が$p$100で売られるか、あるいは個々のものが$p$1で売られるかに関わらず、合計金額は同じであることを意味する。ここでの平均価格は便利な仮想的な代替物であり、個々のものに対して適用しても合計総額に影響を与えない。したがって、100年間にそれぞれ年率1%の増加があったと仮定して、それが元の数値を2倍にするか否かという問いは、類似性のある問いである。そして、先ほど述べた通り、答えは否定的でなければならない。年率1%の増加が100年間続けば、合計はそれ以上に倍増する；実際には約2.7倍となる。真の年間増分は$\sqrt[100]{2}$で測定される；つまり、人口は平均的に年率0.7%で増加したと言える。

これにより、代数の教科書で通常議論される二種類目の平均値の一種である幾何学的な平均が導かれる。僅かな二つの量に関連し、その間に一つの中間値がある場合、この最後の幾何学的平均は最初の二つを比例する平均として最適に説明できる。すなわち、3 : $\sqrt{15}$ :: $\sqrt{15}$ : 5から$\sqrt{15}$は3と5の間の幾何学的な平均である。二つの量間に複数の幾何学的平均を挿入する必要がある場合、各項が全体の連続性において前項に対する一定の比を保つように選ばれるべきである。例えば、最後の段落の例では、1と2の間で99個の中間ステップを挿入し、その結果生成される100つの比率が全て等しいという条件があった。

従って、正確な定量的な結果に関連する場合においては、適切な平均の選択は、その目的に応じた特定の中間値が安全に使用できるかという問いに対する答えによります。つまり、「具体的な目的における正確な目標を達成するために、実際の様々な値をどの程度の中間値で置き換えることができるか」という点は、次の章でより詳しく検討する必要があります。しかし、一般的な比較のために、正確な数値関係が不要な場合においては、ほぼどのような種類の中間値でも私たちの目的を達成できることが安全に言えます。ただし、その種類を調査全体を通じて一貫して使用することです。

例えば、イギリスの異なる郡や地区の住民の身長を比較するか、イギリス人全般とフランス人の身長を比較するか、あるいは特定のクラスまたは地区の身長が増加しているか減少しているかを確認する場合でも、当然、その種類を一貫して使用することにより、どの平均値を選択しても問題ないことが明らかです。平均による作業の大半は、単なる比較的非定量的な説明に過ぎず、またはこれ以上のものは必要ありません。このような状況では、算術平均が最も自然に選ばれます。これは長い間使用されており、広く理解され、使用されるため、また計算が非常に簡単であるからです。

4. 算術平均はほとんどの通常の用途において最も単純で効果的です。実際には、少量の人工的に選択された量に関連する場合、それは唯一考えられる平均であり、例えば数 dozen の長さの測定結果や価格の推定結果に対して他の方法を適用することはありません。

しかし、大量の測定結果を考慮するとき、これらの測定が「確率曲線」のような形でグループ化できる場合、我々は二つの選択肢以上に直面することになります。まず、p.29に示された熟知した曲線に戻りましょう；あるいは、より良いのは次の章（p.476）で最初の形式を示すものを使うことです。そこでは、曲線の頂点を説明する三つの異なる方法を見ることができます。頂点を「最大値」の座標と呼ぶことができます；または、「曲線の中心」と呼ぶこともできます；また（後述される通り）、変数量のすべての異なる値の算術平均が指向する点としても定義することができます。これら三つは全て、位置を説明する異なる方法です；しかし、二項分布や指数形のような対称的な曲線の場合、これらの三つの方法は結果として同じ結果を導き出すことになります：これは現実の例において明らかです。

しかし、非対称な、または偏った曲線の場合に移ると、これら三つの方法による示唆は通常、全く異なるものとなります；したがって、算術平均や一般的な平均とは異なる種類の平均を表すために、前二つは簡潔に説明する価値があります。それぞれについて何か特徴があり、それが直感的に自然で適切であると一般の理解から推奨されることが明らかになります。

5. (3) このこれらの中から、最も頻繁に表される特定の量を選択する最初の方法があります。技術的な名称はまだ獲得されていませんが[1]、そのグラフィカルな表現によって「最大縦座標法」と呼ばれるだけです。しかし私は、このような平均値や基準への依頼が実際には意外と一般的であると考えていますし、もし比較的未開発な人々の判断に潜在する概念を明確にすることができれば、そのような平均値が自然に使用されるいくつかのクラスのケースを見つけることができるでしょう。例えば、魚の大きさは非常に異なるものの、最も一般的な大きさは最大または最小の近くにある漁場があったとします。もし人々が魚を重量で販売していた場合、彼らはおそらくすぐに算術平均や平均という概念についての何らかの理解を得るでしょうし、これは最も適切なテストであることを認識するでしょう。しかし、魚がサイズごとに分類され、それぞれのサイズごとに個数で販売される場合、私は最大縦座標への依頼が始めて他の方法を取ることになると考えています。つまり、最も多く存在するクラスは、同じ漁場や異なる漁場との比較のために一時的に選択される種類として認識されるでしょう。

次の章で見るように、この種の平均が特別なケースにおいて科学的な優先度を持ついくつかの理由もあります；すなわち、扱う量がある大きさに対する「誤差」であり、また正確に正しくまたは非常に近いほど重要であることが、単なる低い平均誤差よりも重要な場合です。

6. (4) 残りの平均は「中央値」と呼ばれるものである。この種の平均は、Mr Galton の著書によって統計学者たちに広く知られるようになり、以下のように説明できる。すべての対象をその大きさの順序で並べるか想像するか、あるいは同じ数の等しい数のクラスに分類されるか想像するか、どちらでも同じことであるが、行の中央にあるもの、または中間クラスの中央にあるものは「中央値」である。私は現在この種の平均が完全には認識されていないとは思わないが、Mr Galton の自然測定法として彼が「パーセントリル」と呼ぶ方法が一般に採用されるようであれば、このようなテストは重要なものになるだろう。この種の平均にはいくつかの顕著な利点がある。まず一点として、ほとんどの統計的な調査では、計算自体が非常に簡単であるという点があり、さらにその過程は別の重要な要素を指定することにも役立つ、すなわち「推定誤差」である。また、Fechner が指摘するところによれば、算術平均の場合には、少数の例外的で極端な値が比較的優勢を占めることによって困惑することがしばしばあるが、中央値の場合（その数だけではなくその極端な大きさも考慮されない）では、このような混乱の重要性は軽減される。

7. 簡単な例を挙げれば、対称的な誤差法則においてはこれらの三種類の平均が一致し、非対称的な法則の場合には完全に異なるものになることが示されるであろう。

[図: 平均の Various 定義]

OBDC上の量を測定する場合、極端な限界がOBとOCであるとする。誤差の法則は三角形BACで表される。このとき、長さODは算術平均、中央値、そして最も頻繁な長さとなる：その頻度は最大縦座標ADで示される。

一方、極端な長さがODとOCである場合、三角形ADCが誤差の法則を表すとする。最も頻繁な長さは以前と同じODであり、これは最大縦座標ADによって示される。しかし、平均値は今度はOXとなる：ここでDX = 1/3 DCである。そして中央値はOYで、DY = (1 - 1/sqrt{2}) DCである。

自然現象から例を挙げると、連続した日における気圧の高さが同一時間に測定されたものがある。4857個のこれらのデータが十分な経験の基礎となるとすれば、確実にその頻度曲線は非対称であるように思われる。ここでの平均高度は29.98であり、中央値は30.01で、最も頻繁な高度は30.05であった。これらの近似がどれほど密接かは、非対称性が僅かなものであることを示している。

8．平均値の種類に関わらず、単にそれが観測された実際の多数の値を一つの代替品として持つことから、重要な情報量を失うことが明白である。実際には、これは経済的な理由で導入されるものである。特定の一目的のために使用されれば、例えば羊の例のように、損失は生じないかもしれないが、一般的な目的では元々の多様性を完全に置き換えることは不可能であり、それらが含んでいたすべての情報も提供できない。この全てを保持するためには別の方法が必要である。

実際的な観点から、我々は通常二つのことを行う：（1）統計表にすべての数字を記録するか、または（2）図形を使用する。後者の案はデータが非常に多く、あるいはその法則性の本質を表示したり発見したりしたい場合に便利である。

平均値を設定することは、ほとんどすべての詳細を省略し、ただ中央値のような何かを示すだけです。それは特定の範囲内の原始的な量がどのように集約されているかについて何も言及しません。たとえば、2つの量がそれぞれ25と27である場合や15と37である場合でも、その算術平均は同じ26になります。

9. この段階で打ち切ることは明らかに問題を不完全な状態に残すことになるでしょう。そこで私たちは自然に、個々の結果が平均値を中心にどの程度集約されているかを示す簡単なテストを求め、一部の情報を回復しようとします。

もし誰もこれが全く新しい問題に取り組む場合、つまり「最小二乗法」理論に関する数学的な要件について何も知らない場合、私は彼はこの業務を行う唯一の方法として以下のように行動するでしょう。彼は言います：「我々がすでに得た平均値は、原始的な量の中間に中央値を設定することで概略的に示しました。もし私たちはこれらの量がどれほど近いかについての概略的な指標を得たいのであれば、その最良の方法は、それらの量から平均値までの差異（技術的には「誤差」と呼ばれるもの）の平均を設定することです。例えば、13人の男性の身長が5フィートから6フィートまで等差に変化している場合、彼らの平均身長は66インチであり、この平均からの平均的な差異は3-3/13インチであると言えるでしょう。

このような視点から見れば、私たちはそれぞれの平均値が目的を果たすかどうか試してみることに進むべきです。その中で2つは完全に目的を達成します：算術平均と中央値です。そしてすぐに明らかになるように、これらはしばしばこの目的のために使用されます。また、我々はもし pleased であれば幾何学的平均値も使用できますが、計算の困難さからそれが退屈であるため、そのような使用は稀です。「最大縦軸」は明らかに適していません、なぜならそれは通常（図443参照）既に得られた平均に戻すための指標となるため、情報を提供しないからです。

ここでの唯一の疑問点は、代数において「誤差の符号」と呼ばれるものが関係するところです。等しい大きさで反対方向にある二つの誤差を代数的に加算すると、相互に打ち消しあいます。しかし、ここで述べるように、誤差が実質的な量として考慮される場合、つまり個々の観察値についてのみ評価される場合、私たちはその実際の大きさだけに注目します。そしてこれらの等しい大きさで反対方向にある二つの誤差は同じ扱いを受けるべきです。

10. これまで議論した様々な平均法の中で、二つが一般的に使用されています。これらの中の一つは天文学や他の計算において「平均誤差」として広く知られており、既に述べた算術平均の原則に基づいて誤差に対して適用されているため、さらなる説明は必要ありません。したがって、先ほどの例では、高さの平均値が66インチであり、誤差の平均値が3-3/13インチでした。

もう一つは「中央値」で、ここでは常に別の名前で知られています、すなわち「推定誤差」として。これは技術的な用語であり、多少混乱を招く可能性があります。この中央値は、我々が等しく超える確率と欠ける確率がある誤差に対応するという意味で定義されます。言い換えると、すべての誤差を大きさ順に並べ替えた場合、その中間にある一つの誤差が対応します。したがって、これは「中央値」の誤差です：もしくは、すべての観測値を連続的に並べ替え、四つの等しい数のクラスに分類する場合、ガルトンによって「四分位」と呼ばれる最初と第三の分類がそれぞれ推定誤差の両側の限界を示し、中央のものは「中央値」を表します。この中央値は、対称的な曲線において算術平均と共鳴します。

受け入れられた用語を守るべきですが、読者にはそのような誤差が厳密な意味での「確率的」なものではないことを理解していただく必要があります。実際、特定の例においてこの誤差を得ることは非常に低い確率です：事実上、我々が正確さを追求し、無限の数の中から一つの正確な値としてそれを考慮するなら、そのような偶然性は無限に低くなります。また、真実の範囲内であるという確率も言えません。定義により、我々はそれ以上になる可能性と不足する可能性が等しいのです。既に述べた通り（441ページ注を参照）、「最大正接」は最も確率的な値として考慮されるべきでしょう。

11. (5) 平方平均誤差．以前の提案によれば、数学の高度な側面に関心がない人なら、目的の指標のために「平均誤差」を選ぶことが自然です。しかし、実際には、分散や乖離量を測定するためのテストとして用いる一般的な種類の平均は非常に異なります．x\_{1}, x\_{2}, ... x\_{n}という大きさがあるとき、これらの通常の平均値は 1/n (x\_{1} + x\_{2} + ... + x\_{n}) であり、その「誤差」はこの値と x\_{1}, x\_{2}, ... x\_{n} の差です。これらを e\_{1}, e\_{2}, ... e\_{n} と呼びます。これらの誤差の算術平均（符号に関係なく）は 1/n (e\_{1} + e\_{2} + ... + e\_{n}) です．一方、平方平均誤差[3]は 1/n (e\_{1}{2} + e\_{2}{2} + ... + e\_{n}{2}) の平方根です。

後者の平均値を他のどの平均値よりも採用する理由は、次の章で説明します。現在のところ、私たちは平均の一般的な論理的性質にのみ関心があり、したがって、どのような中間値でも、我々の様々な尺度が互いにどれほど近いか、そしてその共通の平均からどの程度離れているかをざっくりと示すために十分な役割を果たすことを指摘するだけで充分です。もし「最初の」平均値と「第二の」平均値という用語を使用したとしても、前者は複数の原始的な値を単純な近似値で置き換えるものであり、後者はそのような値が前者からどの程度離れているかを類推的に評価するものです。

12. これまでに述べたのは平均の一般的な性質と、実際には主に使用されている主要な種類の平均に関するものでした。現在は、平均が主にどのような目的で使用されるかについてより詳細に調査しなければなりません。

この点ではまず最初にすべきことは、読者の中に疑問を抱かせることが必要です。彼がこれまで全く疑問を持たなかった問題についてです。誰もが平均値を使用して正確性を確保する実践は多少なりとも知っているでしょう。しかし、我々が最初に行うことは正確性を犠牲することです。実際の結果の複数から一つの結果を得ますが、その結果はおそらくどの一つにも一致しないこともあります。部屋の異なる部分での温度が異なる場合、平均温度が61°Fであると言ったとしても、その正確な温度は部屋内の僅かな部分だけを占めるかもしれません。また、ある小さな男性集団の平均身長が68インチであると言ったとしても、おそらく誰もがこの高さと完全に一致していないでしょう。

このような正確性を確保する主な方法は、我々が真に目指しているものが目の前にある量ではなく、それらが示唆している別の何かである場合にあると言える。もしこれらの量自体が「不正確」であるとする（これについては後で説明する必要がある）、その単純な平均値は個々のどれとも完全には一致しないかもしれないが、我々が必要とするものはそれがもっと近い可能性が高い。この観点を主題に分けて二つの部分に分け、まず量的な考慮が僅かなものである場合と、特定の誤差法則の決定が求められない場合から始め、次にそのような決定を避けることができない場合について述べる。後者の場合についてはここでは一瞥するだけであり、それらはより詳しい検討のために別章で扱われるため、ここで特に詳しく触れることはない。

13. この過程は、どの種類の尺度でも作成する必要がある人々にとって十分に馴染み深いものです。例えば、私はある物体を測定するために銅製の棒を使用しているとします。これは我々が知るところでは、温度によって拡大または収縮します。その結果は僅かに異なるものとなり、時として少し大きすぎ、時として少し小さすぎます。これらの変動はすべて物理学的な事実であり、もし私たちが銅の性質に関心を持っていたら、それらは我々にとって唯一重要な事実となります。しかし、測定された物体の長さに関心を持っている場合、これらの事実は余分で誘導的です。私たちが目指すべきはその影響を排除することであり、これは等しい過剰と欠如が等しく頻繁に起こる限り、算術的な平均値を取ることによって達成できます。そのためには必要なのは、等しい過剰と欠如が等しく普及していることだけです。変動の法則を知る必要はなく、またはそれが特定の一様な種類であることを確認する必要もありません。図29の図において対称的なものであれば、適切で多様な数の測定値の算術的平均値はこの干渉源から解放されます。そして、この変動原因が遵守している同一の一般条件を満たす他のすべての原因についても同じことが言えます。実際、等しい反対の誤差が等しく頻繁に存在することが、一般的な過程である平均値を取り、正確性を確保するための唯一かつ十分な正当化理由となります。

14. これらの主題においてしばしば注意を喚起するために必要となる区別は、特定の客観的な量が存在する場合と存在しない場合との区別です。一般的に「誤差」という言葉を使用することが、この区別の明確さを損なう可能性があります。銅製の棒の場合で述べた過剰と欠如が等しいという点について、我々は正確にその意味を伝えました。すなわち、「真の」長さ（つまり、認定された基準によって決定されるもの）が特定のインチの分数以上となる毎回の状況に対して、対応する欠如がある場合があります。

一方、固定された客観的な基準が存在しない場合、我々が「等しい余剰」や「欠損」という言葉を意味するのは、恒常的な配置の対称性であることが明らかとなる。測定棒の場合には、変動前の何かが存在していたので、それから始めることが可能だったが、多くの場合では我々が見つけるどんな基準点も平均値によって唯一決定される。

たとえば、ある場所での気圧の観測を考えてみよう。その場所で、時間や季節、天候に関係なく大量の観測を行った場合、これらの観測値の平均は「真実」の高さとして一般的に考えられることが多い。我々が本当に意味するのは、瞬間的な気圧はその上にある空気柱の高さによって一部（そして主に）決定されると共に、地元の温度、湿度、風などの他の要因によっても一部決定されることである。これらの要因は時々効果的であり、時々効果的でないが、その範囲は比較的一定であり、その分布はこの範囲を通じて比較的対称であるため、大量の観測値の一組の平均は他のどの一組よりもほとんど同じになる。この平均値の恒常性こそが「真実」である。私は完全に理解しているが、我々はこの恒常性を超える何かがあるべきだと考えるのは難しいかもしれないが、他の類似なケースの例による類推によって誤解される可能性が高い。

15. 実際的な方法として、様々な種類の平均を決定するためのものについてはほとんど述べる必要はありません；算術的な規則は単純で明確であり、それには長い数字列を扱うことによる避けられない労働だけが伴います。おそらくこの部分について最も重要な貢献は、Mr Galton の提案により提供されています。彼は平均値に中位数を代入することを提案し、これによって多くの物体を集めた単なる行為だけで十分な精度で平均を引き出す方法を示しました。したがって、彼はすべての男性を測定せずに多数の男の平均身長を得る巧妙な提案を与えました。

「未開のchieftain はしばしばその部下の兵士たちを身長順に並べたり、またはその能力における人気評価順に並べることを説得されるかもしれません。しかし、それぞれの兵士を個別に測定するためには、通常の強い抵抗を乗り越えるためにも装備と多くの時間を必要とするでしょう（『Phil. Mag.\_』1875年1月）。つまり、広範な経験から、どの程度均質な男性の群れでも身長は対称的に分布することが多いということが知られており、この状況では並べ替えられた行の中間の男性が平均的な男性を代表し、その身長を測定することができます。さらに、中間の身長は極端な身長よりも著しく多く存在するため、長い行の中間の男性を選択時の僅かな誤差は対応する身長を選択時に非常に小さな誤差だけをもたらすでしょう。

16. いまや以前の章で触れられた主題に再び立ち返ることができます。すなわち、平均と平均値を区別しようとすることが時々試みられる試みについてです。前章で述べたように、前者の用語は私たちが自己の構築した仮想的な結果に関連する場合、つまり観察された量から単なる算術的推論を得る場合に限定されるべきだと提案されました。後者については、ある客観的な量が平均を特有に代表すると仮定される場合に適用されるべきだとされています。

第II章 §4で概説した確率に適する三つの主要なクラスに戻りましょう。これらの最初のものは、賭博の結果を含むものでした。サイコロを10回投げます：上側の点数合計は10から60まで変動します。30とする場合、この10個の平均は3と言えます。もう一つの10回の投げ方を取り扱うと、別の平均が得られることもあります。例えば4としましょう。明らかにこれらの平均には何らかの特別な客観的な対応関係はありません。長い間続けていれば、平均は3.5を中心に収束することが分かることでしょう。この最終的な平均はサイコロの構造に関する知識から事前にかなり安定して主張されるかもしれません。しかし、これは型としての性質を持つ他の真実や現実性はありません：単に平均が収束する方向への極限です。

次のクラスは、特に自然種の特性に関して、最も自然な物体のグループのメンバーを占めるものです。ここでも類似した議論が繰り返されるかもしれません。しばしば、個体数が増加するにつれて平均値が近づく「限界」があることが非常に頻繁に見られます。そして、この限界はすべてのものが可能な限りよく似ているように形成された種類の一種であると考える傾向があります。しかし、より詳細に調べてみると、この観点はほとんど正当化されないことがわかりました。安全に断定できるのは、このタイプが一時的に最も多くなる標本を表しているか、または現在の条件下で最も容易に生成される標本を表していることです。

残りのクラスは異なる立場にあります。何か-kindな試みを行う連続的なシリーズにおいて、我々は目標地点からどの程度ずれているかという系列を得ます。これらの数値は算術的に扱い、以前の場合と同様に平均を求めることができます。これらの平均値は仮定であり、つまり、これらは私たち自身の人工的な推論であり、必ずしも客観的な対応があるとは限らないものです。実際、それらが少数のものである場合、おそらくそのような対応関係はありません。何かがタイプに一致するものは、これらの平均値が最終的に向かう「限界」においてしか求めることができません。なぜなら、この限界は目標地点と一致する固定点または目指す対象と一致します。

17. クェテレールの研究における大きな価値と興味を十分に認めつつも、彼は確かに最初に公衆の注意を自然種類が同じ特性を持つことを示す多くのクラスがあるという事実に向けたものであることに注目した第一人者でした。しかし、これらの技術用語を使用してそのような区別を試みることは望ましくないかもしれません。主な反対理由は二つあります。

まず第一に、平均と中間値との単純な反対性は、自然現象における実際の状況よりも非常にシンプルな状態を示唆するように思われます。先ほど述べた三つの類別物についての参照と、それらそれぞれが含む広範で多様な内容を考慮することで、私たちは「虚構」として見做されるものから「実在」であると主張されるものの非常に徐々かつ不感的な進展を思い出すだけでなく、実在型は時として変動的で瞬間的な性質を持つ一方で、虚構は（例えば賭博ゲームにおいて）永遠に固定されると見なされる可能性があるという重要な事実も思い出させます。生産条件が安定している限り、大量の平均値は常に大まかな特性を維持します。もっと重要な区別は、少数の平均値とその変動的な値および不完全かつ偶然に最終目標に達する性質と、多数の平均値とその徐々に近づく最終的価値との間にあるものであると言えます。すなわち、もし最終目標が存在すれば、その客観的な目的点に近づくものです。

また、この章で述べた考慮事項によれば、平均自体の範囲内には、クェテレールが認識したよりもはるかに多様性があることが明らかになります。彼は確かにこの多様性を完全には無視しなかったものの、主として三つの主要な平均値が一つになるような対称的な配置に限定してほとんどその範囲を制限しました。より単純なケース以外では彼の区別を実施するのは困難でしょう。例えば、多少の非対称性がある場合、最大縦軸（最大縦座標）は他の平均値を除き、「平均」として考慮されるべき点となります。算術的な平均に頼ってこの点を見つけることは不可能であり、しかし、最終型の位置を特定するためにはこの点が重要であることを認識すべきです。

18. 我々は幾度となく指摘したように、確率に関心を持つ対象の特徴は、長期的には一貫性が強まる傾向を示すという点です。そしてこれはしばしば「平均的な」こととして説明されています。しかし、このような順序が無秩序から生まれる配列について特に注意すべきであるとする反論が立つかもしれません。その理由は、算術的平均の本質からして、これが避けられない結果であると考えられるからです。平均値を求める過程は、扱う量に対して等化の傾向を保証するものだと主張できます。例えば、十人の男性がいて、そのうち四人は高身長で四人は低身長とします。そして、彼らの五人の中から任意に選んだ平均値を考えましょう。この数には高身長の人だけや低身長の人だけは含まれないということが予想されますので、各平均値が単独の測定値ほど大きく異なることは理にかなわないでしょう。では、増加する観察範囲における等化過程が明らかに算術的必然性から導かれるものであり、従って事前に断言できるものではないかと問われるのは、正当な疑問でしょうか？

上記の反論には何らかの力強さがあるとしても、その主な理由は選択された例の制限によるところが大きいです。この例では、選ばれた数が全体の大きな割合を占め、極端なケースのみが含まれる可能性を完全に排除しました。ここでの混乱はしばしば必要性と経験的な事実との間で生じますが、平均化過程の真正なる結果を正確に定めるために、より精密な例を検討することが適切でしょう。

19. では、ロガリズムの表から十桁をランダムに選びます。同じ桁が連続して十回も現れるような極めて稀な場合以外は、その十桁の平均は必ず最大値と最小値の中間にあるべきです。どのような種類の平均を考えるにしても、それが意味するのは、極端値の中間に何かを取ることです。したがって、十桁の平均は4.5（極端値の平均）よりも近いはずです。

次に、1000桁を選んでみましょう。それらについてより多くのことは、以前の10桁の場合と同じ程度しか断言できません。もし最初から異なる値であれば（つまり、すべて同じでない場合）、その平均は必ず中間にあるべきですが、これ以上のことは算術的に証明することはできません。純粋な算術的な考慮とは比較して、より大きな数の平均が安定化するという信頼に基礎となる物理的事実があります。それは、平均を導き出す要素自体が徐々に均一性を示す傾向があるということです：異なる桁の割合が、それらを選択する個数が増えるにつれて、より近い等しいものになるはずです。

もし1000桁の分布割合が10桁の場合と同じであれば、平均も同じになります。したがって、算術的な過程は、我々が実際に見出すような驚くべき均一性を確保するためには、ほんの少しの助けしかないことが明らかです。

20. 同じことを別の方法で述べることができる。しばしば、元の要素の配置がどのようなものであれ、継続的な平均化プロセスは必ず特殊な二項または指数分布法則に至るという主張がある。この声明は完全に真実である（特定の注意点を除く）が、それ自体は上記で述べたこととは対立しない。私たちが考慮に入れる例を取り上げてみよう。長期的には個々の桁の配置は最も単純な形となる。これは図において直線ではなく、曲線ではなく表されるべきであり、各桁は他のどの桁よりも頻繁に出現するためである。これで全ての「配置」が網羅されていると言える。次に、10つのような桁を平均化した結果を考えると、より広範な配置の余地があることが明らかになる。合計値は0から100までの範囲であり、従って平均値は0から9までの100の異なる値を持つ。そして我々が見るのは、これらの数値の頻度が二項法則または指数分布法則に従って決定されることである。最も頻繁な結果は真の平均値、すなわち4.5であり、この値から左右に減少し、0と10はそれぞれ平均的な状況下で約1010回の機会の中で一度ずつしか出現しないことがわかる。

ここでの説明は、以前のケースと同様です。平均に関して言及される結果の配置は、「必要」であるという点で限られた意味を持ちます。これは特定の物理的仮定や経験の結果として必然的なものだけです。
すべての桁が等しい頻度で出現することが傾向にある場合、そしてそれらが「独立」である（つまり、各桁が他のどの桁とも平等に関連付けられている）場合、その大きさと頻度に関して配置された平均値は、上記で示した利便法則を算術的結果として表示します。経験が提供できる範囲では、選択された桁の真のランダム性、すなわち等しい頻度での再出現と組み合わせの公正さは、実際的に非常に適切に確保されていることが示されます。したがって、個々の結果の原始的な利便法則に関係なく、平均値に関して我々は常に馴染みのある指数分布法則が主張されるという理論的推論は、このような場合において経験により相当に正当化されています。

次章で更なる修正と精密化に関する議論を保留します。

21. 三種類の平均値（つまり、平均誤差、確率誤差、平方平均誤差）を使用して分散量を検討する際には、二つの重要な考慮事項に注意を払う必要があります。これらは次章で方法最小二乗法について触れることになる際により詳細に議論され、正当化されるが、論理的な目的に対するその重要性は非常に大きく、現在において完全に無視することは適切ではない。

まず最初に、任意のケースにおける誤差の実際的な値を知るためには、厳密には極限または最終平均位置が何であるかを知らなければならないことを指摘しておくべきです。誤差の量は常にこの点から理論的に測定されます。しかし、これが我々が常に持っている情報ではありません。

三つの主要な事件のクラスについて再び考えると、ゲームの運命の場合には、この知識を持っていることが多いことがすぐにわかります。限界を経験に訴える代わりに、実際には単純な機械的または算術的な考慮事項によって導き出し、それから個々のケースや一連のケースにおける「誤差」は明らかに、理論が長期的には得られる結果と比較することで見つかります。

第三のクラスである意図的な目標への努力の場合には、我々が正確にこの目標の価値または位置を知っているかどうかはわかりません。天文学的観測ではそれを知らないため、最小二乗法はその値を可能な限り正確に特定する方法であり、射撃のような実験結果ではそれを知っているため、直接経験によって失敗の性質と量をテストすることができます。

残りのケースについて言えば、我々が自然種類または物事のグループとして呼んだものについては、最終的な限界は存在しないだけでなく、その存在自体は常に疑問視されるべきであり、多くの場合では自信を持って否定することができます。限界が存在する場合、つまり実用上永久に固定されたタイプが見られる限り、それは統計学的な労働によってしか特定できません。これを行った後で、それを使って小規模の観測結果をテストすることができます。たとえば、最終的な男性と女性の出生比が約106対100であることが分かれば、特定の地域や町の統計を比較し、その地域での一般的な平均値からの「誤差」について話すことができます。

したがって、実際の多くの場合において我々に必要な作業は、手元にある有限数の測定値や観察値の平均を取ることである。そしてこの平均値を始発点として誤差を測定する。実際の誤差は確実には分かっていないが、確率的に計算されるだけである。しかし、これ自体は最初に見えたように理論的な欠陥とは必ずしも言えず、我々がこれらの方法を使用すること自体があまり稀であり、計算のためにそれらを使用するよりもむしろ最終的な平均値を発見するために使用される場合が多いからである。つまり、その平均値を既に確保していると仮定するのは一種の「論理的循环」（petitio principii）である。

しかし、「誤差」という用語を、確実に分かっている誤差とその量が確定的に指定できるもの、そして確率的に分かっている誤差とその量が確率的に指定できるものとに区別するためにそれぞれ「誤差」および「残差」として使用することを考える価値があるかどうか考えるべきである。実際には、「誤差」と「残差」の2つの用語を、それらによって決定された（一般的に未知の）実際の誤差と観測誤差との間で区別するために提案されている。

22. (2) 他方の問題は、最初の2つの検定法（第446,7頁）が個々の結果が平均値を中心にどの程度集団化しているかをどれだけ信頼的または完全に評価できるかに関わる。答えは、これらは必ずしも完全ではないということである。単一の推定値や量は、さまざまな量の数について十分な説明を与えることは不可能である。この点は非常に重要な問題であり、私はそれが充分に考慮されていないと考えている。その結果として、後述するように、「平均平方誤差」が最小となる方法が必ずしもすべての場合で最良の結果を提供することと仮定されることが過度に簡略化されている。しかし、これによって特定の目的に対する最も良いテストがすべての目的でも最適であるとは必ずしも明らかではない。

したがって、一つの目的に対する最も良いテストが他の全ての目的でも最適であるとは必ずしも明確ではないという点を考慮することが重要である。

これらの検定はすべて平均的なものであり、そして平均というものは必ずしも詳細な情報全てを反映するわけではありません。多くの身長について考えてみましょう：60インチのものがいくつか、61インチのものがいくつか、その他のものもいくつかありました。これらを68インチという平均値に置き換えれば、多くの詳細な情報を失うことになります。そして、これらの誤差や平均からの逸脱量の一部を取り戻すために、それらの身長から平均値への逸脱量を再考することによってその一部を取り戻そうとします。しかし、前と同じく、完全な詳細ではなく、その平均値からの誤差の平均値を代入します。唯一異なる点は、算術平均ではなく、「平方平均誤差」と呼ばれる平均を使用することがしばしば好まれることです。

23. ここで疑問が生じることがあります。それは十分に重要で、短い議論が必要な問題です。測定結果が手元にあるとき、なぜ単に（1）平均値と（2）この平均からの平均誤差を指定するだけで十分であると考えられているのでしょうか？答えは、当然ながら、粗略な近似が求められているという事実によって一部説明されるでしょうが、これだけでは不十分です。さらに重要なのは、私たちは単一の誤差法則を考慮する必要があると仮定していることや、少なくとも馴染みのある法則からの逸脱が微々たるものであるという事実があります。つまり、29ページの図に戻ると、未知の量または調整可能な定数は2つだけであると仮定します：まず第一に、中心の位置、そして第二に、曲線の歪みを表現するための度合いです。平均値の決定は直接かつ即座に前者を指定し、平均誤差（既に述べた方法のいずれかで）の決定は後者を限定的に一つの可能性のある曲線に制限することによって间接的に指定します。

誤差の法則の一つだけを除き、平均誤差の決定は、私たちの機会曲線の概形について僅かな示唆しか与えません。そのような場合、私たちは三つ目や四つ目の「平均」を加えることで、逐次近似法を採用することが便利だったかもしれません。我々が量の平均値とその平均からの偏りを定義するように、この平均誤差（いかなる方法で決定しても）も新たな出発点として取り入れ、それからその平均からの偏りを定義することができます。もしその点がさらに議論に価すると考えれば、図を使ってこのような示唆が最終的な機会曲線または誤差の法則の形態についてどのように逐次近似されるかを容易に説明できるでしょう。

*        *        *        *        *

この本は主に、論理的問題に関心を持つ人々向けに書かれていますので、実際の計算プロセスへの導入としてはそれほど詳しくないため、可能な限り数学的な詳細は避けてきました。そのため、誤差の法則の指数方程式や対応する「確率積分」についての参照も少ないです。ただし、これらの特定のトピックに関連して、多くの学生が困難を感じるか、あるいは感じるべきである二点については、ここで注意を払うことが適切でしょう。

（1）誤差の通常の代数的表現、すなわち \( y = \frac{h}{\sqrt{\pi}} e^{-h^2 x^2} \) について言及した際には、私は常に \( y \) が特定の大きさ \( x \) の誤差の数に比例すると述べてきました。絶対的に \( y \) がその数を表すと言えるのは正確ではないでしょう。なぜなら、可能性の連続性を仮定した場合、どの具体的な大きさに対する実際の誤差の数は無限に小さいからです。したがって、尺度の2つの連続的な分割間で実際に存在する誤差の数を特定するために、連続から離散へ移行したい場合は、通常の微分記号に従って、一つの分割（\( dx \)）における誤差の数は \( dx \left( \frac{h}{\sqrt{\pi}} e^{-h^2 x^2} \right) \) であると言わなければなりません。ここで \( dx \) は長さの単位であり、両方とも \( h^{-1} \) と \( x \) の値を測定する必要があります。

多くの学生が感じる困難は、この公式を実際の統計に適用すること、すなわち適切な単位を代入することです。具体的な数値例として、1460人の男性について身長を「最も近いインチ」まで測定したとします。そして、ここでのモジュールは3.6インチであることが分かっています。すると \( dx = 1 \)（インチ）で、\( h^{-1} = 3.6 \)インチとなります。このとき、\(\sum \left( \frac{h}{\sqrt{\pi}} e^{-h^2 x^2} \right) dx = 1\) であり、つまり全ての連続的な可能性の値の合計は1に等しいです。したがって、ここで求めたい合計が1460である場合、公式を次のように表す必要があります：\( y = \frac{1460}{\sqrt{\pi} \times 3.6} e^{-(x/3.6)^2} \)、または \( y = 228 e^{-(x/3.6)^2} \)。

ここで \( x \) は中央値または平均身長から測定したインチ数を表し、\( y \) は統計テーブルにおけるその身長に該当する男性の数を表します。（各 \( t \) の値に対する \( e^{-t^2} \) の値は手引きで与えられています。）

例として、この公式によって計算した \( x \) から0インチまで8インチまでの数値と、最近マガジンスで実施されたギャロット先生による測定における実際に観察された数値を示します。

インチ  計算値   実測値
x = 0    y = 228   = 231
x = 1    y = 212   = 218
x = 2    y = 166   = 170
x = 3    y = 111   = 110
x = 4    y =  82   =  66
x = 5    y =  32   =  31
x = 6    y =  11   =  10
x = 7    y =   4   =   6
x = 8    y =   1   =   3

ここでは平均身長が69インチでした：dx、つまり1インチとして述べられています。'x = 0'と述べたときには、68.5から69.5の範囲に属する男性の数を計算することを意味します。一方で、'x = 4'と述べたときには、64.5から65.5の範囲または72.5から73.5の範囲に属する男性の数を計算することを意味します。実測結果は、計算値に非常に近いことがわかります：前者の場合においては、平均からの等しい反対方向への離差の平均が採用され、実際の結果は常に反対方向とは限らないことに注意してください。

(2) 他には、馴染みのある確率積分（2/√π）∫\_{0}{t} e{-t2} dt の解釈に関する点があります。多くの手引き書で提供されているこの積分の表を使用して事件の確率を計算した経験がある人なら、任意に選んだ数値 t に対して、上記式の対応する値は、ランダムに取り出した誤差がその限界内にある確率を示します。すなわち、t です。たとえば、t = 1.5 を代入すると結果は0.96となり、つまり誤差の4パーセントしか '1.5' よりも大きいものではないということです。しかし、「1.5」とは何を指すのかという質問に対しては、答えは必ずしもすぐに得られるわけではありません。通常通り、初心者の主な難点は公式の操作ではなく、単位が明確であるかどうかにあります。

この例では、直感的な単位が選ばれ、それに基づいて計算を行うことがわかります：ここでは一つの変数（t）しか存在せず、インチやフィートなどの任意の単位を選べる場合でも、異なる結果を得ることがあります。したがって、我々はその自然な単位で計算を実行する必要があります、これは角度の円周上の測定と度数による測定を区別するときと同じです。

単位は「モジュラス」であり、't = 1.5' を置くことは、「誤差がモジュラスの半分より約二倍大きい」という意味です。モジュラス自体は、測定や観察の性質によって異なる一定の大きさの誤差です。この点を理解するには、以下の積分を形として置くのが良いでしょう：
\[
\left(\frac{2}{\sqrt{\pi}}\right) \int_{0}^{hx} e^{-h^2 x^2} d(hx)
\]
これは、定積分の値は使用される特定の変数に依存しないため、等価です。ここで hx は x : 1/h と同義であり、すなわち x が 1/h の単位で測られた比を表します。しかし、1/h は誤差の法則
\[
y = \left(\frac{h}{\sqrt{\pi}}\right) e^{-h^2 x^2}
\]
におけるモジュラスです。つまり、この式における誤差の数値的な値は、その中に含まれるモジュラスの倍数である整数または分数の回数を表します。

1. この種の平均はフェンシャーなどによって「\_dichteste Werth\_」と呼ばれます。私が見た中で最も適切な引用は、レックスス教授（『Massenerscheinungen』、第42ページ）によるもので、彼はこの平均が約70年という人間の寿命の標準的な長さを明確に示すことを示しています。これは算術平均を使用するときにはほとんど隠れてしまう結果です。

この平均は「可能性のある」値と呼ばれるべきであり（ただし、その名前は他の用途で使用されています）、それは最も頻繁に発生する点を示すからです。つまり、無限に小さな等しい単位の変動を比較したとき、それに対応するものが最も頻繁に表される傾向があることを意味します。

2. この数値結果を説明する図は『Nature』（1887年9月1日）で示されています。上記のように異なる平均を計算する際には、最初の結果が3桁まで与えられていたことに注意してください；しかし、分類する際に、只1桁だけが考慮されました。つまり、29.9は29.900から29.999までのすべての値を含みます。したがって、私のテーブルで最も頻繁に記載された値は30.0ですが、通常の補間原則によればこれは30.05と算定されます。

3. この文脈での表現には多少の曖昧さがあります。例えば、Airyは一般的に「平均平方誤差」を表すために 'Error of Mean Square' という表現を使用しますが、ここでは sqrt{(sum e{2})/n} を示しています。Gallowayはしばしば「誤差の平均平方」という表現を使用し、(sum e{2})/n を示します。私は前者の使用を採用し、それを簡略化して E.M.S. と表記します。Mr Merrimanや他の人々がこの E.M.S. の代わりに '平均誤差' (より自然な意味での広く使用されている表現) を使用することは、私には非常に不運だと考えられます。

4. 実際的にはもちろん、拡大または収縮を考慮に入れる必要があります。しかし論理的な説明の目的では、この変動を平均化によって中和されるような干渉の一例として取り扱うことができます。

5. 厳密には多項式：異なる数値の相対頻度は x のべき乗の係数によって (1 + x + x{2} + ... + x{9}){10} の展開で示されます。

6. Mr Merrimanの「最小二乗法」に関する著作においてです。

\newpage

\chapter{真理への近似手段としての平均理論（Theory of the Average as Means of Approximation）}
\label{ch:xix}

\begin{center}
\large\textit{誤差曲線の縮小（逆問題と順問題）・標本平均の極限分布・天文学的観測の総合}
\end{center}
\vspace{4mm}

第十九章。

平均値の理論は真実への近似手段としての側面を考察する

1. 前章では、主に平均値の質的側面よりも量的側面について論じました。つまり、一般的な性質、主要な種類、そして日常生活や非常に正確ではない推論プロセスでの主な用途について議論しました。現在は、確率に特有の方法で平均値を利用する特定の問題に入り込んで詳細に考察する時です。すなわち、私たちは一定数の測定結果が与えられており、その平均を取る理由や信頼度について何らかの質問をする準備ができていると仮定しなければなりません。つまり、常に同じ種類の平均を取る必要があるのか、またそれが必ずしも必要な場合でも、どのような場合に異なる種類の平均を使用するべきなのか。

このように入ることになる主題は、その最も一般的な理論的な扱いにおいて、おそらく他の数学上の問題よりも深い調査、より多様な意見、そして結果として広範で詳細な歴史と文献を生み出したものかもしれません。しかし、それでも私は、これらの方法やプロセスの基礎となる主な論理原則は特に難しくないと思うのです。ただし、主題についての議論すらも非常に技術的な表現スタイルが採用されるため、数学的リソースに中程度の掌握力しかない人々にとっては、これらの原理をその記号から解き放つことは容易ではありません。本章では、この全体的な理解のためにこれらの困難を排除する試みを行います。

このような扱いは多くの分岐を伴うため、読者はおそらくこの巻の冒頭にある目次を時々参照することが便利であると思われます。

2. 当然の形式で議論するにあたって、主題は平均値に限定される。これは平均値に関する比較的単純さと広く普及している性質により、問題のこの側面を考慮することによる。しかし、その問題は数学的な文献以外でもしばしば「最小二乗法則」または「最小二乗法」として参照され、我々が扱うような場合においては、最小二乗法則はより単純で一般的な算術平均の過程に帰着される。 Herschel によって提示された非常に単純な例を用いて、広い範囲での任務の一般的な性質を説明し、その馴染み深い名称が正当化されることを示す。

ある男が壁に張られた小さな標的、例えばワッフルを射撃するという想定を考える。私たちは射撃痕が中心であるワッフルを中心に集積することを仮定してもよいと考えられる。その密度は中心から距離によって一定の比率で変化していると仮定できる。しかし、現在、中心を示すワッフルが取り除かれ、壁に散在する射撃痕のみを見る状況となったとする。そして、このワッフルの位置を推測することを求められたとする。

もし射撃が一度だけ行われた場合、直感はその標的が実際の中心であることを仮定することを提案するだろう（もちろん非常に慎重に）。二回目の射撃が行われた場合、直感は両点の中間点を選択することを提案するだろう。しかし、三回以上の射撃が行われた場合、直感は迷う。中心の一部であるような点を選ぶべきだと感じながらも、より正確な決定をする方法を見出せない。なぜなら、その慣習的な依頼、つまり算術平均がここには用いられないからだ。この問題に対する規則は我々にどのように進むべきかを教えてくれる。それは各射撃痕からの距離の平方和が最小となる点を選択することを指示する。

この部分は例示のためであり、規則に一般的な呼称を正当化するものである。我々が専ら扱うケースは、比較的単純なもので、その中では線形量またはそのような量を十分に表現できる量のみが考慮されるべき対象となる。これらの点に関して言えば、最小二乗法の規則は最も一般的な意味での平均値の計算へと簡素化され、算術平均を取ることによって実現される。また、誤差の単純な法則、またはそのグラフィカルな等価物である確率曲線が十分に異なる量の頻度を示すのに役立つ。

3. ここで以前の章で既に指摘された誤解や混乱について再度注意を喚起することができる。それは誤差の法則と最小二乗法の混同である。これらは全く異なる種類のものである。前者は物理的な事実の性質であり、その生成は多くの場合我々の制御を超えたものである。後者は、あるいはその簡略化された応用形態、例えば算術平均など、物理学的な意味での法則ではない。むしろそれは私たちに導きを与えるための教則や規則である。法則は、特定の場合において、誤差が大きさと頻度に関してどのように発生するかを述べている。一方で、方法は我々が何らかの数の誤差に対処するためにどうすべきかを指示している。両者の間には関係があることは次のページで明らかにされるであろうが、実際には異なる誤差の法則に対して同じ方法を使用したり、同じ法則に対して異なる方法を使用することも全く問題がない。そのような場合においては、明確な正しさや誤りという観点ではほとんど問題が生じないが、より適切さの観点では問題が生じる可能性がある。

4. 読者には、ピストル発射の例で示されたように理解していただきたいのは、我々が直面している最終的な問題は逆問題であるということです。つまり、我々はいくつかの「誤差」を前提とし、それらが発散する中心点の位置について述べることを試みます。これは原因から結果を観察することによってその原因を決定するようなものですが、逆問題の場合のように、私たちは直接的な問題から考慮に入れるべきです。言い換えれば、我々が直面しているケースに関しては、最終的に私たちの目標である真の確率曲線の中心点が既に分かっていると仮定し、その場合、残される作業は平均値を取ることによる結果の研究に集中することになります。

5. 当面は、確率に関する考慮事項から見れば典型的なケースについて述べることに限定します。すなわち、配置または発展の法則が二項分布の種類である場合です。この法則の性質については第2章で説明しました。そこで、それぞれの出現在数の頻度は、二項式（1 + 1）n の展開の連続する項の大きさに従って規制されていることが示されました。また、nが非常に大きい場合、すなわち影響する状況の数が非常に多く、個々の影響力がそれに比例して小さくなる場合、これらの二項式の連続する項を表す縦軸の頂点を通る曲線が方程式 y = Ae{-h{2}x{2}} によって指定される形に近づくことが指摘されました。

従って、実用的な観点から言えば、二項分布の法則や指数分布の法則を区別せずに話すことも可能です。もし一方の法則や他方の法則による現象の配置が、与えられた数値が大きいときには、それらはすぐに区別できなくなるからです。

しかし、この理由だけではありません。現象自体が連続的な量を表しているとしても、私たちが扱える測定結果はすべて離散的です。百万人の成年男性の正確な身長を考えましょう。実用的にはこれらの数値は連続的な量の変動を代表しています。二つの連続的な量の差異は特に平均近傍では非常に小さく、無視できるほど小さいからです。しかし、私たちのテーブルはおそらくインチ単位でしか表示していません。69インチが何人か、70インチが何人かと記録されています。つまり、これらの表記は硬貨を数えるときの「正面」の回数に相当し、二項分布上の離散的な数値を扱っているようなもので、指数分布上の連続的な量とは異なります。

6. したがって、現時点では二項分布や指数分布の法則に関するこの一般的な問題について、生成条件に関する知識があるかどうかによって、変動する量の2つの異なる場合に区別しなければなりません。

(1) 選択肢の最初は、問題の条件が事前に決定可能である場合です。つまり、特定の経験以前に、長期的にはどの組み合わせがどれほど頻繁に起こるかを言える場合です。この場合、平均値を使用する主な目的、すなわち真の平均値を見つけるという目的は代替されます。私たちは長期的に見ればどのような中央値になるかを言えますから、その平均値を求めることはありませんし、少量の観測結果からそれを決定することも必要ありません。しかし、この場合について慎重に議論する必要があります。なぜなら、他の場合での論理において重要な前提条件となるからです。

この比較的して先験的な知識は、その完全性に関して二つの異なる程度で現れることがあります。まず最初に、関心事に関連する状況においては絶対的に完全である場合があります。たとえば、10ペンスの硬貨を何度も投げ上げる結果について考えます。ここでは平均値が正確に分かっています：表と裏の均等な割合です。また6つの表と4つの裏が出る確率も知っています。つまり、各組み合わせの相対頻度を示す図を作成する必要がある場合でも、経験に頼ることなくそれを描くことができます。

しかし今、長くて細長い目標板について考えてみましょう。その板上の一地点が射撃の中心点としてマークされています[3]。ここで（永久的な偏りを生むような原因が働いていないと仮定する場合）、この中心点は平均値に対応することを知っています。また、この中心点から左右にどれだけ分散があるかについては、一般的な観点からしてbinomial法則に従うことを知っています。しかし、前記の例のように頻度の各程度を表す縦線を立てて相対頻度の分布を描こうとする場合、そうすることはできないことがわかります。新たなデータが必要となります。

上手な射手中にも下手な射手中にも、その射撃は同じ一般的な法則に従いますが、中心点から離れるほど射撃が減少する速度は異なるでしょう。頻度曲線を正確に描くためにはもう一つの「定数」が必要となります。

7. (2) 次に考慮される第二の分類は、論理的な観点からは最初のものと同等であり、長期的には値が集積する状況についてあらかじめの知識がない場合でも、十分な経験により実用的な確信を持ってその状況を指定できるケースを含みます。たとえば、身長表を考えてみましょう。これらの表はしばしば非常に広範で、十万や数十万人を含むこともあります。このような場合においても、平均値または中央値はあらかじめのルールによって同様に確実に決定できます。つまり、同じ人口クラスからもう一度百万人分の測定を行ったとしても、測定器具が実用的に認識できる範囲内での変化により平均値が大きく変わることはありません。

8. しかし、平均値または中央値の指定だけでは、前段落の場合と同様に、私たちが必要とする全ての情報を提供しないこともあります。たまたま二つの人口群の平均身長が同じであっても、その周りでの分散の法則は非常に異なる可能性があります。つまり、一つの系列で特異な巨人や矮小の人間であった人が、別の系列では全く目立たないかもしれません。

この分散の実際的な大きさを決定するプロセスを説明すると、過度に数学的な詳細が必要になりますが、一部の指針は提供できます。私たちがすべきことは、方程式[4] \( y = \left(\frac{h}{\sqrt{\pi}}\right) e^{-h^2 x^2} \) の定数 \( h \) を決定することです。技術的な言葉で言えば、私たちはこの方程式の「モジュラス」を決定する必要があります。上記式における \( 1/h \) はモジュラスと呼ばれ、この方程式によって示される平均周りでの収縮または分散の程度を測定します。モジュラスが大きいほど分散は大きくなります；つまり、中心に近い部分には大きな値が密に集まらず、モジュラスが小さいほど中心周辺に多くの値が密集します。標的の中心周囲で砲弾痕が描かれる曲線を表すモジュラスが小さいほど、射撃手は優れています。

9. 模数の決定にはいくつかの方法があります。上記で最初に議論した場合において、理論的な知識が完全であるため、確率に関する知識から事前に計算することが可能です。硬貨を大量に投げる際など、自然にこの計画を選択します。

通常、測定値や観察結果がある場合の後述の事後的な計画は次の通りです：誤差の平方平均を取ってそれを2倍すると、その結果が模数の平方を与えます。例えば、5つの大きさ 4, 5, 6, 7, 8 を考えてみましょう。これらの平均値は 6 であり、「誤差」はそれぞれ 2, 1, 0, 1, 2 です。したがって「模数の平方」は 10/5 に等しく、すなわち模数は \sqrt{2} となります。

もし大きさが 2, 4, 6, 8, 10 で、同じ平均（6）を示しながらより広範囲に分散している場合、模数はより大きくなります。すなわち、\sqrt{8} ではなく \sqrt{2} となります。

ガルトン氏の方法はよりグラフィカルなもので、統計による相互比較に関する論文（\_Phil. Mag.\_ 1875年）や他の場所で説明されています。以下にその概要を示します。大量の人間の身長測定値を考えるとしますし、すべての人々をその高さ順に並べます。平均高さを選択し、行中央の人物が 69インチであると仮定します。この点から刻度を上げ（または下げる）して、平均より上（または下）の半数の男性を含む高さまで立てるまでします。（実際には、これは約1.71インチが必要となることがわかります：つまり、どの大きなグループの男性の四分の一は69から70.71インチの間になります。）

この数値を 0.4769 で割ると模数が得られます。この例では約3.6インチとなります。

私たちは、誤差の法則が一般的な二項形またはこれに近い形で現れるという前提から始めます。この仮定に基づいて、上記で示した三つの方法は結果において一致します。もし疑問がある場合、あるいは事前に分散の速度を計算することができない場合、私たちは第2の計画によりモジュールを決定することに採用しなければなりません。これは唯一の普遍的に適用可能な計算法です：実際、それがモジュールを与えるべきであることは定義上の真実であり、平均平方誤差の誤差を決定することで、私たちは実質的にはモジュールを決定していることを前章で指摘しました。

10. それでは現在到達した立場は次の通りです。誤差の法則が式 y = (h/√π) e{-h2 x2} のようなシンボリックな形に収束すると仮定する場合、h を決定するためのルールが手元にあることになります。したがって、目的に関連しては、私たちは曲線の頻度についてすべてを知っています：紙に描き出すことができます。中心値の一つの値を与えられた場合、この値からどの距離でも他の任意の値を決定することができます。つまり、100万人の中の何人かが69インチの高さであることを知っている場合、直接観察せずに67インチ、68インチ、70インチ、71インチなどそれぞれの人数も決定できます。

現在、私たちが抱える論理的な興味に値する主要な問題について充分に議論できるようになりました；すなわち、なぜ平均値や中央値を取るのか？そのような操作によって得られる正確な性質と利益は何であるか？高等な学生は当然、誤差の法則（最終的な指数形で）を用いてこれらの問題の答えを導き出すことを好むでしょう。しかし、私は、論理的な過程の本質的理解を得たいと考える人々にとって最善の方法は、初等代数で馴染みのある組合せの問題から始めることだと確信しています；つまり、一定数の誤差を取り、それらを平均化した結果は何になるかを見ることです。その後、二つ以上の誤差を組み合わせて新しい系列を得る算術的な結果を計算し、一般的な性質だけでなく具体的なケースにおけるその内容も計算することに進むべきです。

非常に多くの項を持つ系列を取り扱うのは単純化のため避けますが、十分な数の項が必要で、それによって誤差の法則が標準的な指数形に近似されることを確認すべきです。

この目的のために、誤差または乖離の法則を用いることで十分である。これは、我々の努力が10つの原因によって影響を受け、それぞれが等しい誤差を生み出し、その誤差は正と負に等しく確率的に分布している（あるいは、「10つの等しい加減可能な原因」と表現されるべき）という仮定に基づくものである。これは二項法則に従って形成された最低の数であり、目で見ても、極限または指数法則を適切に示す指標となるであろう。ここでの可能性のあるすべてのケースの総数は210，つまり1024である；これは、単に11つの実質的に異なるケースだけではなく、これらのそれぞれが長期的にはどの程度頻繁に起こるかという相対的な頻度も示すために必要な数である。この合計のうち、252は平均値に位置し、「真の」結果を表し、5つの干渉原因が他の5つを完全に中和したときの結果を示している。また、210は平均値から1単位離れた位置にあるもので、6つの原因が4つに対して組み合わさったときに得られる結果を示しており、そのような順序で、平均値から5つの場所離れた極端な距離では、一つの結果しか得られない。つまり、1024のケースの中で、唯一の場合のみにすべての原因が同じ方向に一致するためである。したがって、この1024の努力のセットは、無限個のそのような努力の分布を適切に表している。これらの配置のグラフィカルな表現はここに示されている。

[図：十分乗法による二項分布]

11. この表現が単独の観測や努力の一完全な集合を表している場合、対応する結合または簡略化された観測の集合においては、その数と配置は何となるでしょうか？特に二つずつの観測について考えてみましょう。まず数に関しては、これは同じものを連続して繰り返すことができないものの組合せの問題ではないことに注意しなければなりません。例えば極端な誤差Fについては明らかに一度ではなく二度繰り返すことができます。このような繰り返しは確かに運が悪いことでしょうが、可能性があるためその位置は集合の中に含まれなければなりません。したがって、同じものを連続して繰り返すことができる1024つのものから二つずつ組み合わせる方法の数は1024 × 1024，つまり1,048,576となります。これが結果を二つずつ取り出す一完全なサイクルにおける数です。

12. 数について述べたところまでにして、次にその配置や分布について考えてみましょう。まず確認すべきは、それぞれの可能な観測のペアが何回現れるか、そして新しい観測結果が結合から得られた各ペアに対してどのように配置されるべきかです。最初の調査に関して言えば、一度にはFが二度繰り返されることを見ることができます。また20回ではFとEを組み合わせることになります（Fが最初に現れた場合、それは10個のEからどれか一つ後に続くことができ、それらのどれか一つもFを後尾に持つことができます）。Eは10 × 10，つまり100の方法で繰り返すこともできます。同様にして他のものについても考えます。

これらの各観測値の位置について、その成分要素の相対的な頻度がすでに指摘された。これは容易に決定できる。二つの誤差を取るとき（先に見たように）それらの間の中央点を選択する方法以外はほとんどない。この中央点は当然、二つが一致した場合にそれぞれと同一になる。従ってFは新しい配置において一度現れることがわかる。これは古い配置で二回繰り返されることからである。GはEとFの中間にあるので20回与えられる。Eは新しい配置において二つの方法で得られる。一つには古い配置で二回繰り返される（これが100回起こることを意味する）こと、もう一つにはDとFの中央点として得られる（これは90回起こることを意味する）。したがってEは全体で190回現れる。

このようにすべての可能性のある出現頻度を計算することを選んだ読者は、その方法で計算できる。もし単にその原理を説明するためにこれを示すことが目的であれば、これが最良の方法かもしれない。しかし、彼はすぐに気づくであろうし、数学的な視点から見れば即座に証明できるが、新しい「誤差の容易性法則」は（1 + 1）20 の展開の連続した項を用いて推論的に得られることが早くわかる。それらは476ページ下部の図において与えられている。

13. 旧の単純な観測値の分布と、新しい結合または簡略化された観測値の分布を直接比較する際には、二つの明確な障害があるように思われる。まず第一に、後者の数は著しく多い。しかし、これは両者を同じスケールに揃えることで容易に解決できる。すなわち、それぞれの総数を同じにすることができる。第二に、新しい位置の半分は旧の観測値には代表が存在しない。つまり、FとEの中間、EとDの中間などにある位置である。これは通常の補間法により解決できる。すなわち、欠損点における数を同じスケールで推定し、その点に埋め込むことで対応することができる。A, B, C等の縦軸の頂点を通る曲線を引くと、中間点での縦軸の長さはそれぞれの大きさに対する誤差頻度を非常に適切に表すものとなる。これらの欠損が埋められたとき、そして数値が同じスケールに揃えられたとき、比較の完全な基礎が得られる。

同様に、三つの観測値やそれ以上の数の観測値をグループ化または「簡略化」することもできる。観測値の可能な組み合わせは非常に多くなる。三つずつ取り扱うときには、(1024)3 となる。三つ以上の観測値になると、既に述べたように、いくつかの異なる処理または簡略化の方法が存在し、算術平均を取るという方法はその一つである。

14. 次の図は、このように複数の観測値の算術平均を取ることによって得られる利点の性質を説明するものである。

曲線ABCDは、既に述べた二項法則に基づいて配置された与えられた数の「誤差」を表している。A'CD'は、単純な誤差ではなく、誤差のペアの平均として与えられた同じ数の類似した配置を示している。また、A"BD"は、三つずつの誤差の平均として得られる類似した配置を表している。これらの曲線は、図の小さなサイズに制限されるが、可能な限りスケールに正確に描かれている。

[図：二つまたは三つの観測値に対するガウス分布]

15. その上の図を見ていただければ、読者には、一つの要素だけにとどまらず、二つ以上の測定や観測を組み合わせて平均値を取ることによって大誤差が比較的少ないという文言の完全な意味が、よりよく説明されるでしょう。この利点は、湖で漁をする場合と同じようなもので、同じ数の魚でも大きなものが多く、小さなものが少ないこと、また袋から硬貨を取り出す場合も同様に、金镑が多く、シリングが少ないということが特徴です。しかし、主題に対する完全な理解の極めて重要な意味合いがあるため、この点をより詳細に説明することが望ましいかもしれません。

まず第一に、一連の「平均」結果はただの「誤差」の一組であるということが明確に理解されるべきです。何らかの装置が真実の結果を得ることを確実にすることができるとは思えません。そのような考え方は確率の基礎自体を誤解することになります。すなわち、我々の労働の結果として実際に状況が悪化する可能性さえ排除することはできません。二つ、三つ、またはより多くの単一結果の平均値は、最初に得られた観測結果よりも遠い最終的な平均値を提供することがあり、すなわち、最初の観測結果よりも悪い結果をもたらすことがあります。しかし、長期的には大きな偏差が稀になるという正当化しか残りません。

また、上の調査においては算術平均を一般的に理解し、使用するものとして述べましたが、観測値や誤差を組み合わせる方法としてほとんどどの対称で規則的な方法でも、同じような一般の結果が得られることが指摘できます。二項法則による規則性の主な特徴は以下の2点でした：（1）中心または真実の結果を中心にした最終的な対称性、および（2）この中心に近づくにつれて相対頻度が増加することです。これらのうちの一つを保持し、もう一つを増加させるという点では、算術平均以外にも特別な優越性はありません。しかし、これを述べる際に注意すべき重要な区別があります。その区別について図を参照して説明します。

16. 任意の一連の測定によって目標とした点O（線D'OD上）が考えられ、または現在の目的においては、すべての測定が最終的に収束する平均値として考えることができます。Oに関する値の対称的な配置とは、どの誤差OBに対しても、長期的にはその反対の誤差OB'が存在することを意味します。つまり、Bを生成する頻度を表す縦線BQを引き立てるとき、等しい長さの縦線B'Q'も引き立てなければなりません。したがって、P点を中心に左右に分かれた曲線の2つの半分、すなわちPQとPQ'は完全に同じです。

[図：分布の対称性]

次に、二次曲線について述べるが、これは二つ以上の単純な誤差の平均値の頻度法則を表すものである。この曲線もまた対称的であることが容易に推測できる。任意の三点B, C, Dに対応する別の三点B', C', D'を考える。したがって、BCDのようなすべてのグループ群に対する規則的な対称的な扱いは、中心Oを中心にした対称的な配置を生むことになる。

普通の算術平均はそのような方法の中の一例に過ぎない。この平均値を定める一つの方法は、B, C, Dからその点までの距離の平方和が最小となる点を選ぶことであるという説明を与えることができる。しかし、立方体、または四乗、あるいはそれ以上の冪が最小となるような点を選んでもよい。これらはすべて、図中の点線に似た曲線を生む。

もちろん、このような方法を採用することは実際的な障害が多い。計算の労力は莫大になり、結果は普通の平均を使用した場合よりも劣る可能性が高いからである。しかし、不一致や誤差のある結果に対処する一般的な原理に関しては、馴染み深い算術平均は無数の可能な方法の中の一例に過ぎず、これらすべてが同じような助けを提供することを忘れてはならない。

したがって、これらの一般的な原則に基づいて、不一致な結果や誤った結果に対処する際には、馴染み深い平均値以外にも多くの選択肢があることを念頭に置いておくことが重要である。

17. また、平均値を用いることで、曲線が中心に傾きやすくなるという効果があることを確認した。これは、平均値の誤差はより良い、つまり小さい種類であるという事実を表している。しかし、個々の誤差に対処する他の方法でも、同様の特性が一般的な規則として生じることが注意されるべきである。ここでは厳密な証明を与えることはできないが、組み合わせの一つの有名な結果に参照することで、この傾向がどのように生じるかを示すことができる。

極端な結果は、どのような種類の平均値でも得られる方法は一つだけであり、それは個々の数値から得られた平均値において極端な結果が繰り返される唯一の方法である。一方で、中間的な結果を得るためには二つの方法があり、それは個々の中間的な数値を用いるか、あるいは個々の数値を反対方向に組み合わせることである。

誤差の二項法則においては、この中心に向かって厚みが増す傾向は、平均値を取る前から個々の値において強く存在していた。しかし、組み合わせの特性により、どのような分布法にもどのような種類の平均値を適用しても、その結果は個々の値に対し、点線が黒線に対して持つような一般的な関係を持つと（大まかに言えば）断定できる。

[6] ここで、黒線は個々の値を表し、点線は平均値を示している。

18. このようにすれば、観測値を組み合わせる、または平均化する、または簡略化する一つの方法が、別の方法よりも理論的な利点があること、つまり実用的な便利性に関係なく、その利点は結果が中心に集まる速度の程度にあるのみである。この優位性については、さらに詳細な分析が必要だが、現在の目的には、ある種の平均値が図479の点線の上側を示し、別の種類の平均値が下側を示した場合、前者の方が良いというだけの説明で十分である。この利点は、代数的な計算において速やかに真値に収束する級数が遅く収束する級数よりも真実に近づくという一般的な利点と同種である。どちらの方法でも同じ労働量でより近い結果を得られるか、または労働量を減らして同じ程度の精度を得られる。

ここで個々の観測値が二項法則に従って集約されている場合を考えると、算術平均は単純かつ容易であるだけでなく、上記の意味での最良な方法でもあることが確実である。そしてこの二項法則、またはそれに近いものが非常に広く存在するため、一般的な算術平均の使用に対する強い前提条件が成立している。

19. 数ページ前の分析は、単純な算術の助けを借りてできる限りの結果を得た。さらに進むには、より高度な数学が必要だが、その種類の結果についての概略的な説明は現在の目的に十分である。結局のところ、各ステップは証明のために複雑な論理が求められるが、それらは単純な算術によって確立できるプロセスの一般化に過ぎない。簡潔に言えば、私たちが行うことは次のようなことである：

(1) 首に、二項式形から有限な要素を持つ証明を、極限または指数形式へ拡張します。私たちは小さな数の離散的な誤差に限定せず、どのような大きさの誤差でも何らかの数の誤差がある可能性を認識します。

(2) 次に、平均を考える際に、二つや三つのみに限定するのではなく、既に私たちが見たように非常に煩雑な計算であるにもかかわらず、任意の数n個の平均値の結果を求めます。実際の結果は非常に単純です。一つの誤差のモジュールがcであるとき、それらn個の平均のモジュールはc ÷ sqrt{n}となります。

(3) 最後に、二つの任意の数の平均の和または差について類似した結論を導きます。たとえば、まずm個の誤差を取り、それらを平均化し、次にn個の誤差を取り、同様に平均化します。これらの平均はほぼ等しいですが、完全には等しくありません。その和または差（もちろん、正の誤差と負の誤差は等価で反対方向にあるため、最終的には区別できません）もまた「誤差」であり、どの大きさも特定の確率や発生可能性があります。私たちが行うことは、これらの誤差のモジュールを指定することです。実際の結果はまた単純です。一つの誤差のモジュールがcであるとき、m個とn個の平均の和または差のモジュールはc sqrt{1/m + 1/n}となります。

20. これまでの問題は直接的なものでした。最終的な平均値または中央位置が与えられていると仮定しました；これは先ず知り（多くの賭博ゲームにおいて）、あるいは直近の物理的考慮から（目標への射撃において）、あるいは広範な統計から（身長の表において）です。したがって、このようなすべての場合で主要な目的は既に前提として与えられており、何が残されているのかという質問が正当化されます。答えは様々です。まず一つには、多くの平均値と少数の平均値を比較してその値を推定することもあります。例えば、一人の人が1000回の事例を集計し、もう一人の人は4000回の類似した事例を集計しました。普通の感覚は後者の集計が前者よりも優れていることを認識できますが、それがあくまでどの程度優れているかについては具体的な概念はありません。ここでも他の場面と同様に、定量的な精度は科学の特権です。この分野からの答えは、長期的には、測定または観察の数の平方根に反比例して、モジュラス（そしてこれと共に確率誤差、平均誤差、平均二乗誤差）が減少することです。（これは上記の公式の2番目から導かれます。）したがって、より広範な統計の推定誤差は比較的少ない統計の半分になります。

別の例を挙げましょう。観察によれば、「2315人の犯罪者の平均身長と8585人の一般成人人口の平均身長との差は約2インチである」ということがわかります（Edgeworth, Methods of Statistics: \_Stat. Soc. Journ.\_ 1885）。普通の感覚では、このような差が有意なものであることを疑う余地はありませんが、その具体的な意味合いを数値的に評価することはできません。科学に訴えることで、これは上記の公式の3番目の例を示すものとなります。私たちは実際に知りたいのは、二つの大きな集団の平均値が指定された量だけ異なる確率比です：この場合、その差は変動量のモジュラスの25倍に相当します。このような差がある確率比は数兆分の一程度です。

21. このように解くことができる直接問題の数は非常に多いが、ここでは前段落の議論を前提とする主な逆問題に限定する必要がある。その問題は次の通りである。一つのこれらの測定または観察のグループの中から少数だけが与えられた場合、それらに対して何ができるか？つまり、最終的にはどの平均値の周りに集まることが予想されるかを決定できるか？

大量のデータが与えられれば、その中心位置は確実性とともに常に増大して明らかになる：しかし現在我々は少数しか手元に持っていないと仮定し、たとえば六つ程度であり、追加する手段がない状況を考えている。

言い換えれば、グラフィカルな例を用いて表現すると、初心者や論理的な学生にとって最も良い方法であるが、直接問題では、決定要素から頻度曲線を描くだけである。すなわち、中心の位置とモジュールの数値的な値である。

一方で逆問題では、少なくとも三つの要素を決定しなければならない。なぜなら、我々は（1）以前と同じように、中心がどこに存在すると仮定するかを決定し、（2）以前と同じように、この中心周りの分散度合いであるモジュールの値を決定する必要があるからである。

これで私たちの知識は完璧ではない。これらの二つの要素のいずれも確実には与えられていないため、確率論において常に必要なものは、それらの可能性の真偽に対する推定である。つまり、これらの要素の値についてできるだけ良い推定をした後でも、その推定に犯された「可能性の誤差」を数値的に指定することも必要である。

確率論においてはこれ以上のことは達成できないが、それ以下のものは我々に提示されるべきではない。

22. (1) この最初の問題に対する答えは非常に単純です。測定回数や観察数が少ないか多いかに関わらず、その平均値が私たちが求める点であるという仮定を立てなければなりません；つまり、少数の平均が最終的な平均と一致することを仮定します。これが実際的に唯一の可能な仮定です。もちろん、観測値が一つしかない極端な場合でも、その一つを単純に採用するべきであり、観測値が増えれば信頼度は徐々に高まります。したがって、このようなデータに基づく知識と完全なデータに基づく知識との間の唯一の違いは、得られた結果自体ではなく、それを抱き込む信頼度にあると言えます。

23. (2) 二つ目の問題に関しましては、つまり平均周りのモジュールまたは分散度を決定する問題についても、ほぼ同じことが言えます。すなわち、完全情報を持っている場合と同じルールで平均平方誤差（E.M.S.）を定めるためのモジュールを指定するために採用します。あるいは正確には、473ページに示された一つのルール、つまり第二のルールに限定されるべきです。これは、仮定により、第一の前提知識がなければならず、第三の条件が満たされないからです。すなわち、平均からの誤差を計算し、その平均平方を求めることが必要であり、この二倍は確率曲線のモジュールの平方に等しいとされます。

24. (3) 第三の問題は、その解法に少し高度な数学が必要とされるが、結果を示すのはそれほど難しくない。私たちが受け入れた少数の平均値が真の平均値と一致する確率について述べる一般的な方法としては、「私たちは、推定された平均値と真の平均値との差異が超える可能性が高い限界をどのように設定すべきか」を知りたいという言い方がある。しかし、これは正確な一致とは当然無数に大きな確率で一致しないので、言葉遣いとしては不正確である。実際に行うことは、「推定誤差」を設定することである。すなわち、推定された平均値と真の平均値との差異が超える可能性が高い限界を設定することである。具体的な数値例として、壁の測定でテープを使用し、その平均値が150フィートであったとする。慎重な測量士は、このような補正を加えてこの結果を提供するだろう。「推定誤差3インチ」。これは意味しているのは、真の値が150フィートであると仮定しても、我々がそのような場合の半数（この種類の場合）で実際には三インチ以内に正確な値に近いという信頼性があるということである。

推定誤差を表す式は、モジュラスの単純な倍数である。それはモジュラスと0.4769...を乗じたものである。それがモジュラスや平均平方誤差（E.M.S.）の関数であることは、当然考えられる。つまり、測定値の誤差が大きいほど、すなわち観測された測定値間の不一致が広いほど、平均値を決定する信頼性は低くなるからである。しかし、もちろん数学がないと、このような数値的な設定は不可能である。

25. したがって、結論としては、少数の観測結果しかない場合と無限の観測結果がある場合では、確率分布曲線を決定する方法は同じであるが、その受け入れられる信頼性は低い。

ある種類の測定を少量行った場合に得られる知識は、無限に多くの測定を行った場合に得られる知識と、程度の違いしかないほど異なっていません。同じ種類の事実は知られていますが、それらについての確信度は低いだけです。しかし一方で、少量の測定値の平均によって得られる知識は、単一の測定から得られる知識とは性質が異なるものです。私たちが想定した射手の場合、彼の的を射抜いた位置が後に取り除かれました。もし彼がただ一度だけ撃った場合、私たちは彼が狙っていた点についての確信度は低くなりますが、彼の射撃技術の質や、次に射る弾丸が前回よりもどの程度遠く飛ぶかを推測する手段も全くないでしょう。しかし複数の撃ち方がある場合、私たちは中心点の位置についてより確信を持ちつつ、将来的な撃ち方がその示唆された場所周辺にどのように集まるかについても知識を得ることができます。彼の射撃技術は結果からすぐに明らかになります。

26. これまで述べてきたように、便利法則が二項型であると仮定してきました。このタイプを詳しく議論する理由はいくつかあります。まず第一に、これは広範な現象において実際に一般的であり、またはそれに近いものである唯一のタイプです。また、その表面上の複雑さにもかかわらず、実際には最も扱いやすいタイプの一つでもあります。これは、便利法則から得られる平均値によって導き出される曲線が常に同じ型を繰り返すためです。平均値の曲線は単一要素と異なるのは、そのモジュールが小さいことだけです；そしてそのモジュールは、一つの要素のモジュール（c）からn個の単一要素を平均化したものを得たもの（c/√n）で与えられる比率が非常に簡単であることに注意してください。

しかし、平均の理論を理解するためには、他のケースも考慮しなければなりません。それでは、最も単純なケースを考えましょう。特定の限界内にあるすべての値が等確率で起こる場合です。これは抽象的な確率論においてはよく知られたケースですが、自然現象においてはそう頻繁に見られません。これは選択肢の対象に対して直接ランダムに行動するときの状況であり、例えば、 logarithm の表から桁をランダムに選ぶときなどに見られます。

読者がそれを好むのであれば、労力をかけずに、p. 476で採用した方法と同じように、二つまたは三つの結果の平均値を計算することができます。この場合、我々が開始しなければならない「確率曲線」は当然、単なる有限な直線となります。質問を単純な組合せ問題として扱うと、その直線を等間隔の点で分割し、それらを二つずつすべて可能な方法で組み合わせることができます。これは数ページ前に議論したケースと同じようにです。

もしそうすれば、我々が見出すものは次のようになります。二つの結果の平均値を考えた場合、「平均の確率曲線」は初期の直線を底とする三角形となります。つまり、この平均化プロセスの最初から、最終的な平均値や中央点が最も確実な結果となることがわかります。三つ、四つ、そしてさらに多くの結果の平均を考えた場合、二項分布法が現れることが見られます。指数曲線の熟知した鐘型は徐々に近づき、最終的にはその形からほとんど区別がつきません。

27. したがって、大量のデータを扱う平均値の場合、原始的なデータにおいて二項分布法を前提とすること自体は一般的には必要ないという結論に達する。平均値に対応する派生曲線の配置の法則、すなわち、その平均値に対する派生曲線の配置は、それによって著しく影響されることはないと考えられる。従って、以前の場合と同様の計算装置を用いることが正当化される。最初の平均値を真の中心または最終的な平均値の確率的な位置として取り、その真実から特定の距離以内である確率を推定するために「平方平均誤差」を計算し、この同じ要素を使って派生曲線のモジュール、つまり収縮量や分散量を決定する。

他の多くの種類の機会法則についても同様の一般的な考慮事項が適用される。大まかに言って、すぐに特定の例外について検討に入ることになるが、原始的な配置（すなわち単一結果の配置）に関係なく、派生結果（すなわち平均値の配置）は中心周辺に密に並ぶ傾向がある。これは既に指摘した組合せの特性から導かれるものであり、極端な値を得るには複数の極端値の繰り返しが必要である一方で、中間値を得るには中間値の繰り返しまたは相反する極端値の相互作用が可能であるからである。原始的な分布が中心を中心に対称的であり、誤差の限界が有限であるか、無限である場合でも平均からの頻度の減少が非常に速い場合、平均値を取ることによる結果は単一結果に頼るよりも良いものとなる。

28. 今、一例を注意深く見ていきます。これは頻繁に実際に起こる可能性が高いものではなく、その考慮は、平均値の結果が個々の値の結果よりも「良い」と呼ぶ意味について、より詳細に問い直すことを強制するためです。図表はこの困難な点を言葉や記号による説明よりもよく理解させます。

[図：非ガウス分布からの二つのサンプルの分布]

黒線は容易に文で述べられる誤差の法則を表し、その後で実際的に起こり得ると見なされるものとして考えることができます。これは、Oから一定距離以内、つまりAとBまで、誤差の確率がOからの距離とともに一様に減少する状態を表しています。しかし、これらの点を超えてEとFまでの間では、誤差の確率は定数で保たれます。点線は、前記二つの平均値を合わせて得られる結果の法則を表します。この後者の法則が前者よりも「良い」かどうかです。後者では確かに大きな誤差は稀になり、中間的な誤差は頻繁になりますが、一方で小さな誤差も稀になります。これが全体として改善であるかどうかは、再考が必要です。

29. 前節で議論したすべての例において、平均値の曲線が個々の結果よりも優れていることは、各点で明らかでした。大きな誤差は少ないし、小さな誤差は多い；ただ一つの中間的な点では二つが等しくなるだけであり、その点には特別な意味や重要性はありませんでした。したがって、一般的な関係に含まれる様々な例を分析する必要はありませんでした。一方が他方よりも良いというだけで十分で、目的のために「モジュール」をこの優越性の尺度として採用すれば良かったのです。実際、平均結果の平均値は個々の結果の平均値よりも良いと単純に述べることができます。

しかし、ヒュームが「精査の精神」と呼ぶところを進めて、「なぜ」平均値に信頼を置くことが十分であるかを問い始めるとき、これまで提示された理由以外にも、誤差の結果についての一部の前提条件が必要であることを発見する。これは、誤差の「損失」または「害」と呼ばれるものを評価することを含んでいる。それは大きな誤差と小さな誤差が、その結果として害悪な影響を持つという点で大まかな立場に等しいと仮定し、しかしその悪影響は誤差の大きさよりも速い比率で増加すると考えていた。

30. 比較のために、誤差の重要性が誤差の大きさに直接比例する場合（もちろん、長期的には正の誤差と負の誤差が相殺されることを仮定します）においては、平均値を取ることで何らの利点も得られないことが明らかになります。このような状況は単なる購入や売買の場合によく見られます。

例えば、ある人が一日に5回、連日多数回、一定価格で布を大量に買うとします。それぞれの回で粗略な測定を行い、結果として生じる誤差が過大か欠如する可能性があると想像してみてください。各個々の布片について別々に支払いを行う場合と、5つの推定された長さを合算し、その平均値に基づいて各個々の布片に対して料金を請求する場合とでは、わずかな違いがあるでしょうか？

後者の場合、それぞれの布片に対する誤差は長期的には前者よりも小さくなることが明らかです。しかし、これは重要な影響を与えるのでしょうか？答えは、取引双方にとって全く影響がないということでしょう。長期的には同じ人々が関与しているため、中間の誤差が大きいか小かに関わらず、結果には何ら影響を及ぼさないからです。

もちろん、これほどことは一般のルールとして考慮するべきではありません。測定を行うほとんどの場合においては、大きな誤差が小さな誤差よりも著しく有害であることが多く、その比率は誤差の単なる大きさよりも大きいことがわかります。すなわち、小さな誤差より大きな誤差の方が、その影響は大きく、その影響度は誤差の大きさに比例します。

すなわち、購入や販売においても、異なる購入者が関わる場合でも、このことは成り立ちます。過剰に測定された方の喜びが、欠測された方の不満を上回ることはほとんどありません。また、科学的な測定の場合には、大きな誤差は単純に致命的であり、その回避が合理的な見込みがない場合、測定自体を行う価値すらないことがあります。

31. 実用的な観点だけを考えるなら、ここまでの議論で十分かもしれません。しかし、平均値の取り扱いが特定の目的（ある指定された量を推定すること）に対する完全な論理的意味を実現するためには、図493に示したような例外的なケースについてより詳しく検討する必要があります。そこで我々が仮定したのは、非常に小さな誤差が頻繁である状況であり、しかし一度一定の小さな範囲を超えると、他のすべての誤差は相当な限界内で等しく起こる可能性が高いという状況でした。

想像するには難しくない例を挙げて、その事態を適切に説明することはできる。最悪の場合でも、多少の遠虑が欠けるように思われるかもしれない。

そこで、イギリスにある会社が急な注文を受け、遠く離れた場所の顧客向けに機械の一部、例えば蒸気エンジンを供給することを考えよう。そして、その仕事は寸分の差で正確であることが絶対的に必要であり、それ以外では全く役立たないものだとする。しかし、二つの仕様書が送られてきたことに注意を向けよ。これらの仕様書は異なる測定に基づいており、一つには必要な部品の長さが60インチと記載され、もう一つには61インチと記載されている。

普通の誤差法、つまり二項分布型であれ何であれ、その会社は非常に悪い仕事を最善に生かすために60インチ半の部品を作ることを選ぶだろう。つまり、この方法で必要な寸分の差を満たす確率が、どちらの仕様書も正確に作成した場合よりも高いと考えられる。しかし、我々の図表で示されたような法則であれば[11]、そうすることにより必要な狭い範囲内に収まる確率はより低くなると推測される。

単なる確率的な観点から言えば、つまりこのような見積もりが何度も実行されることを想定すれば、二つの相反する測定値のいずれかを選んで正確に作成することの方が、その平均値を使用した場合よりも失敗する回数は少ないだろう。

この事例からさらに考えさせられることがある。平均値の取り扱いについて更なる反省が可能である。次に別の例外的なケースを考えてみようが、それには先ほど述べたような異なる考慮点が含まれている。また、例を用いて導入するのが最も便利である。

32. そこで、二つのスカウトが敵の要塞で砲弾の口径を測定するために送られたと想像してみよう。明日占領される予定であり、敵に対して使用されるため、砲弾や-shellの供給が必要だと考えられる。そして結果として、一つのスカウトは8インチ、もう一つは9インチと報告した。

これらの二つの値の平均である8-1/2インチが個々よりも確率的に高い値であると考えるべきだろうか？

答えはこれのように思われます。もし、可能な口径が連続的な量の性質を持つと仮定する理由がある場合、すなわち、ある限界内でのすべての値が適切なものとして考慮されるべきであると仭する理由がある場合（これは、観測や量から観察されたものや測定されたものへ通常の逆手の推移において常に行う仮定です）であれば、平均を最も可能性が高い値として選択することが正当化されます。しかし、反対に、整数インチが常にまたは一般的に好まれると仭する理由がある場合（実際には現在重い砲弾でそのような状況が存在します）、二つの推定のいずれかを無作為に選び、それだけを信頼することがより良い選択となるでしょう。

33. ここでの行動原理は次のように述べることができます。直接的な計算または表示プロセスにおいて、「誤差」の可能性が連続的であると仮定する（つまり、すべての中間値が可能であるという）アルゴリズム式や図表の場合と同じく、逆手のプロセスで原点の確率的な位置を決定する際にも、その位置は特定の限界間のどの点でも成立することが可能なと仮定します。このような例（上記の例）では、位置の可能性が不連続であることが分かっている場合、平均値の価値は完全に破壊されるかもしれません。

上記の例では、砲弾の口径が英インチや他の認知された単位で推移すると考えられることが想定されていました。しかし、砲兵隊が中国または日本にあり、その地での長さ基準について何も知らない場合、この原理を援用することはできなくなります。おそらくこれらの口径は連続的に変動する量の性質ではないと確実であることは高い確率でありますが、実際の採用された基準について全く不知の場合、その口径が連続的な量の性質であるかのように扱うことに実質的な違いはありません。このような状況では、平均値に頼る objections はもう通用しなくなるでしょう。もし我々がたった一つの機会や非常に少ない機会しかない場合、慣習的な方法を遵守することが最も良い選択となるでしょう。

34. しかし、さまざまな対象物の大量の測定値を集積し比較することができるとき、その測定対象が不連続であるという考慮は再び重要となります。すなわち、これらの対象物が有限個の異なる量に限定される傾向があることを示すよりも、すべての中間的な値に対等に適応することはありません。実際、十分な数の測定可能な対象物があれば、その対象物が作られた基準を非常に確かな確率で導き出すことができます。

この問題は、Flinders Petrie 氏が『推論的メトロロジー』という著作において鋭敏さと努力を持って攻めているものであり、実際、その学術的な興味から見ても、確率の研究者にとって推奨されるべき作品です。この論理に基づく主な原理はこれら二つです：（1）すべての技術者は、測定単位を基にした整数または単純分数の作品を作成する傾向がある；そして（2）これを確保するために、誤差法則に従ってそれから逸脱することが許容範囲内であるとします。これらの二つの仮定の結果は、例えばエジプト神殿などの古代建築物の異なる部分や比率の大量の測定値を集積した場合、特定の基準尺度またはその倍数・分数に一定の点で測定値が集積する傾向があることを示します。これらの点は古い建設者が使用していた基準長さまたはその倍数・分数を表しています。

この方法が信頼性のある結果を導き出すためには、多くの実用的な考慮事項を考慮に入れる必要がありますが、その基礎となる主な原理は比較的単純です。

35. これまで考慮したケースは、実際には異なる適用形態における以前に注目した問題の再発現に過ぎません。第2章で述べたように、一般的な誤差法に対応する頂点を持つ曲線ではなく、二つの峰または頂点を示す単一の頂点を持つ曲線が可能なことを指摘しました。この古代建築物の測定に関する問題は、平均関数に関連して同じ問いを再び提起し、やや複雑な形で展開されたものに過ぎません。

簡単な例を取りましょう。大きな誤差が一定の近似値を持つ場合、そのような誤差が小さな誤差よりも明確に多く犯されることがより確率的に高いと仮定します。このとき、曲線の柔軟性は一般的な誤差法の曲線のように中心に向かって一つのピークを形成するのではなく、そこには谷が現れます。実際、二つの二項分布曲線を考え、その中心間に短い間隔があるとしましょう。ここで平均値を計算すると、この結果は直ちに谷を埋める方向へと傾きます；そして十分な数の平均値を取る限り、中心にピークが形成されることが予想されます。実際には、一般的な単一の二項分布曲線が再び現れるようになります。

36. この問いはすぐに浮かび上がります：そのようなことをすべきでしょうか？平均値が、その割り当てられた役割であるように自由に動き、中心に向かって物事を押し寄せることを許すべきでしょうか？この問いに対する答えを得るためには、区別を導入する必要があります。もし、その特徴的な二峰型の曲線が、実際には真の誤差曲線だったとしたら—つまり、それが実際の中心への射撃偏差を表していたとしたら—結果は私たちが望むものになります。それは、本来不利な状況でも平均値を取り用いることで得られる利点の一例となるでしょう。適切な例を提示することは容易です。

例えば、ある人が庇のある場所から標的に向けて射撃する場合を考えましょう。しかし、その射程が広い露出した谷の上か下を通る強風が常に吹いているような状況だとします。もし射撃結果が観測されたら、我々は中心線の左右に二つの中心点を中心に群れて配置されていることを発見するでしょう。そして、その結果を曲線としてプロットした場合、我々が説明したような二峰型の曲線が得られるはずです。しかし、風が等しく強力で反対方向から吹いていると仮定すると、平均化過程によって生じた混乱が修正されることになります。

もし曲線が、自然現象の結果として描かれる場合、それが明確にそれほどよく起こる傾向があるように見えるなら、それは異なる可能性を示すでしょう。たとえば、二つの非常に異質な種族からなる人々の多数を測定した結果について考えてみましょう。そのような場合、曲線は45ページで示されているような形になります。もし二つの混血種族の人数が等しいなら、それは双子の峰を示すでしょう。

また、「一括して平均化すべき範囲」に全体を含めるべきかどうかという問題が生じます。答えは、その目的によります。我々が一貫して比較したい異質な種族全体と他の何か（またはそれ自体の別の時間）との間で比較する場合、分析せずに平均化するのが良いでしょう。既に述べたように（47ページ参照）、人口統計学全般は多くの区別可能な特性を無視しなければなりません。しかし、この異常な誤差曲線の原因を解明したいという目的があるなら、統計を対応する部分に分割し、個々の部分について分析を行う方が良いでしょう。

同様に、古代建築物の測定値を考えます。この場合、すべての「誤差」を集めた一つの統計グループを作ると、結果的な曲線は二つの峰だけでなく、それらが頻度に応じて異なる大きさを持つ一連の峰を示します。それぞれの特定の測定値が発生する頻度に対応したこれらの峰があります。全体として平均化することは可能ですが、そのような作業はほとんど何の合理的な目的も達成しないでしょう；一方で、各々の最大頻度の点には何か重要なことを教えてくれるはずです。

37. 最後に、一つの特殊なケースを挙げます。明確に非対称的または偏った曲線の可能性を考えましょう：

[図：非対称（偏った）分布]

誤差の法則、これはそのグラフィカルな表現であるが、私はこのことは稀ではないと考える。問題となっている曲線は、前章で言及した気圧高度の曲線を僅かに過大評価したものに過ぎない；そこで説明したように、平均値、中央値、最大縦座標は互いに異なることを示す。ここでの疑問は、前の例と異なり、単一の平均値が取られるべきかどうかではなく、どのような種類の平均値を選択すべきかである。前と同じく、答えは我々が目指す特別な目的により依存するものである。すべての一般的な比較のための目的においては、どの平均値でも適切であり、それらの中で算術平均や中央値を選ぶのが自然である。

38. しかし、他の種類の平均値が正当化される可能性がある状況も存在し、これらについての一時的な説明で本稿を締めくくることができる。例えば、ここで問題となっていることが気候の適度性に関するものだと仮定する。通常の算術平均は極端な値の数と大きさに大きく依存しているため、単純に最も頻繁に現れる値で判断することの方がより適切であると考えられる：つまり最大縦座標による判断である。また、この特定の値が最も適切かつ意味深いものとなる視点は、他にも提案できる。

前記の場合、すなわち天候曲線の場合には、客観的な「真実」の値を目標としていなかった。しかし、気象学者ガロント氏によって示された特定の評価クラスを代表する近似した曲線が存在し、彼はこの場合、幾何平均が唯一適切なものであると指摘している。この場合、容易性の曲線はO点で突然終了し、一般的な指数形の短縮版に類似している。その特性については、すでに引用されたマカリスト博士の論文において議論されているが、ここでのその性質を調査しようとすることは過度に複雑な詳細に陥るだろう。

39. このすべての結論から、前章で述べた平均値の性質と機能とは非常に一致しているように思われる。各平均値は、実際の複数の値を一つの代表的な中間値として置き換えるものである。したがって、これらの後者の情報の大半を放棄しなければならない。しばしば、特に一般的な測定では、その一つのことを表すものは明らかに求めているものであり、その場合、問題はどの平均値が求めている「真実」の値と最も一致するかにあるだけである。しかし、統計学の一般的な応用において、現象の後ろには本当に「真実」の値という客観的なものが存在しない場合、問題はさまざまな方向に分岐する可能性がある。目的が多様であればあるほど、どの平均値が最も適切であるかについての判断が必要となる。したがって、我々が手元にある具体的な目的に対して、馴染み深い算術的平均が適切かどうかを疑問視する場合は、その目的自体を明確に定義することが最初のステップとなる。

1. マンフィールド・メリミン氏は1877年（『コンネチカット学術雑誌』）で、最小二乗法に関する408件の著作リストを発表した。

2. つまり、射撃マークの「重心」を「等しい重さ」として考慮し、その中心を取る。実際には、点群の重心を得るための基本的な幾何学的構築が示すように、これはすべてのマークの平均値であるが、一般的にはそう考えられていません。当然ながら、測量におけるような場合で、三つ以上の不一致な方程式を組み合わせたり整合させたりする必要があるとき、最小二乗法のような規則は不可欠となる。単なる平均化では困難から脱出することはできない。

3. この例外的な形状の仮定の唯一の理由は簡潔さを確保することである。二次元での誤差を許容する一般的なターゲットでは、やや複雑な結果が得られる可能性がある。

4. この方程式の一般形は最初に述べられ、（v. p. 29）。ここでは、特殊な形式が与えられており、ここで A として (h/√π) を代入することで、確率論において一般的に使用される。これは、y dx の積分が +∞ から -∞ までで 1 になるためである。つまり、互いに排他的な可能性の総和は通常通り 1 で表される。この表現形式では h は x(-1) 程度の量であり、hx が数値的な指標として存在することからである。その絶対値（モジュラス）はこれの逆数なので、誤差自体と同じ程度の量となる。実際、それを 0.4769... に掛けると、「確率的誤差」と呼ばれるものになる。

5. この説明とグラフィカルな方法については、p. 29 の注を参照してください。

6. 幾乎すべての場合において、上記の議論は実際的な誤差の限界があるどんな誤差頻度の法則にも適用されます。もし誤差の可能性が制限がない場合、平均値の特徴である中心に結果を集積する傾向は状況によって異なります。指数曲線の場合のように、基底への漸近的近似が非常に速いとき、つまり極端な誤差が比較的小さいときには、その特性は依然として成り立ちます。しかし、デ・モガンが示唆しエドゲワース氏が指摘した方程式（y = π/(1 + x2)）のような法則を平均化の結果は中心に集積する傾向を減少させることがあります。（\_Camb. Phil. Trans.\_ 第 10 卷 p. 184, 第 14 卷 p. 160）

7. この公式群その他の類似な公式の証明については、Galloway の『確率論』および Airy の『誤差論』を参照してください。

8. この場合で一般的に使用されるE.M.S.の公式は、$\frac{\sum e^{2}}{n - 1}$であり、$\frac{\sum e^{2}}{n}$ではない。この差異は微々たるもので、$n$が小さい限り無視できるほどである。その正当化は、真の中心からの平方和が最小となる（その中心が最終的な算術平均である）から、誤った中心からの平方和は多少大きくなることにある。

9. 私には、厳格な論理的正確性において、前二節で述べた両方の中心の指定における可能性のある誤差を知りたいと思われます。しかし、この質問について議論し、それに適切に対処できる深い数学的な学者たちは、主に天文学の実用的な要請のために書かれたものが多く、そのためには真の値を求めることのみが重要な要素であることが十分です。従って、一般的に与えられる唯一のルールは平均の可能性のある誤差に関するものです。

10. すなわち、直接それらを使用するのではなく、間接的に使用することを指します。この後者の方法は、「ランダム性」の章で説明されているように、非一様分布を与えることがあります。

11. この法則に非常に類似した状況を想像することは難しくありません。たとえば、二つの測定のうち一つが精密な技術者によって行われ、もう一つは手間を省くために乱数で決定されたもの（一定の範囲内）であると仮定し、その会社はこの事実を知っていますが、当然著者は不明であるとすると、上記のような誤差分布を得ることができます。

インデックス。

事故 342

アイリー、G.B. 447, 484

暗黙の予想 287

アーバンスナット 258

アリストテレス 205, 307

平均
  算術的 437
  几何学的 439
  中央値 442
  平均に関する結果 482
  必ずしも得られない結果 457
  平均の使用 439, 489

バベージ 343

袋と玉 180, 411

信頼性
  正確性 125, 131, 178
  信頼性の段階 139
  信頼性の成長 199
  言葉による表現 143
  量測法 119, 125, 146
  量 133
  信頼性のテスト 140, 149, 294
  不正な信頼性 129
  悩ましさ 127

ベントハム 319, 323

ベルヌーイ 91, 117, 389

ベルティリオン 435

出生、男性と女性 90, 258, 263

ボートレース、オックスフォードとケンブリッジ 339

ボール 183

バクエル 237

ブフン 153, 205, 352, 389

ブルーグァルディック 311

バトラー 209, 281, 333, 366

カリスラーティの表 169

偶然、その意味 245

因果関係
必要性 237
証明 244

重心 467

確実性、法において 324
合理的な確実性 327
仮定的な確実性 210

偶然
と因果関係 244
創造 258
設計 256
天才 353
小さなものを無視する 363
選択 338

シャーヴェンエット 352

分類、数値的なスキーム 48

偶然の一致 245

組合せと順列 87

共産主義 375, 392

概念論 275

偶然性の衝突 418

消耗者保険 227

コルヌ 245, 255, 338

クラッカントホープ 312, 320

クレッグ、J. 192

クロフトン、M.W. 61, 101, 104

但丁 285

偏位
原因 57
射撃からの偏位 38

デ・モラン 83, 106, 119, 122, 135, 177, 179, 197, 236, 247, 296, 308,
350, 379, 382, 483

デ・ロスの裁判 255

桁、ランダム 111, 114

不連続性 116

ランダムな分布 106

図表 29, 45, 118, 443, 476, 481, 493, 501

辩证法 302, 320

ドンキン 123, 188, 283

寿命の持続時間 15, 441

デュシング 259

エーベリングハウス 199

エッジワース、F.Y. 34, 119, 256, 339, 393, 435, 483

エリス、L. 9

流行病 62

誤差、法則 29
非対称的な誤差 34, 441, 443
二項分布の誤差 37, 457, 469, 480
幾何学的誤差 34, 502
混合した誤差 45
生産された誤差 36

誤差
平均的な誤差 446
可能性のある誤差 446, 472, 488
平方平均誤差 447, 488

逃走、狭い 341

期待値、道徳的 388

経験と確率 74

指数曲線 29

異常なこと
その感覚 159, 423
物語 407, 421

論理と確率における誤謬 367

宿命論 243

フェッナー 34, 389, 435, 441

振動 448
無制限の振動 73

フォーブス、J.D. 188, 262

形式的論理 123

形式的なと実質的な扱い 86

自由意志 240

ガロウェイ 248, 448, 484

ギャルトン、F. 33, 50, 70, 318, 442, 451, 473, 502

賭博
と保険 370
不利性 384
最終的な結果 385, 391

ゴッドレイ、H. 99

グロート、G. 307

ガイ 6

ハンモント、W. 266, 297

人類の幸福 382

表と裏 77

遺伝 50, 357

Herschel 30, 466

ホーディン 361

休谟 236, 419, 433

仮説 268

直感的な推論 121

独立な事象 175, 246

帰納法
と確率 194, 201, 208, 233, 358
困難性 213
純粋な帰納法 200

財産の不平等 382

推論、ルール 167

接種 374

保険
  理由の正当化 149
  難しさ 221
  生命保険 151
  特別なケース 224
  理論 372
  品種 374

逆確率 179, 196, 249

不規則性、絶対的および相対的 6

ヤコブス、J. 199

ジャクソン、J.G. 253

ジェボンズ 37, 83, 136, 198, 201, 209, 247

カント 310, 317

ケッコルマン 298, 316

種類、自然の 55

クルグ 324

ラムベルト 309

確率の言語 159

ラプラス 89, 120, 197, 237, 424

法則
  存在しない法則 101
  実験的法則 160
  原因の法則 206

最小二乗法 41, 467

ライプニッツ 309, 320

手紙
  遺失した手紙 162, 368
  誤送された手紙 67, 237, 241

レックス、W. 263, 441

限界
  概念 18, 109, 164
  可能な振れ幅の限界 32

ランダムな線 113

リストラの方法 187

lotteries 128

ルン、J.R. 248

等確率 77, 183

マキャリスター、D. 34, 187, 502

マンセル、H.L. 299, 301, 320

マルティネール 343

実質論と形式論理 265

最大縦軸 441, 455

測定
  信頼性 119
  記憶 192

知覚的特性の測定 49

メリマン、M. 352, 448, 460, 465

ミル、J.S. 131, 207, 266, 282, 402

メリートン、偶然の生成 353

奇跡 428

ミッチェル、J. 260

モダリティ 295
  分類 307
  誤ったもの 297
  形式的なもの 298
  法律におけるもの 319

モジュラス 464, 472, 484

モンロー、C.J. 325, 416

名前、参照 270

民族の比較 51

自然種類 55, 63, 71

必然的で不可能な事柄 310

対象と手段 53

オッカム 314

ペイリ、F. 498

ピースバーグの問題 19, 154

ポアソン 405

プラント 311

法的推定 329

プレヴォスト 348

確率
  定義 165
  相対的なもの 290
  積分的なもの 463

可能性の値 441
  エラー 446, 472

問題、三点型 104

プロクトール、R.A. 262, 378

占い、自殺的 226

神の裁き 89, 431

命題、比例的なもの 2

心理的研究 256

大ピラミッド 251

π、桁数 111, 247

四分位数 446

クエテレット 23, 30, 43, 91, 259, 330, 348, 454

ランダム性
  語源 96
  射撃におけるもの 98
  証明 107

稀な事象 349

現実主義 92

十分な理由 82

残差 460

ローブ
  ランダムに折られた棒 98
  ランダムに投げられた棒 103

規則
  归納的と推論的なもの 176
  続きの規則 191
  冲突する規則 222
  複数の規則 217

系列
  決定的な比 11
  固定的と変動的なもの 16
  理想的なもの 95
  特別なもの 12

シャンクス 248

スキート、W.W. 96

シムゲリウス 306, 316

スミス、P. 251

社会主義 392

スピリチュアルズム 365

星々、108260のランダム配置

統計学
  類推による 473
  不意識的な参照 400

統計誌 6

身長
  人間 25, 471
  フランスとイギリス 44

ステファン、J.F. 282, 323, 326

スターティ、D. 209, 237

主観的・客観的な用語 160

順序
  長期的な 360
  法則 190, 362

スフィールド、G. 248

自殺 67, 237

姓の絶滅 387

驚き、感情 157

三段論法、純粋と模態的 316

テイラー 329

証言
  単独の 411
  結合した 426
  二種類の 409
  値打ちがない 416

トムソン、W. 153, 314, 419

時間
  影響 191
  確率における 279

トドハンター 415

トントネ 380

三角形、ランダムな 103

タッカー、A. 127

種類
  存在する 42, 60, 453
  固定と変動する 64, 93

ウエルベーグ 311

生活における不確実性 370

均一性 240

計算の単位 464

自主的な行為 65, 68, 85

ワットフォード 374

ウォリス、J. 312

ウォーソン、H.W. 387

ホーレイ 297, 307

ウィス 401

ウィトウースト、W.A. 87, 183, 384

ウィルソン、J.M. 104

証人、独立した 405

ウォフ 309

ウールハウス 101

ケンブリッジ：カイ・ジェイ・クラーク、M.A.と息子たちによって大学印刷所で印刷された。

プロジェクト・ガーデンバーグの第3版「確率論的論理」終わり

*** このファイルは57359-0.txtまたは57359-0.zipの名前で保存されます ***
この文書と関連する各種形式のファイルは次の場所にあります：
        http://www.gutenberg.org/5/7/3/5/57359/

更新された版が発行されると、古い版は名前を変更します。

印刷版の著作権が米国で保護されていない作品を作成すると、これらの作品には米国の著作権はなく、因此に財団（そしてあなたも）は、許可なくこれらの作品を複製・配布でき、著作権料を支払う必要はありません。このライセンスの「一般利用条件」部分に記された特別な規定は、Project Gutenberg-tmの概念と商標を保護するために適用されます。Project Gutenbergは登録商標であり、著作料で電子書籍を販売することは許可されていませんが、特定の許可を得た場合を除きます。この電子書籍を有償で提供しない限り、規則に従うことは非常に簡単です。この電子書籍を使用して、派生作品の作成や報告書の作成、パフォーマンスや研究など、ほぼどのような目的でも利用できます。米国で著作権法によって保護されていない電子書籍は、印刷し、配布し、無料で与えることができます——アメリカ合衆国内では電子書籍に対してほとんど何でも行うことができます。再配布は商標ライセンス、特に商業的な再配布に特に注意が必要です。

START: 全体の著作権許諾

Project Gutenberg-tm著作権許諾全文
この作品を配布または使用する前に必ずお読みください

Project Gutenberg-tmの著作物の無料配布を促進することを目的とする財団の使命を保護するために、この作品（または「Project Gutenberg」に関連するいかなる作品も）を使用したり配布したりした場合、あなたは本ファイルやwww.gutenberg.org/licenseで利用可能な全文のProject Gutenberg-tm著作権許諾のすべての条項に準拠することに同意します。

第1条. 一般利用条件とProject Gutenberg-tm電子著作物の再配布

1.A. このProject Gutenberg-tm電子作品の一部を読んだり使用したりすることにより、あなたはその著作権および知的財産（商標・著作権）に関する契約のすべての条項に同意し、理解し、受け入れたことを示します。このライセンスおよび知的財産（商標・著作権）に関する契約のすべての条項に同意しない場合、あなたはProject Gutenberg-tm電子作品の所有を終了し、所持している全てのコピーを返却または破棄する必要があります。もし著作料を支払ってProject Gutenberg-tm電子作品のコピーを取得したか、アクセス権を得たが、この契約の条項に従うことを同意しない場合、あなたは第1.E.8条に定められたように支払いを行った人物または機関から返金を求めることができます。

1.B. "Project Gutenberg"は登録商標です。この商標は、本条約のすべての条項に同意した人々によって電子作品と関連付けられることのみが許可されています。多くのProject Gutenberg-tm電子作品に対しては、この全体的な契約を遵守せずにいくつかの行動を行うことができます。第1.C条をご覧ください。この契約の条項を遵守し、Project Gutenberg-tm電子作品への自由なアクセスを保護する手助けをする場合、Project Gutenberg-tm電子作品に対する多くの活動が可能です。第1.E条をご覧ください。

1.C. ジェンキンス文学アーカイブ財団（以下、「財団」または「PGLAF」という）は、プロジェクト・ゲッutenberg-tm電子作品の集約著作権を所有しています。集合に含まれるほとんどの個々の作品は、アメリカ合衆国において公的使用期間が切れています。アメリカ合衆国に所在している場合、個々の作品がアメリカ合衆国の著作権法によって保護されていないなら、あなたがその作品を複製、配布、上演、表示、または派生著作物を作成することを阻止する権利は主張しておりません。もちろん、財団としては、プロジェクト・ゲッutenberg-tmの目的である電子作品への無料アクセスを促進するために、この協定の条件に従ってプロジェクト・ゲッutenberg-tmの作品を自由に共有することで、プロジェクト・ゲッutenberg-tmという名前が関連付けられるように希望しています。あなたは、この作品を無償で他者と共有する際には、常にその作品と共に付属している完全なプロジェクト・ゲッutenberg-tmライセンスを保持することにより、この協定の条件に従うことができます。

1.D. あなたが所在する場所の著作権法も、この作品に対するあなたの行為を制限します。多くの国における著作権法は常に変化しています。アメリカ合衆国以外に所在している場合、この協定と共にダウンロード、複製、表示、上演、配布、またはこの作品やプロジェクト・ゲッutenberg-tmの他の作品に基づいて派生著作物を作成する前に、その国の著作権法を確認することをお勧めします。財団は、アメリカ合衆国以外のどの国においても、プロジェクト・ゲッutenberg-tmの作品の著作権状況について一切の表明を行っておりません。

1.E. 「ジェンキンス文学アーカイブ」への参照を全て除去していない限り：

1.E.1. 以下の文は、「プロジェクト・ゲッutenberg」という言葉が現れるか、または「プロジェクト・ゲッutenberg」に関連付けられているすべてのプロジェクト・ゲッutenberg-tm作品（そのような言葉が含まれるすべての作品）がアクセスされ、表示され、上演され、視聴され、複製され、配布されるたびに、目立つ位置で表示する必要があります。これは、完全なプロジェクト・ゲッutenberg-tmライセンスへのアクティブリンクを含むか、またはそのライセンスへの即時アクセスが提供されている場合があります：

この電子書籍は、アメリカ合衆国およびその大部分の地域で誰でも無料で利用でき、制限がほとんどありません。これをコピーしたり、無償で配布したり再利用したりする際には、この電子書籍に添付されているプロジェクト・ゲッutenbergライセンスまたはwww.gutenberg.orgでオンラインで見つけることができるライセンスの条件に従ってください。もし所在国がアメリカ合衆国以外の場合、この電子書籍を使用する前に所在国の法規を確認してください。

1.E.2. 个别なプロジェクト・ゲッutenberg-tm電子書籍が米国著作権法によって保護されていないテキスト（著作権所有者の許可を得て掲載されている通知が含まれていない）から派生した場合、その作品はアメリカ合衆国の誰でも無償でコピーしたり配布したりすることができます。もしこの電子書籍を再配布するか、「プロジェクト・ゲッutenberg」の文言が関連付けられているか表示されている場合、以下の1.E.1から1.E.7までの条件に従うか、または1.E.8または1.E.9で定められたプロジェクト・ゲッutenberg-tm商標を使用するための許可を得る必要があります。

1.E.3. 个别なプロジェクト・ゲッutenberg-tm電子書籍が著作権所有者の許可を得て掲載された場合、その利用と配布は1.E.1から1.E.7までの条件に加えて、著作権所有者がこの作品と共に最初に示すプロジェクト・ゲッutenberg-tmライセンスで定められた追加の条件に従う必要があります。これらの追加の条件は、著作権所有者の許可を得て掲載されたすべての作品の最初に見つかるプロジェクト・ゲッutenberg-tmライセンスに記載されています。

1.E.4. この電子書籍またはその一部を含むファイルから、全般的なプロジェクト・ゲッutenberg-tmライセンス条項を削除したり分離したり取り外すことはできません。

1.E.5. この電子書籍の全文またはその一部をコピーしたり表示したり実行したり配布したり再配布したりする際には、段落1.E.1に記載された文を目立つ位置で表示し、プロジェクト・ゲッutenberg-tmライセンスの全文へのリンクや即時アクセスが提供されることを確認してください。

1.E.6. この作品をバイナリ形式、圧縮形式、マークアップされた形式、非所有権的または所有権的な形式（包括的に或はハイパーテキスト形式）に変換し、配布することは可能です。ただし、著作権表示が付されていないProject Gutenberg-tmの作品を「 Plain Vanilla ASCII 」形式以外の形式で提供したり配布する場合、ユーザーに対して追加料金を請求せず、その費用も負担しないで、元の「 Plain Vanilla ASCII 」形式または他の形式の複写本を提供するか、複製または配布の方法を提供する必要があります。代替形式は、第1.E.1条に規定されるProject Gutenberg-tmライセンスを含むべきです。

1.E.7. Project Gutenberg-tm作品へのアクセス、閲覧、表示、演奏、複製、配布に対して料金を徴収することは、第1.E.8条または第1.E.9条に準拠しない限り禁止されています。

1.E.8. Project Gutenberg-tm電子作品の複製本を提供するか、アクセスを提供するか、配布する場合、以下の条件を満たす必要があります：

* あなたが既に適用税金を計算するために使用している方法を使用して、Project Gutenberg-tm作品から得られる総収益の20%の版税料を支払う。この料金は、Project Gutenberg-tm商標の所有者に負担されるべきですが、彼はこの条項に基づいて著作権料をProject Gutenberg文学アーカイブ財団に寄付することに同意しています。版税料は、あなたが（または法律により準備する必要がある場合）定期的な税金申告書を作成した日から60日以内に支払う必要があります。版税料は明確にそのようなものとして記載され、Project Gutenberg文学アーカイブ財団の指定された住所に送るべきです（第4条「Project Gutenberg文学アーカイブ財団への寄付に関する情報」参照）。

* 30日以内に、著作権の全条項に同意しない旨を書面（または電子メール）で通知したユーザーに対して、支払った金銭全額を返還する。あなたはそのようなユーザーに対し、物理的媒体において所持している作品のすべてのコピーを返却または破棄させ、Project Gutenberg-tm作品への使用および他のコピーへのアクセスを停止させるよう要求しなければならない。

* 90日以内に電子作品に缺陷があることが判明し、あなたに通知された場合、第1.F.3条に基づき、支払った金銭全額または代替コピーの返還を行う。

* この契約の無料配布条件を遵守する。

1.E.9. 異なる条件でProject Gutenberg-tm電子作品またはそのグループを料金を徴収して提供したい場合や、配布したい場合は、米国著作権法によって保護されていない作品を作成するためにProject Gutenberg Literary Archive Foundationの志愿者和員工投入了大量精力来識别、進行版权研究、転録和校対。尽管如此，Project Gutenberg-tm電子作品及其存儲介質可能仍可能存在“缺陷”，例如但不限于不完整、不准确或損坏的数据、転録錯誤、版权或其他知識産权侵权、有缺陷或損坏的磁盤或其他介質、計算机病毒或无法被您的設備読取的計算机代碼。

これらの条件を遵守するためには、Project Gutenberg Literary Archive FoundationとProject Gutenberg Trademark LLC（Project Gutenberg-tm商標的所有者）から書面での許可が必要である。Foundationの連絡先は以下のセクション3に記載されている通りである。

1.F.2 保証の制限、損失に関する明示的または黙示的な責任の放棄 - 第1.F.3条で説明される「交換または返金権」以外については、プロジェクト・ゲッutenberg文学アーカイブ財団は、プロジェクト・ゲッutenberg-tm商標を所有し、本契約に基づき電子書籍を配布する他のいかなる第三者も、あなたに対して損失、費用、包括的な弁護士報酬を含む一切の損害に対する責任を放棄します。あなたは、過失、厳格な責任、保証違反、または契約違反による補償の権利については、第1.F.3条に規定されているもの以外には何らの補償を受けないことに同意します。本契約に基づき財団、商標所有者、または配布者があなたに対して実質的な、直接的な、間接的な、結果的、懲罰的、または偶発的な損害を負う責任は一切ありません。

1.F.3 交換または返金の制限的権利 - この電子書籍を受け取って90日以内に欠陥を見つけた場合、あなたはその作品を受け取った人物に対して、作品の欠陥についての説明を記した文書を送ることにより、支払った金銭（あれば）を返還してもらうことができます。物理媒体で受け取った作品の場合、あなたは説明と共にその媒体を返却する必要があります。作品を受け渡した人物または法人は、代替版の提供を選択することができます。電子的に受け取った作品の場合、作品を提供した人物または法人は、代替的な電子配信を選択することができます。もし二度目の配布が欠陥である場合でも、問題を修正する機会を与えずに書面による返還請求を行うことができます。

1.F.4 保証の制限 - 第1.F.3条で定める交換または返還の権利以外については、この作品は「現状」で提供され、いかなる明示的または黙示的な保証も含め、商品性や特定目的に対する適合性を含む一切の保証も提供されません。

1.F.5. 一部の州では、特定の明示的保証の放棄や、特定の種類の損害の除外または制限を認めないことを許可していません。本契約において定められた放棄または制限が適用される州の法に違反する場合、本契約は、該当する州の法律により可能な最大限の放棄または制限と解釈されます。本契約のどの規定も無効であっても、その他の規定は依然として有効です。

1.F.6. 被告人化 - あなたは、以下のいずれかが発生した場合に、被告から一切の責任、費用、および経費（法的費用を含む）を免れるよう、基金会、商標権者、基金会の代理人または従業員、Project Gutenberg-tm電子作品を本契約に基づいて提供する者、そしてProject Gutenberg-tm電子作品の生産、推進及び配布に関連するボランティアを被告から一切の責任、費用、および経費（法的費用を含む）を免れるよう保証します：(a) この著作物またはProject Gutenberg-tmの他の著作物の配布、(b) Project Gutenberg-tmの著作物への変更、修正、追加または削除、(c) あなたが引き起こした欠陥。

第2条 Project Gutenberg-tmの使命についての情報

Project Gutenberg-tmは、最も多様な種類のコンピュータで読み取れる電子作品の無料配布を意味します。そのコンピュータには古式化したものから新しいものまでが含まれます。それは数百人のボランティアと、様々な分野の人々からの寄付によって存在しています。

ボランティアの支援と資金提供は、必要な援助を得られるようにするためのボランティアを支援し、Project Gutenberg-tmの目標を達成し、そしてProject Gutenberg-tmのコレクションが今後も自由に利用可能であることを確保するために不可欠です。2001年には、Project Gutenberg文学アーカイブ財団が設立され、Project Gutenberg-tmとその将来世代のために安定した永久的な未来を提供するための体制が整えられました。Project Gutenberg Literary Archive Foundationについてさらに詳しく知るためには、www.gutenberg.orgのFoundation情報ページをご覧ください。

第3節 Project Gutenberg文学アーカイブ財団に関する情報

Project Gutenberg文学アーカイブ財団は、非営利の501(c)(3)教育法人で、ミシシッピ州の法律に基づいて組織され、国税庁から免税ステータスが与えられています。財団のEIN（連邦税番号）は64-6221541です。Project Gutenberg文学アーカイブ財団への寄付は、米国の連邦法と州法により許容される最大限の範囲で税金が控除されます。

財団の本社事務所はアラスカ州フェアバンクスにあり、郵便宛先はPO Box 750175, Fairbanks, AK 99775ですが、ボランティアと従業員は多くの場所に散在しています。財団の事務局はユタ州サalt lake cityの809 North 1500 Westにあり、電話番号は(801) 596-1887です。連絡先リンクと最新の連絡先情報はwww.gutenberg.org/contact の財団ウェブサイトおよび公式ページで見つけることができます。

追加の連絡先情報：

ドクター・グレゴリー・B・ニューバイ
    会長兼理事
    gbnewby@pglaf.org

第4節 Project Gutenberg文学アーカイブ財団への寄付に関する情報

Project Gutenberg-tmは広範な公衆の支援と寄付に依存し、その使命を果たすことができません。それは機械読取可能な形式で自由に配布できる著作権フリーおよびライセンスされた作品の数を増やすことです。多くの小さな寄付（$1から$5,000まで）はIRSによる税制優遇ステータス維持にとって特に重要です。

財団は、アメリカ合衆国の全50州における慈善活動と寄付に関する法律に準拠することをコミットしています。遵守要件は一様ではなく、これらを満たし続けるには相当の努力、多くの書類作成と手数料が必要となります。特定の場所での寄付募集を行わないように、IRSによる合意確認が得られていない州では寄付を求めません。特定の州の寄付状況を確認するためにはwww.gutenberg.org/donate を参照してください。

我々は、規約要件を満たしていない州からの寄付を求めるものではありませんが、そのような州の寄付者から無断で寄付を受け入れる禁止事項はないと知っています。

国際的な寄付も心から歓迎しますが、国外からの寄付についての税制上の取り扱いについてはいかなる声明もできません。アメリカ合衆国の法律だけが私たちの小さなスタッフを圧倒しています。

現在の寄付方法と連絡先については、プロジェクト・ゲッutenberg（以下略）ウェブページをご確認ください。支払いは切手、オンライン決済、クレジットカードによる寄付など様々な形で受け付けています。寄付をするには、www.gutenberg.org/donate を訪問してください。

第5条 プロジェクト・ゲッutenberg（以下略）電子作品に関する一般情報

プロジェクト・ゲッutenbergの創始者であるマイケル・S・ハート博士は、誰でも自由に共有できる電子版書籍で構成されるライブラリを提唱しました。40年以上にわたって、彼は規約に基づくボランティアのネットワークと共に電子ブックを作成し、配布してきました。

プロジェクト・ゲッutenberg電子ブックは、アメリカ合衆国では著作権保護されていないことを確認した複数の印刷版から作られることが多く、そのため必ずしもどの特定の紙面版と一致しているとは限りません。

多くの人々はまず当ウェブサイトにアクセスします。これは主なPG検索機能が含まれています：www.gutenberg.org

このウェブサイトにはプロジェクト・ゲッutenbergに関する情報が掲載されており、寄付を支援する方法や新しい電子版の作成についての情報を提供しています。また、新刊情報を受け取るためのメールニュース購読も可能です。


\backmatter

\chapter{訳者・数理校閲後記}

本書は、近代頻度論確率論の不滅の金字塔であるジョン・ヴェンの主著『偶然の論理』（\textit{The Logic of Chance}, 3rd edition, 1888）の全19章を、Project Gutenberg（\#57359）の底本に基づき完全翻訳・数理校閲したものである。

ヴェンが19世紀後半に激しく主張した「客観的頻度論」は、単なる数学的モデルの提唱にとどまらず、因果性、帰納推論、言語の様相、人間の認知バイアス、保険とギャンブルの境界、そして標本平均の極限挙動に至るまで、極めて広大な科学哲学とデータサイエンスのフロンティアを切り拓いた。

ラプラスの古典的・主観的決定論と、ヴェンの客観的・経験的頻度論という二大巨星の全容を並び立つ全訳学術書として手元に置くことは、21世紀の現代においてデータとAI、そして確率的推論に向き合うすべての探求者にとって無上の知的財産となるであろう。

\begin{flushright}
2026年10月\\
教育情報分析研究会
\end{flushright}

\end{document}
